%!TEX TS-program = xelatex -output-driver="xdvipdfmx -q -E" "$1"
%!TEX encoding = UTF-8 Unicode
% \documentclass[a4paper,10pt]{memoir}
\documentclass[a4paper,10pt,showtrims]{memoir} % This is for testing on A4 paper
% \documentclass[b5paper,10pt,showtrims]{memoir}
% \documentclass[royalvopaper,10pt,showtrims]{memoir} % this is going to be the final
\raggedbottom
\usepackage[xetex]{graphicx,color}
\usepackage[]{hologo}
\usepackage[EU1]{fontenc}
% \usepackage{lmodern}
\usepackage{fontspec}
\usepackage[]{polyglossia}
\setdefaultlanguage{american} % apperantly for polyglossia
\setotherlanguage{sanskrit}   % ditto
\setotherlanguage{english}   % ditto
% \newfontfamily\sanskritfont[Scale=1.00, Mapping=tex-text]{Devanagari MT} % the scaling does not have effect currently
% \newfontfamily\devanagarifont[Scale=1.00, Mapping=tex-text]{Devanagari MT} % the scaling does not have effect currently
\usepackage{natbib}
%\usepackage[para*,Bplain,Cpara,Dpara,Eplain,addpageno,switch,pagewise]{ednotes}
\usepackage[para*,Bplain,Cplain,Dpara,Eplain,addpageno,switch,pagewise]{ednotes}
% \usepackage[para*,Bplain,Cplain,Dplain,Eplain,addpageno,switch,pagewise]{ednotes}
\usepackage{utf8dev} % I load this for backward compatibility.
% \usepackage{utf8cm}
\usepackage{utf8dingbats}
\usepackage{utf8bodoni}
\usepackage{utf8stix}
\usepackage{utf8applesymbols}
\usepackage{yvied}
\usepackage{ifthen}
\usepackage{mparhack} % to have \marginpar at the correct side.
\usepackage[pagebackref,xetex]{hyperref}
\usepackage{ellipsis}
% \renewcommand{\ellipsisgap}{0.4em}
\usepackage{url}
\usepackage{longtable}
\usepackage{namesandterms}
%%%
%%% END OF LOADING PACKAGES
%%%

\hyphenation{man-u-script man-u-scripts sov-er-eign-ty}

%%% LAYOUTS
% \overfullrule=5pt % be sure to comment out this line
% A4 has 297x210
% \settrimmedsize{277mm}{*}{0.666}
% \setstocksize{9.5in}{6.25in}
% \settrimmedsize{257mm}{182mm}{*}
% \settrimmedsize{260mm}{175mm}{*}
\settrimmedsize{9.5in}{6.25in}{*} % The final production value!
%%% The following lines are to have the trimmed pages in the center of the stock.
\setlength{\trimedge}{\stockwidth} % This is for testing. comment out for the final
\addtolength{\trimedge}{-\paperwidth} % This is for testing. comment out for the final
\setlength{\trimtop}{\stockheight} % \trimtop = \stockheight 
\addtolength{\trimtop}{-\paperheight} % \trimtop = \stockheight - \paperheight 
\settrims{0.5\trimtop}{0.5\trimedge} % comment this out for the final
%%% up to here
\setlxvchars
% \settypeblocksize{*}{1.07 \lxvchars}{1.618}
% \setlrmargins{*}{*}{2}
\settypeblocksize{*}{1.09 \lxvchars}{1.7320508}
\setlrmargins{*}{*}{1.7320508}
\setulmargins{*}{*}{1.4142136}
% \settypeblocksize{500pt}{1.11 \lxvchars}{*}
% \setheadfoot{14pt}{36pt}
\checkandfixthelayout
\typeoutlayout
\typeoutstandardlayout
\setmarginnotes{0.5em}{0.6in}{0pt}
%%% END OF LAYOUT

\newcommand{\st}[1]{\textbf{\large#1}}
\newcommand{\bht}[1]{\textbf{#1}}
\newcommand{\fn}{\footnote}
% see TeXBook p.366
\def\myraggedright{%
  %\hyphenpenalty=-100
  \rightskip=0pt plus1em
  \spaceskip=.3333em \xspaceskip=.5em\relax
}

%\newenvironment{uncertain}{\special{x:textcolor=666666}}{\special{x:textcolor=000000}}
\maxsecnumdepth{book} 
\setsecnumdepth{book} 

% changes made on 20140425
% \linenumbersep=.5em % ednotes
\linenumbersep=10pt % ednotes; in fact lineno.sty
\linenumberwidth=5.525pt

% this is used throughout this trl.  I can't bother reinprementing it for now.  So, it's an empty command.
\newcommand{\refm}[1]{}
% \bibliographystyle{plainnat}
\bibliographystyle{kengo2}
\citeindextrue % natbib specific
% \bibliographystyle{plainnat}
\bibpunct[:\,]{(}{)}{,}{a}{}{,}
\title{God, Reason, and Yoga}
\author{Kengo Harimoto}
\makeindex

\makepagestyle{yvibook}
\makepsmarks{yvibook}{% 
  \createmark{chapter}{both}{nonumber}{}{}
  \createmark{section}{right}{nonumber}{}{}
  \createplainmark{toc}{both}{\contentsname}
  \createplainmark{lof}{both}{\listfigurename}
  \createplainmark{lot}{both}{\listtablename}
  \createplainmark{bib}{both}{\bibname}
  \createplainmark{index}{both}{\indexname}
  \createplainmark{glossary}{both}{\glossaryname}
}

% \makeevenhead{yvibook}{\leftmark}{}{}
% \makeoddhead{yvibook}{}{}{\rightmark}
\makeevenhead{yvibook}{\thepage}{}{\leftmark}
\makeoddhead{yvibook}{\rightmark}{}{\thepage}
% \makeevenfoot{yvibook}{}{\thepage}{}
% \makeoddfoot{yvibook}{}{\thepage}{}
% \makeevenfoot{yvibook}{[Draft]}{}{\today}
% \makeoddfoot{yvibook}{\today}{}{[Draft]}
\pagestyle{yvibook}

\makepagestyle{cried}
\makepsmarks{cried}{% 
  \createmark{chapter}{both}{nonumber}{}{}
  \createmark{section}{right}{nonumber}{}{}
  \createplainmark{toc}{both}{\contentsname}
  \createplainmark{lof}{both}{\listfigurename}
  \createplainmark{lot}{both}{\listtablename}
  \createplainmark{bib}{both}{\bibname}
  \createplainmark{index}{both}{\indexname}
  \createplainmark{glossary}{both}{\glossaryname}
}

\makeevenhead{cried}{\thepage}{}{\rightmark}
\makeoddhead{cried}{\rightmark}{}{\thepage}
\makeevenfoot{cried}{[Samādhipāda I]}{}{[Sūtra 23]}
\makeoddfoot{cried}{[Sūtra 23]}{}{[Samādhipāda I]}
% \pagestyle{cried}

\makepagestyle{test}
\makepsmarks{test}{% 
  \createmark{chapter}{both}{nonumber}{}{}
  \createmark{section}{right}{nonumber}{}{}
  \createplainmark{toc}{both}{\contentsname}
  \createplainmark{lof}{both}{\listfigurename}
  \createplainmark{lot}{both}{\listtablename}
  \createplainmark{bib}{both}{\bibname}
  \createplainmark{index}{both}{\indexname}
  \createplainmark{glossary}{both}{\glossaryname}
}
\makeevenhead{test}{}{}{}
\makeoddhead{test}{}{}{}
\makeevenfoot{test}{}{}{}
\makeoddfoot{test}{}{}{}


% so that we don't get 0.1 for tables, etc. 20140406
\makeatletter
  \counterwithout{figure}{chapter}
  \counterwithout{table}{chapter}
  \renewcommand\@memfront@floats{}
  \renewcommand\@memmain@floats{}
  \renewcommand\@memback@floats{}
\makeatother

\begin{document}
% \nouppercaseheads
\frontmatter

\begin{titlingpage}
% % \maketitle
\newcommand{\titlebreak}{\\ \rule{0pt}{12pt}\hfill}
% % \rule{0pt}{10pt}\hfill{\HUGE A Treatise on Īśvara} \titlebreak\rule{0pt}{20pt}
% % \rule{0pt}{10pt}\hfill{\HUGE Arguing for God} \titlebreak\rule{0pt}{20pt}
% \rule{0pt}{100pt}\hfill{\HUGE God, Reason, and Yoga} \titlebreak\rule{0pt}{20pt}
% % \begin{Large}
% % A Critical Edition and Translation
% % of the Commentary Ascribed \titlebreak to Śaṅkara
% % on Pātañjalayogaśāstra 1.23--28
% % with Appendices \titlebreak Including an Edition and Translation
% % of the \titlebreak Commentary on Pātañjalayogaśāstra 1.1%
% % \end{Large}
%
% % \rule{0pt}{10pt}\hfill{\HUGE \textbdn{A Treatise on Īśvara}}\\
% \vskip 2ex
% \begin{large}
% % \textbdn{
% \hfill A Critical Edition and Translation of the \\
% \rule{0pt}{10pt}
% \hfill Commentary Ascribed to Śaṅkara on \\
% \rule{0pt}{10pt}
% \hfill Pātañjalayogaśāstra 1.23--28%}\\
% \end{large}
%
% \vskip 20ex
% \hfill Kengo Harimoto
% \vfill
% \begin{center}
% Department of Indian and Tibetan Studies\\
% University of Hamburg\\
%
% 2014
% \end{center}
\rule{0pt}{141.421356pt}

\begin{center}
To my parents
\end{center}

\posttitle{A Critical Edition and Translation of the Commentary on the Pātañjalayogaśāstra Ascribed Śaṅkara}

\end{titlingpage}
\firmlists
% \tightlists
\setbeforeparaskip{2.25ex plus .66ex minus .2ex}
\setbeforesubparaskip{2.00ex plus .5ex minus .4ex}
\cleardoublepage
\tableofcontents
% \input{preface}
% \chapter{Acknowledgements}

% Acknowledgments\ldots.

This book owes much to two German scholars who are sadly not with us any more. I am sorry for not being able to present this book in front of them. They are the late Professor Dr.\ Wilhelm Halbfass, and the late Professor Dr.\ Tillman Vetter. 

Professor Halbfass was the patient supervisor of my doctoral thesis in my five and a half years at Penn in Philadelphia. Knowing my interest in the \PYSV, he almost matter-of-factly suggested that I prepare a critical edition of the text as the topic for my doctoral dissertation. I did not have to think much before saying yes. He also made it possible that I had access to manuscripts. The thesis became the basis of this book. I am not happy that I became one of his very last students. 

Professor Vetter practically determined the content of this book. Even before I met him, a brief remark in one of his writings directed me to look at the Īśvara section of the Vivaraṇa. I was fortunate to meet him years later, as a Gonda Fellow in Leiden in 2005 to 2006. He agreed to read my edition with me. In view of the limited time period, I asked him to read the Īśvara section in which I thought he might be interested. In the end, it is not a coincidence that the content of this book corresponds to the part we read together, including the commentary on sūtra 1.1. It was also he who insisted that I have a translation to read the edition with him. 

I also owe much to Professor Albrecht Wezler. He agreed that I work on a critical edition of the \PYSV\ and let me use his copies of the Lahore manuscript which has always been difficult to gain access to. After I came to Hamburg in 2000, he arranged weekly reading sessions for the edition I prepared for my thesis. Much of the improvements in the text come from those reading sessions between 2000 and 2002.

I am thankful to all my teachers. Professor Osamu Hayashima first taught me Sanskrit. It might be that I am the only person who still reads Sanskrit among his students. The way he taught me Sanskrit was very efficient; he told me to learn it by myself. Professor Hiromasa Tosaki imparted to us what it is to be a scholar. Professor Akihiko Akamatsu introduced me to the concept of philology. It took some time for me to digest the concept but I hope now I get it. I have learned many many things from Dr.\ Futoshi Omae when he was an assistant at the Kyushu University. It might have been during the daily chatting over many things with him that I learned most about our field of studies.

When I was at Penn, it was a golden age for classical indological studies there. We had Professors Halbfass, George Cardona, Ludo and Rosane Rocher and many inspiring fellow students. I learned various things from all of them. I have so many fond memories of my time in Philadelphia. One of them is the memory of reading the Vākyapadīya with Prof.\ Cardona for five years.

Working for the Skandapurāṇa project, led by Professor Hans Bakker, in Groningen, the Netherlands, was a new experience for me. Exposure to the Dutch Indology had an effect on me. I had to reshape my self-perception as an indologist. I also thank Prof.\ Bakker's support/insistence to finish this book as soon as possible. %At the same time, I feel apologetic that it took so long.

The same thanks go to Arlo Griffiths whom I first met in 1997 or so at an AOS meeting before seeing each other again in Groningen in 2004. He, too, has been supportive and persuasive. He arranged that I could read my text with Professor Vetter. He also looked at my English many times. I was also fortunate to have had an opportunity to work with Professor Yuko Yokochi in Groningen. In 2012, she arranged an intensive course to read the Vivaraṇa at the Kyoto University. Inputs from scholars and students there were very valuable. The people who were present there include Diwakar Acharya, Somdev Vasudeva, Akihiro Kanabishi, Koreto Ikehata. 

Special mention has to be made to Harunaga Isaacson with whom I seem to share some fate in that we have been moving symmetrically around Hamburg in the past 14 years. Not only has he been a tremendous inspiration, without him, this book would have been impossible. If there is any quality in it, it is largely because of him. I sincerely thank him for his support, not only for this book but also for everything. 

There are many other people to whom thanks are due. Here are some of the names: 
% Without any of the people mentioned below, this work would have been impossible. It took too long for any standard to get this out. I try here to mention those whom this book owes in generally a chronological order.
% Osamu Hayashima
% Hiromasa Tosaki
% Akihiko Akamatsu
% Wilhelm Halbfass
% George Cardona, 
% Ludo Rocher 
% Rosane Rocher
% Hans Bakker
% Albrecht Wezler
% Harunaga Isaacson
% Arlo Griffiths
% Diwakar Acharya
% Friends at PENN
% Friends at RuG
% Friends in Fukuoka
% Friends in Hamburg
% Korean scholars
% Tillman Vetter
% Yuko Yokochi
% Diwakar Acharya,
Ashok Aklujkar,
Orna Almogi,
Loriliai Biernacki,
Piotr Balcerowicz,
Benita von Behr,
Bidur Bhattarai,
Peter Bisschop,
Tim Cahill,
Britton Chance\dag{},
Danielle Cuneo,
Jonathan Duquette,
Christina Edingloh,
John Freeman,
Camillo A. Formigatti,
Angelika Frankowski,
Elisa Freschi,
Takamichi Fujii,
Dominic Goodall,
Kunihiro Harikai,
Tomoko Himeno,
%Koreto Ikehata,
%Akihiro Kanabishi,
Kazuo Kano,
Kei Kataoka,
Kenichi Kuranishi,
Andrey Klebanov,
Masato Kobayashi,
Stasinos Konstantopoulos,
Sonja Lakner,
Philipp Maas,
Kristin McGee,
David Mellins,
Mai Moriguchi,
Shin'ya Moriyama,
Yasutaka Muroya,
Hodo Nakamura,
Masanori Nakagawa,
David Nelson, 
Heidi Neubert,
Sachio Nioka,
Shoko Nioka,
Luther Obrock,
Kiyokazu Okita,
Dimitri Pauls,
Kim Plofker,
Akane Saito,
Florinda de Simini,
Barbara Schuler,
Greg Seton,
Francesco Sferra,
Taisei Shida,
Kiyokuni Shiga,
Jae\-kwan Shim,
Iain Sinclair,
Michael Slouber,
Hans-Jörg Stein,
Vidyasankar Sundaresan,
Peter Szanto,
Masaya Takahashi,
Ryugen Tanemura,
Somdev Vasudev,
Dorji Wangchuk,
Eiichi Yamaguchi%
.


Since it took so long to complete this book, practically anyone I knew has been in one way or another influential. I would like to thank all the friends I have ever had. And at last but not least, the biggest thanks are due to my parents for their unwavering support.
% Koh Endo
% Professor Osamu Hayashima first taught me Sanskrit. For a person like myself, his method of teaching Sanskrit turned out to be best; he was very strict. Professor Hiromasa Tosaki taught me what it means to be a researcher. Without Dr.\ Futoshi Omae, I would not have worked on this intriguing text, the Pātañjalayogaśāstravivaraṇa. I still remember the moment he, then the assistant of the department and history of Indian philosophy and buddhism, Kyushu University, asked me what I wanted to work on for my master's thesis. When I mentioned the Yoga school, he entered the library and came back with the edition of ``the Vivaraṇa.'' He said that I could work on it for ten years; it has been more than twenty years and I am not done yet. I owe much to Professor Akihiko Akamatsu for being patient and not giving up on me. I cannot help feel a sense of irony because the topic of this book is essentially the same as the master's thesis I wrote with him. After more than twenty years, I would hope that I learned something more and this book offers something new.

% To this day, I regret the loss of my dissertation advisor, Professor Wilhelm Halbfass. The inspirations he gave me were enormous; his knowledge and intellect were aspiring; and he was kind. He helped me through obtaining the degree at the University of Pennsylvania. This work, nonetheless, is based on my doctoral dissertation. It was one of my happiest moments when he said, ``Now you have convinced me,'' with regard to the authorship when I showed the final draft of my dissertation. It is a pity that he passed away only five months after I had my defense. He is not there any more to testify that he said that.

% I also owe much to Professors George Cardona, Ludo Rocher and Rosane Rocher of the University of Pennsylvania. Having Prof.\ Cardona as a teacher is a great advantage. He makes it clear that he never openly attacks his student (however he may be wrong). It is a great honor that he considers myself as one of his students. 

% \hfuzz=10pt
\mainmatter
\let\oldfootnoterule\footnoterule
\renewcommand*{\footnoterule}{}
\chapter{Abbreviations}
\newlength{\taiubh}
\settowidth{\taiubh}{TaiUBh} % looks like this is the widest abbreviation
\addtolength{\taiubh}{12pt}
\newlength{\rightcolumn}
\setlength{\rightcolumn}{\textwidth}
\addtolength{\rightcolumn}{-\taiubh}

\newcommand{\abbrindex}[4]{%
  \index{#1@#2|see{#3}}%
  #2 & #3 \citep{#4} \\[0.50ex]
}
% \index{BAU@BĀU|see{Bṛhadāraṇyakopaniṣad}}
% \index{BAUBh@BĀUBh|see{Bṛhadāraṇyako\-pa\-ni\-ṣa\-d\-bhā\-ṣya}}
\section{Texts and editions}
\setlength\LTleft{0pt}
\setlength\LTright{0pt}
\begin{longtable}{@{}l p{\rightcolumn}@{}}
\abbrindex{AA}{AA}{Aṣṭādhyāyī of Pāṇini}{vasu1962ashtadhyayi} %vasu...
% ANSS
\abbrindex{BAU}{BĀU}{Bṛhadāraṇyako\-pa\-ni\-ṣa\-d}{bau}
\abbrindex{BAUBh}{BĀUBh}{Bṛ\-ha\-dā\-ra\-ṇya\-ko\-pa\-ni\-ṣa\-d\-bhā\-ṣya ascribed to Śaṅkara}{bau}
\abbrindex{BhG}{BhG}{Bhagavadgītā}{bhgbh}
\abbrindex{BhGBh}{BhGBh}{Bhagavadgītābhāṣya ascribed to Śaṅkara}{bhgbh}
\abbrindex{BS}{BS}{Brahmasūtra}{bsbh}
\abbrindex{BSBh}{BSBh}{Bra\-hma\-sū\-tra\-śā\-ṅka\-ra\-bhā\-ṣya of Śaṅkara}{bsbh}
\abbrindex{BSi}{BSi}{Brahmasiddhi of Maṇḍana Miśra}{bs}
\abbrindex{ChU}{ChU}{Chāndogyopaniṣad}{upa2}
\abbrindex{ChUBh}{ChUBh}{Chā\-ndo\-gyo\-pa\-ni\-ṣa\-dbhā\-ṣya ascribed to Śaṅka\-ra}{upa2}
% CSS
% \abbrindex{IUBh}{ĪUBh}{Īśopaniṣadbhāṣya ascribed to Śaṅkara}{} 
% \abbrindex{IU}{ĪU}{Īśopaniṣad}{}
\abbrindex{JS}{JS}{Jaiminisūtra a.k.a.\ Mī\-māṃ\-sā\-sū\-tra}{frauwallner:karma,js1,js1a,js2,js3,js4,js5,js6,js7}
\abbrindex{KAS}{KAŚ}{Kauṭilyārthaśāstra}{as:kangle_ed1}
\abbrindex{KaU}{KaU}{Kaṭhopaniṣad}{upa1} 
\abbrindex{KaUBh}{KaUBh}{Kaṭhopaniṣadbhāṣya ascribed to Śaṅkara}{upa1} 
\abbrindex{KeU}{KeU}{Kenopaniṣad}{upa1} 
\abbrindex{KeUBh}{KeUBh}{Kenopaniṣadbhāṣya ascribed to Śaṅkara}{upa1} 
% KU
% KUPBh
\abbrindex{Manu}{Manu}{Manusmṛti}{olivelle:manu_ed}% Olivelle 2006 
\abbrindex{MaU}{MāU}{Māṇḍūkyopaniṣad}{upa1}
\abbrindex{MaUBh}{MāUBh}{Mā\-ṇḍū\-kyo\-pa\-ni\-ṣa\-dbhā\-ṣya ascribed to Śa\-ṅka\-ra}{upa1}
\abbrindex{MBh}{MBh}{Mahābhārata}{poona:mbh} 
\abbrindex{MMK}{MMK}{Mū\-la\-ma\-dhya\-mi\-ka\-kā\-ri\-kās of Nāgārjuna}{mmk} 
% \abbrindex{MS}{MS}{Manusmṛti???}{} 
\abbrindex{MuU}{MuU}{Muṇḍakopaniṣad}{upa1}
\abbrindex{MuUBh}{MuUBh}{Mu\-ṇḍa\-ko\-pa\-ni\-ṣa\-dbhā\-ṣya ascribed to Śaṅkara}{upa1} 
\abbrindex{NBh}{NBh}{Nyāyabhāṣya}{ns}
\abbrindex{NBhu}{NBhū}{Nyāyabhūṣaṇa of Bhā\-sa\-rva\-jña}{nbhusana} 
% \abbrindex{NK}{NK}{Nyāya\ldots???}{}
\abbrindex{NKd}{NKd}{Nyāyakandalī}{pdhs}
\abbrindex{Nkn}{NKṇ}{Nyāyakaṇikā of Vācaspati Miśra}{vidhiviveka,vidhiviveka:tara,vidhiv:stern}
\abbrindex{NM}{NM}{Nyāyamañjarī of Bhaṭṭa Jayanta}{nm:kss}
\abbrindex{NS}{NS}{Nyāyasūtra}{ns}
\abbrindex{NV}{NV}{Nyāyavārttika}{ns}
\abbrindex{NVTT}{NVTṬ}{Nyā\-ya\-vā\-rtti\-ka\-tā\-tpa\-rya\-ṭī\-kā of Vācaspati Miśra}{ns}
\abbrindex{PDhS}{PDhS}{Padārthadharmasaṅgraha}{pdhs}
% \abbrindex{PP}{PP}{Prakaraṇapañcikā of Śālikanātha}{pp}
\abbrindex{PaSu}{PāSū}{Pāśupatasūtra}{pasupatasutra}
\abbrindex{PraPa}{PraPa}{Prakaraṇapañcikā}{pp}
\abbrindex{PYS}{PYŚ}{Pātañjalayogaśāstra}{ys:anss,yvi}
\abbrindex{PV}{PV}{Pramāṇavārttika of Dha\-rma\-kīrti}{pvv,pv:pandeya}
\abbrindex{SBh}{ŚBh}{Śābarabhāṣya}{frauwallner:karma,js1,js1a,js2,js3,js4,js5,js6,js7}
\abbrindex{SK}{SK}{Sāṃkhyakārikās}{yd} %yd
\abbrindex{SRTS}{ŚRTS}{Tattvasaṃgraha of Śā\-nta\-ra\-kṣita}{ts:bbh}
\abbrindex{SS}{SS}{Sphoṭasiddhi of Maṇḍana Miśra}{SS}
\abbrindex{SV}{ŚV}{Ślokavārttika of Kumārila}{sv:umbeka,sv:partha,svkasika} 
\abbrindex{TaiU}{TaiU}{Taittirīyopaniṣad}{upa1}
\abbrindex{TaiUBh}{TaiUBh}{Tai\-tti\-rī\-yo\-pa\-ni\-ṣa\-dbhā\-ṣya ascribed to Śaṅkara}{upa1}
\abbrindex{TS}{TS}{Taittirīyasaṃhitā}{weber:tais}
\abbrindex{TSP}{TSP}{Tattvasaṃgrahapañjikā of Kamalaśīla}{ts:bbh}
%\abbrindex{TP}{TP}{???}{}
\abbrindex{TV}{TV}{Tattvavaiśāradī (a commentary on the Yo\-ga\-bhā\-ṣya) of Vā\-ca\-s\-pa\-ti Miśra}{ys:anss,yvi}
\abbrindex{Upad}{Upad}{Upadeśasāhasrī}{mayeda:upad}
\abbrindex{VidhiV}{VidhiV}{Vidhiviveka of Maṇḍana Miśra}{vidhiviveka,vidhiviveka:tara,vidhiv:stern}
\abbrindex{VMBh}{VMBh}{Mahābhāṣya of Patañjali}{mbh1,mbh3}
\abbrindex{VP}{VP}{Vākyapadīya of Bhartṛhari}{rau:text,rau:index,iyer:vp2text} 
\abbrindex{VS}{VS}{Vaiśeṣikasūtra}{vs}
\abbrindex{YBh}{YBh}{Yogabhāṣya}{ys:anss,yvi,maas:2006}
\abbrindex{YD}{YD}{Yuktidīpikā}{wezler:yd}
\abbrindex{YS}{YS}{Yogasūtra}{ys:anss,yvi,maas:2006}
\abbrindex{YV}{YV}{Yogavārttika of Vi\-jñā\-na\-bhi\-kṣu}{rukmani:yvarttika}
\abbrindex{YVi}{YVi}{Pātañjalayogaśāstravivaraṇa}{yvi}
\index{Jayamangala@Jayamaṅgalā|seealso{Sāṃkhyajayamaṅgalā}}
\index{Mimamsa@Mīmāṃsā|seealso{Mīmāṃsaka}}
\index{Samkhyajayamangala@Sāṃkhyajayamaṅgalā|seealso{Jayamaṅgalā}}
% \abbrindex{}{}
% \abbrindex{}{}
% \abbrindex{}{}
% \abbrindex{}{}
% \abbrindex{}{}
% \abbrindex{}{}
% \abbrindex{}{}
% \abbrindex{}{}
% \abbrindex{}{}
% \abbrindex{}{}
% \abbrindex{}{}
\end{longtable}

\section{Sigla and Others}
See pp.~\pageref{sec:symbols} for details.
\setlength\LTleft{0pt}
\setlength\LTright{0pt}
\begin{longtable}{@{}ll@{}}
	MS & manuscript \\
	em. & emendation \\
	conj. & conjecture \\
	\Tm & Trivandrum manuscript \\
	\Tmac & ditto before correction \\
	\Tmpc & ditto after correction \\
	\Td & A transcript of the Trivandrum manuscript \\
	\Tdac & ditto before correction \\
	\Tdpc & ditto after correction \\
	T & A reading common in \Tm\ and \Td \\
	L & Lahore manuscript \\
	\Lac & ditto before correction \\
	\Lpc & ditto after correction \\
	M & Madras Government Oriental Manuscript Library manuscript \\
	\Mac & ditto before correction \\
	\Mpc & ditto after correction \\
	A & Adyar Library manuscript \\
	\Aac & ditto before correction \\
	\Apc & ditto after correction \\
	\Me & 1952 edition of the YVi \\
	$\Sigma$ & all the witnesses, including the 1952 edition \\
	\mms & the following text is lost in \Tm \\
	\kern 5.3pt \mme & the preceding text is lost in \Tm\\
	\raisebox{-0.5ex}{\mds} & the following text is reported missing in \Td \\
	\raisebox{-0.5ex}{\kern 5.3pt \mde} & the preceding text is reported missing in \Td \\
	\dag & the enclosed text is suspected to be corrupt \\
	\begin{uncertain}{\dn क}\end{uncertain} & (text in a lighter shade) uncertain text \\
	\ldots & in critical text, suspected loss of text \\[0pt]
	[\mist{1}] & physically lost sign in a manuscript \\
	\mist{1} & a sign indicated as lost in the exemplar \\
	\hlg{{\dn क}} & highly illegible but ascertainable as {\dn क} \\
	\ilg & an illegible (but not lost) sign \\
	\ccs {\dn क}\cce & cancelled {\dn क} \\
	\ins{\dn क}\ine & inserted {\dn क} \\
	\open & space left by scribe
\end{longtable}

% \begin{description}
% 	\item[MS]{manuscript}
% 	\item[em.]{emendation}
% 	\item[conj.]{conjecture}
% 	\item[\Tm]{Trivandrum manuscript}
% 	\item[\Tmac]{Trivandrum manuscript before correction}
% 	\item[\Tmpc]{Trivandrum manuscript after correction}
% 	\item[Td]{A transcript of the Trivandrum manuscript}
% 	\item[\Tdac]{A transcript of the Trivandrum manuscript before correction}
% 	\item[\Tdpc]{A transcript of the Trivandrum manuscript after correction}
% 	\item[T]{A reading common in \Tm\ and \Td}
% 	\item[L]{Lahore manuscript}
% 	\item[\Lac]{Lahore manuscript before correction}
% 	\item[\Lpc]{Lahore manuscript after correction}
% 	\item[M]{Madras Government Oriental Manuscript Library manuscript}
% 	\item[\Mac]{Madras Government Oriental Manuscript Library manuscript before correction}
% 	\item[\Mpc]{Madras Government Oriental Manuscript Library manuscript after correction}
% 	\item[A]{Adyar Library manuscript}
% 	\item[\Aac]{Adyar Library manuscript before correction}
% 	\item[\Apc]{Adyar Library manuscript after correction}
% 	\item[\Me]{1952 edition of the YVi}
% 	% \item[L]{}
% 	% \item[L]{}
% 	% \item[L]{}
% 	% \item[L]{}
% \end{description}
\firmlists
\pagestyle{yvibook}
% \part{Introduction}
% \renewcommand{\leftmark}{\MakeUppercase{Introduction}}
\renewcommand{\rightmark}{\MakeUppercase{General Introduction}}
% \input{genint}
% \renewcommand{\rightmark}{\MakeUppercase{Introduction to the Edition}}
% \input{intro}

\newcommand{\starteditionchapter}{
  \begin{sanskrit}%
	\begin{Spacing}{1.30}%
}
\newcommand{\finisheditionchapter}{%
  \end{Spacing}%
  \end{sanskrit}%
  \begin{english}% These two lines are silly; polyglossia does something 
  \end{english}%   funny to the toc file.
}

\part{Critical Text}\renewcommand{\leftmark}{\MakeUppercase{[Critical Text}}
% \input{crittext}

\makechapterstyle{crited}{%
  \chapterstyle{default}
  \chapterstyle{article}%
  \renewcommand*{\thechapter}{\Roman{chapter}}
  \renewcommand*{\printchapternum}{% center number/title
    \centering\chapnumfont \thechapter\space\space}%
  \renewcommand*{\printchapternonum}{\centering}
  \renewcommand*{\clearforchapter}{}% no new page
  \aliaspagestyle{chapter}{yvied}}% no special pagestyle


\chapterstyle{crited}
% \pagestyle{cried}
\chapter{Critical Text 1.2}%
\renewcommand{\rightmark}{\MakeUppercase{1.2]}}
% \thispagestyle{test}
\starteditionchapter
\section{\pysv\ 1.2}
\setlength{\footmarkwidth}{1.8em}
\pagewiselinenumbers
\newlength{\dawidth}
\settowidth{\dawidth}{ड}
\addtolength{\dawidth}{0.08em}
\XeTeXglyph 101 \raisebox{0.25ex}{\scalebox{1}[0.8]{ड}}\kern -\dawidth \raisebox{-0.65ex}{\scalebox{1}[0.8]{भ}} \mll{ONETWOONE}\mula{सर्व्ववृत्ति\vrtms{निरोधे}{TM\Me}{\rdg{॰निरोधे प्य}{L}} त्वसम्प्रज्ञात इति}\donote{ONETWOONE}{{\dn सर्व्ववृत्तिनिरोधे त्वसम्प्रज्ञात इति}}।\llbl{SUTR2INTRO} तुशब्दो ऽवधारणार्त्थः। \mll{SUTR2INTRO}\mula{तस्यै}वं\vrtms{विधस्य}{\Me(em.)}{\rdg{॰वंविषयस्य}{TLMA}} समाधेरसम्प्रज्ञातस्यैव केवलस्य \mula{लक्षणाभिधित्सयेदं सूत्रम्प्रववृते}\donote{SUTR2INTRO}{{\dn तस्य लक्षणाभिधित्सयेदं सूत्रम्प्रववृते}} प्रवृत्तम्---
% evaṃviṣayasya is in all the manuscripts
% I consider "eva" to be part of the mūla. "kevalasya" appears a gloss for "eva". Otherwise "tasyaiva kevalasya" appears slightly redundant.
\newfontfamily
\str{योगश्चित्तवृत्तिनि\mms रोध (२)} \mds \vrt{इति}{L}{\rdg{ईति वक्तव्यं}{M}, \rdg{ई[[ती]]ति वक्तव्यम्}{\Me}}॥

ननु च त\mte स्य लक्षणमभिधातुमेतत्सूत्रम्प्रवृत्तम? तथा च सति योगस्य चित्तवृत्तिनिरोध इति वक्तव्यम्। 
सामानाधि\vrtms{करण्यन्न}{LM\Me\vv{॰ण्यं न}{\Me}}{\rdg{॰करण्यान्न}{\Td}} युक्तम्। न हि लक्ष्यमेव लक्षणं स्यात्। चित्तवृत्तिनिरोधलक्षण इति वा वक्तव्यम्॥
% na hi lakṣyam eva lakṣaṇaṃ syāt | cittavṛtti... is awkward.

नैष दोषः। लक्ष्ये ल\mms क्षण\mds ाध्यासात्। यथा\mte यम्पिण्डो देवदत्त इति लक्षये लक्षणम्प्रत्यस्यते॥
% piṇḍa and Devadatta: BAUBh, MBh, Āryabhaṭīya

ननु \mms च स\mme र्व्ववृत्तिनिरोधे त्वसम्प्रज्ञातः। तस्य चे\mms दं लक्षणम्। ततश्च योगः सर्व्वचि\mme त्तवृत्तिनिरोध इति वक्तव्यम। न चोक्तम्। अतो न व्या\mds पि \mms लक्षणम्॥

सर्व्वशब्दाग्रहणे तावत्कारणमुच्यते---यदि सर्व\mte ग्रहणं क्रियेत तदैकदेशनिरोधनिमित्तस्य सम्प्रज्ञातस्य योगत्वन्निवर्तितं स्यात्। तन्मा भूदिति सर्व्वशब्दाग्रहणम्॥

अतस्तर्हि सम्प्रज्ञातासम्प्रज्ञातयोः सामान्यलक्षणम्, विशेषानुपादानात्॥

नैष दोषः, योगे प्रकृते पुनर्य्योगग्रहणात्। वि\mts शेषाविवक्षायां हि यो\mde ग\mme ग्रहणन्नार्त्थवत्स्यात्॥

ननु चानुशासने गुणभूतत्वाद्योगस्य, इह योगग्रहणे ऽसति चित्तवृत्तिनिरोध इत्युक्ते ऽनुशासनार्त्थता भवेत्सूत्रस्य॥

नैवम्, योगानुशासनस्य प्रस्तुतत्वात्। योगानुशासनं हि प्रारब्धन्नायो\mts गानुशासनम्। त\mde ल्ल\mme क्ष\-णाभिधाने 
ह्यतिप्रसंगः पदार्त्थासंगतिश्च। न हि चित्तवृत्तिनिरोधो योगानुशासनमिति पदार्त्थः संगच्छते॥

एवमपि योगग्रहणात्कुतो ऽसम्प्रज्ञातस्यैव लक्षणम्, न पुनः सम्प्रज्ञातस्य लक्षणं स्यात्? योगग्रहणस्याविशिस्ष्टत्वादित्युच्यते---

न, व्यभिचारात्। [यद्]यद्व्यभिचरति न तत्तस्य लक्षणम्। निरोधस्त्वसम्प्रज्ञाते ऽपि वर्त्तत इति सम्प्रज्ञातं व्यभिचरति। न 
हि विषाणित्वं गोर्लिंगम्, महिष्यादिभिर्व्यभिचारात्॥

ननु चासम्प्रज्ञातमपि व्यभिचरति, सम्प्रज्ञातस्यापि हेतुत्वात्?

सत्यमेवम्। किन्तु न निरोधादन्येन लक्षयितुमसम्प्रज्ञातः शक्यते। निरोध एव तु तस्य लक्षणम्, अन्याभावात्, सम्प्रज्ञातस्य त्वसाधारणवितर्कादिलक्षणलक्ष्यत्वात्। यथा `स्पर्शलक्षणमि'त्युक्ते वायोरेवासाधारण्याल्लक्ष्यते स्पर्शवत्त्वेन। अथ च विद्यत एव स्पर्शवत्वन्तेजःप्रभृतिष्वपि। एवञ्च सति योगग्रहणमनुवादमात्रम्। अत एव च सर्वग्रहणन्न कृतम्। सूत्रे निरोधग्रहणेनैवासम्प्रज्ञातलक्षणत्वे सिद्धे यदि सर्व्वग्रहणं क्रियेत तत्सम्प्रज्ञातस्य योगतां निवारयेत्। तस्मादसम्प्रज्ञातस्यैवेदं लक्षणम्॥

तथा चोपसंहरति भाष्यकारः---न तत्र किञ्चित्सम्प्रज्ञायत इत्यसम्प्रज्ञातः। स योगश्चित्तवृत्तिनिरोध इति। 
चित्ततद्वृत्तितन्निरोधव्याचिख्यासया चित्तं हीत्यादि भाष्यम्। किमात्मकम्पुनस्तच्चित्तं यस्य वृत्तिनिरोधो योग इत्यत आह---
चित्तं हि प्रख्याप्रवृ\linelabel{tmrepeats}त्तिस्थितिशीलत्वात्त्रिगुणमिति॥

तत्र `चित्तन्त्रिगुणमि'ति प्रतिज्ञायते। `प्रख्याप्रवृत्तिस्थितिशीलत्वादि'ति हेतुः। प्रख्या प्रख्यानम्प्रकाशनम्। सा(?) हि 
सत्वगुणस्य धर्म्मः। प्रवृत्तिः प्रवर्त्तनं व्यापारः। स हि रजोधर्म्मः। स्थितिः स्थानं वरणम्प्रतिबन्ध इति। स च तमोधर्म्मः। 
एवंशीला सत्वादयो गुणाः। चित्तञ्च प्रख्यादिशीलम्। तस्मात्त्रयाणां गुणानाम् परिणामश्चित्तम्भवितुमर्हति॥ % NEED A NOTE ON SAA OR SA, REFERENCES TO THOSE DHARMAS OF GU.NAS!

तथा प्रख्यादिहेतोरसिद्धतामाशंक्याह---प्रख्यास्वरूपं हि चित्तसत्वमिति। हिशब्दः प्रसिद्धावद्योतनार्त्थः। प्रसिद्धं हि लोके 
शास्त्रे च सर्व्वावभासकत्वञ्चित्तस्य। चित्तमेव सत्वञ्चित्तसत्वम्, सत्वप्रधानगुणपरिणामत्वात्॥ % need a big note on cittasya sarvaavabhaasakatvam!

अथ वान्यसंसर्गवियोगाभ्यामनेकवृत्तिख्यातिमात्रत्वञ्चित्तस्य स्पष्टन्दर्शयितुं शक्यत इति ब्रवीति---प्रख्यास्वरूपं हि 
चित्तसत्वमिति॥

इदानीं वृत्तयो व्याख्यायन्ते। तासाञ्चानेकत्वं रजस्तमसोरुद्भवाभिभवनिमित्तं विरुद्धत्वञ्चेत्याह---रजस्तमोभ्यां संसृष्टं 
सम्प्रधानाभ्यां यदा संसृष्टं तदानीमैश्वर्यविषयप्रियम्भवति। प्रीतिः रागः। ऐश्वर्यविषयरागयुक्ता वृत्तयो भवन्तीत्यर्त्थः। 
तत्तम्सानुविद्धम्। तदेव प्रख्यास्वरूपञ्चित्तं यदा गुणभूतरजस्केनोद्भूतेन तमसानुविद्धम्, तदानीमधर्मादिचतुष्टयविषया वृत्तयः क्
लिष्टा जायन्ते॥

तदेव प्रक्षीणमोहावरणामिति। न्यग्भूततमस्कन्तमःक्षयादेव जलधारापगमादिव रविबिम्बं सर्वतः 
प्रद्योतमानं रजोमात्रया रजोलेशेनानुविद्धन्तदा धर्म्माद्युपगतं, धर्म्मादिचतुष्टयविषया वृत्तयो ऽक्लिष्टा भवन्ति॥

यदा तदेव रजोमलापेतं स्वरूपप्रतिष्ठं केवलेन स्वेन प्रख्यारूपेणावस्थितन्तदा सत्वपुरुषान्यताख्यातिमात्रम्भवति। सत्वञ्चित्तम्, 
पुरुषो भक्ता, तयोरन्यथाविविक्तता, तत्ख्यातिस्तदवगमः। तन्मात्रग्रहणं \kle शाद्यभावप्रदर्शनार्त्थम्। धर्म्ममेघध्यानोपगतम्भवति। 
धर्म्ममेघो नाम समाधिः। तदेवम्प्रसंख्यानबलाद्रजस्तमसी तिरस्कृत्य केवलेन ख्यात्यात्मना पुरुषस्वरूपमात्रदर्शनेनावस्थानध्यानं 
प्रसंख्यानमित्याचक्षते ध्यायिनो योगिनः॥

एवं वृत्तिस्वरूपे व्याख्याते तन्निरोधप्रदर्शनार्त्थमाह---चितिशक्तिरित्यादि। चितिरेव शक्तिः चितिशक्तिः। यथा पच्यादयः 
शक्तिशक्तिमदपेक्षा प्रापणीयजन्मानो न स्वयं शक्तयः। चितिः पुनः स्वयमेव शक्तिर्नार्त्थान्तरमपेक्षते। तेन 
नित्यावस्थितत्वञ्चितिशब्दस्य चिन्मात्राभिधायिनः शक्तिशब्देन सामानाधिकरण्यं, अविकृतरूपया एव 
चितेर्व्विषयित्वप्रदर्शनार्त्थम्।यस्मादेवमतो ऽपरिणामिनि व्यवस्थिते धर्म्मिणि धर्म्मान्तरतिरोभावेन धर्म्मान्तरप्रादुर्भावः 
परिणामः। न परिणन्तुं शीलमस्या इत्यपरिणामिनी॥

अत एवाप्रतिसंक्रमा परिणामिन एव, चित्तादेर्विषयादौ प्रतिसंक्रमदर्शनात्॥

अत एव दर्शितविषया दर्शितो ऽन्तःकरणेन विषयो ऽस्या इति तेनैव शुद्धा॥

अत एवानन्ता देशतः कालतश्च। पूर्व्वम्पूर्वमुत्तरस्योत्तरस्य हेतुत्वेन द्रष्टव्यम्। वैधर्म्म्यदृष्टान्तञ्चित्तमिन्द्रियाणि च। 
सत्वगुणातिमिका सत्वमेव गुणः सत्वगुणः तस्येवात्मा स्वरूपमवभासो ऽस्याः सेयं सत्वगुणात्मकावभासस्वरूपेत्यर्थः। 
वृत्यविशिष्टवृत्तित्वाद्वा सत्वगुणात्मिकेत्युच्यते॥

यदि वा ख्यात्या सम्बध्यते---ख्यातिः सत्वगुणात्मिकेति। अतश्चितिशक्तेरुक्तलक्षणाया विपरीता विलक्षणा विवेकख्यातिः 
परिणामादिमती॥

यस्मात्पुरुषादुत्कृष्टात्परिणामादिगुणविरहितान्निकृष्टा परिणामादिगुणयुक्तातः स्वरूपदोषदर्शनाद्विरक्तमपरक्तञ्चित्तन्तामपि 
ख्यातिमात्मनो निरुणाद्धि। तदवस्थन्निरोधावस्थं संस्कारोपगं संस्कारमात्रावशेषम्। निरुद्धासु वृत्तिषु वृत्तिजनिताः संस्कारा 
एवावशिष्यन्ते। एतस्यान्निरोधभूमौ यः समाधिः स निर्बीजः। निर्गतम्बीजमत्र \kle शादिबीजं सर्व्वमुत्सन्नमस्मिन्निति॥

तस्मादसंप्रज्ञातसमाधिलक्षणार्त्थमेवं सूत्रमित्युपसंहरति---स योगश्चित्तवृत्तिनिरोध इति॥२॥

% testing for Arlo

प्राङ्गणे

प्र\XeTeXglyph 204 णे

क्लेश

ङ्न
\finisheditionchapter

% \chapter{Critical Text 1.3}%
% \renewcommand{\rightmark}{\MakeUppercase{1.3]}}
% \starteditionchapter
% योगश्चित्तवृत्तिनिरोध इत्युक्तम। तत्र बोद्धृत्वेनैव पुरुषसद्भावाधिगमः। विषयभुतवृत्तिनिरोधे च विषयिणोऽपि बोद्धुः पुरुषस्य निरोधं कश्चिदा\vrtms{शंकेतापि}{}{}। तथा च सति तत्कैवल्यप्राप्त्युपायस्यापि विवेकख्यातेरनर्त्थकत्वम् मन्वीत। योगानुशा\vrtms{सनस्य च}{}{} तदर्त्थस्य निष्फलत्व\vrtms{मपि}{}{} प्रतिपद्येत॥

तस्मात्पुरुषस्य वृत्तिनिरोदहादनिरोधन्दर्शयितुकामः ख्यातेश्चापि फलं साक्षाद्दर्शयन्नाह---\bht{तदवस्थे चेतसी}त्यादि॥

तदवस्थे---निरोधावस्थे। \bht{स्वविषयाभावात्}---स्वविषयश्चित्तवृत्तिः। तदभावे \bht{बुद्धिबोधात्मा}---बुद्धिं वृत्तिरूपेण परिणताम्बुध्यत इति बुद्धेर्बोद्धा---\bht{पुरुषः}। तद्बोधनमेव हि पुरुषस्य स्वरूपम्। नान्यो बोद्धा, नान्यद्बोधनम्। बोद्धुश्च बोधनादन्यत्वे सति विक्रियात्मकता स्यात्। ततश्च दर्शितविषयत्वञ्च न स्यात्। तथा च करणान्तरापेक्षित्वम्पुरुषस्य प्रसज्येत। तस्माद्वोधनम्बोद्धृत्वं चोपचरितम्। त्दवृत्तिसारूप्येन। तथा च वक्षयति---द्रष्टा दृशिमात्र इति। तस्मादाह बुद्धिबोधनमेवात्मा स्वरूपं यस्य स बुद्धिबोधात्मा पुरुषः॥

किंस्वभावः किन्निस्वभावः---सद्भावे वा किंस्वभावः कथं वासद्भाव इति---

\str{तदा द्रष्टुः स्वरूपे ऽवस्थानम्॥३॥}

यदा निरुद्धा वृत्तयः॥
\bht{स्वरूपप्रतिष्ठा तदानीञ्चितिशक्तिः}---\vrtsm{स्वरूपप्रतिष्ठैवे}{\em}{}त्यर्त्थः। \bht{यथा कैवल्ये}। \vrt{कैवल्ये}{}{} सद्भावन्तु `न तत्स्वाभासन्दृश्यत्वादि'त्येवमादिना वक्ष्यति। तत्र सिद्धे सद्भावे तत अकृष्य प्राप्ताशंकानिवृत्यर्त्थमुच्यते---`स्वरूपप्रतिष्ठे'ति॥

सिद्धवद्व्याख्यातं सूत्रम। तत्रेदम्प्रसक्तम्---इतरत्र न स्वरूपप्रतिष्ठेति। अन्यथा हि तदेति विशेषणानर्त्थक्यं स्यात्। न चेदन्यत्र स्वरूपप्रतिष्ठा तत्रावस्थान्तरयोगात्परिणामवत्वादिदोषः प्राप्नोतीति तत्परिजिहीर्षयाहा---\bht{व्युत्थानचित्ते तु सति तथापि भवन्ती}ति। अत्रापि स्वरूपप्रतिष्ठत्वन्दर्शयति। \bht{न तथे}त्यनेन `तदे'ति विशेषणस्यार्त्थवत्वं ख्यापयति॥३॥

% \finisheditionchapter
%
% \chapter{Critical Text 1.4}%
% % \chapter{Commentary on \PYS\ 1.25}%
% \renewcommand{\rightmark}{\MakeUppercase{1.4]}}
% \starteditionchapter
% इतरस्तु \bht{तथापि भवन्ती}त्यनेन चेत्स्वरूपप्रतिष्ठत्वमुच्यते, \bht{न तथे}ति च स्वरूपप्रतिष्ठत्वनिषेधश्चेत्, एकस्य वस्तुनः तथाभावश्चातथाभावश्च विरुध्यत इति वीक्षापन्नः पृच्छति---\bht{कथन्तर्ही}ति। इतर अह---

\str{वृत्तिसारूप्यमितरत्र॥४॥}

इति॥

कस्मात्पुनर्वृत्तिसारूप्यम? \bht{दर्शितविषयत्वात्}। उभयत्र स्वरूपप्रतिष्ठत्वाविशेषे ऽपि वृत्तिसारूप्यासारूप्यकृतो विशष इति न तथाभवातथाभावविरोधः॥

ननु च वृत्तिसारूप्ययोगे सत्यवस्थान्तरयोगात्परिणामितवादिदोषः प्राप्नोत्येव॥

न। दर्शितविषयत्वेन परिहृतत्वात्। चित्तवृत्यध्यारोपितं हि तन्न स्वतः, स्फटिकाद्युपधानोपरागवत्॥

\bht{व्युत्थान} इति निरोधादन्यत्र। \bht{याश्चित्त}स्य क्लिष्टादयो \bht{वृत्तयः}, \bht{ता}भिर\bht{विशिष्टा} सदृशी \bht{वृत्ति} स्वरूपमस्य। % check if there was a visarga after v.rtti
चित्तवृत्तिव्यतिरेकेण स्वरूपव्यतिरेकेण च पुरुषस्य वृत्यभावात्। सदाज्ञातविषयत्वे परिणामानुपपत्तेश्च॥

\bht{तथा च} पूर्व्वाचार्य्य\bht{सूत्रम्}--\bht{एकमेव दर्शनम्} बुद्धिपुरुषयोः। किन्तदित्याह---\bht{ख्यातिरेव दर्शनम्}। बुद्धिवृत्तिरेव दर्शनम्। ख्यायते पुरुषेणेति च ख्यातिः। ख्यायते ऽनया च बुद्धिपुरुषयोः स्वरूपमिति ख्यातिः। सैवार्त्थाकारतया गृह्यत इति करणम्। स्वेन च ख्यातिरूपेण गृह्यत इति कर्म्म च भवति। तथा दृश्यत इति दर्शनम्। दृश्यते ऽनेनेति च ज्ञायते ऽनेनेति च सर्व्वमेवमादि द्रष्टव्यम्॥

यत्तु ख्यातिकर्तृत्वे चित्तस्य तदेव ज्ञायते जानाति चेति ज्ञानान्तरकल्पनावैयर्त्थ्यप्रसंग इति तस्य `न तत्स्वाभासन्दृश्यत्वादि'त्यत्रैव परिहारं वक्ष्यामः॥ % Is this sentence coherent? Yes, it is. It refers to YS 4.19.

ननु च यदि चित्तम्परिविजृम्बहितमेवेदं सर्व्वम्पुरुषश्चेदुदास्ते कथमसौ भोक्तेति—न, कारकवैचित्र्यात्। कानिचिद्व्याप्रियमाणानि कारकाणि कानिचिद्व्यापारादृते सन्निधिमात्रेण कर्तृत्वं प्रतिपद्यन्ते। यथा पचतीत्यधिश्रयणादिषु व्याप्रियमाणः पुरुषः कर्तेत्युच्यते। स्थाली सम्भवधारणाभ्यामव्याप्रियमाणैव पचतित्युच्यते। तथा स्थालीयमध्यगतमाकाशमवकाशन्ददाति। न तस्य व्यापारं प्रतीमः। यथा च राजादर्शनमात्रेणैव कर्त्ता। सर्व्वाणि च कारकाणि स्वेषु व्यापारेषु कर्त्तॄणि।





[14,10]	itaras tu tathaa^ api bhavanty anena cet svaruupa-prati.s.thatvam ucyate 
na tathaa^ iti ca svaruupa-prati.s.thatvani.sedha"s ced ekasya vastunas tathaabhaava"s 
ca^ atathaa-bhaava"s ca virudhyata^ iti viik.saapanna.h p.rcchati katha.m tarhi 
iti./ itara^ aaha v.rttisaaruupyam itaratra^ iti./ kasmaat punar v.rttisaaruupyam ? 
dar"sitavi.sayatvaat // ubhayatra svaruupa-prati.s.thatvaavi"se.se^ api 
v.rtti-saaruupyaasaaruupyak.rto vi"se.sa^ iti, na tathaabhaavaatathaabhaava-virodha.h./

[14,15]	nanu ca v.rtti-saaruupya-yoge sati^ avasthaantara-yogaat 
pari.naamitva-aadi-do.sa.h praapnoty eva. na. dar"sitavi.sayatvena parih.rtatvaat. 
citta-v.rtty-adhyaaropita.m hi tat, na svata.h 
spha.tika-aady-upadhaana-uparogavat. vyutthaana^ iti./ nirodhaad anyatra ya"s 
cittasya kli.s.ta-aadayo v.rttaya.h taabhir avi"si.s.taa sad.r"sii v.rtti.h svaruupam 
asya, citta-v.rtti-vyatireke.na svaruupa-vtyatireke.na ca puru.sasya 
v.rtty-abhaavaat sadaa j~naata-vi.sayatve pari.naama-anupapatte"s ca./

[14,21]	tathaa ca puurva-aacaarya-suutram ekam eva dar"sana.m buddhi-puru.sayo.h. 
ki.m tad ity aaha khyaatir eva dar"sanam. buddhi-v.rttir eva dar"sanam. khyaayate 
puru.se.na^ iti ca khyaati.h khyaate 'nayaa ca buddhi-puru.sayo.h svaruupam^ iti 
khyaati.h. saa^ eva-arthaakaaratayaa [arthaakaaro 'nayaa] g.rhyata^ iti kara.na.m svena 
ca khyaati-ruupe.na g.rhyata^ iti karma ca bhavati. tathaa d.r"syata^ iti dar"sana.m, 
d.r"syate 'neneti ca [j~naayata^ iti j~naana.m ] j~naayate 'neneti ca sarvam 
evam-aadi dr.s.tavyam./

[14,27]	yat tu khyaati-kart.rtve cittasya tad eva j~naayate jaanaati ca^ iti 
(j~naanaad iti ca^ iti) j~naanaantarakalpnaa-vaiyarthya-prasa"nga^ iti tasya 'na 
tat-svaabhaasa.m d.r"syatvaat' ity atraiva parihaara.m vak.syaama.h./

[15,4]	nanu ca yadi citta-parvij.rmbhitam eveda.m sarva.m paru.sa"s ced udaaste 
katham asau bhokteti ? na kaarakavaicitryaat. kaani-cid vyaapriyamaa.naani 
kaarakaa.ni, kaanicid vyaapaaraad .rte sannidhi-maatre.na kart.rtva.m pratipadyante./

[15,7]	yathaa pacati^ ity adhi"sraya.na-aadi.su vyaapriyamaa.na.h puru.sa.h kartety 
ucyate. sthaalii sambhava-dhaara.naabhyaam avyaapriyam.naiva pacati^ ity ucyate. 
tathaa sthaalii-madhya-gatam aakaacam avakaaca.m dadaati, na tasya vyaapaara.m 
pratiima.h./

[15,10]	yathaa ca raajaa dar"sana-maatre.na^ eva kartaa. sarvaa.ni ca kaarakaa.ni 
sve.su vyaapaare.su kartY.ni. na ca^ aaditya.h prakaa"sayan kara.naantaram apek.sate 
vyaapriyate vaa / na hi prakaa"sanam avidyamaanam utpaadayati gamana-aadivat. 
prakaa"saatmataa^ eva hi tasya sad-bhaava.h tathaapi d.r"si-maatra.na gha.ta-aadiinaa.m 
prakaaca-ruupe.na^ abhivyakte.h prakaacayati^ ity ucyate, evam iha^ api 
d.r"si-maatre.na puru.se.na d.r"syamaanaa.m citta-v.rttiinaa.m cid-aatmanaa vyaapyamaanatvaat 
dra.s.taa puru.sa tiimam artha.m dar"sayati cittam ayaskaantama.ni-kalpam iti./

[15,16]	yasya^ apy aatma-mans-sa.myogaat j~naanam utpadyate tasya^ api j~naanena 
j~neya.m vyaapnuvan naatmaa j~naatety ucyate, na vyaapriyamaana.h. j~naanasya^ api 
[j~naanam api] gu.natve sati ni.skriyatvaad aatmana^ eva^ artha.m dar"sayati. na^ 
api kriyaan abhyupagamaat jaanaati^ ity-aadi na^ upapadyate. j~naanena hi j~neya.m 
vyaapnivan jaanaati^ ity ucyate./

[15,20]	ki.m ca^ anyat aatma-manas-sa.myogaad apuurva-j~naanotpattau saty 
aatma-manas-taj-j~naana.m na manasa.h. aatmaa^ eva j~nnaati na mana^ iti ko hetur 
bhavet ? atha^ api syaat aatmana.h samavaayi-kaara.natvaad iti. na^ etad eva.m 
asiddhatvaat. yathaa^ eva j~naanam aatmana^ iti saadhya.m evam aaaatmana^ eva 
samavaayi-kaara.nattva.m na manasa^ iti saadhyam anaatma-sva-ruupatve saty 
ubahaya-sa.myogajatva-abhyupagamaat./

[15,25]	atha^ api syaat j~naanaahita-sa.mskaarasyaatma[na.h]sm.rty-ekakart.rtvena 
pratisandhaana-dar"sanaad iti na tasya^ api saadhyatvaat. tatra^ api 
j~naana-sa.mskaara-sm.rti-pratisandhaanaadiiny\label{JSSP04} aatmana^ eva na manasa^ iti 
saadhaniiyaani. phalena samaanaa"srayatvaad icchaa-aadiinaam iti cet na tatra^ api 
tulyatvaat. sa.myogajatve saty aatmana.h phalecchaadiini na manasa^ iti du.hsaadham 
eva./

[16,4]	ki.m ca^ asya j~naanasya^ api j~neyatva-abhyupagamaad 
anavasthaa-prasa"nga.h. na hi j~naanam anyena kara.na-bhuutena j~naanena vij~naatu.m 
"sakya.m tad apy anyena tad apy anyenety ani.s.tha-prasa"ngaat./

[16,6]	athaanavasthaa-prasa"nga-parihaaraaya duuram api gatvaa ki.mcij j~naanam 
aadya.m vaa j~naanaantare.na na g.rhyate svayam aatmanaa^ eva g.rhyata^ iti tathaa 
j~naanavad arthasya^ api svaya.m graha.na-prasa"nga.h./

[16,8]	atha diipavat svayam avabhaasate ca^ avabhaasayati ca 
avabhaasaatmakatvaat arthasyaanavabhaasaatmakatvaad ado.sa^ iti cet tathaapi 
prakaa"sa-svaruupe.sv arthe.su vij~naanenaagraha.na-prasa"nga.h./

[16,11]	atha jjaanasya^ aatma-samavetatvaavabhaasana-ruupaabhyaa.m svaya.m graha.na.m 
bahi.hprakaa"saatmakaanaa.m tu j~naana-vipariita-dharnatvaad ado.sa^ iti cet tac ca na 
yadi ca^ aatma-samavetatva.m prakaa"saatmakatva.m ca j~naanasya svaya.mgraha.ne 
hetu.h syaat sarve sarvajjaa.h praapnuvanti./

[16,14]	atha graahyaakaaratvam api kaara.na.m cet tac ca na^ anyasya 
j~naanasyaanyaakaaratayaa kara.natva.m yathaa tathaa j~naanasya^ api 
graahyasyaanyaakaare.naanyad eva j~naana.m kara.na.m syaat tathaa ca^ anavasthaa 
du.spariharaa^ eva./

[16,17]	ki.mcaanyat na hi ka"scit svaya.mgraahya-graahakatva-prasaadhanaaya 
d.r.s.taanto vidyate pradiipasya^ api cak.sur-aadi-graahyatvaat. na ca^ apy 
a.m"sa-dvayam asti j~naanasya niravayavatvaat. saty api ca^ a.m"sa-dvaye tayo.h 
prakaa"saatmakatvaan naanyonya-graahya-graahaka-bhaava.h pradiipayor iva dvayo.h./

[16,21]	ki.m ca^ anyat j~naanasya graahya-graahakatva-abhyupagame saty 
aatma-sad-bhaavo api dussampaada^ eva. tasmaad 
arthaakaara-j~naana-sm.rti-pratisandhaana-prayotnecchaa-aadiinaam anaatma-dharmatva.m 
graahyatvaat ruupa-aadiinaa.m iva. paraarthatvaac ca sa"nghaat-aa"srayatva.m 
ruupa-aadi-nidar"sanena^ eva. tathaa^ aa"srayavattvaad anityatvaac ca 
prayatra-puurvakatvaac ca^ ity evam-aadibhyo^ anaatma-dharmatva-siddhi.h./

[16,25]	ayaskaanata-ma.ni-kalpam./ ayakaantama.nir eva. yathaa^ ayakaantama.nir 
ayasa.h svaatma-vyaapti.m pratyupakaroti sannidhimaatra.na tathaa cittam api 
cit-svaruupasyaatmana.h svaatma-vyaapti.m pratyupakaroti d.r"syatvena./

[17,9]	nartakiiva dar"sita-vi.sayatvaad anta.hkara.nenety uktam tatra 
dar"sayit.rtva.m cittasya dra.s.t.rtva.m ca puru.sasya ki.mk.rtam ? 
svasvaamitvak.rtam. katha.m puna.h svasvaamitvam ? vastu-svabhaavatvaat. katham 
it ced ukta.m cittam ayaskaantama.nikalpam iti./ atha citta-v.rttiinaa.m bodhe 
ko hetu.h puru.sasya^ iti ? ucyate yasmaad eva.m sannidhi-maatra-upakaari citta.m 
d.r"syatvena tasmaac citta-v.rttiinaa.m bodhe puru.sasya^ aya.m hetu.h bhavati yo 
'sau^ anaadi-sambanda.h./ 4./

% \finisheditionchapter
%
% \chapter{Critical Text 1.5}%
% \renewcommand{\rightmark}{\MakeUppercase{1.5]}}
% \starteditionchapter
% \input{text05}
% \finisheditionchapter
%
% \chapter{Critical Text 1.6}%
% \renewcommand{\rightmark}{\MakeUppercase{1.6]}}
% \starteditionchapter
% \input{text06}
% \finisheditionchapter
%
% \chapter{Critical Text 1.7}%
% \renewcommand{\rightmark}{\MakeUppercase{1.7]}}
% \starteditionchapter
% \chapter{\pysv\ 1.6--7}
\setlength{\footmarkwidth}{1.8em}
\pagewiselinenumbers


%[18,25]

\section{Sūtras 1.6 and 7 plus the Bhāṣya up to the definition of \emph{pratyakṣa}}
% This is found on Tm8r3 (right); L8r8 (L8r has 11 lines!)
\mll{BHINT6}
\bht{kāḥ punas tāḥ kliṣṭākliṣṭāḥ pañcaprakārā vṛttaya} ity āha---\edn{don't forget the reference in BSBh 2.4.12: nanv atrāpi śrotrādinirapekṣā bhūtabhaviṣyadādiviṣayā aparā manaso vṛttir astīti samānaḥ pañcasaṃkhyātirekaḥ; evaṃ tarhi 'paramatam apratiṣiddham anumataṃ bhavati' iti nyāyāt ihāpi yogaśāstraprasiddhā manasaḥ pañca vṛttayaḥ parigṛhyante -- 'pramāṇaviparyayavikalpanidrāsmṛtayaḥ' nāma;}
\donote{BHINT6}{kāḥ punas tāḥ kliṣṭākliṣṭāḥ pañcaprakārā vṛttayaḥ}

% ya iti starts Tm8r4
\mll{STR6}%
\str{pramāṇaviparyyayavikalpanidrāsmṛtaya} iti||6||%
\donote{STR6}{\textbf{pramāṇaviparyyaya\- vikalpa\-nidrāsmṛtayaḥ||6||}}

%[18,27]
\vrt{etāvatya eva}{\Tm\Me}{\rdg{etāvaty eva}{L}} vṛttayaḥ| tatra—\edn{should this tatra be part of the sūtra or the bhāṣya?}

\mll{STR7}%
\str{pratyakṣānumānā\vrtms{gamāḥ}{\Tmpc L\Me}{\rdg{°gamā`ḥ'}{\Tm}} pramāṇāni||7||}%
\donote{STR7}{\textbf{pratyakṣānumānāgamāḥ pramāṇāni||7||}}%

\vrt{pramāṇākhyā}{\eme}{\rdg{pramākhyā}{\Tm L\Me}} \mds vṛttis tridhaiva bhidyate| tatra pramā\mde ṇākhyāyāś \vrt{citta\-vṛ\-tteḥ}{\Tm\Lpc\Me}{\rdg{cittavṛtte`ḥ'}{L}} prathamo bhedaḥ %(19,5)
pratyakṣam| atas tallakṣaṇam evābhidhīyate tatpūrvvaka\-tvād bhedadvayasya||
\edn{That all the other pramāṇas presuppose pratyakṣa is not particularly special but it is interesting nonetheless that our author makes an explicit statement.}%

\mll{PRTYDEF}
\bht{indriya\vrtms{pranāḍikay\mds ā}{L\Me}{%
	\rdg{pranāḍikāy\lost{1}}{\Tmac},
	\rdg{pranāḍikay\lost{1}}{\Tmpc}}}---%
indriyam iti \mde \vrt{śrotrādi}{\Tm\Me}{%
	\rdg{śrotrābhi°}{L}}
pañcakam, na karmme\-ndriyam, tasya cittavṛtya\vrtmm{hetuka}{\Tm L}{%
	\rdg{°hetu°}{\Me}}%
\mds tvāt|
\vrtsm{tat}{L}{%
	\rdg{\om}{\Me}}%
kriyārtthaṃ hi tat|
\edn{Perhaps tatkriyārthaṃ tat is a good reading. Cittavṛttyahetuka should refer to śrotrādi (they do not act because of cittavṛtti); "they (śrotrādi) exist for the sake of that (cittavṛtti)" remove "ka"! remove the first tat!}%
tasmād indriya\-śabdasya sā\-māny\mde ārtthatve ’pi \vrt{cittavṛtter}{\Me}{%
	\rdg{ucitavṛtter}{\Tm L}%,
	% \rdg{cittavṛtter}{\Me}%
	}\edn{tu and u are hardly distinguishable. This reading might be justified. However, it would be a little peculiar syntax...}
upa\mds vyācikhyā\mde \vrtmm{sitatvā}{\Tmpc L\Me}{%
	\rdg{°sisatvā°}{\Tmac}}%
c chrotrādi\vrtms{viṣayataiva}{\Tm\Me}{%
	\rdg{°viṣatayaiva}{L}}\edn{The verb upavyākhyā appears in the ChU. Due probably to this, Śaṅkara uses this verb in a technical sense, it seems. Now I have a note on this peculiar verb.}%
| vṛtyartthatvāc ca śrotrādī\mds nām|\edn{again, here appears to be a support that śrotrādi are there for the sake of cittavṛtti}
indriyam eva pranāḍikā dvāram---%
śabdādyā\mde kāra\-vṛttirūpeṇa pariṇamam\mds ānasya \bht{cittasya}|
\vrt{atas}{L}{%
	\rdg{atas tat}{\Me}}
tenendriyadvāre\-ṇa \bht{sāmānya\-vi\-śeṣ\mde ātmakabā\mds hyava\mde stv}ā\mds kāratayā pariṇamamā\mde nam upa\-ra\-jyate|
tasya tad\-\bht{upa\-rāgād} dhetoś cittasya yā % L8v1 starts after this
mudrā\mds pratimudrāvad \bht{vṛttiḥ}| sāmānya\-vi\-śeṣātmaka\-va\mde stūparāge 'pi
\bht{%
	\vrtsm{viśeṣāva\mds dhāraṇa}{L\Me}{%
		\rdg{vi\wmhl{1}eṣāva\lost{4}}{\Tm}}%
	pradhānā%
} saiva \bht{pratyakṣam pramāṇām}||% remove danda before?
\donote{PRTYDEF}{indriyapranāḍikayā cittasya sāmānyaviśeṣātmakabāhyavastūparāgād  viśeṣāvadhāraṇapradhānā vṛttiḥ pratyakṣam pramāṇām||}
%Bhāṣya: indriyapranāḍikayā cittasya sāmānyaviśeṣātmakabāhyavastūparāgād  viśeṣāvadhāraṇapradhānā vṛttiḥ pratyakṣam pramāṇām

\section{\emph{Pratyakṣa} and \emph{viśeṣa}}
\subsection{Why the Bhāṣya says \emph{pratyakṣa} is mainly concerned with difference}
%[19,13]
atha svaviṣayasāmānyaviśeṣāvabhāsakasāma\mde rtthye śrotrādīnāñ cittasya ca
\vrt{kiṃ\-kṛtaṃ}{L\Me}{%
	\rdg{kiṃ\wmhl{1}tam}{\Tm}}
viśeṣāvadhāra\mds ṇaprādhānyam iti---\edn{The choice of the verb avabhās and "ka"?}

ucyate---viśeṣapratibandhatvād\mde\ dhāno\-pādān\mds ādī\mde nāṃ vi\mds
śeṣāvadhāraṇa\vrtms{pra\-dhā\-nety| āha---}{L}{%
	\rdg{pradhānateti (tetyāha)}{\Me}}%
\edn{Something is lost, around here, I think. Something like "(viśeṣāvadhāraṇapradhāneti) yuktam| tadā kiṃ sāmānyan nāvadhāryata ity (āha)" Or simply replace āha with uktam? without the question mark; just divide by a comma? āha part of the previous sentence?}%
na sāmānyaṃ nāvadhāryyate? guṇabhūtan tu gṛhyata eva| yathā nīlaṃ rūpam iti
nīl\mde āvadhāraṇa\vrtms{prādhānye}{\Tmpc L\Me}{%
	\rdg{°pradhānye}{\Tmac}}
saty api rūpam eva \vrt{nīlam}{L\Me}{%
	\rdg{\wmhl{1}lam}{\Tm}}
iti guṇabhūtam a\mds vadhārayati tathā duḥkhaśabdaḥ sukhaśabdo mūḍhaśabda i\mde ti ca| \mds
% Tm8r left is over. This was the last line (line 9).
tathā ca viśeṣāvadhāraṇapradhānā vṛttir iti hy upapadyate||

%[19,20]
kiñ ca saṃśayaviparyyayanimittatvāc ca sāmānyopalabdheḥ| % L8v4 starts after tvā
\edn{"kiñ ca" and "...tvāc ca" are redundant.}%
\vrtsm{sāmānye hy upa}{\Tmac\Tmpc L\Me}{%
	\rdg{sāmānye \ccs hy u\cce `hy u'pa°}{\Tm}}%
labdhe bubhutsitaviśeṣasya saṃśayo jā\mds yate|
% Tm8v1 after jā
samupalabdhasāmānyasmṛter vvā viparyyayaḥ| samadhigataviśeṣasya tu \vrt{saṃśayaviparyayau}{\Me}{\rdg{saṃśayaviparyayā}{L}} na sthitim
% L8v5
upā\-śnuvāte iti viśeṣāva\mde dhāraṇapradh\mds ānatocyate||
\edn{Sāmānya being the origin of saṃśaya and viparyaya could be a unique point of view. At least our author is following something. I should figure this out.}%

%[19,24]
yatrāpi sāmānyam evāvadidhārayiṣitaṃ gaur iti vāśva \vrt{iti vā}{\Me}{%
	\rdg{iti vātra}{L}}%
, tatrāpy anyata\-ra\-parityāgenānya\vrtms{tarāvadhāraṇeti}{\Me}{%
	\rdg{°tarāva\ilg dhāraṇeti}{L}} % L8v6 starts after vi next line
viśeṣāvadhāraṇapradhānataiva| sarvam e\-va hi \vrt{vastu}{\Lpc \Me}{%
	\rdg{vastu \ccs ta\cce}{L}}
vastvantarāpekṣayā sāmānyaṃ viśeṣaś ca||

%[19,27] A new view
\subsection{Criticism of the view that \emph{sāmānya} is unreal (Buddhist)}
yasya punaḥ sāmānyam avastv adhyāropitam evoṣarāmbuvat nendriyaviṣayam iti|
% L8v7 after ta next line
\edn{Certainly a Buddhist but who? Dignāga?}%
\edn{The matter is a bit more complicated. See the Tarkabhāṣā of Mokṣākaragupta toward the end of pratyakṣa and Kajiyama's notes 131ff (section 7.1; pp. 56ff.). According to Kajiyama, Dharmottara and Ratnakīrti contributed to Mokṣākara's saying indriya also takes sāmānya as object but I don't see those two clearly saying it in the passages cited by Kajiyama. The criticism that Buddhist won't acquire sāmānya to form inferences is already in Kumārila.}%
tasya pratiniyata\-sāmānyā\-dhyā\-ropaṇā\-nupapattiḥ| na hy anupala\mde bdhapū\-rvvamu\mds khye \mde sambhava\mds ty adhyāropaṇā, smṛtyasambhavāt|
\edn{The meaning is clear but the use of the compound anupalabdhapūrvamukhya and its association of smṛti might be of interest}%
na kadācid adṛṣṭamukhyodakasya salilā\-dhyā\-ropaṇam ūṣareṣu vidyate||
\edn{Is he talking about some sort of saltlake like places? The example used by the Buddhist means something like positing (salty) water in salt rock desert?}%

%[p]
athāsmarann apy adhyāropayet||
\edn{This, too, must have some background...}%

%[s]
rūpasvalakṣaṇe śabdatvādy apy adhyasyet||
\edn{Again the meaning is more or less clear but apparently this is a very abridged form of a more lengthy discussion.}%

athendriyāntara\mde viṣayatvād \vrt{rūpasvalakṣa\mds ṇe}{\Me}{%
	\rdg{rūpasvalakṣa\lost{1}}{\Tm},
	\rdg{rūpasvalakṣaṇa°}{L}}
\vrt{śabdatvādyadhyāso}{\Lpc\Me}{%
	\rdg{śabdatvā\ccs du\cce dyabhyāso}{L}}
viruddha iti cet---na, avastutve % L8v9 from saty
saty aviśeṣāt| \vrt{viśeṣābhyupagame}{\Lpc\Me}{\rdg{viśeṣā\ccs bhyu\cce bhyupagame}{L}} ca vastutvaprasaṃgaḥ| śa\-bdasvalakṣaṇād iva rūpasvalakṣa\mde ṇasya||
\edn{OK, here is more detail. Not necessarily an abliged discussion. Do we see something like this elsewhere?}%

%[20,6]
ki\mds ñ cānyat---\vrtsm{deśakālānyatva}{\eme}{\rdg{deśakālānya°}{L\Me}}viśi\mde ṣṭa\vrtms{pratyakṣābhyupagame}{\Tm L}{\rdg{pratyakṣā[na]abhyupagame}{\Me}}
\edn{There are only deśakālānyathā or deśakālānyatva in use it should be deśakālādi. See the next sentence. deśakālādyaviśiṣṭa...!}%
samastaloka\mds vya\-va\-hāra\vrtms{vilopa}{\Me(°paḥ)}{%
	\rdg{°vilopā}{L}}
syāt|
%1 ) L8v10 after hī below
na hīdam aham amuṣminn avakāśe kāle cāsminn adrākṣam etasmān nimittād evaṃviśeṣaṇam iti ca \vrt{smṛtyupapattiḥ|
% 2)
na}{\Me}{%
	\rdg{smṛtyupapattin na}{L}}
ca \mde pratyakṣā\vrtms{dṛṣṭe\-ṣu}{L\Me}{\rdg{°dṛ\wmhl{1}eṣu}{\Tm}} viśeṣyavad vi\vrtms{śeṣaṇeṣu}{L\Me}{\rdg{°śeṣaṇe\ccs nu\cce ṣu}{\Tm}} smṛtiḥ saṃbhavati|
% 3)
na ca kalpa\mds nāpoḍhena vya\-va\-hāra\-\vrtms{gocarātītena}{\Me}{%
	\rdg{gocarātī\wmhl{1}e\wmhl{1}a}{L}} % L8v11 at tī
vyavaharamāṇo dṛśyate|
\edn{conclusion p. 273 of Prasannapadā vol. 1.}%
tasmāl lokaprasiddhapratyakṣānu\-saraṇa\-m eva nyāyyam||

\section{\emph{Yogipratyakṣa} and \emph{mānasapratyakṣa}}
% [20,11]
\subsection{Question: Are they separate \emph{pramāṇa}(s) or kinds of \emph{pratyakṣa}?}
nanu ca yoginān ni\mde rvvikalpasamādhijan
\edn{This terminology reflects that of the Nyāyabhāṣya ./6\_sastra/3\_phil/nyaya/nyayabhasya.txt:ubhayathā darśanam, ``asty ātmā'' ity āptopadeśāt pratīyate, tatrānumānam --- ``icchādveṣaprayatnasukhaduḥkhajñānāny ātmano liṅgam'' iti, pratyakṣaṃ --- yuñjānasya yogasamādhijam ``ātmamanasoḥ saṃyogaviśeṣād ātmā pratyakṣa'' iti/ agnir āptopadeśāt pratīyate ``atrāgniḥ''iti, pratyāsīdatā dhūmadarśanenānumīyate, pratyāsannena ca pratyakṣata upalabhyate/ vyavasthā punaḥ --- ``agnihotra juhuyāt svargakāmaḥ'' iti, laukikasya svarge na liṅgadarśanam, na pratyakṣam/ stanayitnuśabde śrūyamāṇe śabdahetor anumānam, tatra na pratyakṣam, nāgamaḥ/ pāṇau pratyakṣata upalabhyamāne nānumānam, nāgama iti/}%
darśanan tasyānevaṃlakṣaṇatvāt tathā sukharāgādi%
\vrtmm{vijñāna\mds sya cānindriya}{L\Me}{%
	\rdg{vijñāna\lost{7}}{\Tm}}%
\vrtmm{pranāḍikā}{L}{%
	\rdg{pranāḍī°}{\Me}}%
\mde pūrvvatvāt \vrt{pramāṇāntaratvam}{\Tm\Me}{\rdg{pramāṇāntaram}{L}} a\-bhy\-upetavyam pratyakṣatvaṃ vā vaktavyam|

\subsection{Answer: \emph{Pratyakṣa}} % Change this at least!
% Consider not emending!
ucyate---\vrt{na| pratyakṣasya}{\conj}{\rdg{na pramāṇasya}{\Tm L}, \rdg{pratyakṣasya (na pramāṇasya)}{\Me}} puruṣapratyayāpekṣatvena \vrt{pari\-hṛta\-tvāt|
na}{\Tmpc L\Me}{\rdg{parihṛtatvena}{\Tmac}} hy apratyayā \vrtsm{vṛttiḥ pra}{\Tm L\Me}{\rdg{vṛttiḥ\ittspc{3}pra°}{\Tm}}tyakṣam|\edn{This part should mean "na pramāṇāntaram". Does this mean what may be called yogipratyakṣa or sukhaduḥkha, etc., are not pramāṇa because they do not [generate] a notion in people? And by puruṣapratyaya, does he mean something like hānopādāna? A tentative translation of this sentence would be "No. For, since pramāṇa is expected to [produce] notions [of hāna and upādān] in man, [they---yogins' darśana and sensations such as sukha-duḥkha] do not qualify [as pramāṇa]." Pratyaya in Śaṅkara is always in the form of a sentence? Not really. Even in the Vivaraṇa itself, there are usages like ahampratyaya or ihapratyaya.}\edn{20161014: The above was probably wrong. The answer "na" is probably OK. This should mean that there is no need to tell that those two are part of pratyakṣa or there is no need to postulate (an)other pramāṇa(s) to incorporate those. "parihṛtatvāt" means that those two possiblities are precluded because of the definition(?) that follows}\edn{Also, this part appears to tell a rathe unique view on pratyakṣa. Do the followings help? BSBh 1.3.33 pratyayāpratyayau hi sadbhāvāsadbhāvayoḥ kāraṇam; nānyārthatvam ananyārthatvaṃ vā; BSBh 2.2.25, 29?}
\edn{Does he mean anumāna, śabda/āptavacana, etc., to be apratyaya?}%
caitanyodaya\vrtms{hetor eva}{\Tm\Lpc\Me}{%
	\rdg{he\ccs tu\cce tor eva}{L}}
sapratya\-yāyāḥ pra\-tya\mds kṣatva\mde m| % L9r2 starts at yāḥ pra
\edn{This sentence basically says the same thing as the previous one.}%
tathā cāha—%
\mll{PRMPHL}%
\bht{phalan tadaviśiṣṭaḥ pauruṣeyaś citta\-vṛtti\-bodha} iti||%
\donote{PRMPHL}{phalan tadaviśiṣṭaḥ pauruṣeyaś citta\-vṛttibodhaḥ}

evañ ca sati sukharāgādijñānasya
\edn{We could emend this to sukharāgādivijñānasya but may not be necessary. The two expressions mean the same thing.}%
\vrt{kliṣṭākliṣṭarūpasya}{\Tm\Me}{\rdg{kliṣṭaṃ kliṣṭaṃ rūpasya}{L}}
\edn{kliṣṭākliṣṭa is an expression to qualify all the vṛttis. See YBh 1.3(?).}%
\vrtsm{tadaviśiṣṭa}{\Tmpc\Me L}{\rdg{tadaviśiṣṭa\ccs viśeṣa\cce}{\Tm}}puruṣapra\-tyayaphalāvasānatvāt pratyakṣatā \vrtsm{siddhai}{L\Me}{\rdg{siddhe°}{\Tm}}va||
\edn{Is the argument as follows?: The Bhāṣya says the result of pratyakṣa is tadaviśiṣṭaḥ pauruṣeyaś cittavṛttibodhaḥ; any cittavṛtti that produces that sort of result qualify as pratyakṣa; the cognition of sukha, rāga, etc., produce such results; therefore they are pratyakṣa? They are not another kind of pramāṇa because...they can be categorized in pratyakṣa?}%
% L9r3 ends with siddhaive


%[21,1]
indriyapranāḍikāgrahaṇaṃ laukika\vrtms{pratyakṣodbhāvanānuvāda}{\Tm}{%
%	\rdg{pratyakṣotbhāvanānuvāda}{\Tm},
	\rdg{pratyakṣotbhāvānuvāda}{L},
	\rdg{pratyakṣodbhavānuvāda}{\Me}}
eva| anyathā hī\-ndriya\vrtms{nirapekṣasya}{\Tmpc L\Me}{%
	\rdg{nirapekṣa\ccs tva\cce sya}{\Tm}}
\vrtsm{yogīśvara}{L\Me}{%
	\rdg{\wmhl{2}gīśvara°}{\Tm}}%
\vrtmm{yor jñāna}{\Tmpc L\Me}{%
	\rdg{yor jjana°}{\Tmac}}%
syāpratyakṣatā syāt|
\edn{So, the first category in quesion, yogipratyakṣa, indeed is pratyakṣa.}%
% Tm9r starts at ādipratibandha
asti hi kleśādi\vrtms{pratibandhābhāve}{\Tm\Lpc\Me}{%
	\rdg{°pratibandhābhāva}{\Lac}}
sarvvārtthena cittasatvena sarvvaviṣayeṇa % L9r4 after satvena sa
yugapatsūkṣma\-vyavahi\-tā\-di\-grahaṇam| tathā cāha---\vrt{tārakaṃ}{\Me}{\rdg{kārakam}{\Tm L}} sarvaviṣayaṃ \vrtsm{sarvvathāviṣayam a}{\Tmpc\Me L}{\rdg{sarvvathāviṣaya`ma'°}{\Tm}}kramañ \vrt{ceti}{L\Me}{\rdg{ce\wmhl{1}}{\Tm}} \vrt{vivekajañ}{L\Me}{\rdg{vi\wmhl{1}ekajaṃ}{\Tm}} jñānam iti| \vrt{tasmān}{L\Me}{\rdg{\wmhl{1}smān}{\Tm}} nāvadhāryyata indriyapranāḍikayai\-veti||
\edn{The use of iti at the end of this paragraph appears interesting.}%

%[21,6] L9r5 starts after bāhya
bāhyavastūparāgād iti viparyyaya\vrtmm{nivṛtyarttham| viśeṣā}{L\Me}{\rdg{°nivṛtyartthaviśeṣā°}{\Tm}}vadhāraṇapradhāna\-tvena saṃśayanivṛttir iti||
\edn{again iti at the end}%

\section{\emph{Pramāṇaphala}}
\subsection{Two views: \emph{pramāṇa} and its \emph{phala} are identical; there is an independent \emph{phala}}
%[21,8]
tatra \vrt{keṣāñ cit}{L\Me(°āṃcit)}{\rdg{ke\wmhl{2}ñcit}{\Tm}} phalan tad eva pramāṇam|\edn{I wonder if the author of the Bhāṣya in fact had Dignāga in mind. The order of the whole commentary (Vivaraṇa) on pratyakṣa reflects that of the Pramāṇasamuccaya. I wonder if this indeed reflects the thought of the Bhāṣyakāra. This seems unlikely. It was the Vivaraṇakāra's idea that somehow the definition in the Bhāṣya should cover yogipratyakṣa and mānasapratyakṣa. The Bhāṣyakāra apparently had no idea his definition would become problematic.}
\edn{Is he basically following the Pramāṇasamuccaya? PS 1.8}%
apareṣān tu pramāṇād \vrt{anyatra}{\Tm L}{\rdg{anyat(tra)}{\Me}} \vrt{pramāṇavyaktyantare}{\Tm L}{\rdg{pramāṇavyaktyantaram(re)}{\Me}}|
\edn{The above sentence looks really weird. I should be looking at the Tantravārttika, Prakaraṇapañcikā, NVTT, NBh, NMañ for hānopāānopekṣā. For pramāṇavyakti, see the TSP, NVTTP.}%
% L9r6 after gha below
\edn{Our author does not like either view.}%
yathā ghaṭavijñānasya pramāṇāsya
\vrt{hānopādānopekṣā}{\Me}{\rdg{hānopādānāpekṣā}{\Tm}, \rdg{hāno\ccs cā\cce pāadānopekṣā}{L}} buddhayaḥ phalam iti
\edn{This is a quote from the Nyāyabhāṣya! No emendation???}%
\vrt{guṇadoṣatadubhayaśūnyatvaviṣayatvād}{\Me(\eme)}{%
  \rdg{guṇadoṣavadubhayaśūnyatra`tva'viṣaya\-tvād}{\Tm},
  \rdg{guṇadoṣavadabhavaśūnyatvaviṣayatvād}{L},
  \rdg{guṇadoṣatadubhaya(vadabhava)śūnyatvaviṣaya\-tvād}{\Me}%
  }
upādānādibuddhīnām
\vrtsm{pramāṇabhūta}{\Tmpc L\Me}{\rdg{brahmāṇaprata°}{\Tmac}}ghaṭavi\mds jñānād
\vrt{a\mde tyantabhedaḥ}{\Lpc\Me}{%
	\rdg{\lost{1}bhyanta\wmhl{1}e\wmhl{2}}{\Tm},
	\rdg{atyanta\ccs meva\cce bhedaḥ}{L}}%
||
\edn{He is in fact following this Naiyāyika idea of hānopādānādibuddhis being the phala of pramāṇa? Most likely not. At the end, he says "nāpi tadvyatiriktapramāṇāntare."}%

\subsection{There is no independent \emph{pramāṇaphala}}
\vrtsm{ataś ca nānya}{\eme}{\rdg{\wmhl{1}taś cānya°}{\Tm}, \rdg{tataś cānya°}{L\Me}}%
% L9r6 after vi below
\vrtms{viṣayasya}{\Tmpc L\Me}{\rdg{°viṣayīsya}{\Tmac}}\edn{This should refer to another reason that is going to be stated why pramāṇa and pramāṇaphala are not two different things. Therefore, tataś ca (for the same reason) does not seem right. L9r6 after vi below}
\vrt{pramāṇasyānyaprameyasya}{}{%
  \rdg{pramāṇasyānyaprameyasya viṣaye}{\Tm L},
  \rdg{pramāṇasya anya(prameyasya)viṣaya°}{\Me}
}
\edn{Could anyaprameyasya an interpolation from a marginal comment?}%
\vrt{pramāṇavyaktyantare}{\Tmpc}{%
  \rdg{pramāṇavyaktyantara°}{\Tmac},
  \rdg{pramāṇe vyaktyantare}{L},
  \rdg{°pramāṇa(ṇe)vyaktyantaraṃ (re)}{\Me}%
} phalam iti|
\edn{The whole of the above is hopelessly corrupt... but at least vyaktyantare phalam or pramāṇavykatyantare phalam appears consistent. So, that reading might be OK.}%
na hi khadiracchedanasya \vrt{khadiropādānan}{\Lpc\Me}{%
	\rdg{khadiropādānān}{\Tmac},
	\rdg{khadiropādā\ccs na\cce nan}{L}} \vrt{tatparityāgo}{\Tm L\Me}{\rdg{tatparityāgo\ittspc{6}}{\Tm}} vā
\edn{The space left in Tm is curious. There could have been something written and erased although doing so without visible traces in a Malayalam manuscript should not be easy because writing literally leaves the surface scratched. A possible candidate that could come there is something like upekṣaṇa.}%
phalaṃ \vrt{bhavituṃ}{\Me}{%
	\rdg{\ilg bhavitu`ṃ'}{\Tm},
	\rdg{bhavatim}{\Lac}}
\edn{I cannot read what the letter before bha is in Tm; Td ignores it.}%
yuktam|
\vrtsm{viyoga}{\Me(\eme)}{%
	\rdg{api yoga°}{\Tmac},
	\rdg{aviyoga°}{\Tmpc L},
	\rdg{vi(api)yoga°}{\Me}%
	}%
sambandhaphalatvāt tyāgo\mds pādānayoḥ|
% L9r8 after tya
\edn{This sentence should explain why tyāga and upādāna are not the result of cutting of khadira tree but the result of ...yogasambandha. This might be simply viyogasambandhaphalatvāt, meaning that a relationship that is detached from khadiracchedana. The decision of whether taking or abondoning the resulting khadira wood comes not from the cutting but from somewhere else.}%
na ca phalasya phalaṃ sambha\mde vati| \mds dvai\mde dhībhāvaphalavyavahitatvāc ca||
\edn{NS/NBh/NV etc., BAUBh, ĀŚV, Bhāmatī, Upad etc. dvaidhībhāva means "making something into two." Something appears to have gone horribly wrong with these paragraphs. Cutting khadira is a Śābarabhāṣya thing. So, the translation might be something like "[Abondoning or taking of the khadira wood] is detached from the result that is making [the tree] into two."}%

% [21,16]
kiñ ca phala\mds sya \mde cet phalam, kriyānapekṣāprasaṅgaḥ| phalārtthaṃ hi
\vrt{kriyā\-pekṣyate}{\Tmpc L\Me}{%
	\rdg{kriyāpekṣyakṣyate}{\Tmac (\emph{kṣya} appearing twice at the end of a line and at the beginnning of the next line)}}%
| tac ce\mds t pha\mde laṃ vināpi kriyayā phalād api
% L9r9 after kriyayā
\edn{phalād api part should mean something like "simply from the [first] result."}%
\vrtsm{sidhyati, ko 'py a}{\conj}{%
	\rdg{sidhyati ko hy a°}{\Tm},
	\rdg{siddhya ko hy a°}{L}
	\rdg{(ato hy a) tato na hy a°}{\Me}}%
nekāy\mds āsā\-pekṣitaniṣpattikriyām anutiṣṭhet|\edn{No need for conjecture???} \vrtsm{sādhanānve\mde ṣaṇā}{\Me}{\rdg{sādhanānveṣaṇaṣa•ā°}{L}}nupapattiś ca, kri\mds yāyā abhā\-vā\mde t||

%[22,4]
nāpi kriyaiva phalam, \vrtsm{kriyānupa}{L\Me}{\rdg{\wmhl{1}yānupa°}{\Tm}}rama\vrtms{prasaṃgāt}{\Tmpc L\Me}{\rdg{°prasaṃgā}{\Tmac}}|
\vrt{phalāvasānā hi}{\eme}{\rdg{\ccs rtthapā\cce `\ilg\ilg'vasānāpi}{\Tm}, \rdg{phalāvasānāpi}{L}, \rdg{phalāvasanāpi [hi]}{\Me}} \mds kriyā| kri\-yāyāś ca
% L9r10 after phala below
phalatve phalā\mde nabhilāṣitva\vrtms{prasaṅgaḥ}{L\Me}{\rdg{°prasaṃ\wmhl{1}ḥ}{\Tm}}|
\mds tasyā duḥkhatvāt||
\edn{Here it is again. The action being pain.}%

\subsection{\emph{Siddhānta}s}
%[22,6]
\subsubsection{No independent \emph{phala} nor are \emph{pramāṇa} and \emph{phala}  identical}
tasmān na pramāṇasvarūpam eva phalam| \vrt{nāpi}{\Me}{\rdg{nādhi}{L}} tadvya\mde ti\mds rikta\vrtms{pramāṇāntara}{L}{\rdg{pramāṇāntaram}{\Me}} ity āha---\bht{phala\mde n tadaviśiṣṭa} ityādi| \vrt{tadaviśiṣṭas}{L\Me}{\rdg{\wmhl{1}daviśiṣṭas}{\Tm}} tayā \vrt{pramāṇākhyayā}{\Tm\Lpc\Me}{\rdg{pramāṇākhyāyā}{\Lac}} vṛ\mds tyā\-viśi\-ṣṭa\-s tulyas tadākāratvena|
% L9r10 after tadākā
\bht{pau\mde ruṣeyaḥ} puruṣasyāyam \mds bodha iti| puruṣasya hi vṛtyākāratā|
\edn{So far, it appears that our author is saying that in pratyakṣa (or in pramāṇa in general?), bodha and pramāṇa (a vṛtti) share the same ākāra, and the bodha and puruṣa share the same ākāra.}%

\subsubsection{The \emph{phala} is not the resulting \emph{dravya}, either}
tasmān \vrt{na kriyājan}{\eme}{\rdg{nākriyājaṃ}{L\Me}} dravyasvarūpamātraṃ phalaṃ|

\subsubsection{The phala is a special state of the \emph{dravya}}

kin tarhi dravyasyai\mde
vā\vrtms{vasthāntaraviśeṣaḥ kriyājaḥ kriyāparyavasānakāle}{\eme}{%
	\rdg{°vasthāntaravi\wmhl{2}ṣa\ccs kriyā\cce\ins(kāryyā)\ine jaḥ \wmhl{1}ryyavasānakālaḥ}{\Tm},
	\rdg{°vasthāntaraviśeṣakriyājākriyāparyāvasānakālaḥ}{L}, %
	\rdg{°vasthāntaraviśeṣaḥ (ṣa)kriyājaḥ(jā)kriyāparyāvasānakālaḥ}{\Me}%
	}
kri\mds yo\vrtmm{parāmopa}{L}{%
	\rdg{°paramopa°}{\Me}}%
patteḥ phalam|\edn{Is this a parallel? ./6\_sastra/3\_phil/vedanta/atmasiddhi.text:yastu sugatamatāvalambī vijñānābhinnahetujatayā tayorapi tadantarbhāvam abhimanyate; kaṇabhakṣapakṣāśrayaṇena vā tayorātmaviśeṣaguṇatvam;tābhyāṃ sukhaduḥkhādhikaraṇaṃ vyācakṣīta; svataḥsukhītyetadvimarśaṃ vātratyam / rāgadveṣādayastu caitanyasyaivāvasthāviśeṣāstadvadeva pratyakṣībhavantīti na tannidarśanenānumānodayaḥ / sukhaprayuktaviṣayīkāracaitanyaṃ rāgaḥ, tadvirodhaprayuktaviṣayīkāraṃ tadeva dveṣaḥ, bhūtaduḥkhajñānena cetaścalanaṃ śokaḥ, āgāmitajjñānena cetaścalanaṃ bhayamityādi lakṣaṇagranthādevāvagantavyamityalaṃ pravistareṇa / Or in BSBh, look for avasthāntara. It is the cause of bhedābhedabhāva Grammarians? About odana and phala} % L9v1 before tteḥ
phalaniṣpattāv eva \mde hi \vrt{kriyoparāma}{L}{%
	\rdg{ṣi`(kri)'yoparāma}{\Tm},
	\rdg{kriyoparama}{\Me}}
\mds upapadyate|
% bodhaḥ should be typeset in bold

\subsubsection{Either \emph{kriyā} or \emph{dravya} is figuratively the \emph{phala} of the other}
evaṃbhūtasya mukhyaphalatve
\vrtms{pradhānakriyāvyavahitaviprakṛṣṭāyāḥ}{\Lpc}{%
	\rdg{pra\ccs mā\cce dhānakriyāvyavahitaviprakṛṣṭāyāḥ}{L},
	\rdg{pradhānakriyāvyavahita(viprakṛṣṭāyāḥ)°}{\Me}}
kriyāyā dravyan dravyasya vā kriyā phalam ity \vrtsm{a\mde nyonya}{L\Me}{%
	\rdg{\lost{1}nyo\ins nya\ine °}{\Tm}}% L9v2 after sthāpa
pratyupasthā\vrtms{pakatvena}{L\Me}{%
	\rdg{°pakatv\wmhl{1}na}{\Tm}}
pra\-dhāna\-phalābhimukhī\vrtms{karaṇāt}{\Tmpc L\Me}{%
	\rdg{karaṇā\ins t\ine}{\Tm}}
pha\mds latvopacāro na virotsyate||
\edn{So, basically the three things are not the phala: kriyā pramāṇa dravya. The last part I cannot say I fully understand but one can figuratively say either dravya or kriyā is the phala? I still don´t understand in what context our author is talking about the phala. Regular kriyā or specifically pramāṇa? Is there the picture of pramāṇa to kriyā to the change in dravya? Looks like pradhānakriyā is pramāṇa and pradhānakriyāvyavahitaviprakṛṣṭakriyā is the secondary;}%

\paragraph{Criticism: the \emph{phala} is only figurative?}
%[22,14] The following criticism appears reasonable. At any rate, what is not pramāṇaphala is more or less over.
nanu ca bhavato ’py upacaritam eva phalam, vṛtyaviśiṣṭāyāḥ
\edn{buddheḥ? The bhāṣya actually reads "pauruṣeyo bodhaḥ." So, it should be masculine. Or he means vṛttyaviśiṣṭāyā avasthāyāḥ? <- That should be it!}%
\edn{Here is a nice one from the Upad: ./upad.txt:atrāha codakaḥ---avagatiḥ pramāṇānāṃ phalaṃ kūṭasthanityātmajyotiḥsvarūpeti ca vipratiṣiddham| ity uktavantam āha---na vipratiṣiddham| kathaṃ tarhi| kūṭasthanityāpi satī pratyakṣādipratyayānte lakṣyate, tādarthyāt| pratyakṣādipratyayasyānityatve 'nityeva bhavati| tena pramāṇānāṃ phalam ity upacaryate|108|}%
phalatvenābhyupagamāt| puruṣasya ca \vrtsm{śuddhatvāt, tadākāratāyā anṛta}{\Me}{%
	\rdg{śuddhatvād akāratāyānṛta°}{L}}%
tvam| na hi sphaṭikamaṇer % L9v3 after ma
a\mde laktakopadhānā\vrtms{kāratā}{L\Me}{\rdg{°k\wmhl{1}ratā}{\Tm}} satyā||
\edn{upādhāna?}%

\paragraph{Answer: yes}
%[22,17]
bāḍham| tathāpi tv etad eva %Tm9v1 starts here
phalam mukhyam, dṛṣṭatvāt, mukhyaphalāntarādarśanāt, sarvvasyaiva kriyākārakajātasyaitatpha\-lārtthatvāt| etadavasānaṃ hi tat sarvvam|
% L9v4
tathā ca pradarśitam||

\paragraph{Criticism: then, the result being more or less erroneous, no reason to pursue \emph{pramāṇa}}
%[22,20]
nanu ca \vrtsm{mithyā}{\Me}{\rdg{mitthyā°}{\Tm L}}bhūtasyaiva mukhyaphalatvaṃ vṛtyāk\mde ārasya,\edn{This notion that pramāṇa is anyway wrong is present in the intro to the BSBh.}
\edn{vṛtyākāra was puruṣa. This sentence does not make much sense.}%
tathā ca tadabhilāṣānupapattiḥ|
\edn{A tentative translation: For [puruṣa] who takes the shape of the [citta]vṛtti, the main result [then] is nothing but unreal. That way, seeking [the main result of pratyakṣa/pramāṇa] will be impossible. doesn't tad refer to puruṣa?}%

\paragraph{Answer: correct}
bāḍham| \vrtsm{mithyā}{\Me}{\rdg{mitthyā°}{\Tm L}}bhūta % Tm9v2
\mds eva vṛtyanugato bhogaḥ| tathā cāha---satvapuruṣayo\mde r atyantāsaṃkīrṇa\mds yoḥ pratyayāviśeṣo bhoga % L9v5 starts at rṇṇa
\edn{YS 3.35}%
 iti| bhā\-ṣyam api---tayoḥ svarūpopalambhe sati kuto bhoga iti||
\edn{What he is doing in these paragraphs is that pramāṇa is essentially mithyā for mixing up puruṣa with cittavṛttis. Pratyakṣa/pramāṇa does not lead to mokṣa. It rather leads to saṃsāra.}%

\paragraph{\emph{samyagdarśana} (rather than \emph{pramāṇa}/\emph{pratyakṣa}) brings about \emph{nivṛtti}}
ata eva \vrt{ca}{\Me}{\rdg{\om}{\Tm L}} samyagdarśanād \mde ātyantikī \vrt{nivṛttir}{L\Me}{\rdg{ni\wmhl{1}ttir}{\Tm}} upapadyate||
\edn{If we adopt the reading na, samyagdarśana means pramāṇa, not vivekakhyāti. But the following suggest ca is better.}%

\paragraph{\emph{Mokṣa} is brought about only by \emph{samyagdarśana} (not by \emph{pramāṇa} that leads to action)}
%[22,25]
yeṣām \vrt{amithyā}{L\Me}{\rdg{amitthyā}{\Tm}} phalan \vrt{te\mds ṣāṃ}{L\Me}{%
	\rdg{te\lost{1}e\lost{1}}{\Tmac}, % L9v6 before vairāgya
	\rdg{te\lost{3}}{\Tmac}}
vairāgyakāraṇābhāvāt, samyagdarśanā\mde nveṣaṇānupapatte\mds r mmokṣābhāvaḥ| na ca kriyāpariniṣpādyo mo\-kṣaḥ, anityatvaprasaṅgāt||
\edn{This is such a typical Śaṅkara thing to say!}%

\subsection{The reason why \emph{puruṣa}’s nature, being aware of \emph{buddhi}, is not established here/Explanation of the intent of the last statement with regard to \emph{pratyakṣa} in the Bhāṣya}
%[23,4]
atha buddhiboddhṛtvam puruṣasya sa\mde tvaṃ cāsiddham iti cet,
% L9v7 after si above
āha---%
\mll{PRTSMVDBU}%
\bht{pratisaṃvedī buddheḥ %Tm9v4
pu\mds ruṣa ity upariṣṭāt pravedayiṣyāma} iti||\mde%
\donote{PRTSMVDBU}{pratisaṃvedī buddheḥ puruṣa ity upariṣṭāt pravedayiṣyāmaḥ}
\edn{This ends the discussion on pramāṇaphala. He does not believe in pramāṇa or pratyakṣa to lead to mokṣa.}%
\edn{The part being referred to in the Bhāṣya is probably not YBh 2.17 but YS 2.20. The Vivaraṇakāra refers to this very passage of YBh 1.7 when commenting on the sūtra (2.20) as the one being referred to be established in the future. YBh 2.20 also has "sa puruṣo buddheḥ pratisaṃvedī." In that sense, whether the reading was pravedayiṣyāmaḥ or upapādayiṣyāmaḥ may be determined by looking at MSS for the part where they record the Vivaraṇa on 2.20. The edition reads upapādayiṣyāmaḥ in YVi 2.20 (p. 189)... So, Td reads upapādayiṣyāma(ḥ) in YVi 2.20 (p. 161/PDF p. 61). Both readings are attested in different manuscripts. Does this mean that our author had access to more than one manuscript of the Yogabhāṣya?}%

\subsubsection{Criticism of the Naiyāyika position of \emph{pratyakṣa}, including their position on \emph{pramāṇaphala}}
% [23,6]
ye \edn{Naiyāyika definition of pratyakṣa; still inside the discussion on pramāṇaphala... Better consult NBh 1.1.4 or thereabout}cāpīndriyārtthasannika\mds rṣasya pratyakṣatvam abhyupayanti,
% L9v8 after kriya
na teṣām api samyag abhyupetam, sannikarṣasya kriyāphalatvāt|
sannikarṣa\mde ś ca saṃyogaḥ|
\vrt{sa ca}{L\Me}{\rdg{\wmhl{1}ca}{\Tm}} kriyājas tair iṣyate|
yato yad upajā\mds yate \vrt{tat tasya}{\Me}{\rdg{tat tasya \ilg}{L}} pramāṇam| \vrt{ta\mde c ca}{L\Me}{\rdg{ta\ittspc{4}c ca}{L}} tasya phalaṃ yad \vrt{upajāyate}{\Tm\Lpc\Me}{\rdg{upajā\ccs\ilg\cce yate}{L}}||
\edn{sannikarṣa = pratyakṣa sannikarṣa <- kriyā; sannikarṣa = kriyāphala sannikarṣa = saṃyoga saṃyoga = kriyāja (kriyāphala) i.e., sannikarṣa/saṃyoga comes from kriyā or kriyāphala, i.e., pratyakṣa is kriyāphala (the result of an action) A is born from B then B is A's pramāṇa; A that is born is the result (phala) of B B -> A then B is pramāṇa and A is phala This is a general principle that both sides should accept but according to the Naiyāyikas saṃyoga = sannikarṣa = pratyakṣa = pramāṇa is a phala of an action not the source (pramāṇa); here is a contradiction}%

% [23,9] L9v9 after iti
phalam api \mds sannikarṣo jñānam prati pramāṇam iti cet---
\edn{Then the Naiyāyikas say sannikarṣa = saṃyoga = pratyakṣa = pramāṇa is a phala (kriyāja) but toward jñāna, it is pramāṇa (source), generator}%

tac ca na,\edn{[Siddhāntin] The key here is the distinction between utpādaka and jñāpaka. Utpādaka should refer to sannikarṣa's being a kāraka} jñ\mde āpakapra\mds māṇa\mde pra\mds stāve kārakapramā\mde ṇopanyāsasyā\vrtms{yuktatvāt}{L\Me}{\rdg{°yu\wmhl{1}\ccs sy\cce āt}{\Tm}}| sannikarṣo hi jñānasyotpāda\mds ko \mde na tasya jñāpakaḥ|
% tasya should be understood as jñānasya; so, "na jñānasya jñāpakaḥ" L9v10 after kāraṇa
prameyādhigamāya hi pra\-tyakṣādīny upa\mds nyastāni, na jñānotpattikāraṇam| \vrt{tatprakriyā\mde yāṃ hy aprakṛtam}{\conj}{%
	\rdg{\lost{4}yāṃ hy aprakṛtaṃ}{\Tm},
	\rdg{tatprakriyāyānyaprakṛtam}{L},
	\rdg{tatprakriyāyām a(yānya)prakṛtaṃ}{\Me}} %
prakriyeta||
\edn{prakriyā is something like a themed discussion? Cf. prakaraṇa? The verb pra-kṛ means something like "to talk about something"?}%

%[23,13]
athā\mds pi syāt\mde ---jñānasyāpi prameyatvāt tac cāpy utpadyamānam pramīyata iti tadutpattikāraṇasya sannikarṣasya jñāpakatvam eveti||
\edn{The point of this opponent's argument is jñānotpattikāraṇa is jñāpaka, and hence sannikarṣa is pramāṇa; no need to distinguish utpādaka and jñāpaka; makes sense, no?}%

tac ca na, jñāna\mds ṃ jñanenaiva jñāyeta\mde, teṣān tathābhyupagamāt|
% L9v11 after jñā of jñāyate above
\edn{Where do the Naiyāyikas teach this? The point should be that one cannot identify jñāpaka and jñānasyotpādaka because only another(?) jñāna is the jñāpaka of jñāna?}%
%
\edn{./6\_sastra/3\_phil/buddh/durvekamisra\_dharmottarapradipa.text:nanu kiṃ ghaḍādidṛṣṭāntabalāt prakāśāntaraprakāśyatā jñānasyāstām, āhosvit pradīpadṛṣṭāntasāmarthyāt prakāśāntaranirapekṣatayā svaprakāśatā? atrocyate | jñānasya svaprakāśarūpatvābhāve ghaṭādeḥ svarūpaprakāśanānupapatteḥ svaprakāśataiveṣṭavyā | atra ca prayogaḥ -- yad avyaktavyaktikaṃ na tad vyaktam yathā kiñcit kadācit kathañcid avyaktavyaktikam | avyaktavyaktikaś cāyaṃ ghaṭādiḥ jñānaparokṣatva iti vyāpakānupalabdhiprasaṅgaḥ | jñānasya jñānāntareṇa vyaktau hetur ayam asiddha iti cen na, nīlajñānodayakāle 'siddhatvād hetor nīlasya parokṣatvaprasañjanāt | na ca bhavatām api mate sarvaṃ vijñānam ekārthasamavāyinā jñānena jñāyate | bubhutsābhāve tadabhāvāt, yathopekṣaṇīyaviṣayā saṃvit | tat upekṣaṇīyam eva tāvad avyaktavyaktikatvād avyaktaṃ prasajyeta | na ceyam ./1\_veda/4\_upa/agamasastra\_w\_sankara.text:sopāyaṃ paramārthatattvaṃ laukikaṃ śuddhalaukikaṃ lokottaraṃ ca krameṇa yena jñānena jñāyate tajjñānam / ./1\_veda/4\_upa/bausbh.text:nanvanyajñānenānyanna jñāyata iti / ./6\_sastra/3\_phil/vedanta/bsbh.text:vyatireke hi satyekavijñānena sarvaṃ vijñāyata itīyaṃ pratijñā hrīyeta  /}%
na ca sannikarṣo jñānam| \vrtsm{jñānotpādaka}{L\Me}{\rdg{jñāno\wmhl{1}ādaka°}{\Tm}}tvenābhyupagamāt||
\edn{So, jñāna, not being a jñāna, it cannot be known}%

% [23,16]
atha \edn{This one sounds alright, again}jñānasya jñāne sannikarṣo \vrt{nimittatvāt pramāṇam}{\Me(conj.)}{\rdg{nimittapramāṇatvam}{L\Tm}} iti \vrt{cet---%
adṛṣṭe}{\Tmpc(ced adṛṣṭe°)\Me}{%
	\rdg{ce`d a'dṛṣṭe°}{\Tm}, %
	\rdg{cedaṣṭe}{L}}%
śvare\-cchāviṣayanimeṣonmeṣādīnām\edn{Perhaps both jñānasya and jñāne are necessary in this context. No need to emend, probably} % L10r1 after ādīnām
api nimi\mds ttatvāt \mde pramāṇatve kaḥ pradveṣaḥ? na ca jñānasya jñeyatva ātmanaś ca svabhāvacidrūpatvānabhyupagame śakyānavasthā parihartum||
\edn{"It is not possible to get rid of the anavasthā if jñāna is to be known and as long as one does not accept ātman being essentially the consciousness."}%

% [23,19]
athātmasamavetamātratayaiva % L10r2 after jñā next
jñānasya jñeyatā| na jñānāntaram apekṣyate yenāti\vrtms{prasajyeteti}{\Tm\Lpc\Me}{%
	\rdg{°prasajyete\ccs te\cce ti}{L}} %
cet---

\vrt{tasyāpy}{\eme}{\rdg{asyāpy}{\Tm L\Me}} anyat pratyakṣalakṣaṇaṃ vācyam| \vrtsm{anindriyasannikarṣādilakṣaṇa}{L\Me}{\rdg{anindriyasa\wmhl{1}ka\wmhl{5}ṇa}{\Tm}}tvād \vrt{ātmani}{L\Me}{\rdg{ā\wmhl{2}}{\Tm}} jñānasamavāyatvasya| sannikarṣasya ca jñānotpatti% This was the end of Tm9v L10r2 starts after jñāno
phalatvāt tadu\-tpattivyavadhāne sati tasya \vrt{jñānasya}{\eme}{\rdg{jñānena}{\Tm L\Me}} \vrt{nānantaryyam}{\Tm\Me}{\rdg{\om}{L}}||
% Was there the end of Tm9v somewhere in this paragraph?

% [23,23]
\vrtsm{tasmāt tadā}{\Tm\Me}{%
	\rdg{\om}{L}}%
nan\vrtms{taryyeṇāpi}{\Tm\Lpc\Me}{%
	\rdg{°taryyaṇāpi}{\Lac}}
na sannikrṣaphalatvaṃ jñānasya| % Were these omissions supplied somewhere in L?
\edn{This ends the discussion on sannikarṣa}%

\paragraph{Examinations of the view that ideas of taking or removing are the result of \emph{pramāṇa}}
tathā hānādibuddhīnām api \vrt{rūpādijñānaphalatvānupapattir uktā}{L\Me}{\rdg{rūpādijñāa\wmhl{8}ā}{\Tm}}|
\edn{Well, the earlier discussion.}%
%[23,25] L10r4 after ca ya
vyabhicārāc ca| \vrt{yathaikasminn}{\Tm\Lpc\Me}{%
	\rdg{yathaivasminn}{\Lac}}
\vrtsm{aṃganādi}{\Tm\Lpc\Me}{%
	\rdg{aṃgānādi°}{\Lac}}%
piṇḍaviṣaye jñāne samāne bahūnām \vrt{utpanne}{L\Me}{%
	\rdg{ulpa\ccs t\cce nn`e'}{\Tm}}
kasyacid upādānabuddhiḥ, anyasya heyabuddhiḥ, kasyacid ubhayābh\mds āvaḥ\mde | tatra hānopādānabuddhir vvyabhicarati|
\vrt{anyataratraikam eva jñānan tad}{L\Me}{% L10r5 before kam eva
	\rdg{anyatara\wmhl{1}ai\wmhl{6}ānan tad}{\Tm}}
\vrt{evobhayaṃ}{L\Me}{%
	\rdg{evo\wmhl{1}yaṃ}{\Tm}}
vyabhicaraty upekṣakasya| yac ca yad vyabhicarati na tat tasya phalam||
% t tasya phalaṃ starts Tm10r3.
\edn{There might be still textual problems but overall the idea is clear.}%

%[24,1]
kiñ ca, hānopādānābhāvamātratvād abhāva evopekṣā| na cābhāvaḥ \mds pramā\-ṇa\-sya va\mde stunaḥ \mds % L10r6 after pha next line
phalam \mde bhavitum arhati| rāgadveṣābhyāñ copādānaparityāga\-bu\-ddhī|
\vrt{tadabhāvāt}{L\Me}{%
	\rdg{tadabhā\wmhl{1}t}{\Tm}}
\vrt{kāraṇābhāve}{\Tm\Lpc\Me}{%
	\rdg{kāraṇābhā\ccs bā\cce ve}{L}}
kāryyābhāva iti pramāṇavijñāne saty % Tm10r4 after sa
api hā\-no\-pādānabuddhyabhāva evopekṣeti
% L10r6 after evope
vyabhicārān na hānādibu\mds ddhiphalatvam||
\edn{"...the notions of hāna, etc., are not phala" Again, the text looks more or less alright. There might be a significance with regard to our author's position on abhāva.... rāga -> upādāna(buddhi); dveṣa -> parityāga(buddhi)/hāna(buddhi)}%

%[24,5]
kiñ ca, rāgadveṣatadabhāvā\mde nām apy upādānādibuddhiphalatve pramāṇatva\-prasaṅgaḥ|
\edn{Err, "There would be the undesirable consequence that rāga, dveṣa, and upekṣā are the pramāṇa with regard to the notions of upādāna, etc., being the result"? The idea seems OK given the diagram above.}%
pauruṣeyacittavṛttibodhaphalatve tu na % Tm10r5 after ki L10r8 after kiñci
kiñcid viru\mds ddhyata iti||
\edn{"But nothing contradicts if [the result/phala is] pauruṣeyacittavṛttibodha."??? iti is again significant here}%

\subsubsection{The problem of \emph{saṃvedyatva} of \emph{cittavṛtti}}
%[24,7] % XYZ
\paragraph{Is \emph{pramāṇa} a \emph{karaṇa}, one of \emph{kāraka}s?}

\mde nanu ca cittavṛtteḥ saṃvedyatvam asiddham| na hi ka\mds raṇam eva karma bhavati| na hi dātram eva lūyate||
\edn{Cittavṛtti and saṃvedyatvam; a new topic So, here cittavṛtti = pramāṇa = karaṇa; cittavṛtteḥ saṃvedyatvam is a problem because of the statement "phalaṃ tad(pramāṇākhyacittavṛtti)aviśiṣṭaḥ pauruṣeyaś cittavṛttibodhaḥ| pratisaṃvedī puruṣa ity upariṣṭāt pravedayiṣyāmaḥ" We have come back to the interpretation of the Bhāṣya For this opponent, the above implies that cittavṛtti is a karaṇa and it cannot be a karman}%

nai\mde ṣa doṣaḥ| kriyākāraka\mds yor bbhinnadha\mde rmmatvāt| kriyā\vrtms{nirvṛttim}{\Me}{\rdg{nivṛttiṃ}{\Tm}}
\edn{kriyānivṛtti or nirvṛtti? Rather nirvṛtti, no? The edition reads that so it is probably my typo when I made the e-text}%
prati sa\-m\-bhū\-ya % L10r9 after sambhūya
kārakāṇi \vrt{vyāpriyante| % Tm10r6 after priya
na \mds parasparaṃ}{L}{%
	\rdg{vyāpriyante na \lost{5}}{\Tm},
	\rdg{vyāpriyante parasparaṃ}{\Me}}
viṣayaviṣayibhāvena|  kri\mde yā tu \vrt{sattā\-mā\-treṇaiva}{\Me}{%
	\rdg{sattāmā\wmhl{1}eṇ\wmhl{2}va}{\Tm},
	\rdg{sattātreṇaiva}{L}}
\mds phalahetur iti karaṇatvam asyām upacaryyate| na
\vrt{kārakavat| kurvvad dhi kārakam| yadi kriyāpi kārakaṃ syā\mde n na}{\conj}{%
	\rdg{[\mist{5}sarvva\mist{2}raka\mist{8}]\ccs sva\cce ān na}{\Tm},
	\rdg{kārakavat kurvvad dhi kāraka yadi kriyādhikākrakaṃ syān na}{L},
	\rdg{kārakavat| kurvad dhi kārayed iti kriyādhikārakaṃ syān na}{\Me}}
bhūtam bhavyāyopadiśyeta|| % L10r10 after bhavyā
\edn{This paragraph is totally obscure... Well, this is from the Śābarabhāṣya! Shows up in the ŚBh, NMañ, Bhāmatī, ŚV commentaries, etc. In essence, the author is saying that pramāṇa is kriyā and therefore not a kāraka and therefore not a karaṇa?}%

%[24,12]
\vrtsm{kiñ cānavasthā}{\Tmpc L\Me}{%
	\rdg{kiñcanavasthā°}{\Tmac}}%
prasaṃga\mds ś ca| % Tm10r7
kriyā cet karaṇaṃ syāt, tannirvvarttyayā \mde
\vrt{cānya\-yā kriyayā}{L\Me}{%
	\rdg{cā\wmhl{5}ā}{\Tm}} %
bhavitavya\mds m| na hy anyāṃ kriyām antareṇa tasyāḥ
\vrt{karaṇatvaṃ}{\eme}{%
	\rdg{kāraṇatvaṃ}{L\Me}} %
ghaṭate| yāpi cānyā sāpi
\vrt{karaṇaṃ}{\eme}{%
%	\rdg{karaṇatvaṃ}{},
	\rdg{kāraṇaṃ}{L\Me}} %
syād anyasyāḥ kriy\mde āyāḥ| sāpy anyasyā ity anavasthā| % L10r11 after s of sāpy
tataś ca kriyāsantānānupara\mds māt % Tm10r8
phalānabhinirvvṛttiḥ|
\vrt{tatrānyā}{L}{%
	\rdg{tatrāntyā}{\Me}}
cet kriyānta\mde reṇa vinā phalahe\mds tur bhavet, prathamasya tathātve kaḥ pradveṣaḥ|
\edn{This sentence does not look good, The last sentence should contribute to the point that kriyā is not a kāraka; and it appears that it is still in the context of anavasthā the reading antya, in contrast to prathama, might be good. These two paragraphs are really hard to figure out}%

\paragraph{\emph{Siddhānta}: \emph{cittavṛtti} is \emph{saṃvedya} as the revealer (\emph{abhivyañjaka})}

cittavṛttes tv abhivyañjakatvāt samvedyatety adoṣaḥ|\edn{Does this part have anything to do with Nyāyasūtra Nyāyabhāṣya NS NBh 1.1.22?} % L10v1 starts after ado
\edn{"tu" is here. As opposed to what? Probably kriyā}%
%[24,18]
dṛśyate cā\mde bhivyañjakā\-nāṃ
\vrtsm{svarūpa}{L\Me}{%
	\rdg{svarū\wmhl{1}°}{\Tm}}%
saṃvedyatayaiva karaṇatvam pra\mds dīpādīnām|
% Tm10r9 after dīpādīnām
\edn{pradīpa, etc., are abhivyañjaka; they are saṃvedya; by being naturally known, they are karaṇa; so is the cittavṛtti...}%
tathā cittavṛtter api|

\subparagraph{Why is the eyesight not grasped while it is an \emph{abhivyañjaka}?}

cakṣurādiṣv abhi\mde vyañjakatve \mds 'py agrahaṇād vyabhicāra
\edn{The sight, etc., do reveal things but they are not grasped, i.e., asaṃvedya, therefore there is a vyabhicāra: not everything that reveals does so by being saṃvedya.}%
iti \vrtsm{cen na, samāna}{\conj}{%
	\rdg{cetanasamāna°}{L}
	\rdg{cet, samāna (cetanasamāna)}{\Me}}%
jātīyasaṃkarāt||

katham?
\edn{The following should explain what he means by samānajātīyasaṃkarāt; why the sight, etc., are not saṃvedya "due to the confusion with things of the same kind"(?)}%

% \subparagraph{According to the theory that sense faculties are material}
%[24,21] L10v2 after yeṣa
yeṣān tāvad bhautikānīndriyāṇi,\edn{position A (Naiyāyikas)} teṣā\mde ñ cakṣuṣas taijasatvād bāhyenālokena \vrtsm{tulya}{\Tmpc L\Me}{\rdg{`tu'lya°}{\Tm}}jātīyena
% Tm10v1 after jā
vyatibhedād gṛhyamāṇasyaivāvivekaḥ| yathā pradīpānāṃ bahūnāṃ prabhāvyatibhedāt kasya kā prabheti % L10v3 afte prabheti
na vivicyate| tathā śrotrādiṣv api||
\edn{This paragraph appears to be relatively well-preserved}%

%[24,24]
% \subparagraph{According to the theory that sense faculties are not material}
yeṣām apy abhautikāni,\edn{Who?}
\vrt{teṣām}{\Me}{%
	\rdg{keṣām}{L}}
apy a\mde bhivyañjakatvena samānajātīyatā||
\edn{Thus far is an explanation why cakṣur etc., are not saṃvedya even though they are abhivyañjaka}%

vṛtter
\vrtsm{apy a\mds vyañjaka}{\conj}{%
	\rdg{apy abhivyañjaka°}{L},
	\rdg{abhivyañjaka°}{\Me}}%
tvād agrahaṇam
\vrt{iti cet, tan na}{L}{%
	\rdg{iti| tan na}{\Me}}%
|| % surviving part of Tm10v2 is below
\vrtsm{prāptacaita\mde nyopagraha}{L\Me}{%
	\rdg{\lost{7}pagrahaṇa°}{\Tm}}%
tay\mds ā
% L10v4 starts after vi of vibhāgya
\edn{The expression prāptacaitanyopagraharūpa is from YBh 2.20, 4.22. This adjective modifies buddhivṛtti The point the S has to prove is that a) cittavṛtti is abhivyañjaka and that b) it is not a grāhya...? It is not grāhya because it is a kriyā...? Shouldn't he be arguing that cittavṛtti is saṃvedya and it is a abhivyañjaka? The opponet here is saying "cittavṛtti is not a vyañjaka and that's why it is not cognized." Our author should counter on both counts, no?}%
\vrt{vyañjakatvād}{\conj}{%
	\rdg{vibhāgyatvād}{L},
	\rdg{'bhibhāvyatvād}{\Me}}
avyavahitatvāc ca| sarvvagatatve ’py ātmanaḥ sūkṣmatvena svacchataratvena ca
\vrtsm{cittenā}{\Me}{%
	\rdg{citte tā}{L}}%
vyavadhānāc ca grahaṇaṃ citta\mde vṛtteḥ||
\edn{Several things to note here. Puruṣa is equated as ātman, it is considered to be sarvagata and sūkṣma. This sentence should be explaninng how the grahaṇa (saṃvid) of cittavṛtti occurs and it appears as though it is not happening.}%

%%%% Conclusion!
\subsection{Conclusion: \emph{pramāṇaphala} is what the Bhāṣya says it is}
tasmāc cittavṛtti\vrtms{pramāṇenaiva}{\Tmpc L\Me}{\rdg{°pramāṇainaiva}{\Tmac}} tadākāratay\mds ā pauruṣeyo ’pi bodhaḥ % L10v5 after bodhaḥ; Tm10v3 after taya
saṃvedya\-māna eva phalam iti siddham ||
\edn{This must be the conclusion of the discussion on pramāṇaphala!}%

\subsubsection{How is \emph{cittavṛtti} object of awareness (\emph{saṃvedya})? Not as \emph{pratyakṣa}}
%[25,4]
yeṣāṃ punaḥ
\vrt{pratyakṣeṇa saṃvedyā}{\eme}{%
	\rdg{pratyakṣeṇāsaṃvedyā}{\all}}
cittavṛttis
\vrtsm{teṣām phala}{\Me}{%
	\rdg{teṣām aphala}{L}}%
grahaṇapramāṇā\-bhāvaḥ|
\edn{Who are these people? What's the point of bringing this up? This discussion in fact is related to the one that starts at XYZ Now I'm pretty sure that the reading should be pratyakṣeṇa saṃvedyā. The difference between the view being attacked and our author's is whether cittavṛtti is saṃvedya by being abhivyañjaka or by being pratyakṣa! Since it is hard to make sense out of this discussion, should we emend the text to "pratyakṣeṇa saṃvedyā"? That would be the Buddhist and Naiyāyika position. The last compound should mean the abhāva of both phalagrahaṇa and pramāṇa Should the last part be pramāṇavyāpārābhāvaḥ? So the whole compound is phalagrahaṇapramāṇavyāpārābhāvaḥ?}%
% This is Tm10v and L10v
phale ca gṛhīte \vrtsm{pramāṇāstitva}{\eme}{%
	\rdg{jñānāstitva}{L},
	\rdg{jñānāsiddhatva}{\Me}}%
\mde syārtthāpattiḥ|
\edn{Is he talking about paracittavṛtti? Most unlikely. Arthāpatti suddenly appearing here is also suspicious. Maybe the argument is if phala is grasped, the existence of pramāṇa/jñāna is inevitable. But...}%
na ca phalaṃ gṛhya\-te % L10v6
\vrtsm{pramāṇavyā\-pārāna}{\Tmpc L\Me}{%
	\rdg{pramāṇavyāpārād a}{\Tmac}}%
bhyupagamāt|
\edn{So, "ca" here means "but," stating what happens if we accept the opponent's position. The phala is not grasped and therefore the knowledge or pramāṇa does not exist. Because the opponent cannot accept the operation of pramāṇa because doing so would result in anavasthā situation.}%
%%% THE FOLLOWING IS IRRELEVANT This statement is not incompatible with the idea "pratyakṣeṇa saṃvedyā" because with that view, the final function of pramāṇa is not specified. cittavṛtti is still known by a pramāṇa but what happens then? A translation could be something like "Also, the effect is not grasped because [according to the theory] they do not accept the function of pramāṇa (for its effect being end of up again as an object of a pramāṇa, i.e., pratyakṣa)."
%%% THE FOLLOWING IS IRRELEVANT The big problem I have is pramāṇavyāpārānabhyupagamāt. Does this refer to a doctorinal position that does not accept the utility of pramāṇa or a consequence of the doctorinal position being discussed? It is hard to conceive anyone who does not accept the working of pramāṇa. It should be then that if we accept certain premises held by the opponent, it becomes inevitable that they also do not accept the utility of pramāṇa.
tadabhyupagame cānavasthāprasaṃgaḥ|
yat phalāvagrahaṇāya pramāṇam abhyupagamyate tasyāpy anyat phalaṃ tatpramāṇasyāpy anyad iti na kvacid avatiṣṭheta||
\edn{tad here is pramāṇavyāpāra tasyāpīti tatpramāṇasyāpi so, pramāṇa P1 -> phala Ph1-> vyāpāra V1 is the basic. If phala is cittavṛtti grasped by pratyakṣa then pramāṇa P1 -> (cittavṛttibodha Ph1 = Ph2 <- pratyakṣa P2) -> vyāpāra V1 then pramāṇa P1 -> ((cittavṛttibodha Ph1 = Ph2 <- pratyakṣa P2) = Ph3 <- pratyakṣa P3) -> vyāpāra V1 then}%
%%% THE FOLLOWING IS IRRELEVANT % From this, it does appear that the second alternative above was the case. The opponent did not say they do not accept the utility of pramāṇa.
%%% THE FOLLOWING IS IRRELEVANT % The logic goes something like, if we accept the premise of the opponent (cittavṛtti is a) saṃvedya or b) asaṃvedya by pratyakṣa),
%%% THE FOLLOWING IS IRRELEVANT % I have the feeling this is OK by our author
%%% THE FOLLOWING IS IRRELEVANT % Oh, this is not about pramāṇaphala being the pauruṣeyaś cittavṛttibodhaḥ but about sukha, dveṣa, etc., being pratyakṣa. Our aurhot is attacking the view that does not accept those as pratyakṣa. However, I have difficulty understanding the relationship between that part (sukha, dveṣa, etc., being pratyakṣa) and the pramāṇaphala being cittavṛttibodha part.
%%% THE FOLLOWING IS IRRELEVANT % The above is probably wrong. This is not about sukhaduḥkhadveṣa, etc., but is about pramāṇa, viparyaya, smṛti, nidrā, etc.
\edn{Is Sureśvara dealing with the same type of person? ./6\_sastra/3\_phil/vedanta/naiskarmyasiddhi\_incomplete.txt:na ca sukhaduḥkhādisaṃbandho .avagatyātmanaḥ pratyakṣādipramāṇair gṛhyate yena virodhaḥ pratyakṣādipramāṇair udchāṭyate| katham| śṛṇu| EVen Dharmakīrti explicitly(?) says: rāgadveṣasukhaduḥkhādīnāṃ sarvacittacaittānām ātmasaṃvedanaṃ pratyakṣam avikalpatvāt | Cf. this from the Brahmasiddhi: tathā hi---pramāṇasya phalam arthāntaram anarthāntaraṃ vā sarvaparīkṣakaiḥ pratijñāyate prajñāyate ca| na ca tad asaṃvedyam, tadasaṃvedyatve sarvāsaṃvedyatvaprasaṅgāt, tatkṛtatvāt saṃvedyabhāvasya bhāvānām| na ca saṃvedyam, phalāntarānupalabdher anavasthāprasaṅgāc ca| tasmāt saṃvedyam, ātmaprakāśatvāt; asaṃvedyaṃ ca, viṣayavat karmabhāvābhāvāt| Maṇḍana's position is that pramāṇaphala is both saṃvedya and asaṃvedya.}%

%[25,8]
atha na
\vrt{\dag pramāṇavyāp\mde āraḥ ta\ldots syāt\dag}{L}{%
	\rdg{\lost{5}āraḥ \ccs ta\cce syāt}{\Tm},
	\rdg{pramāṇavyāpāras tasya}{\Me}}% L10v7 after atha na pramāṇa
||
\edn{The original reading in the common ancestor was "pramāṇavyāpāraḥ tasyāttadā." L actually reads: L10v7 starts vyāpāraḥ ta syāttadāpramāṇābhāvāt phalasyāsatvaṃ prasaktaṃ It is probably reasonable to assume something is missing between the ta and syāt. Since vyāpāraḥ ends with a visarga, there was an awareness that it ended an utterance. Then what would follow is tasya...āt. Then the phrase may be followed by another sentence There is probably a significant portion missing!}%

tadā pramāṇābhāvāt phalasyāsattvaṃ prasaktam| yatra hi sadupalambha\-kāni pramāṇāni na vyāpāraṃ gacchanti, tan nāstīty
\vrt{abhyupagacchāmas teṣām}{L}{%
	\rdg{abhyupagacchāmaḥ (teṣām)}{\Me}}%
||
\edn{tad here must be phala, teṣāṃ pramāṇānām. The basic flow of events appears to be pramāṇa -> (pramāṇa)phala -> (pramāṇa)vyāpāra. The removal of teṣām from the text might be preferable but then where did it come from?}%

%[25,11]
athāgṛhyamāṇam api tad eva jñā\mde nam, pra\mds māṇaṃ phalasyāpīti cet---%
% L10v8 after pramā
\edn{This does sound strange. "phalam api"? tad should be phala. "Even when not being grasped, the same [phala] is the knowledge; the phala, too, has a pramāṇa" or rather the second part should be "the pramāṇa is the evidence (pramāṇa) also for the [existence of] the phala [i.e., there is no need for the phala/jñāna to be grasped. We know that it exists because vyāpāra is being produced]" This could be an answer to the anavasthā problem in the sense there is no need for the phala to be grasped Looks like pramāṇa -> phala = jñāna -> vyāpāra}%
na, a\mde bhivyañjakadharmavyatikramāt|
\edn{Be reminded that abhivyañjakadharma belongs to cittavṛtti according to our author}%
\vrt{na hy abhivyañjakam agṛhyamāṇam}{\Tmpc}{
	\rdg{na hy ativyañjakam agṛhyamāṇam}{\Tmac}, % Tm10v6 starts at ṇam
	\rdg{na hy a`bhivyañjakam agṛyamāṇam'}{L},
	\rdg{na hy abhivyañjakaṃ gṛhyamāṇam}{\Me}}
\edn{If bhi and ti looked similar, this is another indication that the text was from a northern script origin.}%
\vrt{abhi\-vyañjayad}{\Lpc\Me}{%
	\rdg{abhivya[ñja\mist{1}]d}{\Tm},
	\rdg{abhivya\ccs ñjaka\cce yad}{L}}
dṛṣṭam|
\edn{Revealer is always grasped. See above! Reminds me of the lamp being a bad example of. But isn´t this complete opposite of what our author argued for? cittavṛtti is saṃvedya/gṛhya for being abhivyañjaka? This is true if we follow the reading of Lpc (and M, A, and E)}%
na hy \vrt{ulkā}{\Me}{\rdg{utkā}{\Tm L}} giriśikhare ’pi nijan dharmmaṃ atikrāmati||
\edn{Perhaps the most plausible understanding of this example is that fire does what it does with or without anyone watching.}%

%[25,14]
athābhivyañjayad eva phalaṃ praty api pramāṇam iti ce\mde t
\edn{"But the same [pramāṇa] revealing [the object] is the pramāṇa also toward the phala." "abhivyañjayat" should modify "pramāṇam" that should be the pramāṇa for phala, too -> the pramāṇa for the phala was the problem}%
siddhaṃ \vrtsm{tarhīpsita}{\Tmpc L\Me}{\rdg{tarhi smita°}{\Tmac}}m asmāka\mds m|
\vrt{rāgādīn\mde āṃ}{\Me}{% L10v9 after dī
	\rdg{ragādīnāṃ}{L}
}
vṛttiviśeṣaṇatvena grahaṇopapatteḥ|
\edn{Are rāga, etc., here refer to the phala? The word grahaṇa/graha is being used for the phala in this context. Be reminded that rāga and dveṣa give rise to upādāna and hāna/tyāga Need to read this a bit more carefully but looks OK in essence.}%
rakto ’ha\mde n dviṣṭo ’ham
\vrt{iti ca}{\Tm\Lpc}{%
	\rdg{iti ce}{\Lac},
	\rdg{iti}{\Me}}%
| yady
\vrtsm{ahampratyaya}{\Me}{%
	\rdg{ahaṃpratyayi}{\Tm L}}%
viśeṣaṇaṃ,
\vrt{tadāpi}{\Me}{%
	\rdg{tathāpi}{\Tm L}}
\vrtsm{vṛttibhāvabhāvi}{\eme}{%
	\rdg{vṛttibhāvābhāvi}{\Tm L\Me}}%
\mds tvenai\-va grahaṇād vṛttiviśeṣaṇ\mde ānā\mds ṃ
\vrt{rāg\mde ādīnāṃ}{\Tm\Me}{%
	\rdg{\om}{L}}\edn{This is weird. Where did the edition get the reading?}
grahaṇam upapadyate| \vrt{nānyathā}{\Tm L}{% L10v10 after nyatha
	\rdg{a(tenā)nyathā}{\Me}}%
| rūpādi\vrtmm{vat pramānā}{\Tm L}{\rdg{van mānā}{\Me}}ntarāpekṣatvaprasaṅgāt| tataś cānavastheti||
\edn{Here again, iti is used very effectively.}%

%[25,19]
\subsubsection{Final word on \emph{pramāṇaphala}}
\tstwll{PNDT}{tasmāt\ldots upekṣaṇīyam}
tasmāt pramāṇa\vrtmm{phalābhijñānābhi}{L}{%
	\rdg{°phalajñānābhi°}{\Tmac},
	\rdg{°phalānabhijñānām abhi°}{\Me}}%
mānamātram eva kevalaṃ paṇḍitam\vrtms{manyā\-nām}{\Tm\Me}{%
%	\rdg{manyānām}{\Tm\Me},
	\rdg{manyamānām}{L}}
ity upe\mds kṣaṇī\mde yam iti||\edn{Tm has a section marker here. So, this was a major end of a discussion.}
\donote{PNDT}{Cf.\ BSSBh intro.: kathaṃ punar avidyāvadviṣayāṇi pratyakṣādīni pramāṇāni śāstrāṇi ceti|
%
ucyate---dehendriyādiṣv ahaṃmamābhimānarahitasya pramātṛtvānupapattau pramāṇapravṛttyanupapatteḥ| nahīndriyāṇy anupādāya pratyakṣādivyavahāraḥ saṃbhavati| na cādhiṣṭhānam antareṇendriyāṇāṃ vyavahāraḥ saṃbhavati| na cānadhyastātmabhāvena dehena kaścid vyāpriyate| na caitasmin sarvasminn asati asaṅgasyātmanaḥ pramātṛtvam upapadyate| na ca pramātṛtvam antareṇa pramāṇapravṛttir asti| tasmād avidyāvadviṣayāṇy eva pratyakṣādīni pramāṇāni śāstrāṇi ca|
%
paśvādibhiś cāviśeṣāt| yathā hi paśvādayaḥ śabdādibhiḥ śrotrādīnāṃ saṃbandhe sati, śabdādivijñāne pratikūle jāte tato nivartante, anukūle ca pravartante; yathā daṇḍodyatakaraṃ puruṣam abhimukham upalabhya māṃ hantum ayam icchatīti palāyitum ārabhante, haritatṛṇapūrṇapāṇim upalabhya taṃ praty abhimukhībhavanti; evaṃ puruṣā api vyutpannacittāḥ krūradṛṣṭīn ākrośataḥ khaḍgodyatakarān balavata upalabhya tato nivartante, tadviparītān prati pravartante, ataḥ samānaḥ paśvādibhiḥ puruṣāṇāṃ pramāṇaprameyavyavahāraḥ| paśvādīnāṃ ca prasiddho 'vivekapuraḥsaraḥ pratyakṣādivyavahāraḥ, tatsāmānyadarśanād vyutpattimatām api puruṣāṇāṃ pratyakṣādivyavahāras tatkālaḥ samāna iti niścīyate| (BSSBh NSP ed pp. 40–43)}
\edn{Again, the combination of tasmāt and iti So, after yeṣām was the guys who do not accept pramāṇaphala according to the 1952 edition but not according to the 1999 edition...}%
\edn{tatkāla BAUBh yathā vijñātadigvibhāgasyāpyakasmāddigviparyayavibhramaḥ. samyagjñānavato'pi cet pūrvavadviparītapratyaya utpadyate, samyagjñāne'pyavisrambhācchāstrārthavijñānādau pravṛttirasamañjasā syāt, sarvaṃ ca pramāṇamapramāṇaṃ saṃpadyeta, pramāṇāpramāṇayorviśeṣānupapatteḥ. etena samyagjñānānantarameva śarīrapātābhāvaḥ kasmādityetatparihṛtam. jñānotpatteḥ prāk ūrdhvaṃ tatkālajanmāntarasaṃcitānāṃ ca karmaṇāmapravṛttaphalānāṃ vināśaḥ siddho bhavati phalaprāptivighnaniṣedhaśrutereva; 'kṣīyante cāsya karmāṇi' 'tasya tāvadeva ciram' 'sarve pāpmānaḥ pradūyante' 'taṃ viditvā na lipyate karmaṇā pāpakena' 'etamu haivaite na tarataḥ' 'nainaṃ kṛtākṛte tapataḥ' 'etaṁ ha vāva na tapati' 'na bibheti kutaścana' ityādiśrutibhyaśca; 'jñānāgniḥ sarvakarmāṇi bhasmasātkurute' ityādismṛtibhyaśca.}%
\edn{An interesting tidbit from the Pramāṇasamuccaya 1.8 of Dignāga: tām upādāya pramāṇatvam upacaryate nirvyāpāram api sat. tad yathā phalaṃ hetvanurūpam utpadyamānaṃ heturūpaṃ gṛhṇātīty kathyate nirvyāpāram api, tadvad atrāpi. This is about pramāṇaphala. Steinkellner's edition pp. 3–4.}%

\end{document}
%Tm10v
[25,21]
anumānam athocyate 'numeyasya
% Tm: anumānam iti sunocyate a`nu'meyasya
tulyajātīyeṣv anuvṛtta iti| anumeyasya dharmaviśiṣṭasya| tad dharmava%Tmv9/10
ttayā tulyajātīyeṣv anuvṛttaḥ sadbhāvamātreṇa, tadviruddham dhamrmakebhyo vyāvṛttas teṣv asann ity arthaḥ| sambandha iti| kaḥ sambandhaḥ? kayor vā? na hi sambandhināv anāpekṣya sambandhaḥ sambandhatām iyānn asau prasiddha%Tm11r1/p41 of pdf1
vac cocyate ya iti||

[25,26] nanu cānumānasya prakṛtatvāl liṅgaliṅginoḥ sambandhaḥ
prasiddha eva nāgamakatvātn na kevalaḥ sambandho% Tmac sambandhe/Tmpc sambandho
gamaka ucyate liṅgasyaiva sambadho vyāvṛttyanuvṛttyoḥ| kaḥ punar asau? liṅgam eva| na hi liṅgād anya%Tm10/11r2
d anuvṛttaṃ vyāvṛttañ ca ta%ñca Tmac
c ca gamakan nanu liṅgam api naiva gamakaṃ 26,5 liṅgisambandhānapekṣaṃ yadi cānapekṣita%<<yadi cānupekṣā>> yadi cānepekṣita° Tm
liṅgisambandhaṃ gamakaṃ yat kiñcid gamakaṃ syāt, tasmāt liṅgaliṅgisambandho vaktavyaḥ||

[26,7] noktatvād anu%Tm10/11r3
meyasya tulyajātīyeṣv anuvṛttaḥ ity ādinaivoktatvān na hy anumeyenāsaṃ%\mds
baddhaṃ liṅ%\mde
ginā tat sadṛśajātīyeṣv anuvartate%anupravarttate tm
bhinnajātīyebhyo vā nivartate| tasmād ya iti ca prasiddham eva saṃbandham anuvadati

[26,10] tadviṣa%Tm10/11r4
yā evam anvaya% yā eva samanvaya° Tm
vyatirekaviśiṣṭasambandiviṣayā vṛttiḥ %\mds
sāmānyāvadhāraṇapradhānā liṅ%\mde
ga% gi Tm
sāmānyāvadhāraṇapradhānānu% °pradhānā anu° Tm
mānan

[26,12] yathā% tathā tm
deśāntaraprāpter gatimac candratārakañ caitravad iti %Tm10/11r5
yady api la%\mds
kṣaṇā%\mde
bhidhāna%\mds
mā%\mde
treṇa tad abhivyaktis tathāpy udāharaṇapratyudāharaṇa%\mds
pradarśanād atitarām ity udāharaṇa%\mde
pratipādanan tatra gatimac ca%\mds
ndratāraka%\mde
m iti candratārakasya gatimattvam anumeyan tatsamāna%Tm10/11r6
jātī%\mds
yeṣv anuvṛttā deśāntarapr%\mde
āptir gatimatsu caitrā% hard to read Tm
diṣ%\mds
u, bhinnajātīyebhyaś cāgatibhyo vindhyādibhyo vyāvṛttety avyabhica%\mde
rantī %\mds ...
%\mde
% Tm tvañ jātīyakam avinābhāvena sambandhaḥ samba %Tm10/11r7 \mds
% \mde dikāryyakāraṇādibhava%\mds　...
nityaṃ gatimatā sambandhāt||

[26,19] bhavatu tāvad aviruddhānāṃ nityasambandhād deśāntaraprāptyādīnāṃ gamyagamakatā, vadhyaghātakādīnāṃ tu sam%\mde
bandhābhāve katham ity ucyate gamyagamakabhāvbhāvena

[26,21] tatrā%Tm10/11r8\mds
pi nityasambandhād adoṣaḥ| na hy atrāpi prāpt%\mde
ilakṣaṇaḥ sambandho%\mds
'bhipreyate| sāmānyatodṛṣṭā(nu)bhāvaprasaṅgāt| na hi candratārakādiṣu deśāntaraprāpteḥ pr%\mde
āptilakṣāṇasambandhaḥ pratyakṣeṇa gṛhyate agṛhyamā%Tm10/11r9 \mds
ṇaś ced gamakaḥ sarvaṃ sarvasya gamakaṃ syāt| na hi kiṃcit kenacid asambaddham, ākāśādīnāṃ sarvatra bhāvāt||

[26,25] tasmād ganyagamakabuddhyutpattihetur eva sambandho anumānavisaye %\mde
tathā ca badhyaghātakayor api anyatarābhāvada%Tm10/11v1 \mds
rśane 'nyatarabhāvo gamyata ity anumānam eva| yasyāpy ucite deśe abhāvadarśanād anyatra bhāvo gamyate| yathā ghātakadarśanena badhyābhāvo 'numīyate| caitrasya jīvato %\mde Tm10/11v1
gṛhābhāvadarśane bahirbhavaś ceti

[27,1] nanu cābhāvārthā%Tm10/11v3 \mds
pattī ete, nānumāne, (na) lakṣaṇāsaṅgateṅ gavādivat||

[27,2] na ya%\mde
thoktalakṣāno%\mds
papatteḥ badhyaghātakayor virodhalakṣaṇaḥ sambandhaḥ| yatra yatra ghātakabhāvas tatra badhyābhāva evety avinābhāvaḥ%\mde
na hi badhyasyābhāvaṃ ghātakabhāvo%ghātaka(kā)bhāvo E
 vyabhicarati badhyasya vā%Tm10/11v3 \mds
bhāvo ghātakā(ka)bhāvam, agnibhāvam iva dhūmabhāvaḥ||

[27,6] %\mde
kathaṃ punar avagamyate %\mds
ghātakabhāvadarśanad eva tadgocare badhyasyābhāvabuddhir iti ? na punaḥ sadupalambhakapramāṇānāṃ gh%\mde
ātaka iva badhye vyāpārā%vyāpāraḥ Tm
bhāvād ity

[27,8] ucyate ghātakabh%Tm10/11v4 \mds
āve vyāpṛ(pyāvṛ)tapramāṇasyaiva badhyābhāva%\mde
buddhiupapattir dhūmabhāve vy%\mds
āpṛ(vṛ)tatapramāṇasya badhyavi%\mde
śeṣābhāvabuddhyutpattir ddhūmaviśeṣavyāpṛtapramā%Tm10/11v5
ṇasyai%mds
vāgniviśeṣabuddhir i%\mde
ty

[27,12] abhāvapramāṇavādinas tu sarvatra% (tu sa)rvv<sma> Tm
pramāṇābhāvasyāviśiṣṭatvād virodhiviśeṣadarśanānuvidhānaṃ virodhyantarābhāvadarśanasyānupapannaṃ, liṅgaviśeṣānuvudhānam iva liṅgiviśeṣadarśanasya| na hi sadupalambhakapramāṇavyāpārābhāvasya kaścid viśeṣaḥ yena viśeṣabuddhis tad viśeṣam anuvidhīyate||

[27,16] atha jajñāsite viśeṣe sadupalambhakapramāṇaābhāvād iti cet tan nāsataḥ smṛtīcchādisahabhāvāpekṣanupapatteḥ| na hy asat sahāyāpekṣaya kāryaṃ nirvartayad dṛṣṭaṃ śaśa-_Viṣāṇādi| dṛṣṭam iheti cet na vipratipattisthānatve dṛṣṭāntābhāvāt| tulyam iti cet na ghaṭādeḥ sato 'pekṣādarśanāt||

[27,20] na cāpi nirvyāpārasyāvastutvād abhāvasya pramaṇabhāvo yoktaḥ| na hi nirastavyāpārakaṃ kiṃcit kāryārambhakaṃ dṛśyate loke| tasmād vyāpārarūpapramāṇajanitā prameyābhāvabuddhiḥ prameyaviṣayatvat ghaṭādibuddhivat| tathā ca vastutvāt pravṛttinivṛttihetutvāt| vipakṣaḥ śaśaviṣāṇādiḥ| tathābhāva[pramā]pramāṇā(ṇa)bhāvajanitā na bhavati, yathā 'bhihitahetu-dṛṣṭāntadarśanāt(-nena)||

[27,25] ghātakadarśanaṃ badhyābhāvabuddher liṅgaṃ tad bhāvitvād agnyādibuddhinimittadhūmādibuddhivat saṃśayanivartakatvāc ca| badhyābhāvabuddhir vā liṅgajanitā niścayarūpatvād agnibuddhivad iti||

[28,1] iha gaṭo nāstīti kila nānumānaṃ liṅgaliṅgisambandhāgrahaṇāt| na hi yathāgnyādibuddheer dhūmadiliṅgānveṣaṇaṃ tathā iha nāsti ghaṭa ity atra liṅgānveṣaṇam asti| prāg eva liṅgānveṣaṇāt buddhyutpatteḥ||

[28,4] naiṣa doṣaḥ ihapratyayaākārasyaiva liṅgatvāt| na hi prathamam eva liṅgabuddhyupajanane sati liṅgānveṣaṇaṃ bhavati| ihetyākāroparañjitā buddhir liṅgam umuparañjakavyatiriktaghaṭādyabhāvabudder ghātakabuddhivad eva| virodhalakṣaṇa eva cātrāpi sambandhaḥ||

[28,8] yatra punar aviruddhāny uparañjakāni tatrāstitvabuddheḥ etenaivārthāpatter anumānatvaṃ vyākhgyātam avinābhāvena sambandhena| tatrāpi jīvato hi devadattasya gṛhābhāve bahir-bhāvaḥ, yadā gṛhe bhāvaḥ na tadā bahirbhāva ity anyatareṇānumānam | sākṣāt dṛśyamānabahir-bhāvābhyantarābhāvatvaṃ dṛṣṭāntaḥ| vṛttyuparañjakānuparañjakatvena hy abhāvagrahaṇasyābhihito hetuḥ||

[28,13] agnyādāv api dhanādiśaktigrahaṇam anumānam eva| kāryanirvartakayor avinā-bhāvasambandhasyātmani dṛṣṭatvād bhujyādau| dṛśyate hi pratyagātmani śakto 'ham idaṃ kāryaṃ kartum, aśakto 'ham ity ātmavad ātmaviśeṣaṇasya śakter grahaṇam| śaktinirvartyam eva dahanādikāryam apīty anumānasiddhiḥ||

[28,17] nanu cābhāvasya pratyakṣatām upagāyanti kecit| yathā kilādarśakena pradīpādinā sati gṛhyamāṇe vastuni na tad avayavini yad upalabhyate tan nāstīty evaṃ sataḥ prakarśakaṃ pramāṇam asad api prakāśayati| tathā codāharanty anaṅkitaṃ vāsaānayety praviśya bhāṇḍāgāram aṅkitānaṅkitavāsas sannidhāv aṅkitavāsograhaṇenaiva
pratyakṣeṇāṅkitābhāvaviśiṣṭaṃ vāsa upalabhyanayati| tasmāt pratyakṣa evābhāvo lakṣaṇāntarānupapatter iti||

[28,23] tan na buddhyākāreṇārthāvadhāraṇānṛtatvaprasaṅgāt| yadākārā hi buddhis tasya pratyakṣaāvadhriyamāṇātmatābhyupagamyate| yathā nīlaṃ vastu pītam ity atrāpi viśeṣaṇaviśeṣyasambandhapratyayo mithyā prāpnoti| tadyatheha ghaṭo nāstītīhapradeśasya ghaṭābhāvenāsati sambandhe sambandhapratyayah, na hi ghaṭābhāvasyehapradeśasya cā[prāpta]prāptilakṣaṇaḥ samavāyābhidhānovā sambandhas tair abhyyatheha devadatta iti| tasmād abhāva(ve)pratyayavā(yā)n nīlaṃ vastv iti viśeṣaṇaviśeṣyasambandhapratyayo api sambandhapratyayatvān mithyā prasajyeta||

[29,4] yathā cāsty apīhadeśaghaṭābhāvasambandhe sabandhapratītis tathaivānīle api vastuny asaty eva nīlatvasambandge nīlatvasambandhapratyayaḥ prāpnoti, tathā sambandhiṣv api| athāsty eva sambandhaḥ dviṣṭhatā viruddhyeta sambandhasya| abhāvasyāvastutvābhyupagamāt||

[29,7] athāpy ekaśataṃ ṣaṣṭhyarthā iti yāvat ṣaṣthyarthāḥ sambandha iṣyeta, tadā dvividha eva sambandhaḥ samavāyākhyaḥ samyogākhyaś ceti na yujyate||

[29,9] athai(tathai)kaśatasyāpi ṣaṣṭhyarthānāṃ ubhayāntaḥpātitvam ucyeta nābhāve 'anyatarasyāpy asambhavāt||

[29,11] kiñcā(ntva)nyat| iha nabhasi kusumaṃ nāstīti sambandhino nabhaso apy apratyakṣatvāt gaganakusumābhāvasyāpratyakṣatvaṃ
prāpnoti| athāpratyakṣatvam eveti ced āpatitaṃ gaganakusumābhāvavad apratyakṣatvam iha bhū-pradeśe ghaṭābhāvasyāpi| viśeṣo vaktavyaḥ||

[29,14] atha vyomno 'pratyakṣatvād iti tan na puṣpasya patato nabhasy astitvadarśanād vyomāpratyakṣatvam ahetuḥ| yathaiva nipatatkusumasadbhāve vyomāpratyakṣatvam ahetus tathā sadbhāvaviruddhagrahaṇe apy ahetuḥ syāt| athākāśāpratyakṣatvād eveti cet, asambandham api yat kiṃcid vyavahitāpratyakṣatvādi so api hetuḥ syāt| atha nabhaḥ-kusumābhāvaḥ pratyakṣa iti cet indriyārthasannikarṣānarthakyaṃ sarvatra prasajyate| tathā ca saty asannikarṣṭopalabdheḥ sarveṣāṃ sarvajñatvaṃ prasaktam||

[29,20] athānumeyaṃ tatreti cet nabhaḥ kusumābhāvavad iha ghaṭābhāvasyāviśeṣād anumānikatva(kīya)m eva siddham| evaṃ ca saty abhāvapramāṇavādino api vyavahitaviprakṛṣṭodakādāv analādyabhāvasyānumeyatvād iha ghaṭābhāvo api tatsāmānyād anumānagrāhya eveti niścetuṃ yuktam||

[29,24] vindhyāś cāprāpter agatir iti ca prayogāntaram| agatir vindhya iti pratijñānam| aprāpter iti hetuḥ caitravad iti kākalocanavad ubhayatra sambadhyate| gatimac candratārakam ity anena sādharmyeṇa, vindhyāś cāaprāpter ity anena vaidharmyeṇa sambandhaḥ evaṃ ca saty anvayavyatirekā(ka)darśanakṛto na doṣaḥ||

[30,7] atha vaitasmin prayoge vyatirekaviparyayadarśanena doṣābhāva ity upanyastam| lakṣaṇasya ca lakṣyaviśeṣeṇohyatvād adoṣaḥ||

[30,9] āptaḥ parānujighṛkṣāprayukto doṣarahito dṛṣṭānumitārthaḥ| tena dṛṣṭo vā anumito vārthaḥ paratra śrotṛviśeṣe svabodhasaṅkrāntaye svādhigamasaṃkramaṇāya śabdena śabda iti vākyaṃ tato viśiṣṭārthapratipatteḥ| tasmāt tadupadiṣṭād vākyāt| [tadarthaviṣayā] vākyārthaviṣayā yā cittavṛttiḥ sānumānavad eva sāmānyāvadhāraṇapradhānā||

[30,14] ittham āptyādidharmagrahaṇapūrvakā śrotur āgamaḥ na vaktuḥ| dṛṣṭānumitārtho hi vaktā tasya laiṅgakam aindriyakaṃ vā taj-jñānam| śrotur eva tu śravaṇapūrvakaṃ niścitaṃ vijñānaṃ laiṅgikaindriyakavyatiriktam| sa āgamaḥ [na] pratyakṣam indriyā[viṣaya]viṣayatvāt yathā-abhihitaliṅgaliṅgisambandhanirapekṣatvāc ca nānumānam| na hi svargāpūrvadevatādiśabdā(bde ')arthānvitaṃ jñānam āsāditaliṅgaliṅgisambandhapurassaram upajāyate| phalaṃ tu sarvatra
pauruṣeyo vṛttibodha eva||

[30,20] yasyāgamasya vaktāgamākhyavṛttyutpattinimitte śabdoccāraṇāt āgamasya prayogopanyāsena pratipāditaṃ phalaṃ tu pūrvavad eva| yadāpy eka eva prayogas ta(thā)dāpi sādhyasādhanābhāvamātrasya vivakṣitatvān na viparītavyatirekadarśanād ity ucyate||

[30,23] aśraddheyārthavaktrabhihitatvenāgama eva viśiṣyate| katham asāv aśraddheyārtho vaktety ata āha na dṛṣṭo 'numito vārthaḥ tena vaktā| śiṣṭalekaprasiddhānyathā pratipādanāt buddhārhatādiḥ| na hi rāgādidoṣānupahatāntaḥkaraṇaḥ prasiddham anyathā nirvarṇayati| sa āgamaḥ plavata ābhāsī-bhavati||

[30,27] yady api pratyakṣānumānaviṣayābhāsadarśanaṃ svaśabdena na kṛtaṃ tathā apy āgamābhāsopadarśanenaiva pūrvatrāpi yathābhihitalakṣaṇaviparītatvena pratyakṣānumānābhāsādi draṣṭavyam arthād uktam||

[31,1] mūlavaktari tv ādyavaktarīśvare tann nimitto ya āgamaḥ sa viplavanimittābhāvānnirviplavaḥ syāt| yathā-artha(ya)vaktṛvat||

[31,3] upamānaṃ (mānaṃ) śabdapūrvakaṃ tadāgamāntaḥpātitvān na
pramṇāntaram| yathā gaur iva gavaya iti vanecarabhihatād eva vākyād avadhṛtasādṛśya punar vane gavayadarśane api
pūrvaśrutavākyopasthāpitasādṛśyasmaraṇamātram eva nādhikam| sādṛśyaṃ tu vākyād eva pūrvapratipannam||

[31,7] nāpi saṃjñāsaṃjñisambandhapratipattita eva mānaṃ tasyā api tadvākyād evāvagatatvāt| na hi gaur ity etāvatānavagatagavayanāmakasya vanagatagavayadarśino api gavayo 'yam iti saṃjñāsaṃjñisambandhāvabodho bhavati||

[31,10] tasmād gaur iva gavaya iti gavayaśabdo api tadvākye 'pekṣitavyaḥ| tathā ca sati pūrvaśrutād eva sarvaṃ pratītam| athāpi syād ayaṃ(dasya)gavaya iti na pūrvabuddhiviṣayam iti tad api tatrabhavatām api paricodanaparihāram ava[pa]śyatām api
pīḍitakaivale[kaule]yakavṛttikoṇo'vasthānam iva lakṣyate||

[31,14] tathā hi deśakālasannikarṣaviprakarṣādibhedenanekopamānabhedaḥ prasajyeta| tathāstv iti cet tat-pūrvadhiyā eātatvam antar-nihitasamastaviśeṣam eva hi vā neyau pūrvam upapāditam| na hi gaur ity ukte deśakālādiviśeṣābhāve na go-śabdād arthapratipattiḥ||

[31,17] kiṃ cānyat na hi saṃjñāsaṃjñisambandhapratipattir upamānaṃ yena hy apamīyate tadupamānaṃ sādṛśyajñānaṃ hi tat| na hi nāmināmasambandhapratipattyā kiṃcid upamīyata uttarakālatvāt| avadhṛtasādṛśyasya hi paścād ā(sya)yaṃ gavaya ity evaṃ so api nipatati||

[31,21] evam tarhi syād aśabdapūrvakaṃ pramāṇāntaraṃ yathā gavayadarśanād etat sadṛśo gaur iti grāmīṇasya buddhiḥ tad api na
prameyābhāvāt| sati hi prameye sadupalmbhakatvābhimatasya pramāṇasya vyāpāraḥ na hi go-gavayayor vyāvṛttam anuvṛttaṃ ca vastutaḥ sādṛśyaṃ nāmaikam asti||

[31,25] nanu ca sādṛśyabuddhir asty eva nāsau nirālambanā bhavitum arhati| satyam evaṃ tatraiva vicāraṇā kim asau vanayūthādipratyayavat ? atha ghaṭādibuddhivad iti| kim cātaḥ yadi vanapratyayavad bhavet tadā yathaiva viśiṣṭasanniveśānugatā mahīruhā evātmavyatirekeṇāvidyamānavanapratyayahetavas tathaiva go-gavayāvayavā eva dvayadarśanasmṛtyapekṣayā svātmavyatirekeṇāsataḥ sādṛśyasya
pratītau kāraṇam iti prāptam avastutvam||

[32,1] atha punar ghaṭādipratyayavad vastutvaṃ pramāṇatvaṃ ca tadā syād iti||

[32,2] tatra kīdṛśaṃ vastv iti vicāryate 32,3  kiṃ go-gavayayor anuvṛttam ekaṃ sādṛśyaṃ gotvagavayatvābhyām anyat ? te eva vānyonyaviśiṣṭe yadi vā tābhyāṃ bhinne go-gavayādhāre dve sādṛśye parasparānupekḍe, sāpekṣe vā ? 32,6  kiṃ vyaktir eva sānyatarajātiviśiṣṭā vā kiṃ jātiḥ sācānyataravyaktiviśiṣṭā vyaktidvayaviśiṣṭā vā sambandho vaiteṣāṃ yathā-vikalpayogam ity ekaikānyonyaviśiṣṭaviśeṣaṇatvādinā bhūyāṃso 'nye api vikalpā udbhāvayitavyāḥ| tatra yady ekaṃ sādṛśyam ubhayagataṃ tadā nopameyatvam ekatvāt yathā gaur iva gaur iti||

[32,10] atha tenānyataravyaktir upamīyate naitad api svātmagatatvāt gaur iva gotvena| atha tadvatyā vyaktyā vyaktyantaram upamīyate naitad api na hi gotvvaṃ khaṇḍamuṇḍopameyatvakāraṇam| atha dharmāntarāpekṣayā sādṛśyaṃ tadbuddheḥ kāraṇam iti naṣṭaṃ tarhi sādṛśyaṃ sa evāstu dharmaḥ tadbuddheḥ kiṃ te sādṛśyahatakena kalpitena||

[32,14] kiṃ ca sa cāpy apekṣyamāṇo dharmaḥ kim ekaḥ sādṛśyavad ubhayatrātha bhinna ity evam ādivikalpāḥ syur ity asambhavaḥ pūrvavad eva| atha bhinnam\ṃ sādṛśyam ucyate tac ca na bhedād eva na hi bhavati gaur ivāśva iti| yathaivaikaikatvād aghaṭamānatā tathā bahubhir apekṣyamāṇaiḥ| etena jātivyaktisambandhādayaḥ pūrvavikalpitāḥ
pratyuktāḥ||

[32,19] atha buddhyutpattyanyathānupapattyā sādṛśyakalpanā saty evaṃ vanavivarapūrvadakṣiṇottarādibuddhiyogānyathānubhavāt [nyathānupapattyā] vastutvaṃ kiṃ na kalpyate| syād eṣā tava buddhir vṛkṣādaya eva tatra tadbuddhijanmahetava iti| yady evaṃ mithyāpratyayo gavayena sadṛśī gaur iti vanabuddhivat| atha ye vyāvṛttāḥ ka ihopamārthaḥ gaur ivāśva iti| athobhaye api te 'vayavā vyāvṛttanuvṛttātmānas tadbudhdihetavaḥ tathāpi vanapratyayavad eva
prāptam||

[32,25] tasmāt parikṣyamāṇaṃ nātivimardanakṣamam iti ta evāvayavā ubhayatropalabhyamānā bhūyāṃsaḥ sādṛśyabuddheḥ kāraṇaṃ vanabuddher iva
pādapāḥ [iti] nopamānaṃ nāma prmāṇāntaram iti siddham||

[32,27] nāpi premeyadvitvāt pramāṇadvitvakalpanam| kutaḥ pramāṇavaśād arthapratīteḥ| na hy arthavaśāt pramāṇakalpanā| yadi cārthavaśāt
pramāṇakalpanā syāt adhunā 33,4       tan na śakyaṃ jānānām asarvajñatvāt caityavandanāt svargo bhavatīty atra caityavandanasya svargahetutvabuddheḥ pratyakṣanumānāviṣayatvān mithyātvaṃ syāt| tad arthaṃ vā pramāṇāntarakalpanam||

[33,7] atha sā buddhir anmānaphalam iti cet pramāṇād anyatra
phalaprasaṅgaḥ evaṃ canirmokṣaḥ(ka)raṇatrayādivākyārthaviṣyabuddhir apīti||

[33,9] nāpi cātra pramāṇalakṣaṇanirūpaṇatātparyaṃ
prasiddhapramāṇādivṛttinirodhābhyupāyābhipradarśanaparatvāt| ata eva na sūtrakāraḥ sūtrayāñcakāra| bhāsyakāreṇa tu kiñcit
prasiddhānuvādamātram avalambitam iti tad anusareṇena kim api jalpitam iti||6-7||

% \finisheditionchapter
% Upad. 14.4 buddhi as light!
% Upad 14 in general as dream being sm.rti!
% {\dn chaayaa madhura"siitalam} % refer to Upad. 18.34!!! {\dn vastu
% cchaayaa sm.rter anyan maadhuryaadi ca kaara.nam}

% \rmfamily
\nolinenumbers
\chapterstyle{default}
\openright
\part{Translation}
\renewcommand{\leftmark}{\MakeUppercase{[Translation}}
% \flushbottom
% \chapter{Translation}
\input{trl02}
% \input{trl03}
% \input{trl04}
% \input{trl05}
% \input{trl06}
% \input{trl07}

% \appendix
% \renewcommand{\leftmark}{\MakeUppercase{Appendices}}
% \appendixpage
% \pagestyle{yvibook}
%
% \chapter{Critical Text 1.1}%
%   \renewcommand{\rightmark}{\MakeUppercase{Critical Text 1.1}}
% \starteditionchapter
% \label{CTEXT01}%
% \pagewiselinenumbers
\setlength{\footmarkwidth}{1.8em}
\label{edition01}
%\tstl{hari}
\var{$\ast$}{\cm Before the beginning: \dn [ह]रि श्रीगणपतये नमः [अविघ्न]मस्तु॥ \Tm, \dn पतञ्जलये नमः \LMA}%\donote{hari}{Cf.\ the beginnig of the manuscript of the \Svaditamkarani\ by \Paramesvara~I on the \Nyayakanika: \dn हरिः श्रीगणपतये नमः गणपतिरभिमतमनिशं दिशतु\cm\ \citep[59]{vidhiv:stern} with \Tm. In that sense, the statement made by Sarma 1991 about the ``generic statement'' in the manuscript of Yuktibhāṣā should also be looked at.}

\begin{verse}
\mds%\com{zero}{यस्मिन्न स्तः \ldots\ प्रणमामि}
\llbl{p01_1}\labelh{FVFPEDN}%
\tstwll{FRSTVRS1}{यस्मिन्न स्तः \ldots\ कैटभशत्रुम्प्रणमामि}%
\tstwll{YTASTM1}{यस्मिन्न स्तः \ldots\ निखिलानाम्}%
\varl{2}य\mms स्मिन्न\mde\ स्तः\mme\ कर्म्मविपाकौ %$\dagger$
यत \mms आस्तां%$\dagger$
\donote{2}{%
  \vnote{L\-\Mpc A\-\Me}{%
    \rdg{[य\mist{3}कर्म्मवि\ccs वा\cce]पा[कौ यत\mist{3}]}{\Tm},
    \rdg{\mist{3}स्तः कर्मविपाको यत आस्तां}{\Td},
    \rdg{यस्मिन्नन्तः कर्मविपाकेन यत आस्तां}{\Mac}}} 
\\ \hspace{60pt} %
% \vrt{\XeTeXglyph 175 \XeTeXglyph 349\mme शा य\mms स्मै\mme\ \mds ना\mms लमलंघ्या\mde\ निखिला\mme नाम्।}%
\vrt{\kle\mme शा य\mms स्मै\mme\ \mds ना\mms लमलंघ्या\mde\ निखिला\mme नाम्।}%
  {\LMA\-\Me}{%
    \rdg{[\mist{2}ा य\mist{7}ा]नाम्}{\Tm}, 
    \rdg{\kle शा यस्मै \mist{5} नखिलानाम्}{\Td}}%
\donote{YTASTM1}{See \rYS{1.24}{1.24}: {\dn \kle शकर्मविपाकाशयैरपरामृष्टः \barfu पुरुषविशेष \barfu \barfu ईश्वरः॥}; Cf.\ \SVSAP{069--072}{kk.\ 69--72}: {\dn सर्वेषां तु फलापेतं न स्थानमुपपद्यते। न \barfu चाप्यनुपभोगो ऽसौ कस्यचित्कर्मणः फलम्॥७०॥ अशेषकर्मनाशे वा पुनःसृष्टिर्न युज्यते। कर्मणां \barfu वाप्यभिव्यक्तौ किं निमित्तं \barfu तदा भवेत्॥७१॥ ईश्वरेच्छा यदीष्येत सैव स्याल्लोककारणम्। ईश्वरेच्छावशित्वे हि निष्फला \barfu कर्मकल्पना॥७२॥}, 75: {\dn कस्यचिद्धेतुमात्रत्वं यद्यधिष्ठातृतेष्यते। कर्मभिः सर्वजीवानां तत्सिद्धेः सिद्धसाधनम्॥}.}%
\\
\tstwll{KLYT1}{नावच्छिन्नः \barfu कालदृशा \barfu यः \barfu कलयन्त्या}नावच्छिन्नः %
\vrt{कालदृशा}{LMA\Me}{%
    \rdg{कालदिशा}{T}%
	%\rdg{कालदृशा}{LMA\Me}%
} %
यः कलयन्त्या%
\donote{KLYT1}{See \rYS{1.26}{1.26}: {\dn स पूर्वेषाम\-पि गुरुः कालेनानवच्छेदात्॥}.  Cf.\ \rBhG{10.30b}{10.30b}: {\dn कालः कलयतामहम्} and \rnnBhGBh{10.30}{10.30} on it: {\dn कालः कलयतां कलनं गणनं कुर्वतामहम्}.} %
\\ \hspace{60pt} \vrt{लोकेशन्तं}{TL}{\rdg{लोकेशस्तं}{MA\Me}}
%\resume{one}
कैटभशत्रु\mts म्प्रणमामि॥१॥%
\donote{FRSTVRS1}{Cf.\ the salutatory stanza of the \NKn: %
    {\dn परामृष्टः \kle शैः कथमपि न यो जातु भगवान्न धर्माधर्माभ्यां त्रिभिरपि विपाकैर्न च \barfu तयोः। परं वाचां तत्त्वं यमधिगमयत्योमिति पदं  न\-म\-स्या\-मो विष्णुं तममरगुरूणामपि गुरुम्॥}. %
	See also \Paramesvara's commentaries (the \Jusadhvamkarani\ and the \Svaditamkarani, especially the latter) on the \NKn\ where the YVi is referred to \citep[106--25]{vidhiv:stern}. %
	Cf.\ also \index{Halbfass, Wilhelm}\citet[207]{halbfass:sankara}.%
}%\donote{zero}{This verse is in the Mattamayūra meter.}
%\bigskip
% \newpage
%\com{versetwo}{यस्सर्व्ववित् \ldots गुरोरपि}
% \tstwll{SRVVVS1}{यस्सव्ववित् सर्व्वविभूतिशक्तिः}यस्स\mme र्व्व\mde वि\mts त्\mte\  सर्व्वविभूतिशक्ति%
\\ \tstwll{SRVVVS1}{यस्सव्ववित् सर्व्वविभूतिश\kti ः}\llbl{p01_2}यस्स\mme र्व्व\mde वि\mts त्\mte\  सर्व्वविभूतिश\kti%
\donote{SRVVVS1}{See \pageref{ANUMANAONE},\lineref{ANUMANAONE}--\pageref{thrprfe},\lineref{thrprfe}.}%
\tstwll{VIHIND1}{विहीनदोषोपहितक्रियाफलः}%
र्व्विहीन%
\vrtmm{दोषोपहित}{\Tmpc LMA\Me}{%
    \rdg{॰दोषो\ins प\ine हित॰}{\Tm}, 
	\rdg{॰दोषो ऽपि हि तत्}{\Td}}%
क्रिया%
\vrtms{फलः}{M(em.?)A\Me}{%
    \rdg{॰फलम्}{TL}}%
\donote{VIHIND1}{See %
    \pageref{PASRR25B},\lineref{PASRR25B}--\pageref{PASRR25E},\lineref{PASRR25E};  %
	Cf.\ \pageref{DSSR25B},\lineref{DSSR25B}--\pageref{SRRTV25E},\lineref{SRRTV25E}.}।%
%{conj.}{%
%    \rdg{॰र्व्वीहीनदोषो\ins प\ine हितक्रियाफलम्}{\Tm},
%    \rdg{॰र्विहीनदोषो ऽपि हि तत्क्रियाफलम्}{\Td},
%    \rdg{॰र्विहीनदोषोपहितक्रियाफलम्}{L},
%    \rdg{र्वि\vv{॰: वि॰}{E}हीनदोषोपहितक्रियाफलः}{MA\Me}}% 
 \\
विश्वोद्भवा\vrtmm{न्तस्थिति}{\Tm LMA\Me}{%
  \rdg{॰न्तः स्थिति॰}{\Td}%
%  \rdg{(विश्वोद्भवान्तस्थिति॰)}{L}%
}हेतुरीशो \tstwll{GRGR1}{नमो ऽस्तु तस्मै गुरवे गुरोरपि}\vrt{नमो ऽस्तु}{T\-L\vv{नमो स्तु}{\Tm L}\-A(em.?)\-\Me(em.)}{\rdg{नम स्तु}{M}} तस्मै %\resume{v2}
\vrt{गुरवे}{\Td LMA\Me}{\rdg{गुरवै}{\Tm}} गुरोरपि\donote{GRGR1}{See \rYS{1.26}{1.26}: {\dn  स पूर्वेषाम\-पि गुरुः कालेना\-नव\-च्छेदात्}.} ॥२॥%\donote{v2}{For b-p\=ada, See \rYS{1.24} and YBh 1.24 ({\dn स हितफलस्य भोक्ता}); d-p\=ada, see \rYS{1.26}.}%\donote{versetwo}{\come{versetwo}}
\end{verse}
\newpage

\llbl{p01_3}\mll{ys1}\str{\vrt{अथ \barfu योगानुशासनम्}{TLMA}{om.\ \Me}॥१॥}\donote{ys1}{{\dn अथ योगानुशासनम्॥}}

% \bigskip
%\bhsy{अथेत्ययमधिकारार्थः। योगानुशासनं शास्त्रमधिकृतं वेदितव्यम्}
\labelh{FIRSTSTC}%
\llbl{p01_4}\tstwll{FSTSTC1}{अथेत्यादि पातञ्जलयोगशास्त्रम्; तस्य विवरणमारभ्यते}%\begin{uncertain}%
अथेत्यादि पातञ्जलयोग%
\vrtmm{शास्त्र}{TLMA}{%
    \rdg{॰शास्त्र\eil सूत्रभाष्य\eir ॰}{\Me}}%
\vrtms{म्; तस्य}{conj.}{%
    \rdg{\om}{\all}} %
विवरणमारभ्यते%\end{uncertain}%
।%
\donote{FSTSTC1}{\labelh{OPENG}%
    Cf.\ the openings of the \rnnBhGBh{00intro}{intro.}:
    {\dn तदिदं गीताशास्त्रं समस्तवेदार्थसारसंग्रहभूतं दुर्विज्ञेयार्थं तदर्थाविष्करणायानेकैर्विवृतपदपदार्थवाक्यवाक्यार्थन्यायमप्यत्यन्तविरुद्धानेकार्थत्वेन लौकिकैर्गृह्यमाणमुपलभ्याहं विवेकतो ऽर्थनिर्धारणार्थं संक्षेपतो विवरणं करिष्यामि।} \citep[5--6]{bhgbh}; %
    \rnnChUBh{00intro.}{intro.}: %
    {\dn ओमित्येतदक्षरमित्याद्यष्टा(॰दिरष्टा॰} or {\dn ॰द्याष्टा॰?)ध्यायी छान्दोग्योपनिषत्। तस्याः सङ्क्षेपतो ऽर्थजिज्ञासुभ्यः ऋजु विवरणमल्पग्रन्थमिदमारभ्यते।} \citep[1]{upa2}; %
    \BAUBh: %
    {\dn `उषा वा अश्वस्य' इत्येवमाद्या वाजसनेयिब्राह्मणोपनिषत्। तस्या इयमल्पग्रन्था वृत्तिरारभ्यते।} \citep[1]{bau}; %
    \KeUBh:
    {\dn `केनेषितमि'त्याद्योपनिषत्परब्रह्मविषया वक्तव्येति \barfu नवमस्याध्यायस्यारम्भः।} \citep[17]{upa1};
    \KaUBh:
    {\dn अथ काठकोपनिषद्वल्लीनां सुखेनार्थप्रबोधनार्थमल्पग्रन्था वृत्तिरारभ्यते।} \citep[55]{upa1};
    \MuUBh:
    {\dn `ओं \barfu ब्रह्मा देवानामि'त्याद्याथर्वणोपनिषत्।} \citep[127]{upa1};
    \MaUBh:
    {\dn `ओमित्येतदक्षरमिदं सर्वं तस्योपव्याख्यानम्।' वेदान्तार्थसारसंग्रहभूतमिदं प्रकरणचतुष्टय`मोमित्येतदक्षरमि'त्याद्यारभ्यते।} \citep[212]{upa1}} %
तत्रानाख्यातसम्बन्धप्रयोजनन्न %\tstwll{PPNN}{पुरुषप्रवृत्तिनिवृत्तिभ्याम्}
पुरुष\-प्रवृत्ति\-निवृत्ति\-भ्या\-%\donote{PPNN}{blah}
म्पर्या\-प्नोतीति सूत्रकाराभिप्रेते पुरुष\-प्रवृत्ति%
\vrtmm{निवृत्ति\-नि\-मि\-त्त}{\conj}{%
    \rdg{॰निमित्त॰}{\all}}%
\vrtms{भूते}{TMA\Me}{%
\rdg{॰भूतैः}{L}} %
\tstwll{SBPY}{सम्बन्धप्रयोजने प्रकटीक्रियेते}सम्बन्धप्रयोजने पूर्व्वम्प्रकटीक्रियेते॥%
\donote{SBPY}{Cf.\ \SVPRA{011--25}{kk.\ 11--25}: {\dn    
	`अथातो धर्मजिज्ञासा' सूत्रमाद्यमिदं कृतम्।     धर्माख्यं विषयं वक्तुं मीमांसायाः प्रयोजनम्॥११॥    
	सर्वस्यैव हि शास्त्रस्य कर्मणो वापि कस्यचित्।     यावत् प्रयोजनं नोक्तं तावत् तत् केन गृह्यते॥१२॥     
	मीमांसाख्या तु विद्येयं \barfu बहुविद्यान्तराश्रिता।     न शुश्रूषयितुं शक्या प्रागनुक्त्वा प्रयोजनम्॥१३॥     
	विद्यान्तरेषु नाप्येतद्यद्यभीष्टं प्रयोजनम्।     अनर्थप्रापणं तावत् तेभ्यो नाशङ्क्यते क्वचित्॥१४॥     
	मीमांसायां त्विहाज्ञाते दुर्ज्ञाते \barfu वाविवेकतः।     न्यायमार्गे महान् दोष इति \barfu यत्नोपचर्यता॥१५॥     
	तस्मात्प्रयोजनं पूर्वमुक्तं सूत्रकृता स्वयम्।     यत्स्वेनोक्तं वदेयुस्तद्भाष्यकारादयः कथम्॥१६॥     
	सिद्धार्थं \barfu ज्ञातसम्बन्धं \barfu श्रोतुं श्रोता प्रवर्तते।     शास्त्रादौ तेन \barfu वक्तव्यः \barfu सम्बन्धः \barfu सप्रयोजनः॥१७॥     %
	शास्त्रं प्रयोजनं चैव सम्बन्धस्याश्रयावुभौ।     तदुक्त्यन्तर्गतस्तस्माद्भिन्नो नोक्तः प्रयोजनात्॥१८॥     
	सिद्धिः श्रोतृप्रवृत्तीनां \barfu सम्बन्धकथनाद्यतः।     तस्मात्सर्वेषु शास्त्रेषु सम्बन्धः पूर्वमुच्यते॥१९॥     
	यावत्प्रयोजनेनास्य सम्बन्धो नाभिधीयते।     \barfu असम्बद्धप्रलापित्वाद्भवेत्तावदसंगतिः॥२०॥     
	इह त्वाक्षिप्य सम्बन्धं भाष्य एवाभिधास्यते।     धर्मप्रसिद्ध्यसिद्धिभ्यां तस्मान्नान्यो ऽभिधीयते॥२१॥     
	न चाप्यत्राथशब्देन शास्त्रसम्बन्ध उच्यते।     सम्बन्धक्रिययोर्ह्येष ब्रूते शास्त्राच्च ते पृथक्॥२२॥     
	यो ऽप्ययं शास्त्रसम्बन्धो वर्ण्यते \barfu कैश्चिदादितः।     क्रियान्तर्यरूपो वा गुरुपर्वक्रमो ऽपि \barfu वा॥२३॥     
	तदतद्भावयोस्तस्य विशेषो नोपलभ्यते।     श्रोतुर्विधौ निषेधे वा ज्ञाने वा शास्त्रगोचरे॥२४॥     
	\barfu तस्माद्व्याख्याङ्गमिच्छद्भिः \barfu सहेतुः \barfu सप्रयोजनः।     %
	शास्त्रावतारसम्बन्धो वाच्यो नान्यस्तु \barfu निष्फलः॥२५॥}; %
    \rnnNBh{1.1.1}/\rnnNV{1.1.1}\ 1.1.1, and the opening discussion of the \VMBh.  Cf.\ also \dnfiller the introduction to the \BhGBh: {\dn विशिष्ट\-प्रयोजन\-संबन्धा\-भिधेय\-वद्गीता\-शास्त्रं यत\-स्तद\-र्थ\-विज्ञानेन समस्त\-पुरुषा\-र्थ\-सिद्धि\-रत\-स्तद्वि\-वरणे यत्नः क्रियते \barfu मया।} \citep[7]{bhgbh}; \rMaUBh{1.01}{1.1}: {\dn वेदान्ता\-र्थ\-सार\-संग्रह\-भूत\-मिदं प्रकरण\-चतुष्टयमो\-मि\-त्ये\-तद\-क्षर\-मि\-त्या\-द्या\-रभ्यते। अत एव न पृथक्सम्बन्धाभिधेयप्रयोजनानि वक्तव्यानि। यान्येव तु वेदान्ते सम्बन्धा\-भिधय\-प्र\-यो\-जना\-नि तान्ये\-वे\-ह भवितुम\-र्ह\-न्ति। तथा\-पि प्र\-करण\-व्याचिख्या\-सुना संक्षेपतो वक्तव्यानि। तत्र प्रयोजन\-वत्साधना\-भिव्यञ्जकत्वेनाभिधेयसम्बद्धं शास्त्रं पारम्पर्येण विशिष्टसम्बन्धाभिधेयप्रयोजनवद्भवति। किं पुन\-स्त\-त्प्रयोजन\-मित्युच्यते। रोगार्तस्येव \barfu रोग\-नि\-वृत्तौ \barfu स्व\-स्थ\-ता। तथा दुःखात्मकस्यात्मनो द्वैतप्रपञ्चोपशमे \barfu स्व\-स्थ\-ता अद्वैत\-भावः प्रयोजनम्।} \citep[212--3]{upa1}; \rKaUBh{1.1.01}{1.1.1}: {\dn प्रयोजनं चास्या उपनिषद \barfu आत्यन्तिकी \barfu संसारनिवृत्तिर्ब्रह्मप्राप्तिलक्षणा। \barfu सम्बन्धश्चैवंभूतप्रयोजनेनोक्तः।} \citep[58]{upa1}; and \rMuUBh{1.1.01}{1.1.1}. Cf.\ also the beginning of the \VMBh\ where the purpose (\prayojana) of the \emph{śabdānuśāsana} and the relationship between the speech (\sabda) and the meaning (\artha) are discussed in succession.}%

% \newpage
\llbl{p01_5}तत्र प्रयोजनन्तावत्---

\tolerance=1000
\emergencystretch=1.5pt

\llbl{p01_6}चिकित्सा\varmm{शास्त्र इव तच्चतुर्व्यूहत्व}{em.,
  {\dn ॰शास्त्रेवतच्चतुर्व्यूहत्व॰} \Tm L,
  \rdg{॰शास्त्रवत्तच्चतुर्व्यूहत्व॰}{\Td},
%  {\dn ॰शास्त्रवतच्चतुर्व्यूहश्च} L,
  {\dn ॰शास्त्रे तच्चतुर्व्यूहश्च} MA,
  {\dn ॰शास्त्रे तच्चतुर्व्यूहत्व॰} \Me%
}प्रदर्शन\varms{द्वारेण}{T\Me,
 {\dn ॰चारेण} LMA%
} व्याख्यातम्। \barfu तद्यथा %
\tst{चिकित्साशास्त्रञ्च\-तु\-र्व्यू\-ह\-म्---रोगो रोगहेतुरारोग्यम्भैषज्यमिति}{\rYBh{2.15}{2.15}.}% \com{vdhpnym}{विधिप्रतिषेधनियमद्वारेण}
---विधिप्रतिषेध%
\varms{नियमद्वारेण}{em.,
% Probably I need to explain this. See YVi p. 184.
  {\dn ॰नियमचारेण} TLMA,
  {\dn ॰नियमद्वारेण \eil च तत्\eir } \Me(conj.)} %\donote{vdhpnym}{\come{vdhpnym}}
चतुर्व्यूह%
\varmm{विषय}{\Td LMA\Me, {\dn ॰विषया॰} \Tm}\-%
व्या\-ख्या\-न\-परम्। \var{एव\-मिहापि}{TMA\Me, {\dn एवमिहावि॰} L}  \tstl{ys2.15}\quoteo परिणाम\-\pause{ys2.15}<\selm>\sf{\dn ताप\-}\varmm{\dn संस्कार\-दुःखै}{TL\Mpc A\Me, {\dn ॰दुःखसंस्कारै॰} \Mac}%
र्गुण\-वृत्ति\-विरोधाच्च दुःखमेव सर्व्वं \resume{ys2.15}विवेकिन\quotec%
\donote{ys2.15}{%
	\rYS{2.15}{2.15}.  See also \citet{wezler:yvi2}\index{Wezler, Albrecht}.%
} %
इत्यारभ्य चतुर्व्यूहत्वं शास्त्रस्य प्रदर्शितम्। \tstl{ybh215a}\pause{ybh215a}<{\dn दुःखप्रचुरः \barfu संसारो}\sel{\dn तद्दृशेः कैवल्यम्}>तद्यथा---%
\tst{दुःख\-%
\varms{प्रचुरः \barfu संसारो}{%
	M(em.)A\Me, 
	{\dn ॰प्र\-चुर\-सं\-सा\-रो} TL%
} %
\barfu हेयः}{%
	Cf.\ \rYS{2.16}{2.16}: {\dn हेयं दुःखमनागतम्॥}.%
}; %
\tstl{ys2.17}%
त\-स्या\-%
\pause{ys2.17}<\selm>विद्या\-\mms नि\mme मित्तो \barfu द्रष्टृ%
\varmm{दृश्य}{%
	TMA\Me, 
	{\dn ॰दृश्यं} L%
}%
\resume{ys2.17}%
संयोगो \barfu  हेतुः%
\donote{ys2.17}{%
	Cf.\ \rYS{2.17}{2.17}: {\dn द्रष्टृदृश्ययोः \barfu संयोगो \barfu हेयहेतुः॥}.%
}%
; %
\tst{विवेकख्यातिरविप्लवा \barfu हानोपायः}{Cf.\ \rYS{2.26}{2.26}: {\dn विवेकख्यातिरविप्लवा \barfu हानोपायः॥}.}%
; \barfu विवेक\-ख्यातौ च सत्यामविद्यानिवृत्तिः, %
\tst{तन्निवृत्तावात्यन्तिको द्रष्टृदृश्यसंयोगोपरमो हानम्; तदेव कैवल्यमिति।} %
    {Cf.\ \rYS{2.25}{2.25}: {\dn तदभावात्संयोगाभावो हानम्, तद्दृशेः कैवल्यम्॥}.}%
%
\donote{ybh215a}{%
	Cf.\ \rYBh{2.15}{2.15}: {\dn तत्र दुःखबहुलः संसारो \barfu हेयः। प्रधान\-पुरुष\-योः \barfu संयोगो \barfu हेयहेतुः। संयोगस्यात्यन्तिकी निवृत्तिर्हानम्। हानोपायः सम्यग्दर्शनम्।}. %
	Cf.\ also \rNBh{1.1.1}: {\dn हेयम्, तस्य निर्वर्तकम्, हानमात्यन्तिकम्, तस्योपायो ऽधिगन्तव्य इत्येतानि चत्वार्य्यर्थपदानि \barfu सम्य\-ग्बुद्ध्वा निःश्रेयस\-म\-धि\-गच्छति।} \citep[33]{ns}; %
	\rNV{1.1.1}: {\dn हेय\-हानो\-पाया\-धि\-गन्तव्य\-भेदा\-च्चत्वा\-र्य्यर्थ\-पदा\-नीति। हेयं दुःखं तद्धेतुश्च। दुःखमुक्तं। हेतुरविद्यातृष्णे धर्माधर्माविति। हानं तत्त्वज्ञानम्। तत् पुनर्यथार्थावस्थितपदार्थाधिगतिस्तच्च प्रमाणम्। उपायः शास्त्रं तदप्युक्तम्। \barfu अधिगन्तव्यो \barfu ऽपवर्गः \barfu स \barfu पुनरात्यन्तिको \barfu दुःखाभावः।} \citep[11--2]{ns}.%
} %
आरोग्यस्थानीयकैवल्य%
\varmm{प्रयुक्तत्वा}{\Td LM\Me, {\dn ॰प्रयुत\-त्वा॰} \Tm, {\dn ॰प्रयुक्त्वा॰} A}%
दस्य तदेव कैवल्यम्प्रयोजनम्॥

\tolerance=200
\emergencystretch=0pt

\llbl{p01_7}ननु च हेयतद्धेतू न %
\var{प्रस्तोतव्यौ}{%
	T\Lpc, %
	{\dn प्रस्तोतव्येन} \Lac MA, %
	{\dn प्रस्तोतव्ये} \Me}%
, निष्प्रयोजनत्वात्। हानप्रयुक्तं हि शास्त्रम्। \barfu तदुपायभूता विवेकख्यातिरेव \barfu वक्तव्या। %
\labelh{NVKANTHA}%
\tstl{nv111}\pause{nv111}<न हि कण्टकविद्ध॰\sel चोद्येते>%
न हि कण्टकविद्धचरणतलस्य %
\vrtsm{तदपनयन}{T\-L\-\Mpc\-\Me}{%
	\rdg{तदु\-प\-नयन॰}{\Mac\-A}%
}%
म्मुक्त्वा दुःख%
\varms{तद्धेतू}{%
	TL\Mpc A\Me, %
	{\dn ॰तद्धेतु} \Mac} %
चोद्येते॥%
\donote{nv111}{%
	Cf.~\rBAUBh{4.3.07}{4.3.7}: {\dn कण्टकविद्धस्य हि कण्टकवेधजनितदुःखनिवृत्तिः फलं, न तु कण्टकविद्धमरणे तद्दुःखनिवृत्तिफलस्याश्रय उपपद्यते।}; %
	\rNV{1.1.1}: {\dn श्रेयः पुनः सुखमहितनिवृत्तिश्च, \barfu तच्छ्रेयो भिद्यमानं द्वेधा व्यवतिष्ठते दृष्टा\-दृष्ट\-भेदेन। दृष्टं \barfu सुखमदृष्टमहितनिवृत्तिः। \barfu अहितनिवृत्तिरप्यात्यन्तिकी \barfu अनात्य\-न्ति\-की च। \barfu अना\-त्यन्ति\-की, कण्टकादेर्दुःखसाधनस्य परिहारेण। \barfu आत्यन्तिकी \barfu पुनरेकविंशतिप्रभेदभिन्नदुःखहान्या।}~\citep[5--6]{ns}.%
}

\llbl{p01_8}\mms नैतदेवम्\mme\kern 2pt । हेयतत्कारणापेक्षत्वाद्धानोपायस्य। यावदिदमनेनोपायेन हातव्यं संसारचक्रम्, अस्य \barfu चाविद्या\-%
\varms{निमित्तो}{%
	LMA\Me, 
	{\dn ॰नि\-मि\-तो} \Tm, 
	{\dn ॰निमित्तौ} \Td%
} %
द्रष्टृ\-दृश्य\-संयोगो हेतुरिति नाख्यायते, तावन्नाविद्याप्रतिपक्ष\-भूता विवेक\-ख्याति\-रा\-ख्यातुं शक्यते। हेयतद्धेतु\vrtmm{मतो ऽपि \barfu हानो}{em.}{\rdg{॰मतो हि हानो॰}{\Tmpc\Td\Lpc}, {\dn ॰मतोर्व्विहानो॰} \Tmac, {\dn ॰मतो हि \ccs ना\cce हानो॰} L, \rdg{॰मतो हानो॰}{MA\Me}}पाया\-र्त्थिनो \var{रोगिणो}{TMA\Me, {\dn रोगिणौ} L} भै\-ष\-ज्या\-र्त्थित्वदर्शनात्। न हि रोगं रोग\varmm{हेतुञ्चा}{\Tdpc\Mpc A\Me\vv{\dn ॰हेतुं चा॰}{\Mpc A\Me}, {\dn ॰हेतुस्त्वा॰} \Tm, {\dn ॰हेतुश्चा॰} \Tdac, {\dn ॰हेतुत्वा॰} L, {\dn ॰हेतुं त्वा॰} \Mac}न\varms{पेक्ष्य}{\Td LM\Me, {\dn ॰पेक्ष्या} \Tm, {\dn ॰पयदय॰} A} चिकित्साशास्त्रमुपदिश्यते॥

\llbl{p01_9}\llbl{smb}सम्बन्धो ऽपि---\llbl{smb_end}

\llbl{p01_10}विवेकख्यातेर्हानमेव फलं साध्यम्, \varl{10}हानस्यापि\pause{10}<\sel> विवेकख्यातिरेव \resume{10}साधनमिति\donote{10}{TMA\Me, {\dn हा\-न\-स्या\lost{14}मिति} L} साध्य\-{}साधन\-{}यो\-{}रित\-{}रेतरनियम एव, %
\varsm{नान्यः \barfu सम्ब}{T, {\dn नान्यसम्ब॰}\vv{॰संब॰}{MA\Me}LMA\Me\@}%
न्धः। तद्यथा---\varsm{भै\-ष\-ज्य}{\Td LMA\Me, {\dn भैषज्या॰} \Tm}स्यारोग्यमेव फलम्, आरोग्यस्यापि भैषज्यमेव \barfu साधनमिती\varms{तरेतरनियमः। स च शास्त्रादिति॥}{TMA\Me, {\dn ॰तरेतर\lost{9}स्त्रादिति} L}

\llbl{p01_11}\varsm{तस्मात्सत्सम्ब}{%
	TLMA, {\dn तस्मात्ससंब॰} \Me%
}%
न्धप्रयोजनं योगानुशासनम्॥
  
% \medskip

\llbl{p01_12}ननु च यदि हानम्प्रयोजनम्, तदुपायश्च विवेकख्यातिरिति, तत्र वक्तव्यम`थ विवेकख्यात्यनुशासनमि'ति। किमर्थम`थ योगानुशासनमि'ति सूत्रितम्?

\llbl{p01_13}तदुपायत्वाद्योगस्य। %\com{upayaeva}{उपाय एव वक्तव्य स्यात्}
उपाय एव वक्तव्य स्यात्%
%\donote{upayaeva}{\come{upayaeva}}
---\vrt{उपेय\-प्र\-दर्शने}{em.}{\rdg{उपेयो\-प\-प्र\-दर्शने}{TL\Mpc A\Me}, \rdg{उपेयोप\ins प्र\ine दर्शने}{M}} \varsm{ह्युपेय}{T, {\dn हि उपेय॰} LMA\Me}मुपायम्प्रति साकांक्षमेवेत्युपायः सांगकलापः पुनरपि 
\var{वक्तव्य स्यात्}{%
	TMA\Me\vv{॰व्यः स्या॰}{M\-A\-\Me}, 
	{\dn वक्तव्यः \lost{1}}{\dn ात्} L%
}\fubar । \varsm{तस्मिंस्त्वभि}{TL, {\dn तस्मिंश्चाभि॰} MA\Me}हिते सर्व्वमभिहितमिति॥

\llbl{p01_14}कथन्तदुपायत्वम्?

\llbl{p01_15}यत आह सूत्रकारः---%
\tst{`योगांगानुष्ठानादशुद्धिक्षये ज्ञानदीप्तिरा विवेकख्यातेः'}{%
	\rYS{2.28}{2.28}.%
}%
\varl{iti1}\pause{iti1}<{\dn ॰ख्यातेः'; `निर्व्विचार॰}>\resume{iti1}%
\donote{iti1}{%
	TLMA, 
	{\dn ॰ख्यातेः' इति। `निर्विचार॰} \Me}%
; %
\tst{`निर्व्विचारवैशारद्ये ऽध्यात्मप्रसादः'}{%
	\rYS{1.47}{1.47}.%
}; %
\tst{`ऋतम्भरा तत्र प्रज्ञे'ति}{%
	\rYS{1.48}{1.48}.%
} %
च। तथा भाष्यकारो ऽप्याह---%
\tst{`यस्त्वेकाग्रे चेतसि स भूतमर्त्थम्प्रद्योतयती'ति}{%
	\rYBh{1.01}{1.1}.  Cf.\ the reading {\dn सद्भूतमर्थम्} instead of {\dn स भूतमर्थम्} \nMaas\citep[5]{maas:2006}. See \pageline{BhT11s}--\pageline{BhT11e} in this edition.%
}%
\fubar ।  भूतार्त्थावगतिश्च \barfu विवेकख्यातिः।
%
त\vrtmm{स्माद्युक्त}{\Td\-M\-A\-\Me}{%
  \rdg{॰स्मात्द्युक्त॰}{\Tm}, 
  \rdg{॰\hlg{स्म}ाद्युक्तर॰}{L}}%
न्त\-दु\-पा\-य\-योगा\-नु\-शा\-स\-न\-मे\-वा\mms दौ\mme\ सूत्रितमिति॥

\llbl{p01_16}ननु च \barfu योगांगानुष्ठानाद्विवेक\varms{ख्यातिः। तत्र}{TM(em.)A\Me\vv{॰तिस्त॰}{TMA}, \vv{॰तिः, त॰}{\Me}, {\dn ख्यातिस्यत्र} L} वक्तव्यम`थ \varsm{योगांगानु}{\Tm LMA\Me, {\dn योगानु॰} \Td}शासनमि'ति॥

\llbl{p01_17}न। फलेनोपक्रमात्। योगांगानुष्ठानस्य हि फलं योग इति \barfu तदुपक्रमो \barfu युक्तः॥

\llbl{p01_18}\varsm{यद्येवं हानेनो}{\Tm LMA\Me, {\dn यद्येवमाह अनेनो॰} \Td}{\dn पक्रम्येत॥
  
\llbl{p01_19}न। तस्योपेयत्वात्। उपेयमेव हि तत्, योगः पुन स्वांगानामुपेयश्चोपायश्च विवेकख्यातेर\-%
\tstwll{animaadi}{अणिमादि॰}%
णिमादि\donote{animaadi}{%
	See \rYS{3.45}{3.45}: {\dn ततो \barfu ऽणिमादिप्रादुर्भावः \barfu कायसंपत्तद्धर्मानभिघातः॥} and the 
	\rnnYBh{3.45}{3.45} on it: {\dn तत्राणिमा \barfu भवत्यणुः। लघिमा लघुर्भवति।} \ldots\  %{\dn महिमा महान्भवति। प्राप्तिरङ्गुल्यग्रेणापि स्पृशति चन्द्रमसम्। प्राकाम्यमिच्छानभिघातः। भूमावुन्मज्जति निमज्जति यथोदके। वशित्वं भूतभौतिकेषु वशी भवत्यवश्यश्चान्येषाम्। ईशितृत्वं तेषां प्रभवाप्ययव्यूहानामीष्टे। यत्र कामावसायित्वं सत्यसंकल्पता यथा संकल्पस्तथा भूतप्रकृतीनामवस्थानम्।} %
	{\dn एतान्यष्टावैश्वर्याणि} \citep[164--5]{ys:anss}.%
}%
प्रा\mms प्ते\mme श्चे%
\varms{त्य\-भ्यर्हिततरः}{%
	\Tm\Tdac \Me(conj.?), %
	{\dn ॰त्यभ्यक्ततरः} \Tdpc, %
	{\dn ॰त्यह्यहिततरः} L, %
	{\dn ॰त्य \ccs इ\cce \ins हि\ine ततरः} M, %
	{\dn ॰त्यहिततरः} A, %
	{\dn ॰\eil त्यभ्य\eir र्हिततरः} \Me%
}%
, फल%
\vrtmm{द्वयनिमित्त}{conj.}{%
	\rdg{॰वन्निमित्त॰}{T}, %
	\rdg{॰(प/व)यनिमित्त॰}{L}, %
	\rdg{॰पयनिमित्त॰}{M ({\dn प} copiest marks as uncertain)}, %
	\rdg{॰पयनिमित्त॰}{A}, %
	\rdg{॰तदुपायरूप॰}{\Me}%
}%
त्वात्॥
  
\llbl{p01_20}उत्तर\varms{सूत्रार्त्थत्वाच्च}{LMA\Me\vv{॰र्थ॰}{MA\Me}, {\dn ॰सूत्रत्वाच्च} T}\fubar । न ह्य`थ योगांगानुशासनमि'त्युक्ते }\tst{`योगश्चित्तवृत्तिनिरोध' इति}{\rYS{1.02}{1.2}.} युक्तं \var{वक्तुम्। \barfu तदा}{M(em.)\vv{वक्तुं तदा}{M}A\Me, {\dn वक्तुन्ता\vv{॰क्तुम्। ता॰}{\Td}दा} TL} हि \tstsm{यमनियमादि सूत्र}{\rYS{2.29}{2.29}: {\dn यमनियमासनप्राणायामप्रत्याहारधारणाध्यानसमाधयो ऽष्टावङ्गानि।}.}मेव पठितव्यं स्यात्॥
  
\llbl{p01_21}अथ तदेव पठितव्यमिति \var{चेत््---न}{\Me(em.), {\dn चेत् न\vv{॰न्न}{MA} हि} TLMA}\fubar । \var{योगोपायाया}{TLMA, {\dn योगोपे\edl पा\edr याया} \Me} \barfu हानोपायभूताया \var{विवेकख्यातेः}{TMA\Me, {\dn विवेक\lost{2}तेः} L} फलस्य च हानस्य \varsm{सम्बन्ध}{T\-M\-A\-\Me, {\dn यम्बन्ध॰} L}\varmm{दर्शना\-}{TLMA, {\dn ॰\eil प्र\eir दर्शना॰} \Me}र्त्थ\varmm{त्वा`द्योग}{TL\Me(em.), {\dn त्वायोग॰} MA}श्चित्तवृत्तिनिरोध' इत्यस्य॥

\llbl{p01_22}\mms\varsm{कथम्पु}{\Td \Me(em.)\vv{\edl स\edr\  कथं पु॰}{\Me}, {\dn स कथम्पु॰} LMA}नरे\mme\varmm{तत्सम्बन्ध}{TMA, {\dn ॰तत्सम\-ब\-न्ध॰} L, {\dn ॰त\-त्सं\-ब\-न्ध\-प्र॰} \Me}दर्शना\vrtms{र्त्थम्}{TLA\Me\vv{र्त्थं}{TL}}{\rdg{॰र्थ॰}{M}}, \vrt{यावता}{LMA\Me}{\rdg{यावतां}{T}} निर्बीजसमाधिलक्षणाभिधानमेतत्॥
 
\llbl{p01_23}सत्यमेवम्। तथापि हानतदुपायसम्बन्ध\varms{म\-व\-द्यो\-त\-य\-दे\-व}{\Td(em.)\-A(em.)\-\Me(em.), {\dn ॰मपद्योतयदेव} \Tm, {\dn ॰मवद्योतयातिव} L, {\dn ॰मवद्योतयादेव} M} \varsm{सूत्र}{T\-L\-\Me(em.), {\dn सून॰} MA}न्निर्बीजसमा\mms धि\mds \vrtsm{लक्षणता}{LM\Apc \Me}{\rdg{लक्षण\ccs भिधानमेतत्सत्यमेवं तथापि\cce ता॰}{A}}मा\-द्य\-ते। न हि निरोध\-लक्षणात्स\-माधेरर्त्थान्तरं हानम्। किन्त्वे\mde ता\mme\varmm{वान्वि}{T, {\dn ॰वात्वि॰} L, {\dn ॰वता वि॰} MA\Me}शेषः---निरोधलक्षण\vrtms{समाधौ}{\Me(em.)}{\rdg{॰समाधौ न}{T}, \rdg{॰समाधेर्न्न\vv{॰र्न॰}{MA}}{L\-M\-A}, \rdg{॰समाधौ\edl धेर्न\edr }{\Me}} पुनःप्रवृत्तिः, हाने त्वात्यन्तिकी निवृ\mms त्ति\mds रिति। समाध्य\mte वस्थायान्तु हानाविशेष एव॥

%\begin{uncertain}% \com{ysbh13}{तथा चाह \ldots\ कैवल्यमि'ति}
\llbl{p01_24}तथा चाह \tstl{ys13}`\vrt{तदा}{\Tmpc \Td LMA\Me}{\rdg{त\ccs था\cce दा}{\Tm}} \var{द्रष्टुः}{\Tmac \Td LMA\Me, {\dn द्रष्टृ॰} \Tmpc} \vrt{स्वरूपे ऽवस्थानम्}{TLMA}{\rdg{स्वरूपेऽवस्थानम्' \eil इति\eir}{\Me}}'%
\donote{ys13}{\rYS{1.03}{1.3}}; %
\tstl{bh334part}`स्वरूपप्रतिष्ठा %
\vrt{वा}{em.}{\rdg{च}{$\Sigma$}} %
चितिशक्तिः' कैवल्यमिति%
\donote{bh334part}{Part of \rYS{4.34}{4.34} (the last sūtra): {\dn पुरुषार्थशून्यानां गुणानां प्रतिप्रसवः कैवल्यं स्वरूपप्रतिष्ठा वा चितिशक्तिरिति}.%  See the note on page~\pageref{ysbh13note}.
}\fubar । स्वरूप\vrtms{प्रतिष्ठत्वं हि}{\Me(em.)}{\rdg{॰प्रतिष्ठं हि}{TLMA}} \vrt{कैवल्यं हानम्}{conj.}{\rdg{कैवल्यमाह}{$\Sigma$}}\fubar ।%\end{uncertain}%\donote{ysbh13}{\come{ysbh13}}
\ ततश्च निर्बीजसमाधिना कैवल्यमेव शास्त्रार्त्थप्रत्य\mts यद्रढि\mde\-म्ने साक्षा\mme त्क्रियते॥
%\special{x:textcolor=000000}

\llbl{p01_25}तस्माद्यदुच्यते \tstwll{KAISC1}{कैश्चिन्निर्बीजसमाधिः कैवल्यसाधनमिति}\varsm{कैश्चिन्निर्बीज}{TL, {\dn कैश्चिद्बीज॰} MA, {\dn कैश्चित् \eil स\eir  बीज॰} \Me}समाधिः कैवल्यसाधनमिति\donote{KAISC1}{Cf.\ \Vacaspati\ on \rYS{1.03}{1.3}: {\dn निःश्रेयसस्य हेतुः समाधिरिति हि श्रुतिस्मृतीतिहासपुराणेषु प्रसिद्धम्॥} \citep[2,4--5]{ys:anss}.}, तन्न। किन्तर्हि ख्यातिरेव साधनमविद्यानिवृत्तिद्वारेण। अविद्यानिमित्तो हि \barfu बन्धः। तस्माद्यद्यपि योगानुशासनं सूत्रितम्, तथापि ख्यात्यर्त्थत्वाद्योगस्य तद्वारेण ख्यातिहानसम्बन्धमपि दर्शयतीति सुष्ठूच्यते। उत्तरसूत्रसम्बन्धार्त्थञ्च योगानुशासनमिति सूत्रयितव्यमिति॥

% \bigskip

\llbl{p01_26}\begin{uncertain}यद्यपि \mms\varsm{फलार्त्थि}{\Td\Lpc MA\Me, {\dn फलाति॰} \Lac}नान्तत्फल\mme साधनम्प्र\varmm{त्युत्प}{\Td LMA\Me, {\dn ॰त्युल्प॰} \Tm}न्नाकांक्षाणां साधनोपदेशो ऽप्यस्ति, \var{तथापि}{TMA\Me, {\dn \lost{1}थापि} L} न तत्प्रयोजनम्। सर्व्वस्य फलोद्देशेन \vrtsm{प्रवृत्ते}{\Tmpc \Td LMA\Me}{\rdg{प्रवृत्ति॰}{\Tmac}}स्तदेव प्रयोजनमिति॥\end{uncertain}

\begin{messwithp}{1000}{2pt}

%\newpage
\llbl{p01_27}\mll{ast1}%\com{anusasana}{योगानुशासनमिति \ldots शास्त्रम्}%
`योगानुशासनमि'ति---%
\fakepar
\tstwll{taiubhanu}{यथा शिष्यो ऽनुशिष्यते \ldots ॰नुशासनमित्युच्यते।}%
\begin{uncertain}यथा शिष्यो ऽनुशिष्यते विशिष्ट%
	\varmm{प्रवृत्तिनिवृत्ति\linelabel{nymdv}नियम}{%
		TMA\Me, 
		{\dn  ॰प्रवृत्तिनियम॰} L%
	}%
	द्वारेण \barfu तथा विशिष्टसाध्यसाधनतदंगनियममात्रसादृश्याद%
	\varmm{न्तेवास्य}{%
		\Td LA\Me, 
		{\dn ॰न्तेवास्या॰} \Tm, 
		{\dn ॰न्तेनास्य॰} M%
	}%
	नु\varmm{शासनवद्योगा}{%
		LMA\Me, 
		{\dn  ॰शासनविद्योगा॰} T%
	}%
	नुशासनमित्युच्यते।
\end{uncertain}%
%\resume{taiubhanu}
\donote{taiubhanu}%
{\labelh{TAIUBH111REF}Cf.\ \rTaiU{1.11.1}{1.11.1}: {\dn वेदमनूच्याचार्यो ऽन्ते\-वासि\-नम\-नुशास्ति} and the \rnnTaiUBh{1.11.1}{1.11.1} on it: {\dn वेदम\-नूच्ये\-त्येवमा\-दि\-कर्तव्यतो\-प\-देशारम्भः प्राग्ब्रह्म\-विज्ञाना\-न्नियमेन कर्त\-व्यानि \barfu श्रौत\-स्मार्त\-कर्माणी\-त्येवम\-र्थः। } {\dn \barfu अनुशासनश्रुतेः पुरुषसंस्कारार्थत्वात्। संस्कृतस्य हि विशुद्धसत्त्वस्यात्मज्ञानमञ्जसैवोत्पद्यते। `तपसा कल्मषं हन्ति विद्ययामृतमश्नुते'} (\rManu{12.104}{12.104}) {\dn इति हि \barfu स्मृतिः। \barfu वक्ष्यति \barfu च---`तप\-सा ब्रह्म विजिज्ञासस्व'} (\rTaiU{3.3.01}{3.3.1}) {\dn इति। \barfu अतो विद्यो\-त्पत्त्यर्थम\-नुष्ठेयानि कर्माणि। \barfu अनुशास्तीत्यनुशासनशब्दात्। अ\-नुशासना\-तिक्रमे हि \barfu दोषो\-त्पत्तिः। प्रागु\-पन्यासा\-च्च कर्मणाम्, केवल\-ब्रह्म\-विद्या\-रम्भाच्च। पूर्वं कर्माण्यु\-पन्यस्तानि। उदितायां च ब्रह्मविद्यायां, `अभयं प्रतिष्ठां विन्दते'} (\rTaiU{3.3.01}{3.3.1}) {\dn `न बिभेति कुतश्चन' `किमहं साधु नाकरवम्'} (\rTaiU{2.9.01}{2.9.1}) {\dn इत्येवमादिना कर्मनैष्किञ्चन्यं दर्शयिष्यतीत्यतो ऽवगम्यते---पूर्वोपचितदुरितक्षयद्वारेण विद्योत्पत्त्यर्थानि कर्माणीति। मन्त्र\-वर्णाच्च---`अविद्यया मृत्युं तीर्त्वा विद्ययामृतमश्नुते'} (\rIU{11}{11}) {\dn इति। ऋतादीनां \barfu पूर्वत्रोपदेशः \barfu आन\-र्थक्य\-परिहारा\-र्थः। इह तु \barfu ज्ञानो\-त्पत्त्यर्थ\-त्वा\-त्कर्तव्य\-नियमा\-र्थः। वेद\-मनूच्या\-ध्याप्या\-चार्यो ऽन्तेवासिनं शिष्यम\-नुशास्ति ग्रन्थ\-ग्रहणादनु पश्चाच्छास्ति तदर्थं \barfu ग्राहयती\-त्यर्थः। अतो ऽवगम्यते ऽधीत\-वेदस्य ध\-र्मजिज्ञासाम\-कृत्वा गुरु\-कुलान्न समा\-वर्तितव्यमिति। `बुद्ध्वा कर्माणि कुर्वीत'} (\Apastambhadharmasutra\ 9.21.5?) {\dn इति स्मृतेश्च।} \citep[420]{upa1}.%
} %
\barfu अ\-नुशिष्टिरनुशासनम्। \barfu यो\-गो ऽनुशिष्यते ऽने\varms{नास्मिन्निति}{\Td LM\Me, {\dn ॰नास्तिन्निति} \Tm, {\dn ॰नान्यस्मिन्निति} A} \barfu वा \mula{यो\-गा\-नुशासनं \var{शास्त्रम्}{\Td LMA\Me, {\dn सस्त्रम्} \Tm}}॥%\donote{anusasana}{\come{anusasana}}
\donote{ast1}{{\dn योगानुशासनं शास्त्रम्।}}

\end{messwithp}

\llbl{p01_28}\mll{athy1}%\com{athad}{अथेत्ययमधिकारा॰ \ldots\ प्रक्रियार्त्थो दृष्ट' इति}
\mula{अथेत्ययम\varmm{धिकारा}{LMA\Me, 
{\dn ॰धिकारो} T}र्त्थः}---अधिकार आरम्भः प्रस्तावः, अर्त्थो ऽभिधेयो ऽस्य। \tstl{sistasmrti}\pause{sistasmrti}<{\dn शिष्टस्मृतिप्रामाण्यात्}>\varsm{शिष्टस्मृति}{TL\Me, {\dn शिष्यस्मृति॰} MA}\varl{maanya}\pause{maanya}<॰माण्यात्॥ ननु>\-प्रामाण्यात्\donote{sistasmrti}{The \emph{śiṣṭasmṛti} here is the \VMBh.  See the beginning of the \VMBh: {\dn अथ शब्दा\-नुशासनम्। अथेत्ययं शब्दो ऽधिकारार्थः प्रयुज्यते।}.}॥\donote{athy1}{{\dn \barfu अथेत्ययमधिकारार्थः।}}\fakepar
%
\llbl{p01_29}ननु\resume{maanya}\donote{maanya}{TM(em.)A\Me\vv{॰ण्यान्ननु}{\Tm MA}, \rdg{॰प्रामण्यन्ननु}{L}} चानन्तर्य्यार्त्थस्याथशब्दस्य स्मरन्ति \barfu शिष्टाः। तथा \vrt{चाह}{\Tmpc \Td LMA\Me}{\rdg{चा\ccs न\cce\ins ह\ine}{\Tm}}---\tst{`वृत्तादनन्तरस्य प्रक्रियार्त्थो दृष्ट' \barfu इति}{\rSBh{1.1.1}. See also \rBSBh{1.1.01}{1.1.1}: {\dn अत्राथशब्द आनन्तर्यार्थः परिगृह्यते, नाधिकारार्थः, ब्रह्मजिज्ञासाया अनधिकारत्वात। मङ्गलस्य च वाक्यार्थे समन्वयाभावात्। अर्थान्तरप्रयुक्त एव ह्यथशब्दः श्रुत्या मङ्गलप्रयोजनो भवति। पूर्वप्रकृतापेक्षायाश्च फलत आनन्तर्याव्यतिरेकात्। सति चानन्तर्यार्थत्वे यथा धर्मजिज्ञासा पूर्ववृत्तं वेदाध्ययनं नियमेनापेक्षते, एवं ब्रह्मजिज्ञासापि यत्पूर्ववृत्तं नियमेनापेक्षते तद्वक्तव्यम्। स्वाध्यायनान्तर्यं तु समानम्। \ldots} \citep[47--51]{bsbh}.}\fubar ॥%\donote{athad}{\come{athad}}

\llbl{p01_30}न। सर्व्वत्रारम्भार्त्थत्वात्, आनन्तर्य्यस्य च गम्यमानत्वात्। \tstwll{PITAPUTRA}{यथा `पुत्र' इत्युक्ते \ldots\  पुत्रशब्दस्यार्त्थः \barfu पिता}%\com{putreti}{यथा `पुत्रे'त्युक्ते \ldots पुत्रस्यार्त्थः पिता}
यथा `पुत्र' \barfu इत्युक्ते \barfu पिता गम्यते, \var{न च}{\Tmpc \Td LMA\Me, {\dn न व} \Tmac} \var{पुत्रशब्दस्यार्त्थः}{LMA\Me\vv{\dn ॰र्थ्ः}{MA\Me}, {\dn पुत्रस्या}[\mist{1}]{\dn र्त्थः} \Tm, {\dn पु\-त्र\-स्या\-र्थः} \Td} \barfu \mms पि\mme ता।%\donote{putreti}{\come{putreti}}
\labelh{PTRPITR}\donote{PITAPUTRA}{Cf.\ \rBSBh{2.2.17}{2.2.17}: {\dn यथैको ऽपि सन्देवदत्तो लोके स्वरूपं संबन्धिरूपं \barfu चापेक्ष्यानेकशब्दप्रत्ययभाग्भवति---मनुष्यो \barfu ब्राह्मणः \barfu श्रोत्रियो \barfu वदान्यो \barfu बालो \barfu युवा \barfu स्थविरः \barfu पिता पुत्रः पौत्रो भ्राता जामातेति, यथा चैकापि सती रेखा स्थानान्यत्वेन निविशमानै\-कदश\-शत\-सहस्रा\-दि\-शब्द\-प्रत्यय\-भेदम\-नुभवति} \citep[521]{bsbh}.  This in turn alludes to \rYBh{3.13}{3.13}: {\dn यथैका रेखा शतस्थाने शतं दशस्थाने दशैका च एकस्थाने। यथैवैकत्वे ऽपि स्त्री माता चोच्यते दुहिता च स्वसा चेति।}.}
 \varsm{तथे}{\Td MA\Me, {\dn तथै॰} \Tm L}हा\mms प्या\mme रम्भ एवाभिधीयते, आनन्तर्य्यन्तु प्रतीयते। \var{अत एव च}{\Td\Lpc MA\Me, {\dn अत वेच} \Tm, {\dn अत एचव} \Lac} `वृत्तादनन्तरस्य प्रक्रियार्त्थो दृष्ट' इत्युच्यते। \barfu \varsm{यदि चा}{\Tmpc \Td\LMA\Me, {\dn यदि वा॰} \Tmac}न\-न्तर्य्या\-र्त्थो \barfu ऽभविष्यत्तदा `वृत्तादा\varms{नन्तर्य्यार्त्थे}{em., \rdg{॰नन्तर्य्ये र्त्थे}{\Tmac}, {\dn ॰नन्तर्यो र्त्थे} \Tmpc, {\dn ॰नन्तर्यो ऽर्थे} \Td, {\dn ॰नन्तर्य्यर्त्थे} L, {\dn ॰नन्तर्यो ऽर्थे} M, \rdg{॰नन्तर्ये ऽर्थे}{A\Me}} \var{प्रक्रियाया}{\Td\LMA\Me\vv{॰यायाः}{\Me}, {\dn प्रकियाया} \Tm}' इत्यवदिष्यत्॥
 
\llbl{p01_31}अनन्तरस्येति चानन्तरभाविन एवाभिधानम्।
अपि चैवं हि स्मृतिः---%
\labelh{VMBH1138_E}%
\tstl{mbh1138}\pause{mbh1138}<{\dn किञ्चिदव्ययं}\sel{\dn विद्यत इति}>%
`किञ्चिदव्ययं वि\-भ\-\mms \ktya \-\mme%
\varms{र्त्थप्रधानं}{%
	\Td\LMA\Me, 
	{\dn ॰र्त्थं प्रधानं} \Tm%
} %
किञ्चित्क्रियाप्रधानम्। उच्चै\-{}र्नीचैरि\-{}त्यादि विभ\ktya र्त्थ\-प्रधानम्। हिरु\-क्पृथ\-गि\-त्यादि %
\var{क्रियाप्रधानम्}{%
	T\LMA, 
	{\dn क्रियाप्रधानम् \eil इति\eir} \Me%
}%
\fubar । \fubar न \fubar चैत\-द्व्य\-ति\-रे\-के\-णा\-%
\varmm{व्य\-या\-ना}{%
	T\-M(em.)\-A\-\Me, 
	{\dn ॰प्रयाना॰} L%
}%
म%
\vrtms{र्त्थो}{\Td\LMA\Me}{%
	\rdg{॰र्त्थ॰}{\Tm}%
	} %
\var{विद्यत' इति}{em. {\dn वि\-द्य\-ते} $\Sigma$.}।\donote{mbh1138}{\rnnVMBh{1.1.38 or 1.2.47}{1.1.38 or 1.2.47} ad \rAAP{1.1.38}{1.1.38} or \rnAAP{1.2.47}{1.2.47}.  Cf.\ the readings in Kiel\-horn's third edition (p.\,95): {\dn किं\-चि\-दव्ययं विभ\ktya र्थप्रधानं किंचित्क्रियाप्रधानम्। उच्चैर्नीचैरिति वि\-भ\-\ktya\-र्थ\-प्र\-धा\-नं हिरु\-क्पृथगि\-ति क्रियाप्रधानम्। त\-द्धितश्चापि कश्चिद्विभ\ktya र्थप्रधानः \barfu कश्चित्क्रियाप्रधानः। तत्र यत्रेति विभ\ktya र्थप्रधानो नाना विनेति \barfu क्रियाप्रधानः। न चैतयोरर्थयोर्लिङ्गसंख्याभ्यां योगो ऽस्ति॥} and (p.~223) 
{\dn अव्ययं हि किंचिद्विभ\ktya र्थप्रधानं किंचित्क्रियाप्रधानम्। उच्चैर्नीचैरिति विभ\ktya र्थप्रधानं हिरुक्पृथगिति क्रियाप्रधानम्। तिबन्तं चापि किंचिद्विभ\ktya र्थप्रधानं किंचित्क्रियाप्रधानम्। काण्डे कुद्ये इति विभ\ktya र्थप्रधानं रमते ब्राह्मणकुलमिति क्रियाप्रधानम्। न चैतयोरर्थयोर्लिङ्गसंख्याभ्यां योगो ऽस्ति।}
} तत्रा\-थ\-शब्द\-स्या\-नन्तर्य्या\-र्त्थ\-त्वे सति \var{वि\-भक्ति\-श्रवणं युक्तम्}{T, {\dn वि\-भक्ति\-श्रय\-ण\-म\-युक्तं} L, {\dn वि\-भक्ति\-श्रवण\-म\-युक्तं} M, {\dn वि\-भक्ति\-श्रवण\-म\-प्यु\-क्तं} A, {\dn विभ\ktya श्रवणम्। विभक्तिश्रवणमयुक्तम्} \Me}, \var{सत्वप्रधानत्वात्}{\Tm L\Me\vv{॰सत्त्व॰}{\Me}, {\dn स त्वप्रधानत्वात्} \Td MA}\fubar । \barfu आरम्भक्रिया\varms{र्त्थत्वे तु सति}{em., {\dn ॰र्त्थत्वे तु स तु} \Tm L, {\dn ॰र्थत्वे तु} \Td, {\dn ॰र्थत्वे तु न} M\-A\-\Me} \vrtsm{वि\-भ\-\ktya}{\Tm\-L}{\rdg{विभ\ktya र्थ॰}{\Td MA}, \rdg{विभ\ktya\edl र्थ\edr ॰}{\Me}}\vrtmm{श्रवणं यु}{T}{\rdg{॰श्रवणमयु॰} L\-M\-A\-\Me\-\vv{॰श्रव\ccs न\cce णमयु॰}{L}}क्तम्, अस\mts त्वप्रधानत्वा\mte दिति॥

\llbl{p01_33}\labelh{TSMDTH}तस्मादथशब्दः प्रस्तावार्त्थ \vrtsm{एवेति युक्त}{TMA\Me}{\rdg{एवेति\lost{2}}{L}}म्॥

\llbl{p01_34}`अ\vrtmm{थेत्यधिका}{TMA\Me}{\rdg{॰[\hlg{थे}]त्यधिका॰}{L}}रार्त्थ' \barfu इति---स्वरूपा\varmm{वद्योतना}{TMA\Me, {\dn ॰वद्योतन\ccs ा\cce ॰} L}र्त्थ इति\mms \var{श\mme ब्दः}{\Td\LMA\Me, {\dn \lost{1}ब्दाः} \Tm}। % need checking (\Tm)!!!
यथा `गौरि'त्याहेति। प्र\vrtmm{सिद्धो ऽप्य\-थ}{TLM\Me\vv{॰द्धो प्य॰}{\Tm L}}{\rdg{॰सिद्ध्येप्यथ॰}{A}}\-शब्दार्त्थो ऽधिक्रियमाण\-\varms{विषये}{\Td\LMA\Me, {\dn विषय॰} \Tm} शिष्यबुद्धिसमाधानार्त्थं कीर्त्त्यते। \barfu सिद्धो \mts\vrt{ह्यय\-न्न्या\-यः \barfu श्रु\-तौ\mme}{\Me(em.)}{\rdg{ह्ययन्न्या\lost{2}तो}{L}, \rdg{ह्ययन्यायः \barfu श्रुतो}{M}, \rdg{\ccs ह्य\cce यं \barfu ना\ccs न्य\cce \ins य\ine ः। \barfu श्रुतो}{A}} % need checking!!!
\tstl{bau244}\pause{bau244}<व्याख्यास्यामि\sel निदिध्यायस्व>`\vrt{व्या\mde\-ख्या\-स्या\-मि}{\Td MA\Me}{\rdg{\lost{3}या\-ख्या\ccs\hlg{य्य}ाव्या\-दि\cce \ins स्या\ine मि}{\Tm}, \rdg{व्याख्यास्यास्यादि}{L}} ते \mds व्या\mms च\mte \vrtms{क्षाणस्य}{T\-\Me}{\rdg{॰क्ष\-ण\-स्य}{L\-M\-A}} तु मे \var{निदिध्यासस्वे\quotec ति}{\Tmpc\Me, {\dn निदिध्यासर्त्थेति} \Tmac, {\dn निदिध्यासस्येति} \Td, {\dn नि\-दि\-ध्या\-स\-ञ्चे\{॰सं चे॰ \textrm A\dn\}ति} \LMA}\donote{bau244}{\rBAU{2.4.04}{2.4.4} (\rnBAU{4.5.05}{4.5.5}).}\fubar ॥
  
\llbl{p01_35}\mll{khp1}कः \barfu पुनर्य्योगो यस्यानुशासनम्प्रस्तुतमित्याह---\linelabel{dpatha}\mula{योगः समाधिरि}ति। \barfu %
\tstwll{TWODP}{`यु\-जि\-{}र्य्यो\-गे' \ldots\ `युज समाधौ'}%
समाधिग्रहणा%
% \tstms{`द्यु\-जि\-{}र्य्यो\-ग' इति}{Dhātupāṭha 7.7.} निवर्त्त्य %
`द्यु\-जि\-{}र्य्यो\-ग' इति निवर्त्त्य %
%\tstl{dhp468}\pause{dhp468}<`युज समाधावि'त्य॰>%
`युज \var{समाधावि'त्ययं}{\Td MA\Me, {\dn समाधामित्ययं} \Tm L}%
%\donote{dhp468}{Dhātupāṭha 4.68.} %
\barfu गृहीतः।%
\donote{TWODP}{Dhātupāṭha 7.7 and 4.68.}
\linelabel{dpathaend} योगः समाधानम्॥\donote{khp1}{{\dn योगः \barfu समाधिः।}}

\llbl{p01_36}ननु च \mts `योगः स\mte माधिरि'\vrtms{त्ये\-व\-म\-न्ते\-न}{TL\Me}{\rdg{॰त्ये\-व\-म\-र्थे\-न}{\Mac}, \rdg{॰त्येव मतं तेन}{\Mpc A}} व्याख्याते सूत्रे `स च सार्व्वभौम' इत्यादि \varsm{भाष्यम\-सम्ब}{T\Mpc A\Me, {\dn भाष्यसम्ब॰} L, {\dn भाष्यं सम्ब॰} \Mac}न्ध\-न्दृश्यते? 

\begin{messwithp}{1000}{2pt}

\llbl{p01_37}नैष \barfu दोषः। इह `समाधिरि'त्युक्ते समाधीयमानापेक्षत्वात्; समाधीयमानानाञ्च बहुत्वात्; किमात्मा समाधीयते, किं \barfu वा
\varsm{शरीर}{\Td\LMA\Me, {\dn शारीर॰} \Tm}म्, \barfu आ\mts होस्विदिन्द्रियाणी\mte ति %\com{bahuvidaam}{बहुविदां विप्रतिपत्तेः}
\vrt{बहुविदां}{T\LMA}{\rdg{बहु\edl विदां\edr धा}{\Me}} \barfu विप्रतिपत्तेः%\donote{bahuvidaam}{\come{bahuvidaam}}
; \barfu स्व\-वि\-शे\-ष\-णाकांक्षत्वाच्च \barfu समाधेः॥

\end{messwithp}

\llbl{p01_38}\mll{ksyy1}कस्यायं किंविशेषणक इत्यनुषक्ते 
\vrt{प्रश्न}{em.}{\rdg{प्रश्ने}{\all}}
इदमाह---\mula{स च सार्व्वभौमश्चित्तस्य धर्म्म} इति। 
चित्तस्य धर्म्मो नात्मादीनाम्। तच्च चित्तं स्वयमेव समाधीयते 
\tstl{trentara}\pause{trentara}<{\dn कर्त्रन्तर॰}>\varsm{क\mms र्त्रन्त\mme र}{\Td, {\dn त्रेन्तर॰} L, {\dn \mist{5}न्तर॰} M, {\dn \mist{2}त्रेन्तर॰} A, {\dn \eil समाधात्र\eir न्तर॰} \Me}\donote{trentara}{Cf.\ note 34 of \citet{wezler:yvi1}\index{Wezler, Albrecht}.}निरपेक्षत्वात्॥\donote{ksyy1}{{\dn स च सार्व्वभौमश्चित्तस्य \barfu धर्म्मः।}}

\llbl{p01_39}\mll{bhmy1}काः पुनर्भूमय इत्यत आह---\mula{क्षिप्तम्मूढमि}त्यादयो \mula{भूमय} इति। 
\tst{क्षिप्तमिति कर्म्मकर्त्तरि \barfu निष्ठा। यथा भिन्नः कुसूलः स्वयमेवेति।}{%
  See \rVMBh{3.1.87, Vārttikas 5--7, esp.\ 7}{3.1.87, Vārttikas 5--7, esp.\ 7}: \ldots {\dn क्त। भिन्नः कुशूल इति कर्म। स  
  यदा स्वातन्त्र्येण विवक्षितस्तदास्य कर्मवद्भावः स्यात्। तस्य \barfu प्रतिषेधो वक्तव्यस्तस्मिन्प्रतिषिद्धे ऽकर्मकाणां भावे क्तो 
  भवतीति भावे क्तो यथा स्यात्। भिन्नं कुशूलेनेति। क्त॥} \citep[68]{mbh3}.%
} 
\mula{क्षिप्तम}निष्टविषया\vrtms{सञ्जनेन}{\Td\-L\-M\-A\-\Me}{\rdg{॰सञ्जनेने}{\Tm}} स्तिमितम्; 
\mula{मूढ}न्निर्व्वि\-{}वेकम्; \mula{विक्षिप्त}न्नानाक्षिप्तं कर्म्मकर्त्तर्य्येव; विवेकाक्षमञ्च वि\-क्षिप्त\-त्वादेव; 
\mula{एकाग्र}न्तुल्य\-प्र\mms त्य\mme\-य\-\vrtms{प्रवाहम्}{\Me(em.)}{\rdg{॰वाहम्}{T\LMA}}; \mula{निरुद्धमिति} प्रत्ययशून्यम्---चित्तम्॥\donote{bhmy1}{{\dn क्षिप्तम्मूढं विक्षिप्तमेकाग्रन्निरुद्धमिति \barfu भूमयः।}}

\llbl{p01_40}ननु च भूमि\mms षु\mme\ धर्म्मेषु विवक्षितेषु किमर्त्थं क्षिप्तमित्यादिना धर्म्म्युच्यते?

\llbl{p01_41}नैष \barfu दोषः। धर्म्मिणा धर्म्म एवोपदिश्यते। 
धर्म्माणान्धर्म्मिविषय\varms{त्वात्। \barfu यथा}{\Td\LMA\-\Me\vv{\dn त्वाद्यथा}{\Td\LMA}, {\dn त्वात् \barfu द्यथा} \Tm} गोत्वे किं लिंगमिति पृष्टे विषाणी ककुद्मान्प्रान्तवालधिरिति धर्म्मिणा धर्म्म 
एवोपदिश्यते। तस्माद्विक्षेपादयश्चित्तस्य \var{भूमयो}{\Tmpc\Td\Lpc MA\Me, {\dn भूमय} \Tmac, {\dn भूयो} \Lac} धर्म्मा \barfu इत्यर्त्थः॥ 
  
\llbl{p01_42}ननु च भूमीनाञ्चित्तधर्म्मत्वे समाधेश्च, कथम्भूमिभिः समाधिराधारत्वेन विशेष्यते 
\var{सार्व्वभौम}{\Me(em.),  {\dn सार्व्वभूम} T\LMA} \barfu इति---\var{न}{TLM, om.\ A\Me}\fubar । 
\barfu सामान्य\-भूत\-त्वात्समाधेः, भूमीनाञ्च विशेषत्वात्---\barfu यथा 
\vrtsm{क्षिप्त}{\Me(em.)}{\rdg{विक्षिप्त॰}{\Td MA}, \rdg{विक्षिप्ता॰}{\Tm L}}न्तिष्ठति, मूढन्तिष्ठति, \vrtsm{विक्षिप्त}{\Me(em.)}{\rdg{क्षिप्त॰}{T\LMA}}मेकाग्र\vrtmm{न्निरुद्ध}{conj.}{\rdg{\om}{$\Sigma$}}ञ्चेति क्षिप्तादिभूमिषु 
स्थितिरनुवर्त्तते सामान्यम्। %
\linelabel{smdhsth}%
\barfu स्थितिश्च \barfu समाधिः। तस्मात्सामान्यरूपः सर्व्वासु भूमि\mms षु\mme\ 
\vrt{साधारण्येन}{em.}{\rdg{साधान्येन}{\Tmac}, \rdg{प्राधान्येन}{\Tmpc \Td\LMA\Me}} वर्त्तत इति॥

%\comml{aa5141}
\llbl{p01_43}\tst{`सर्व्वभूमिपृथिवी'त्यैश्वर्य्यार्त्थीयो \barfu ऽण्प्रत्ययः। `अनुशतिकादि'त्वादुभयपदवृद्धिः}{See \rAAP{5.1.41--2}{5.1.41--2}: {\dn सर्वभूमिपृथिवीभ्यामणञौ। \barfu तस्ये\-श्व\-रः।} and  \rAAP{7.3.20}{7.3.20}: {\dn अनुशतिकादीनां च}.}%\donote{aa5141}{\come{aa5141}}
॥ 
% the note should be more elaborate!!!

\llbl{p01_44}\tstwll{bhuumis}{अन्ये पुनर्बाह्याध्यात्मिकान्संयमविषयान्भूमय
इत्याचक्षते।}अन्ये\labelh{ANYEPUNARedn}%\com{samayama}{अन्ये \ldots संयमो भवति}
\ पुनर्बाह्याध्यात्मिकान् 
\vrtsm{संयम}{T\Me}{\rdg{समयम॰}{\LMA}}% CHECK!
विषयान्भूमय इत्याचक्षते।\donote{bhuumis}{Cf.\ \rYS{3.06}{3.6}: {\dn तस्य भूमिषु \barfu विनियोगः} and the \rnnYBh{3.06}{3.6} on it. The objects (\visaya) of \samyama\ are described in \rYS{3.16}{3.16}\,ff.}
तेषां \tst{`विक्षिप्ते चेतसी'ति}{See\ p.~\pageref{VKSPTCTS1}} सामानाधिकरण्यानुपपत्तिः, स्ववचनविरोधश्च। 
कथं? क्षिप्यते ऽस्मिन्निति क्षिप्तम्, मुह्यते ऽस्मिन्निति मूढम्, \vrt{विक्षिप्यते ऽस्मिन्निति \barfu च विक्षिप्तम्}{conj.}{\rdg{\om}{$\Sigma$}}\fubar । यदा चैवन्न तदा \barfu \vrtsm{संयम}{\Tmpc\Td\Me}{\rdg{समयम॰}{\Tmac LMA}}\-%
विषयता, तत्र क्षेपाद्यसम्भवात्। एकाग्रावस्थायां हि \vrt{संयमो}{T\Me}{\rdg{समयमो}{LMA}} भवति॥%\donote{samayama}{\come{samayama}}

\llbl{p01_45}\vrt{किञ्चान्यत्}{T\Me(em.)}{{\dn किञ्चान्य} LMA}। \tst{क्षिपेश्च ध्रौव्याद्यर्त्थाभावादधिकरणे \mts निष्ठा\vrt{प्रत्यया\mte भावः}{LMA\Me}{{\dn \lost{3}भावः} \Tm, {\dn \mist{3}मा\-नः} \Td}}{See \rAAP{3.4.76}{3.4.76} {\dn क्तो ऽधिकरणे च \barfu ध्रौव्य\-गति\-प्रत्य\-व\-साना\-र्थेभ्यः॥}}॥

\llbl{p01_46}किञ्च निरुद्धे चेतसि \vrtsm{संयम}{T\Me(em.)}{\rdg{समयम॰}{LMA}}स्यापि \barfu तत्राभावाद्वृत्तिशून्य\vrtmm{त्वात्संयम}{T\Me(em.)}{\rdg{॰त्वात्समयम॰}{LMA}}विषयाभावः। न हि निरोधे विषयविशेषे चित्तन्निरुध्यते। तदा हि विषयित्वान्निवर्त्तते। परिगणनानुपपत्तिश्च। न हि \vrt{क्षि\-प्ता\-दयः}{\Me(em.)}{\rdg{विक्षिप्तादयः}{TLMA}} पञ्चैव \vrt{संयमस्य}{T\Me(em.)}{\rdg{समयस्य}{LMA}} विषयास्तेषामानन्त्यात्॥

\llbl{p01_47}ननु \vrt{चान्यथा}{T\-M\-A\-\Me}{\rdg{चान्यधा}{L}} योगं केचिदिच्छन्ति? तथा चाहुः---
\begin{quote}\firmlist
%\com{vs}{इन्द्रियमनो॰ \ldots\ संयोगः}%
\llbl{p01_48}%
\tstwll{VSFIVETWO}{इन्द्रियमनोर्त्थसन्निकर्षात्सुखदुःखे। तदनारम्भ आत्मस्थे मनसि सशरीरस्य सुखदुःखाभावः प्राणमनोविनिग्रहापेक्षः \barfu संयोगः}%
\str{इन्द्रियमनोर्त्थसन्निकर्षात्सुखदुःखे। तदनारम्भ आत्मस्थे मनसि %
\vrt{स\-श\-री\-र\mms स्य\mme}{\Td\Me(em.)}{%
	\rdg{स \barfu सशरीर[\mist{1}]}{\Tm}, 
	\rdg{शरीरस्य}{LMA}, 
	\rdg{\eil स\eir शरीरस्य}{\Me}%
} 
सुखदुःखा\-%
\vrtms{भावः}{TL\Me}{%
	\rdg{भावं}{MA}%
} %
\barfu प्राण\-मनो\-वि\-निग्रहा\-पेक्षः \vrt{संयोगः}{TLMA}{%
	\rdg{संयोगो \eil योग\eir\  इति}{\Me}%
}%
॥%
}\donote{VSFIVETWO}{\rVS{5.2.16--7}.  See \citet{wezler:yviandvs}\index{Wezler, Albrecht}.  Cf.\ \Carakasamhita\ \citep{caraka} Śarīra\-sthā\-na 2.138--139:
{\dn
    आत्मे\-न्द्रिय\-मनोऽर्थानां सन्नि\-कर्षा\-त्प्रवर्तते।
    सुख\-दुःख\-मनारम्भादा\-त्म\-स्थे मनसि स्थिरे॥
    नि\-वर्तते तदु\-भयं वशि\-त्वं चो\-पजायते।
    स\-शरीर\-स्य योग\-ज्ञाः तं \barfu योग\-मृषयो \barfu विदुः॥}; % \end{verse}
\rNBh{4.2.38}: {\dn स तु प्रत्याहृत\-स्ये\-न्द्रिये\-भ्यो मन\-सो धारके\-ण प्र\-यत्नेन धार्य\-माण\-स्या\-त्मना \barfu संयोग\-स्तत्त्व\-बुभुत्सा\-विशिष्टः। सति हि तस्मि\-न्निन्द्रिया\-र्थेषु बुद्ध\-यो नो\-त्पद्यन्ते, तद\-भ्यास\-वशा\-त्तत्त्व\-बुद्धि\-रुत्पद्य\-ते॥}; \PDhS\ \citep[187]{pdhs}: {\dn अस्म\-द्विशिष्टानान्तु योगि\-नां युक्तानां योगज\-धर्मा\-नु\-गृही\-तेन मनसा स्वात्मा\-न्तरा\-काश\-दिक्काल\-परमाणु\-वायु\-मनस्सु तत्समवेत\-गुण\-कर्म\-सामान्य\-विशेषेषु सम\-वाये चा\-वितथं स्वरूप\-दर्शन\-मुत्पद्यते। वियुक्ता\-नाम्पुन\-श्चतुष्टय\-सन्निकर्षा\-द्योगज\-धर्मा\-नुग्रह\-सा\-मर्थ्यात्सूक्ष्म\-व्यवहित\-विप्रकृष्टेषु प्रत्य\-क्षमुत्पद्यते। तत्र सामान्य\-विशेषेषु स्वरूपा\-लोचन\-मात्र\-म्प्रत्यक्ष\-म्प्रमाण\-म्प्रमेया द्रव्यादयः पदार्थाः प्रमाना\-त्मा प्रमिति\-र्द्रव्यादि\-विषय\-ञ्ज्ञानम्। सामान्य\-विशेष\-ज्ञानो\-त्पत्तावविभक्त\-मालोचन\-मात्र\-म्प्रत्यक्ष\-म्प्रमाण\-मस्मिन्नान्यत्प्रमाणा\-न्तरम\-स्त्य\-फलरूपत्वात्। अथ वा सर्वेषु पदार्थेषु चतुष्टय\-सन्निकर्षा\-दवि\-तथम\-व्यपदेश्यं यज्ज्ञान\-मुत्पद्यते तत्प्रत्यक्ष\-म्प्रमाणम्प्रमेया द्रव्यादयः पदार्थाः प्रमाता\-त्मा प्रमिति\-र्गुण\-दोष\-माध्यस्थ्य\-दर्शन\-मिति॥}; \PaSu\ \citep[6]{pasupatasutra}: {\dn किं तदि\-त्यु\-च्यते योगम्। \barfu अत्रा\-त्मेश्वर\-संयोगो \barfu योगः।}}%\donote{vs}{\come{vs}}

\llbl{p01_49}\linelabel{vscomstart}\mula{तदि}ति प्रकृतापेक्षम्। \linelabel{yosau}यो \barfu ऽसौ \mula{सुखदुःख}योरात्मे\mula{न्द्रियमनोर्त्थसन्निकर्षो} हेतुः, \barfu \linelabel{tasya}%
त\-स्या\-\mula{नारम्भ}स्तदनुत्पत्तिः। स कथम्भवति? \mula{आत्मस्थे मनसि} %
\vrt{नेन्द्रियस्थे}{TMA\Me}{%
    \rdg{ने[\mist{4}]}{L}}%
। \mula{स\-श\-री\-र\-स्या}\-विशीर्ण्णशरीरस्य। \barfu तदा %
\tst{`कारणाभावात्कार्याभाव' इति}{\rVS{1.2.1}.} %
सन्निकर्षा\-{}भावे \barfu \mula{सु\-ख\-दुः\-ख}\-%
\vrtmm{योरप्य\mula{भाव}}{T\Lpc MA\Me}{%
    \rdg{॰योरप्यवा॰}{\Lac}%
}ः। \vrt{तस्याम\-वस्था\-यां}{MA\Me}{\rdg{तस्यावस्थायां}{\Tm}, \rdg{तद\-व\-स्था\-यां}{\Td}, \rdg{तस्यामस्थायां}{L}} \barfu यो \barfu ऽसौ \barfu विभोरात्मनो \barfu मनसा \vrt{\mula{संयोगः}}{TMA\Me}{\rdg{सायोगः}{L}}, \linelabel{vscomsa}स \mula{प्राण\-मनो\-विनिग्रहा\vrtms{पेक्षः}{MA\Me}{\rdg{॰पेक्ष॰}{TL}}} \barfu संयोगविशेषो \barfu योगः\linelabel{vscomend}
\end{quote}

%\newpage
---\varl{ntoc1}\pause{ntoc1}<{\dn इति॥ अत्रो॰}>{इति॥%
\fakepar
\llbl{p01_50}अत्रो}\donote{ntoc1}{LMA\Me\vv{इत्यत्रो॰}{LMA}, \rdg{इत्यन्तो॰}{T}}च्यते---


\llbl{p01_51}\linelabel{aatmasthe}`आत्मस्थे मनसी'त्ययुक्तम्, \barfu \vrtsm{सर्व्वदात्म}{TLMA}{\rdg{सर्वदा आत्म॰}{\Me}}स्थत्वा\mms न्मन\mme सः॥

\llbl{p01_52}इन्द्रियादिसन्निकर्षा\vrtmm{नारम्भा}{\Tmpc\Td\Lpc MA\Me}{\rdg{॰दारम्भा॰}{\Tmac}, \rdg{॰\ccs द\cce नारम्भा॰}{L}}पे\vrtmm{क्षया`त्म}{LMA}{\rdg{॰क्षया आत्म॰}{T\Me}}\labelh{TEcorr1}\vrtms{स्थ' इति---इति}{em.}{\rdg{॰स्थ इति}{TLMA}, \rdg{॰स्थे इति}{\Me}} चेत्---`तदनारम्भ' इत्येव सिद्ध\-{}त्वा\-{}दा`त्मस्थे मनसी'त्यनर्त्थकं स्यात्॥

\llbl{p01_53}कि\mts ञ्चान्यत्---मु\mte क्तस्यापि मनस आत्मस्थत्वादिन्द्रियसन्निकर्षाभावाच्च योग\mms ः\mme\ प्राप्नोति।
\tstwll{atmpm1}{तस्यापि हि \barfu सर्व्वगतत्वात्}तस्यापि हि \barfu सर्व्वगतत्वा\donote{atmpm1}{See \PDhS\ \citep[70]{pdhs}: {\dn तथा चात्मेतिवचना\-त्परम\-मह\-त्परि\-माणम्}.  Cf.\ \rVS{7.1.28--9}: {\dn \barfu विभवान्महानाकाशः। \barfu तथाचात्मा।}.}\tstwll{mnsnt1}{मनसश्च नित्यत्वात्}न्मनसश्च नित्यत्वा\donote{mnsnt1}{See \rVS{3.2.2}: {\dn द्रव्यत्वनित्यत्वे वायुना व्याख्याते।}; \PDhS\ \citep[130]{pdhs}: {\dn नित्यं परमाणुमनस्सु तत्पारिमाण्डलुयम्}.}त्स\vrtmm{र्व्वदात्म}{LMA\Me}{\rdg{॰र्व्वन्त\vv{॰र्वं त}{\Td}दात्म॰}{T}}स्थत्वमेव॥ % VS 3.2.3 manas' nityatva; PDhS 70 tathaa caatmetivacanaatparamamahatparimaa.nam (VS 7.1.29 tathaacaatmaa)

\llbl{p01_54}\tstwll{ATMPRDS1}{किञ्चात्मनः \ldots\ मिथ्यात्वात्}किञ्चात्मनः प्रदेशाभावादात्मस्थ \vrt{इत्ययुक्तम्}{TMA\Me}{\rdg{इत्य[\mist3]}{L}}। %\ldots\ % ref?
न \barfu चा\mts प्युपचरितात्म\mte प्रदेशसंयोगः परमार्त्थस्य योगस्य हेतु \vrtsm{स्यात्। उप}{\Tm\Lpc MA\Me\vv{स्यादुप॰}{\Tm MA}}{\rdg{स्यात् अप॰}{\Td}, \rdg{॰त्वादुप॰}{\Lac}}\-च\-रित\-स्य मिथ्यात्वात्॥\donote{ATMPRDS1}{\labelh{BSBHTHREETHREE}Cf.\ \rnnNS{2.2.13–17}{2.2.13–17}/\rnnNBh{2.2.13–17}\ 2.2.13–17: {\dn अनित्यः शब्द इत्युत्तरम्। कथम्?---{\dnb आदिमत्त्वादैन्द्रियकत्वात्कृतकवदुपचाराच्च॥} \ldots\ इतश्च शब्द उत्पद्यते नाभिव्यज्यते---कृतकवदुपचारात्। तीव्रं मन्दमिति कृतकमुपचर्यते, तीव्रं सुखं मन्दं सुखं तीव्रं दुःखं मन्दं दुःखमिति, उपचर्यते च तीव्रः शब्दो मन्दः शब्द इति। \ldots\ 
% p.603
%vyañjakasya tathābhāvād grahaṇasya tīvramandatā rūpavaditi cenna abhibhavopapatteḥ---saṃyogasya vyañjakasya tīvramandatayā śabdagrahaṇasya tīvramandatā bhavati na tu śabdo bhidyate yathā prakāśasya tīvramandatayā rūpagrahaṇasyeti, tac ca na, evam abhibhavopapatteḥ---tīvro bherīśabdo mandaṃ tantrīśabdam abhibhavati na mandaḥ/ na ca śabdagrahaṇam abhibhāvakam, śabdaś ca na bhidyate, śabde tu bhidyamāne yukto 'bhibhavaḥ/ 
तस्मादुत्पद्यते शब्दो नाभिव्यज्यत इति। \ldots\ 
%
% p.604
% abhibhavānupapattiś ca vyañjakasamānadeśasyābhivyaktau prāptyabhāvāt---vyañjakena samānadeśe 'bhivyajyate śabda ity etasmin pakṣe nopapadyate 'bhibhavaḥ/ na hi bherīśabdena tantrīsvanaḥ prāpta iti/
%
% p.605
%aprāpte 'bhibhava iti cet śabdamātrābhibhavaprasaṅgaḥ---atha manyeta asatyām prāptāv abhibhavo bhavatīti, evaṃ sati yathā bherīśabdaḥ kañcit tantrīsvanam abhibhavati evam antikasthopādānam iva davīyaḥstho pādānān api tantrīsvanān abhibhaved aprāpter aviśeṣāt/ tatra kvacid eva bheryāṃ praṇāditāyāṃ sarvalokeṣu samānakālās tantrīsvanā na śrūyeran iti/ nānābhūteṣu śabdasantāneṣu satsu śrotrapratyāsattibhāvena kasyacic chabdasya tīvreṇa mandasyābhibhavo yukta iti/
%
% mandasyābhibhavo yukta iti/] p.606
%kaḥ punar ayam abhibhavo nāma? grāhyasamānajātīyagrahaṇakṛtam agrahaṇam abhibhavaḥ; yatholkāprakāśasya grahaṇārhasyādityaprakāśeneti
॥१३॥
%
%
%_____________________________________________________________________
%
%
%     ********************** NySBh_2,2.14 **********************
%
%
% p.607
%
%     NyS_2,2.14: 
{\dnb न, घटाभावसामान्यनित्यत्वान्नित्येष्वप्यनित्यवदुपचाराच्च॥१४॥}
%
\ldots\ 
%na khalu ādimattvād anityaḥ śabdaḥ/ kasmāt? vyabhicārāt/ ādimataḥ khalu ghaṭābhāvasya dṛṣṭaṃ nityatvam/ katham ādimān? kāraṇavibhāgebhyo hi ghaṭo na bhavati/ katham asya nityatvam? yo 'sau kāraṇavibhāgebhyo na bhavati na tasyābhāvo bhāvena kadācin nivartyata iti/ yad apy aindriyakatvāt, tad api vyabhicarati, aindriyakaṃ ca sāmānyaṃ nityaṃ ceti/ 
यदपि कृतकवदुपचारादिति, एतदपि व्यभिचरति; नित्येष्वनित्यवदुपचारो दृष्टः---तथा हि भवति वृक्षस्य प्रदेशः कम्बलस्य प्रदेशः, एवमाकाशस्य प्रदेशः आत्मनः प्रदेश इति भवतीति॥१४॥
\ldots\ 
%
%_____________________________________________________________________
%
%
%     ********************** NySBh_2,2.15 **********************
%
%
%% p.608
%
%     NyS_2,2.15: tattvabhāktayor nānātvasya vibhāgād avyabhicāraḥ ||
%
%nityam ity atra kiṃ tāvat tattvam? arthāntarasyānutpattidharmakasyātmahānānupapattir nityatvam, tac cābhāve nopapadyate/ bhāktaṃ tu bhavati yat tatrātmānam ahāsīd yad bhūtvā na bhavati na jātu tat punar bhavati, tatra nitya iva nityo ghaṭābhāva ity ayaṃ padārtha iti/ tatra yathājātīyakaḥ śabdo na tathājātīyakaṃ kāryaṃ kiṃcin nityaṃ dṛśyata ity avyabhicāraḥ//
%
%
%_____________________________________________________________________
%
%
%     ********************** NySBh_2,2.16 **********************
%
%
%% p.609
%yad api sāmānyanityatvād itīndriyapratyāsattigrāhyam aindriyakam iti---
%
%     NyS_2,2.16: santānānumānaviśeṣaṇāt ||
%
%nityeṣv avyabhicāra iti prakṛtam/ nendriyagrahaṇasāmarthyāc chabdasyānityatvam/ kiṃ tarhi indriyapratyāsattigrāhyatvāt santānānumānaṃ tenānityatvam iti //
%
%
%_____________________________________________________________________
%
%
%     ********************** NySBh_2,2.17 **********************
%
यदपि नित्येष्वप्यनित्यवदुपचारादिति। न---%
%
%     NyS_2,2.17: 
{\dnb कारणद्रव्यस्य प्रदेशशब्देनाभिधानात्॥१७॥}
%
नित्येष्वप्यव्यभिचार इति। एवमाकाशप्रदेश आत्मप्रदेश इति नात्राकाशात्मनोः कारणद्रव्यमभिधीयते यथा कृतकस्य। कथं ह्यविद्यमानमभिधीयते? \barfu अविद्यमानता च \barfu प्रमाणतो \barfu ऽनुपलब्धेः। किंतर्हि तत्राभिधीयते संयोगस्याव्याप्यवृत्तित्वम्। \ldots\ %
%परिच्छिन्नेन द्रव्येनाकाशस्य संयोगो नाकाशं व्याप्नोति, अव्याप्य वर्तत इति, तदस्य कृतकेन द्रव्येण सामान्यम्। न \barfu ह्यामलकयोः संयोग आश्रयं व्याप्नोति। %
\barfu सामान्यकृता च भक्तिराकाशस्य प्रदेश इति। \barfu अनेनात्मप्रदेशो \barfu व्याख्यातः।} \ldots\ \citep[594–613]{ns}; \SplitNote%
%\rBSBh{2.2.17}{2.2.17}: {\dn तथाण्वात्ममनसामप्रदेशत्वान्न \barfu संयोगः संभवति। \barfu प्रदेशवतो द्रव्यस्य \barfu प्रदेशवता द्रव्यान्तरेण संयोगदर्शनात्। कल्पिताः प्रदेशा अण्वात्ममनसां भविष्यन्तीति चेत्, न, अविद्यमानार्थकल्पनायां सर्वार्थसिद्धिप्रसंगात्।} \ldots\ \citep[521]{bsbh}; 2.3.53: {\dn अथोच्येत---विभुत्वे ऽप्यात्मनः शरीरप्रतिष्ठेन मनसा संयोगः शरीरावच्छिन्न एव आत्मप्रदेशे भविष्यति, अतः \barfu प्रदेशकृता व्यवस्थाभिसंध्यादीनामदृष्टस्य सुखदुःखयोश्च भविष्यतीति। तदपि नोपपद्यते। कस्मात्? अन्तर्भावात्। विभुत्वाविशेषाद्धि सर्व एवात्मानः सर्वशरीरेष्वन्तर्भवन्ति। तत्र न वैशेषिकैः शरीरावच्छिन्नो \barfu ऽप्यात्मनः प्रदेशः कल्पयितुं \barfu शक्यः। कल्प्यमानो ऽप्ययं निष्प्रदेशस्यात्मनः प्रदेशः \barfu काल्पनिकत्वादेव न पारमार्थिकं कार्यं नियन्तुं शक्नोति। \ldots} \citep[627]{bsbh}.%
}

\llbl{p01_55}\linelabel{indriyamano}किञ्चान्यत्। मनोर्त्थसन्नि\vrtmm{कर्षादि}{em.}{\rdg{॰कर्षाभाव इ॰}{$\Sigma$}}त्येव \vrtsm{सिद्धेरि}{\Me(em.)}{\rdg{सिद्धमि॰}{TLMA}, \rdg{सिद्धेः\edl द्धं\edr\ इ॰}{\Me}}न्द्रियग्रहणमनर्त्थकं \vrtsm{स्या}{TLA\Me}{\rdg{॰त्वा॰}{M}}त्॥\FootnotetextC{}{Cf.\ also \rBSBh{2.2.17}{2.2.17}: {\dn तथाण्वात्ममनसामप्रदेशत्वान्न \barfu संयोगः संभवति। \barfu प्रदेशवतो द्रव्यस्य \barfu प्रदेशवता द्रव्यान्तरेण संयोगदर्शनात्। कल्पिताः प्रदेशा अण्वात्ममनसां भविष्यन्तीति चेत्, न, अविद्यमानार्थकल्पनायां सर्वार्थसिद्धिप्रसंगात्।} \ldots\ \citep[521]{bsbh}; 2.3.53: {\dn अथोच्येत---विभुत्वे ऽप्यात्मनः शरीरप्रतिष्ठेन मनसा संयोगः शरीरावच्छिन्न एव आत्मप्रदेशे भविष्यति, \barfu अतः \barfu प्रदेशकृता व्यवस्थाभिसंध्यादीनामदृष्टस्य सुखदुःखयोश्च भविष्यतीति। तदपि नोपपद्यते। कस्मात्? अन्तर्भावात्। विभुत्वाविशेषाद्धि सर्व एवात्मानः सर्वशरीरेष्वन्तर्भवन्ति। तत्र न \barfu वैशेषिकैः \barfu शरीरावच्छिन्नो \barfu ऽप्यात्मनः प्रदेशः कल्पयितुं \barfu शक्यः। कल्प्यमानो ऽप्ययं निष्प्रदेशस्यात्मनः \barfu प्रदेशः \barfu काल्पनिकत्वादेव न पारमार्थिकं कार्यं नियन्तुं शक्नोति। \ldots} (ibid.:\ 627)%
%\citep[627]{bsbh}%
.}

\llbl{p01_56}अथेन्द्रियैर्व्विना मनोर्त्थसन्निकर्षा\mts भावात्तदना\mte रम्भ इति प्राप्त्य\vrtms{भावे}{\Td\LMA\Me}{\rdg{॰भावो}{\Tm}} प्रतिषेधो न युक्त इति चेत्---नैतदेवम्। विनेन्द्रियैर्म्मनोर्त्थसन्निकर्षो नास्तीत्यर्त्थात्सन्निकर्षप्रतिषेधं कुर्व्वता विषयप्राप्तिरिन्द्रियद्वारा \barfu प्रतिपादिता %
%\rule{0pt}{10pt} %
\mms 
स्यात्॥\linelabel{indriyamanoend}

\llbl{p01_57}\linelabel{sukhaduhkhe}\vrtsm{किञ्चा\mds न्य}{LMA\Me}{\rdg{किंत्वा\mist{1}}{\Td}}त्---सर्व्वप्राण\mde भा\mme जामिन्द्रियमनोर्त्थसन्निकर्षा\vrtmm{त्सुखदुःख}{\Tmpc \Td\LMA\Me}{\rdg{॰त्सुख॰}{\Tmac}}प्रतिलम्भस्य \vrtms{प्रसिद्ध\-त्वा}{\Td\LMA\-\Me}{\rdg{प्रसिद्धत्वप्रसिद्धत्वा॰}{\Tm}}\-त्, % ref. 
 आ\-\vrtmm{त्ममनः}{em.}{\rdg{॰त्मनः॰}{$\Sigma$}}सम्बन्धस्य च नि\-त्य\-त्वात्, सुखदुःखाभावो योग इत्येव \vrt{सिद्धेः \barfu सूत्रशेषो}{\conj}{\rdg{\ccs सि सू\cce सिद्धे सूत्रविशेषो }{\Tm}, \rdg{सिद्धे सूत्रविशेषो}{ \Td}, \rdg{सिद्धेस्तत्र \ilg षो}{L}, \rdg{सिद्धेस्तत्र शेषो }{MA}, \rdg{सिद्धेः, तत्र \eil सूत्र\eir शेषो}{\Me}} \barfu ऽनर्त्थक स्या\-\mms त्॥

\llbl{p01_58}अथापि स्यादि\mms न्द्रि\mme य\mms मनो\mds र्त्थसन्नि\vrtmm{कर्षग्रहणा}{L\Mpc A\Me}{\rdg{॰कर्ष\il ग्रहणा\ir ॰}{M}}दृते\mde\  सुखदुःखाभाव\mme\ इत्येतावति सति मुक्तस्या\-{}पि \barfu सुखदुःखाभावाद्यो\vrtmm{गित्वप्र}{\Me(em.)}{\rdg{॰ग\ccs स्य\cce त्वप्र॰}{\Tm}, \rdg{॰गस्य त्वप्र॰}{\Td}, \rdg{॰गत्वप्र॰}{LMA}}संग इति तन्निवृत्यर्त्थं सूत्र\vrtmm{शेषग्रहण}{\Td\LMA\Me}{\rdg{॰शेषप्रग्रहण॰}{\Tm}}मिति---

\llbl{p01_59}तच्च न, मुक्तस्य सुखदुःखप्राप्त्यभावात्। प्राप्तिपूर्व्वप्रतिषेधस्य  न्याय्यत्वात्। \mts यस्य सुख\-दुः\-ख\-प्राप्तिः\mte, त\vrtmm{द्विषयः सुख}{\Me(em.)}{\rdg{॰द्विषयसुख॰}{TMA}, \rdg{॰द्विषये सुख॰}{L}}दुःखा\-भाव एव योग इति सूत्रशेषप्रत्याख्यानं युज्यते॥\linelabel{sukhaduhkheend}

\llbl{p01_60}\linelabel{sasariirasya}किञ्चान्यत्---उपात्ते ऽपि सशरीर\vrtms{ग्रहणे}{\Td MA\Me}{\rdg{ग्रहणो}{\Tm}} मुक्ताद्विनिवृत्तिर्न्न \vrt{शक्यते}{\Td\LMA\Me}{\rdg{युज्यते शक्यते}{\Tm}} कर्त्तुम्, निष्प्रयोजन\-{}त्वात्। \barfu सुखदुःखाभावश्चेद्यो\mts गः विद्यमान\mte मपि शरीरं \vrtsm{कार्य्याकरणा}{\Me(em.?)}{\rdg{कार्य्या\vv{॰र्या॰}{MA}कारणा॰}{\Tm LMA}, \rdg{कार्यकारणा॰}{\Td}}दनर्त्थकं स्यात्। कार्य्यान\-{}पेक्षायां हि मुक्तामुक्तयोरविशेष एव॥\linelabel{sasariirasyaend}

\linelabel{vinigraha}%
\llbl{p01_61}किञ्चान्यत्---प्राणमनो%
\vrtmm{विनिग्रहा}{T\Me}{%
	\rdg{॰विग्रहा॰}{L\Mpc A}, 
	\rdg{॰ग्रहा॰}{\Mac}%
}%
पेक्ष इत्यपि न %
\vrt{घटते}{TMA\Me}{%
	\rdg{घ\ilg ते}{L}%
}%
, विधारकस्य %
\tstsm{प्रयत्न}{%
	See \PDhS\ \citep[70]{pdhs}: {\dn तस्य गुणाः %बुद्धि\-सुख\-दुःखे\-च्छा\-द्वेष\-%
	\ldots ॰प्रयत्न॰%
	\ldots %
	%\-धर्मा\-धर्म\-संस्कार\-संख्या\-परिमाण\-पृथक्त्व\-संयोग\-विभागाः%
	}. Cf.\ also \rVS{3.2.4}.%: {\dn %
	%प्राणा\-पान\-निमेषो\-न्मेष\-जीवन\-मनो\-गती\-न्द्रिया\-न्तरविकाराः %
%	\ldots सुख\-दुःखे इच्छा\-द्वेषौ प्रयत्न\-श्चेत्या\-त्म\-लिङ्गानि}.%
}%
स्यात्मसमवेत\-{}\varl{tvaatvi}\pause{tvaatvi}<{\dn ॰त्वात्। वि॰}>त्वात्। \tstwll{manasahkryvtvm11}{विद्यमाने ऽपि मनसः क्रियावत्त्वे}वि\donote{tvaatvi}{em.\ \rdg{॰त्वा\-द\-वि॰}{\all}}द्यमाने ऽपि म\mms न\mds सः \barfu क्रिया\vrtms{व\mte त्वे}{LM\Me}{\rdg{[\mist{1}त्वे]}{\Tm}, \rdg{\mist{1}त्व}{\Td}}\donote{manasahkryvtvm11}{Cf.\ \PDhS\ \citep[89]{pdhs}: {\dn क्रियावत्त्वान्मूर्तत्वम्}.}, % this is accidentally corect.
 \vrtsm{इन्द्रियेणा}{T\Lpc MA}{\rdg{इन्द्रि\ccs या\cce येणा}{\Lac}}संयुज्य \vrtsm{मनः}{TLM\Me}{\rdg{मतः}{A}}\vrtms{प्राणनिग्रहो}{TLMA}{\rdg{प्राण\eil मनो\eir निग्रहो}{\Me}} न कर्त्तुम्पर्य्याप्यते॥\linelabel{vinigrahaend}

\linelabel{tadanaarambha}%
\llbl{p01_62}इन्द्रिय%
\vrtms{संयोगे}{T\Lpc MA\Me}{%
	\rdg{॰संहोगे}{\Lac}%
} %
च मनसस्तदनारम्भ \mms इत्य\mme युक्तम्। %
\vrt{मनसश्च}{TL\Mpc A\Me}{%
	\rdg{मनसश्च \ccs प्रनसश्च\cce}{M}%
} %
प्राणनिग्रहे %
\vrtsm{व्यापृत}{TL\Me(em.)}{%
	\rdg{व्या\ins प्या\ine वृत॰}{M}, 
	\rdg{व्या\edl प्या\edr वृत॰}{A}, 
	\rdg{व्यापृ(वृ)त॰}{\Me}%
}% vyaap.rta is rather frequently used by Sankara
त्वादनार\-म्भा\-नु\-%
\vrtms{प\-{}पत्तिः}{TMA\Me}{%
	\rdg{॰पपतिः}{L}%
}%
। न चासतीन्द्रियवायु\-सन्निकर्षे \barfu प्रा\mms ण\mds %
\vrtms{निग्रहः}{\LMA\Me}{%
	\rdg{[\mist{3}]ा}{\Tm}, 
	\rdg{\mist{3}ा}{\Td}%
}%
\mte । न \vrtsm{च नि}{em.}{\rdg{चानि॰}{$\Sigma$}}गृहीते प्राणे योगः, आ\-त्म\-%
\vrtmm{स्थ\-विधारक}{T\Me(em.)}{%
	\rdg{॰स्वविकधारक॰}{L}, 
	\rdg{॰स्वविधारक॰}{MA}%
}%
\-प्रयत्न\-%
\vrtmm{द्वारमनो}{T}{\rdg{॰द्वारा मनो॰}{\LMA\Me}}गतक्रियासम्बन्धहेतुत्वात्प्राणनिग्रहस्य॥\linelabel{tadanaarambhaend}

\llbl{p01_63}\linelabel{vyogasamadhi}अथापि स्यात्---योगः समाधिरिति॥ 

तच्च न, \vrtsm{निष्क्रिय}{TMA\Me}{\rdg{निष्कर्य्य॰}{L}}त्वात्। आ\vrtmm{त्ममनसोर्न्नि}{em.}{\rdg{॰त्मनो नि॰}{$\Sigma$}}त्यसमाधानमेव। \tst{स्थितिश्च समाधिरित्युच्यते}{See \pageref{smdhsth},\lineref{smdhsth}.}॥% ref?
\linelabel{vyogasamadhiend}

\llbl{p01_65}तस्माद्यु\vrtmm{क्तम}{TLM\Me}{\rdg{॰क्तसम॰}{A}}भिधीयते `स च सार्व्व\vrtms{भौमश्चित्तस्य}{\Td\LMA\Me}{\rdg{॰भौमश्चित्तत्तस्य}{\Tm\ (the second {\dn त्त} beginning of a line)}} धर्म्म' इति॥


%\newpage
\llbl{p01_66}यद्येवं समाधिश्चेद्योगः, स च सार्व्वभौमश्चित्तधर्म्मः, स च सर्व्वजनानामयत्नसिद्धः, क्षिप्ताद्यनुगतत्वात्; तथा च सति यथैव श्वास\vrtmm{प्रश्वासा}{TMA\Me}{\rdg{प्रश्वसा॰}{L}}दीनाम्प्रयत्नादृते सिद्ध\vrtmm{त्वाल्लक्षणम}{TLMA}{\rdg{॰त्वात्तत्करणम\edl ल्लक्षणम\edr ॰}{\Me}}\vrtmm{नर्त्थक}{\Tmpc\Td MA\Me}{\rdg{॰नर्त्थ॰}{\Tmac L}}\-म्, एवं योगानुशासनमप्यनर्त्थकम्प्राप्तमित्यत आह---

\llbl{p01_67}\mll{atrvks1}\mula{\vrt{अत्र विक्षिप्ते}{TL}{\rdg{अनविक्षिप्ते}{MA}, \rdg{\eil तत्र\eir\  विक्षिप्ते \edl अनविक्षिप्ते\edr}{\Me}} चेतसि\labelh{VKSPTCTS1} विक्षेपोपसर्ज्जनीभूतः समाधिर्न्न योगपक्षे वर्त्तत} इति।\donote{atrvks1}{{\dn अत्र विक्षिप्ते चितसि विक्षेपोपसर्ज्जनीभूतः समाधिर्न्न योगपक्षे वर्त्तते।}} क्षिप्ताद्य\-{}नुगतस्य योगपक्षे ऽनभिप्रेतत्वात्। न हि \vrt{क्षेपाद्यनुगतः}{\Tm\Lpc MA\Me}{\rdg{विक्षेपाद्यनुगतः}{\Td}, \rdg{क्षेप\ccs त्व\cce ाद्यनुगतः}{L}} समाधिर्भूतार्त्थावद्योतनादि\vrtms{क्षमः}{TL\Mpc A\Me}{\rdg{क्षमं}{\Mac}}, क्षेपादिप्रधानत्वात्। \vrtsm{विक्षेपनि}{\LMA\Me}{\rdg{विक्षेपन्नि\vv{॰पं नि॰}{\Td}॰}{T}}\vrtms{षेधेनैव}{\Tmpc\Td\LMA\Me}{\rdg{॰षधेनैव}{\Tmac}}%\com{nsehn}{निषेधेनैव}\donote{nsehn}{\come{nsehn}}
\ हि \tstwll{pradhaanamalla}{प्रधानमल्लनिबर्हणन्यायेन}प्रधानमल्ल\vrtmm{निबर्हण}{TLM\Me}{\rdg{॰निर्बहण॰}{A}}न्याये\-न\donote{pradhaanamalla}{See \rBSBh{1.4.28}{1.4.28}; \rnBSBh{2.1.12}{2.1.12} for the use of this maxim.} क्षिप्तमूढयोरपि \barfu निषेधः \barfu कृतः। \barfu प्रधानता च विक्षेपसमाधेः, अधिकारयोग्यत्वात्। विक्षेपस्थं हि \barfu चित्तमपक्षपाता\vrtms{दिष्टं}{T\Me(em.?)}{\rdg{॰दृष्टं}{LMA}} \vrtsm{विषय}{TMA\Me}{\rdg{विषमय॰}{L}}मु\-प\-नेतुं शक्यते। न \vrt{तु}{em.}{\rdg{हि}{$\Sigma$}} \barfu विषयासञ्जनादिना \vrtsm{मूढ}{TLMA}{\rdg{क्षिप्त॰}{\Me(em.?)}}\vrtmm{मनिष्ट}{em.}{\rdg{॰मिष्ट॰}{$\Sigma$}}विषयवियोगादिना \barfu वा \vrtsm{क्षिप्त}{TLMA}{\rdg{मूढ॰}{\Me}}म\vrtmm{न्यत्रो}{TLMA}{\rdg{॰न्यथो॰}{\Me(em.?)}}पनेतुम् %\com{vryt}{वार्य्यते}%
\vrt{पार्यते}{\Me(em.)}{\rdg{[\mist{2} े]त}{\Tm}, \rdg{\mist{3}त}{\Td}, \rdg{वार्य्यते}{LMA}}%\donote{vryt}{\come{vryt}}
। `योगपक्ष' इति---समाधित्वे सत्यपि %\com{krykrn1}{कार्य्याकरणात् \ldots कार्याकरणा॰}
\vrtsm{कार्य्याकरणात्प}{\Me(em.)}{%
  \rdg{कार्य्या\vv{॰र्या॰}{M}कारणात्प॰}{\Tm LM}, %
  \rdg{कार्यकरणात्प॰}{\Td}, %
  \rdg{कार्यकारणात्प॰}{A}}क्ष इत्युच्यते। \tstwll{ngrjn}{यथा गच्छतः \ldots\ न स्थितिरित्युच्यते}यथा \vrt{गच्छतः}{TLM\Me}{\rdg{गच्छेतः}{A}} प्रतिपदं विद्यमानापि स्थितिः स्थिति%
\vrtmm{कार्य्याकरणा}{T\Me(em.)}{%
  \rdg{॰कार्य्या\vv{॰र्या॰}{MA}कारण॰}{LMA}}%\donote{krykrn1}{\come{krykrn1}}
  न्न स्थिति\-रित्युच्यते॥\donote{ngrjn}{Cf.\ \MMK\ 2.1--25.}% definitely deserves mention of Nagarjuna

%\pagebreak

\llbl{p01_68}`विक्षेपोपसर्ज्जनीभूत'त्वं विक्षेपस्थैकसमाधिव्य\ktya भिप्रायेण। `सार्व्वभौम' इति \vrt{तु}{TL\Mpc A\Me}{\rdg{{\cm om.}}{\Mac}} प्राधान्यं सामान्येन सर्व्वभूमिषु \barfu वृत्तेः॥

\llbl{p01_69}\linelabel{BhT11s}यद्युप\mts सर्ज्ज\mte नीभावः \vrtsm{योगपक्षा}{conj.}{\rdg{समाधिपक्षा॰}{$\Sigma$}}वृत्तौ कारणम्, इहाप्येकाग्र \mts एकाग्र\mte तो\vrtms{पसर्ज्जनत्वं}{\LMA\Me}{\rdg{॰पसर्ज्ज\vv{॰र्ज॰}{\Td}नं}{T}} स्यात्। ततश्च योगपक्षा\-वृत्तिरिति चेत्, अत आह---\mll{ystkgr1}\mula{यस्त्वेकाग्र} इति।
%
नैकाग्रतायाम्भूम्युपसर्ज्जनत्वम्। \barfu कुतः? 
  %\com{klskrmbl1}{\kle शकर्म्माप्रबलत्वात्। \kle शकर्म्मप्रबलत्वेन हि}%
\vrtsm{\kle शकर्म्माप्रबलत्वात्। \kle शकर्म्मप्र}{T}{
  \rdg{\ins \kle शकर्म्माप्रबलत्वात्\ine \kle शकर्म्माप्र॰}{L},
  \rdg{\kle श\-कर्मा\-प्तत्वात् \kle श\-कर्माप्र॰}{MA},
  \rdg{\kle शकर्माप्रबलत्वात्। \kle शकर्माप्र॰}{\Me}}%
बल\vrtms{त्वेन हि}{TLMA}{%
  \rdg{त्वे न हि}{\Me}}
  %\donote{klskrmbl1}{\come{klskrmbl1}} 
\barfu विक्षेपादीनामुदयः। 
\vrtsm{एका}{T\Lpc MA\Me}{%
  \rdg{एवका॰}{\Lac}%
}\-ग्रता\-भूमौ \mts स\-मा\mte धि\-र्य्यः \mula{स भूतमर्त्थं} यथात्मानम् \mula{प्रद्योतयत्य}वगमयति। 
\begin{uncertain}
अयोग्यर्त्थज्ञानमयथाभूतत्व\vrtmm{गन्धि\-त}{T\-\Lpc\-M\-A\-\Me}{%
%  \rdg{॰गन्धित॰}{T\Lpc MA\Me},
  \rdg{॰गन्थित॰}{\Lac}%
}मेवेति भूतग्रहणम्। %
\end{uncertain}\barfu
\mula{क्षिणोति} \vrt{क्षपयति}{T\Me(em.?)}{%
  \rdg{पक्षयति}{LMA}%
} \labelh{AVGRH1}\tst{पञ्चपर्व्वणो \barfu ऽविद्यादीन् \barfu \mula{\kle शान्}}{See \rYS{2.03}{2.3}: {\dn अविद्यास्मितारागद्वेषाभिनिवेशाः \barfu \kle शाः॥}.}\fubar । % ref!
\mula{कर्म्मबन्धनानि}---%\com{krmbndh1}{कर्म्माण्येव बन्धनानि}
\vrt{कर्म्माण्येव बन्धनानि}{T\Lpc \Mpc A\Me}{%
  \rdg{\om}{\Lac},
  \rdg{\ccs कर्म धर्माधर्म\cce कर्माण्येव बन्धनानि}{M}%
}%\donote{krmbndh1}{\come{krmbndh1}}
, धर्म्माधर्म्मवि%
\vrtms{मिश्राणि; यदि \barfu वा कर्म्मजानि तानि}{em.}{%
	\rdg{॰मिश्राणि यदि कर्म्मजानि तानि}{T}, 
	\rdg{॰मिश्राणि यदि कर्मजाति तानि}{\LMA}, 
	\rdg{॰विमिश्रकर्मजातानि}{\Me(em.?)}%
} %
\vrtsm{जन्म\mds ादी\mde नि ब}{\Tm\LMA}{%
	\rdg{जन्मादिनिब॰}{\Me}%
}\-न्ध\-ना\-नि---\mula{श्लथयति} शिथिलीकरोति। \mula{निरोधमामुखी\-करोत्य}भिमुखीकरोति।%
\donote{ystkgr1}{%
	{\dn यस्त्वेकाग्रे, स भूतमर्त्थम्प्रद्योतयति, क्षिणोति \kle शान्, कर्म्मबन्धनानि श्लथयति, निरोधमामुखीकरोति।}%
}%
\labelh{SSMPRJN} %
\mll{SSMPRJ1}\mula{स सम्प्रज्ञात \barfu इत्या}%
\vrtms{\mula{ख्यायत}}{T}{%
	\rdg{॰ख्यायते}{\LMA\Me}%
}%
\donote{SSMPRJ1}{%
	{\dn  स सम्प्रज्ञात इत्याख्यायते।}%
} \barfu आचार्य्यैः॥\linelabel{BhT11e} % This aacaarya must be the author of the YBh. -> YVi-kaara uses plural to refer to a single author.

\llbl{p01_70}\mll{vtrknugt1}\mula{स च वितर्क्कानु\vrtmm{गतो विचारा}{\Td \Me(em.)}{\rdg{॰गतो पि\vv{॰गतो ऽपि}{MA} विचारा॰}{TLMA}}नुगत} इत्यादि भाष्यम्---सम्प्र\vrtmm{ज्ञातस्य प्रा}{TL\Mpc A\Me}{\rdg{॰ज्ञातस्य \ccs सा\cce प्रा॰}{M}}थम्यात्। लक्ष्णा\-{}भिधाने \vrt{तु}{conj.}{\rdg{\om}{$\Sigma$}} प्राप्ते, \barfu अभ्यर्हितत्वादसम्प्रज्ञा\vrtms{त\-\mds स्यै\mde\mms व लक्ष\mme णं}{\LMA\Me}{\rdg{॰तस्ये[\mist{3}]णं}{\Tm}, \rdg{॰त\mist{1}पलक्षणं}{\Td}} युक्तमिह वक्तुम्---\mula{इति} \tstwll{refys117}{उपरिष्टात्प्रवेदयिष्यामः}सम्प्रज्ञातमु\-\mula{प\-रिष्टा\-त्प्रवेदयिष्याम}\donote{refys117}{See \rYS{1.17}{1.17}: {\dn वितर्कविचारानन्दास्मितारूपानुगमात् \barfu संप्रज्ञातः॥}.} इत्युच्यते॥\donote{vtrknugt1}{{\dn स च \barfu वितर्क्कानुगतो विचारानुगत \barfu इत्युपरिष्टात्प्रवेदयिष्यामः।}}%
\fakepar
\llbl{p01_71}इतश्चासम्प्रज्ञातस्य लक्षणमिहैव व\mts क्तव्यम््---सम्प्र\mte ज्ञातनिरपेक्षो \barfu ऽप्यसम्प्रज्ञातः \barfu  \vrtsm{प्रकृष्ट}{TLMA}{\rdg{पर\edl प्रकृष्ट\edr ॰}{\Me}}\-वै\-रा\-ग्य\-विरामप्रत्यया\vrtms{भ्यासाभ्यां}{\Td\LMA\Me}{\rdg{॰भ्यासाभ्यासाभ्यां}{\Tm}} सिध्यतीत्येतत्प्रदर्श\mms ना\mme र्त्थम्। इह \vrt{तु}{T}{\rdg{त}{L}, \rdg{\om}{MA\Me}} सम्प्रज्ञातलक्षणाभिधाने तदुत्तरकाले चासम्प्रज्ञातलक्ष\vrtmm{णाभिधाने सम्प्र}{T\-\Me(em.)}{\rdg{॰णमभिधानम्प्र\vv{॰नं प्र॰}{MA}॰}{LMA}}\labelh{TEcorr2}\mts ज्ञातापेक्षयै\vrtmm{वासम्प्र}{T}{\rdg{॰व असम्प्र\vv{॰संप्र॰}{MA}॰}{\LMA\Me}}\-ज्ञात\vrtmm{स\-मा\-धा\mde व\-धि}{LM\Me}{\rdg{॰समानाधि॰}{A}}\mme कार \barfu इत्याशंका स्या\-त्। तस्मादु`परिष्टात्प्रवेदयिष्याम' इत्याह॥१॥

% \finisheditionchapter
%
% \chapter{Translation of the Commentary on \PYS\ 1.1}
% \renewcommand{\rightmark}{\MakeUppercase{Translation of YVi 1.1}}
% \thispagestyle{plain}
% \section{Opening stanzas}
\begin{verse}
	\labelh{OSTZ1}
  In\refm{2}%\pratika{yasmin na}
\pratikap{p01_1}\labelh{FVFPTRL} Him \karman\ and its ripening do not exist,%
  \footnote{\label{NONFSFP}This is a reference to \rYS{1.24}{1.24}: \emph{kleśakarmavipākāśayair aparāmṛṣṭaḥ}.  Impurities (\klesa\/s) are mentioned in the next pāda.%
  }
  From Him they originated,%
  \footnote{%The verb \emph{āstām} may be read in several ways: as the third person dual, imperfect of the root \emph{as-} (2nd class) in the sense ``the two (\karman\/s and their ripening) originated from/because of Him''; if we read the phrase as \emph{yataḥ} (ablative) + \emph{ā} + \emph{astām}, even the interpretation, ``may the two exist up to Him (viz., function except for Him),'' is possible.  %
	The part \emph{[karmavipākau] yata āstām} may be interpreted in two ways.  One way is to understand it to mean that the teachings on \karman\/s and their results, i.e., the \Veda s, were given by \Isvara.  That \karman\/s and their ripening, in the context of \dharma, or in the sense of Vedic rituals and their fruition, were taught by \Isvara\ is expressed a few times in this text.  See syllogisms following the one starting with \emph{anekakartṛ\/°} \reftoeditionpar{p25_39} translated on p.~\pageref{VRNASRM25}.  Also, in one interpretation of the part of \rYBh{1.24}{1.24} where it asserts that \Isvara\ is the source (\nimitta) of \sastra\ and that \sastra\ is the source of \Isvara\ as well (starting with the paragraph \emph{atro\-cyate} \reftoeditionpar{p24_13} translated on p.~\pageref{SSTRISVR24}), the author interprets the word \nimitta\ to mean the cause.  The \sastra\ in this case most likely includes the Vedas.  Another indication that the author had the Vedas in mind in the expression \emph{[karmavipākau] yata āstām} is his marked hostility towards the \Mimamsaka s. \Sabara's interpretation of the word \atha\ is examined shortly (see pp.~\pageref{SBRATH}\,ff.);  he alludes to \Kumarila\ many times as an opponent in the \Isvara\ section. \Kumarila\ champions the view that the Vedas have no personal author and strongly skeptical of the creator god.  Our author even cites and criticizes \Kumarila's view on \sphota\ in the commentary on \rYS{3.17}{3.17}. It is natural for such an author to mention the relationship between \Isvara\ and the Vedas in the opening stanza.%

	Another likely meaning of \emph{[karmavipākau] yata āstām} is that the rules regarding \karman\/s (in the sense of deeds) and their ripening were set by \Isvara\ and he oversees their operation.  This old idea is mentioned several times in the \BhGBh\ in the form of \emph{sarvasya karma\-phala\-jātasya dhātāraṃ vidhātāraṃ} (\rBhGBh{08.9}{8.9}), \emph{dhātā karma\-phalasya prāṇibhyo vidhātā} (\rBhGBh{09.17}{9.17}), \emph{dhātāhaṃ karma\-phalasya vidhātā} (\rBhGBh{10.33}{10.33}), etc.  Since the author shows affiliation to the \BhG\ in the third pāda of this stanza by using the verb \kal\ in relation to \kala, it is conceivable that he had this old idea in mind.  Also, one area of \Kumarila's criticism of \Isvara\ involves the relationship between the law of \karman\/s (or \dharma) and \Isvara\ (\SVSAP{043b, 50, 69--72, 75}{kk.\ 43b, 50, 69--72, 75}).  One common theme that appears in \SVSAP{072 and 75}{kk.\ 72 and 75} is that the involvement of \Isvara\ in the workings of \karman\/s and their fruits is redundant.  It is also conceivable that the author wanted to give an answer to this problem.%
	}\\
  Impurities (\klesa\/s) are not enough [to affect] him; but not surmountable for everybody [else], \\%
  He cannot be fixed by the compelling eye of \inidx{Kāla},%
  \footnote{The readings in two primary manuscripts differ. \Tm\ reads \emph{kāladiśā} and L reads \emph{kāladṛśā}.  Both seem possible.  Even the singular for the dual \DD\ compound \emph{kāladiś} may still be acceptable because its etymologically related synonym \emph{deśakāla} is allowed to be used as singular.  However, I slightly favor the reading \emph{kāladṛśā} because it is less likely for it to change to \emph{kāladiśā} (more familiar sounding and graphically simpler) than the other way around.  This part refers to \rYS{1.26}{1.26}: \emph{sa pūrveṣām api guruḥ kālenānavacchedāt}.%
  } \\
  I salute Him, the Lord of the world, the enemy of \Kaitabha.%
  \footnote{\label{VSNSM1}This stanza pays homage to \Visnu\ (the enemy, i.e., the slayer, of \Kaitabha).  On our author's affiliation to the \Vaisnavism, see pp.\ \pageref{NRYNN25}, \pageref{NNRYN26}, \pageref{VISNU27}, etc.  Despite the common view that associates \Sankara\ with \Saivism, it has been pointed out that \Sankara\ had a religious background closer to \Vaisnavism.  See \citet{hacker:relation}.  Cf.\ \citet[4]{mayeda:thousand}.}%  On the other hand, however, we may practice some caution assuming the genuinness of the verse since we have evidence not all the closing verses are genuine.  See p.~\pageref{CLSVRS} in the introduction.}
  
  He\refm{3}%\pratika{yas sarvavit}
\pratikap{p01_2} is omniscient, ruling all, and omnipotent,\footnote{Various schools or authors attribute certain sets of characteristics to \Isvara. In our author's case, it is typically three: he often refers to \Isvara\ as all-knowing (\sarvajna), having all the power or capable of everything (\sarvasakti), and the \inidx{lord} of all (\sarvesvara) or the \inidx{supreme lord} (\paramesvara).  See for example, p.~\pageref{TTHSKT25} where the author constructs two \anumana\/s (inferences) that assert that \Isvara\ is \sarvasakti\ and \paramesvara, taking the one on \jnana\ in the \YBh\ as a model.}\\
  {}[He is] free from defects, [but] actions and rewards are superimposed on Him.\footnote{\label{OSTWO}The question whether \Isvara\ is capable of actions is discussed in conjunction with the question whether \Isvara\ has a body or not.  First the opponent asks the question in two paragraphs starting with \emph{tathāsarvajño} \reftoeditionpar{p25_58} (trl.\ p.~\pageref{TTHASRV25}), and the answer is given starting with the paragraph \emph{athāśarīratvād asarvajño} \reftoeditionpar{p25_130} translated on p.~\pageref{ATHSRRTVS}. There again the intended opponent is \Kumarila.

In this stanza the use of the word \emph{upahita}, deriving from the root \emph{upa-dhā-}, suggests the author's view that \kriya s and \phala s attributed to him are not real.  There is some ambiguity in the \phala\ part of this compound.  When the word appears in conjunction with the word \kriya\ (action), the most natural way to understand it is that it is the effect of an action.  Even though that may be figurative, it is hard to conceive that \Isvara\ receives a result of his action.  Another possible interpretation, adopted here, is that the \phala\ is the reward for actions of others.  The difficulty with this interpretation is that such concept is not mentioned in the body of the text.  Another problem with this compound, and in an effect, with the whole stanza is that the mention of defects, actions, and rewards with regard to \Isvara\ is redundant; it has already been mentioned in the previous stanza.  Note that this compound is not without textual problems.  The primary manuscripts, \Tm\ and L, both read \emph{°phalam} at the end of this line.  This reading is certainly not acceptable, and therefore a simple emendation, first introduced by the copyist of M, \emph{phalaḥ} is adopted here.  There remains some possibility that the reading \emph{phalam} was a result of further corruption.}%\\
\\  The Lord who is the cause of creation, end, and maintenance of everything,\\
  Reverence be to Him, the teacher even of the teacher.
\end{verse}\label{OSTZ2E}
%\section{First {Pāda}}

\section{Commentary on Pātañjalayogaśāstra 1.1}
\begin{quotation}
\st{Now\refm{4}%\pratika{atha}
\pratikap{p01_3} the instruction of yoga.~(\rYS{1.01}{1.1})}
\end{quotation}

\subsubsection{Explanation of the \sutra}

% \paragraph{On the word \emph{atha}}
The%\pratika{athetyādi}
\pratikap{p01_4}\labelh{ATHAATHATRL} \PYS\ starts with [the word] ``now (\atha{})'';  its explanation (\vivarana)\labelh{VIVPLNT} is commenced.\footnote{\label{STRMTSY}I have emended \emph{°śāstravivaraṇam} found in the manuscripts to \emph{°śāstram; tasya vi\-va\-raṇam}, which would be \emph{śāstrantasya vivaraṇam} without punctuation.  Considering that \emph{nta} and \emph{sya} look similar in the Malayalam script, this is a small emendation.  The adjective \emph{anā\-khyā\-ta\-samba\-ndha\-pra\-yo\-janam} in the next sentence refers to the noun \emph{°śāstram} (a neuter noun), rather than \emph{vivaraṇam}; this is indicated by the subsequent discussion on the relationship and the objective of a \emph{śāstra}.  The use of the adjective presupposes that the word modified by it has appeared earlier.  Note also that I read \emph{athetyādi} as a separate word, also modifying a neuter noun \emph{°śāstram}.

In addition, regardless whether this is \Sankara's composition or the one that imitates his style, many commentaries that are probably genuine \Sankara's, particularly those on the Upaniṣads, open up with the same formulation to the proposed reading here, viz., first a compound consisting of the first word(s) of the \emph{mūla} text and the word \emph{ādi}, followed by the title of the \emph{mūla} text, forming a phrase ``The TEXT TITLE, starting with the word(s) SUCH AND SUCH\ldots.''  That phrase is followed by the main clause, ``its commentary (\emph{vivaraṇa} or \emph{vivṛti}) is commenced (\emph{ārabhyate}).''  See the reference register in the edition (p.~\pageref{FIRSTSTC}).  That this beginning of the YVi is clearly intended to have the same basic structure (the first word of the \emph{mūla} text compounded with \emph{ādi}, followed by the title of the \emph{mūla}, and the phrase \emph{vivaraṇam ārabhyate}) is a support for my emendation.}  There (at the beginning of the explanation) if [the \sastra\ (discipline)] does not have the relationship (\sambandha) and the objective
(\prayojana) [explicitly] stated, it cannot [bring] a person to taking an action or backing off [from something].\footnote{The convention that a \sastra\
must have the objective (\prayojana), the relationship (\sambandha), and sometimes the \abhidheya\ is quite old.  Such discussions are found, for example, in \SVPRA{011--25}{kk.\ 11--25},
%\begin{quote}
%  \dnp{%
%    `अथातो धर्मजिज्ञासा' सूत्रमाद्यमिदं कृतम्। 
%    धर्माख्यं विषयं वक्तुं मीमांसायाः प्रयोजनम्॥११॥\\
%    सर्वस्यैव हि शास्त्रस्य कर्मणो वापि कस्यचित्। 
%    यावत् प्रयोजनं नोक्तं तावत् तत् केन गृह्यते॥१२॥ \\
%    मीमांसाख्या तु विद्येयं बहुविद्यान्तराश्रिता। 
%    न शुश्रूषयितुं शक्या प्रागनुक्त्वा प्रयोजनम्॥१३॥ \\
%    विद्यान्तरेषु नाप्येतद्यद्यभीष्टं प्रयोजनम्। 
%    अनर्थप्रापणं तावत् तेभ्यो नाशङ्क्यते क्वचित्॥१४॥ \\
%    मीमांसायां त्विहाज्ञाते दुर्ज्ञाते वाविवेकतः। 
%    न्यायमार्गे महान् दोष इति यत्नोपचर्यता॥१५॥ \\
%    तस्मात्प्रयोजनं पूर्वमुक्तं सूत्रकृता स्वयम्। 
%    यत् स्वेनोक्तं वदेयुस्तद्भाष्यकारादयः कथम्॥१६॥ \\
%    सिद्धार्थं ज्ञातसम्बन्धं श्रोतुं श्रोता प्रवर्तते। 
%    शास्त्रादौ तेन वक्तव्यः सम्बन्धः सप्रयोजनः॥१७॥ \\
%    शास्त्रं प्रयोजनं चैव सम्बन्धस्याश्रयावुभौ। 
%    तदुक्त्यन्तर्गतस्तस्माद्भिन्नो नोक्तः प्रयोजनात्॥१८॥ \\
%    सिद्धिः श्रोतृप्रवृत्तीनां सम्बन्धकथनाद्यतः। 
%    तस्मात्सर्वेषु शास्त्रेषु सम्बन्धः पूर्वमुच्यते॥१९॥ \\
%    यावत्प्रयोजनेनास्य सम्बन्धो नाभिधीयते। 
%    असम्बद्धप्रलापित्वाद्भवेत्तावदसंगतिः॥२०॥ \\
%    इह त्वाक्षिप्य सम्बन्धं भाष्य एवाभिधास्यते। 
%    धर्मप्रसिद्ध्यसिद्धिभ्यां तस्मान्नान्यो ऽभिधीयते॥२१॥ \\
%    न चाप्यत्राथशब्देन शास्त्रसम्बन्ध उच्यते। 
%    सम्बन्धक्रिययोर्ह्येष ब्रूते शास्त्राच्च ते पृथक्॥२२॥ \\
%    यो ऽप्ययं शास्त्रसम्बन्धो वर्ण्यते कैश्चिदादितः। 
%    क्रियान्तर्यरूपो वा गुरुपर्वक्रमो ऽपि वा॥२३॥ \\
%    तदतद्भावयोस्तस्य विशेषो नोपलभ्यते। 
%    श्रोतुर्विधौ निषेधे वा ज्ञाने वा शास्त्रगोचरे॥२४॥ \\
%    तस्माद्व्याख्याङ्गमिच्छद्भिः सहेतुः सप्रयोजनः। 
%    शास्त्रावतारसम्बन्धो वाच्यो नान्यस्तु निष्फलः॥२५॥  }
%\end{quote}
 \rnnNBh{1.1.1}/\rNV{1.1.1}.  The tradition may originally have been inspired by the opening discussion in the \VMBh\ although it discusses \emph{prayojana}, \abhidheya\ and \sambandha\ in a quite different context. \Sankara\ also held that a \sastra\ must have \sambandha\ and \prayojana.  For example, see the
introduction of the \rnnBhGBh{00intro}{intro.}, \rMaUBh{1.01}{1.1}, \rKaUBh{1.1.01}{1.1.1}, and \rMuUBh{1.1.01}{1.1.1}.  On the discussions of \emph{prayojana\/} in various Buddhist texts, see \citet{funayama:prayojana}.

\Sankara, like our author here, presupposes the view that the effect of a \sastra\ are \pravrtti\ (taking an action) and \nivrtti\ (backing off from something).  He, however, is critical of the view along with the related idea that a \sastra\ leads a person to \pravrtti\ and \nivrtti\ by means of \vidhi\ and \pratisedha.  See for example, \rBSBh{1.1.04}{1.1.4} \citep[particularly pp.~111\,ff., 130\,ff., 137\,ff.]{bsbh}.  There, \Sankara\ emphasizes that \sastra s do not have to have them as purpose; telling the truth without leading to any action can be the goal of \sastra s, especially the \Upanisad s.  In \rBSBh{1.1.01}{1.1.1} \citep[41--44]{bsbh} a somewhat related topic that \pramana\ leads to \pravrtti\ and \nivrtti\ is examined.  According to \Sankara, if the effect of \pramana\ is either \pravrtti\ or \nivrtti, human beings are no better than animals.%
% in the last reagrd, it might be helpful to take a close look at the Hetubindu and the Brahmasiddhi. Particularly the latter seems to offer a relevant discussion.
% % must have an objective and a relationship to human activities seems
% % to have been established as early as the time of NV. See NBh/NV
% % 1.1.1, {śV Pratijñāsūtra} kk.~11--25.  Also cf.\ the beginning
% % of the MBh, where \emph{proyojana}s are discussed.  On the
% % discussions of prayojana in various Buddhist texts, see
% % \cite{funayama:prayojana}.
}
Therefore the relationship and the
objective [of the \sastra], which are intended by the \Sutrakara, and which are the cause of commencement and termination of human [activities], are first discussed.

\paragraph{Objective}
Of\refm{6}%\pratika{tatra}
\pratikap{p01_5} those two, first, the \inidx{objective}---\labelh{PRYJN}

[It]%\pratika{cikitsāśāstravat}
\pratikap{p01_6} is explained [by the \Bhasyakara] by means of showing that it (the \sastra)
consists of \inidx{four divisions}, as in the \sastra\/ of healing (\cikitsasastra). %\footnote{On this discussion of
%  quadruple division see \cite{wezler:yvi2}.  Although Wezler ascribes
%  the beginning of the quadruple division to Buddhist, there are
%  evidences that suggest the existence of a text called \emph{%
%    Cikitsāśāstra}.  For example, a fragment (\emph{chāyā
%    madhuraśītalā}) quoted in YVi (p.~70, l.~5 of printed edition)
%  as from the \emph{Cikitsāśāstra} is also found elsewhere (in
%  \emph{Nyāyamañjarī}?  as a half verse and in \emph{%
%    Dvādaśāranayacakra}?  as a complete verse).  On the similar
%  discussion on \emph{heya}, etc., cf.\ NBh/NV 1.1.1.} 
Namely---the \sastra\/ of healing consists of four divisions: the disease,
the cause of the disease, freedom from disease, and the treatment; the goal [of this \sastra] is to explain the domains of the \inidx{four divisions} by means of a limited [set of] \inidx{injunctions} and \inidx{prohibitions}.\footnote{The phrase \emph{cikitsāśāstraṃ caturvyūham, rogo rogahetur ārogyam bhaiṣajyam iti} \index{caturvyuha@\emph{caturvyūha}} (including \emph{iti}) is found in \rYBh{2.15}{2.15}.  The editors of the 1952 edition suggested to emend the text by inserting \emph{tac ca} in between \emph{°dvāreṇa} and \emph{caturvyūhaviṣaya°}, making the sentence into two.  The text certainly reads more smoothly with the change, but it is still an intelligible sentence without the change; the text up to \emph{bhaiṣajyam iti} forms the subject and the rest predicate.  Even though the part \emph{cikitsāśāstram \ldots\ iti} is taken from the \bhasya, the author's intention was not to quote from the \bhasya, but use that phrase as part of his sentence.  The word \emph{iti} at the end of the phrase should be seen as terminating a list of four divisions.

The expression \emph{vidhi\-prati\-ṣe\-dha\-niyama\-dvāreṇa} is a result of an emendation proposed in the 1952 edition and adopted here.  It is a reasonable emendation, given a similar expression \emph{pra\-vṛtti\-ni\-vṛtti\-ni\-yama\-dvā\-re\-ṇa} later in the paragraph starting with \emph{yogānuśāsanam iti} \reftoeditionpar{p01_27} (translated on p.~\pageref{YGNSSM}).  It is a curious expression.  Literally the expression should mean ``by restricting\ldots.''}  In the same manner, here also, it is shown that the \sastra\ [of yoga] consists of four divisions [in the section] beginning with [the {sūtra}], ``For the one who knows the difference (\emph{vivekin}), everything is suffering, because of the sufferings from change, torment, and residues [of past experiences], and because of the conflict among the activities of the [three] \guna s (\rYS{2.15}{2.15}).''  That is:\refm{7} the object to-be-removed (\heya) is reincarnation (\samsara) full of suffering;\footnote{See \rYS{2.16}{2.16}: \emph{heyaṃ duḥkham anāgatam}.} its cause (\emph{hetu}) is the union of the seer and the to-be-seen (\emph{draṣṭṛdṛśyasaṃyoga}) that originates from \ignorance\ (\avidya);\footnote{See \rYS{2.17}{2.17}: \emph{draṣṭṛdṛśyayoḥ saṃyogo heyahetuḥ}, and \rYS{2.24}{2.24}: \emph{tasya hetur avidyā}.} the means of removal (\emph{hānopāya}) is the unwavering discriminative knowledge (\vivekakhyati);\footnote{See \rYS{2.26}{2.26}: \emph{viveka\-khyāti\-r a\-vi\-pla\-vā hāno\-pāyaḥ}.} and when the discriminative knowledge is present, \ignorance\ comes to an end. When it (ignorance) comes to an end, the union between the seer and the to-be-seen ends. This is the removal (\hana).  That verily is emancipation (\kaivalya).\footnote{See \rYS{2.25}{2.25}:
  \emph{tadabhāvāt saṃyogābhāvo hānaṃ tat dṛśeḥ
    kaivalyam}.  The expression in our text: \emph{duḥkha\-pra\-curaḥ saṃsāro heyaḥ; tasyā\-vidyā\-nimitto draṣṭṛ\-dṛśya\-saṃyogo hetuḥ; vi\-ve\-ka\-khyā\-ti\-r a\-vi\-pla\-vā hā\-no\-pā\-yaḥ; vi\-ve\-ka\-khyā\-tau ca sa\-tyām a\-vi\-dyā\-ni\-vṛ\-ttiḥ, ta\-nni\-vṛ\-ttā\-v ā\-tya\-nti\-ko dra\-ṣṭṛ\-dṛśya\-saṃ\-yogo\-pa\-ra\-mo hā\-nam} is a close parallel to that in \rYBh{2.15}{2.15}: \emph{tatra duḥ\-kha\-ba\-hu\-laḥ saṃ\-sā\-ro
    he\-yaḥ\dd\ pra\-dhā\-na\-pu\-ru\-ṣa\-yoḥ saṃ\-yo\-go he\-ya\-he\-tuḥ\dd\
    saṃ\-yo\-ga\-syā\-tya\-nti\-kī ni\-vṛ\-tti\-r hā\-na\-m\dd\ hā\-no\-pā\-yaḥ
    sa\-mya\-gda\-rśa\-nam\dd}.  Notable differences are the mention of \avidya\ as the cause of \samyoga\ and the replacement of \samyagdarsana\ with \vivekakhyati.  The first difference is a result of including the teaching of \rYS{2.24}{2.24} (\emph{tasya hetur avidyā}), and the latter incorporates the expression from \rYS{2.26}{2.26}.  The effect of these differences is that emphasis is put on the view that \avidya\ is the ultimate culprit.} As [this \sastra] aims at emancipation
which is comparable to freedom from disease [of the \sastra\ of healing], that very emancipation is its objective.

[Objection]\refm{8}%\pratika{nanu ca}
\pratikap{p01_7} The to-be-removed and its cause should not be
placed at the beginning, because [doing so is] pointless.  Because the
objective of the \sastra\/ is the removal.  [Thus] discriminative
knowledge, which is the means to that [aim (the removal)], only should be
discussed.  For, nobody enjoins the pain and its cause to someone
whose foot sole is pierced by a thorn, forgetting about the removal of
that [thorn].\footnote{Pain from a thorn is also found in \rNV{1.1.1} as an
  example of suffering (\emph{duḥkha}).  See the reference register accompanying this part of the text on p.~\pageref{NVKANTHA} for the text of the \NV\@.}

[Answer]\refm{9}%\pratika{naitad}
\pratikap{p01_8} That is not correct.  Because the means of removal
presupposes the to-be-removed and its cause.  It is not possible to
discuss discriminative knowledge, which is the opposite of 
\ignorance, until it is stated that this cycle of reincarnation is what 
should be removed by this means of removal, and that its cause is the
union---itself caused by \ignorance---between the seer and the
to-be-seen.  Because we also\footnote{The presence of the word \emph{api} here is conjectural. Seeing various corrections both in \Tm\ and L, I infer that there was a confusing correction already in the archetypal manuscript. The exact situation in the archetype is hard to reconstruct but I consider the presence of \emph{vi} in \Tm\ significant. In the Malayalam script \emph{p} and \emph{v} are often hard to distinguish. The presence of the word \emph{api} fits in the context. The subject matter of this phrasal clause expressing a reason is not identical to the preceding. The reason rather functions by pointing out a real word example. In that sense, qualifying the subject with ``also'' is appropriate.} experience that a patient, who has something
to-be-removed and its cause, and who seeks for the means of removal
[of the disease] goes for treatment.  For, the \sastra\/ of healing is not taught without presupposing disease and the cause of disease.

\paragraph{Relationship}
The\labelh{SMBDHPI}%\pratika{sambandho 'pi}
\pratikap{p01_9} relationship, too [is properly presented in the \sastra]---

The%\pratika{vivekakhyāter}
\pratikap{p01_10} result (\emph{phala}) to be achieved (\emph{sādhya}) by discriminative knowledge is but the removal; also, the means (\emph{sādhana}) of removal is but discriminative knowledge.  Precisely this mutual determination (\emph{itaretaraniyama}) of the aim (\emph{sādhya}) and the means (\emph{sādhana}) [is the relationship]; there is no other relationship.  This is analogous to the mutual determination that the result of the treatment is but the freedom of disease and that the means toward freedom of disease is but the treatment.  %Similarly, the result of the treatment is but the freedom of disease, and the means toward freedom of disease is but the treatment.  Therefore they restrict each other. 
 And this [mutual determination] is based on the \sastra.\footnote{The word \emph{iti} at the end of this paragraph marks the end of the discussion on \prayojana\ and \sambandha\ of this \sastra.  The use of \emph{iti}, followed by a concluding remark that starts with the word \emph{tasmāt} is found frequently in this text.  See another example on p.~\pageref{TSMDTH}.}

In%\pratika{tasmāt}
\pratikap{p01_11} conclusion,\refm{11} the instruction of yoga (\yoganusasana) [that is this \sastra] has proper relationship and objective.\labelh{PSCOMP}

\paragraph{Justifying the first \sutra}\labelh{JSTF}
[Objection]\refm{12}%\pratika{nanu ca}
\pratikap{p01_12}\footnote{Our author starts discussions to justify why the reading \emph{atha yogānuśāsanam} is the first \sutra, raising several other alternatives for the \sutra.  I doubt that there were real criticisms on the wording of the first \sutra\ of the \PYS.  The following arguments rather appear to be purely an intellectual exercise to write a commentary that guards the text being commented.} If the removal is the
objective, and its means is discriminative knowledge, then [the first \sutra] should say ``Now the instruction
of discriminative knowledge (\emph{atha
  vivekakhyātyanuśāsanam}).''  For what purpose does the [first]
\sutra\ say ``Now the instruction of yoga (\emph{atha
  yogānuśāsanam})''?

[Answer]%\pratika{tadupāyatvād}
\pratikap{p01_13} Because\refm{13} yoga is the means of
obtaining it (\vivekakhyati).  Precisely the means should be
stated.  For, if the goal (to be obtained, \emph{upeya}) is [first] explained, then since the
goal presupposes the means (to obtain it), it will furthermore be
necessary to speak of the means along with its divisions and
in its entirety.  On the contrary, when it (the means, yoga) is denoted,
then everything is denoted.

[Objection]\refm{14}%\pratika{kathan tad°}
\pratikap{p01_14} How is [yoga] the means to obtain that (\emph{vivekakhyāti})?

[Answer]%\pratika{yata aaha}
\pratikap{p01_15} Because\refm{15} the \Sutrakara\ says: ``When
impurities are removed as the result of exercising the divisions of
yoga, the intelligence begins to shine until [it becomes]
discriminative knowledge (\rYS{2.28}{2.28})''; ``When skilled in \emph{%
  nirvicāra[samādhi]}, there will be internal tranquility (\rYS{1.47}{1.47})'';
and ``Then there will be an intellect called carrier of the truth (\emph{%
  ṛtambharā}) (\rYS{1.48}{1.48}).''  Also the \Bhasyakara\ says:
``That (\emph{samādhi}), % needs to be reconsidered
when the mind is focused, illuminates the real object \ldots\ (\rYBh{1.01}{1.1}).''  And discriminative knowledge is the understanding
of the real object.  Accordingly, it is appropriate that the instruction of yoga---the means to obtain it (\emph{vivekakhyāti})---has been put in a sūtra at the beginning.

[Objection]\refm{16}%\pratika{nanu ca}
\pratikap{p01_16} Is it not that discriminative
knowledge arises after exercising the divisions of yoga? Then [the first
\sutra] should state ``Now the instruction of the divisions of
yoga (\emph{atha yogāṅgānuśāsanam}).''\footnote{The \PYS\ in fact teaches the eight divisions of yoga (\emph{yogāṅga}s), starting from 2.28.}

[Answer] No.\refm{16+1}%\pratika{na}
\pratikap{p01_17} Because [a \sastra\/
should have] a beginning with the result.\footnote{The answer in this paragraph, which may be paraphrased as ``it is appropriate to start a \sastra\ by stating the result'' may seem contradictory to what has just been argued: the means to obtain a goal, rather than the goal should be taught.  As can be observed in the paragraphs that follow, our author distinguishes a goal (\emph{upeya}) and the result (\emph{phala}).  That seems to be his strategy to avoid such a contradiction.  Still, the arguments appear weak.  This is one reason to suspect that this part was an intellectual exercise to guard the text being commented.}  For, since yoga is the result of
exercising divisions of yoga, it is appropriate to begin with that.

[Objection] If%\pratika{yady evaṃ}
\pratikap{p01_18} so,\refm{17} the beginning should be the removal.

[Answer] No.%\pratika{na}
\pratikap{p01_19}\refm{17+1} Because that is a goal.  Namely, it is nothing but
a goal.  On the other hand, yoga is the goal (\emph{upeya}) for its own divisions, and it is also the means to obtain (\emph{upāya}) discriminative knowledge and [supernatural powers, such as] the ability to become subtle (\emph{aṇimā}),\footnote{In the \PYS\ \rYS{3.44}{3.44} and accompanying \bhasya\ mention the well-known supernatural powers \emph{aṇimā}, \emph{mahimā}, etc.} and so on.  Therefore it is more valuable, because it is the cause of two fruits.\footnote{The reading \emph{phaladvayanimittatvāt} is conjectural.  The readings in two primary manu\-scripts (\Tm\ and L) are unintelligible.  It is not even certain what the scribe of L meant to write because there the letters \emph{pa} and \emph{va} are often indistinguishable.  Of the three possible readings found in the primary manuscripts (\emph{phalavayanimitta°} or \emph{phapapayanimitta°} for L and \emph{phalavannimitta°} for \Tm), \emph{phalavayanimitta°} requires the least change to produce a reasonable reading, \emph{phaladvayanimitta°}, which I have adopted.  We can be reasonably certain that the corrupt reading \emph{phalavayanimitta°} was already present in the common ancestor (possibly the direct exemplar of both \Tm\ and L).  The reading found in \Tm\ was probably a result of the scribe's attempt to produce a meaningful reading.}

[It is\pratikap{p01_20} not\refm{18} appropriate to have the
\sutra, ``Now the instruction of the divisions of yoga,'' at the
beginning]%\pratika{uttarasūtrārtha°}
 also because [the first \sutra] introduces the next \sutra.
For, if the first \sutra\ was to say ``Now the instruction of the
divisions of yoga,'' then it would not be appropriate to say
``[It (\emph{asamprajñāta-samādhi}) is] yoga [and it is] the termination of the activities of the mind (\rYS{1.02}{1.2}).''\footnote{\label{YS12first}The translation here follows the interpretation of the \sutra\ in the YVi where the author states that the \sutra\ is the definition of the \emph{asamprajñāta-samādhi}.  This is somewhat counter-intuitive, but our author is clear about \sutra\ 1.2 being the definition of \emph{asamprajñāta-samādhi}.  See, for example, the last two paragraphs of the commentary of the current \sutra.  See also note \ref{no_yogah}.}
In that case, the [next] \sutra\ should read ``\emph{yama, niyama},''
and so on (\rYS{2.29}{2.29}).

If%\pratika{atha tad eva}
\pratikap{p01_21} you say---[Objection]\refm{19} But then that is exactly how [the \sutra] should read---[Answer] No.  For the next \sutra, ``[It (\emph{asamprajñāta-samādhi}) is] yoga [and it is] the termination of the activities of the mind,'' has the purpose of showing the relationship
between discriminative knowledge---whose means is yoga,\footnote{The editors of the 1952 edition suggested the reading \emph{yogopeyāyā(ḥ)} instead of \emph{yogopāyāyā(ḥ)}, the reading attested in all the manuscripts.  It is understandable that they preferred a \TP\ compound rather than a \BV\ compound, but the following word, another adjective to modify the word \emph{vivekakhyāteḥ}, reads \emph{hānopāyābhūtāyā(ḥ)}.  The author made an extra effort to makes it clear that \emph{hānopāya} is to be read as a \TP\ compound.  I take it that the author was illustrating the position of \emph{vivekakhyāti} using two compounds that both have the word \emph{upāya} as a member.  Even without the use of \emph{°bhūta}, in the second compound the case ending of the word \emph{yogopāyāyāḥ} makes it clear that it is a \BV\ compound.}
and which is the means to obtain the removal---and the removal which is
the result.

[Objection]\refm{20}%\pratika{katham punar}
\pratikap{p01_22} How, then, can it
have the purpose of showing the relationship while it states the
definition (\emph{lakṣaṇa}) of the seedless immersion (\emph{%
  nirbīja\-samādhi})?\footnote{Our author considers \rYS{1.02}{1.2} to be the definition of \emph{nirbīja-samādhi}/\emph{asamprajñāta-sa\-mā\-dhi}.  See note~\ref{no_yogah}. Cf.\ also note \ref{YS12first}.}
  % SUPPLY THE REFERENCE!

[Answer]\refm{21}%\pratika{satyam evam}
\pratikap{p01_23} That is correct.  Even so, the \sutra, which
is [in fact] only setting forth the relationship between the removal and its means,
amounts to the definition of seedless immersion (\emph{%
  nirbīja\-samādhi}).  For, the removal is not different from the
immersion which is characterized as suppression  (\nirodha).
There still is the following distinction: in the case of immersion
characterized as suppression, there will be further
commencement (\emph{pravṛtti}); on the other hand, in the case of removal, cessation (\emph{nivṛtti}) is permanent.  But in the state of immersion [characterized as suppression] there is no distinction from the removal at all.\footnote{This paragraph offers a glimpse of our author's soteriological position, which is more clearly stated in two paragraphs below. While defending the arrangement of the \PYS, he implies that the best yoga can offer, \nirbijasamadhi, is still inferior to removal (\hana).  Discriminative knowledge is, as repeatedly expressed, is the means to achieve the final goal.  What one might consider as the highest goal in the system of the \PYS, \nirbijasamadhi, is just a temporal experience of real salvation, according to our author.}

Namely,\refm{22}%\pratika{tathā cāha}
\pratikap{p01_24} [the \Sutrakara] says [in the immediately following \sutra\ after 1.2], ``At that time the seer stays in his original state (\rYS{1.03}{1.3}),''  [using the same expression as in the last \sutra] ``Or\footnote{I have emended the particle \emph{ca} found in all the manuscripts and in the 1952 edition to \emph{vā}, following the reading of the \sutra.  There is no indication in the YVi manuscripts that our author had the reading \emph{ca} for \sutra\ 4.34 when it is quoted and explained.  The exchange may have been introduced by a scribe who was unaware that the particle \emph{vā} was part of the quotation but thought that there should be \emph{ca} to connect two sentences.

The editors of the 1952 edition proposes to insert \emph{iti} after the first quote (from \rYS{1.03}{1.3}).  I do not find it necessary.} the power to cognize (\citisakti)\footnote{In the \PYS, \purusa\ is identified as \citisakti.  In the YS the word \citisakti\ appears only in the last \sutra\ 4.34, but it already appears in the \bhasya\ on \sutra\ 1.2.  The author of the \bhasya\ continues to paraphrase \drastr\ in \rYS{1.03}{1.3} with \citisakti.  The author of the YVi glosses the word as a \KD\ compound (\emph{citir eva śaktiś citiśaktiḥ}) \citep[11,26]{yvi}.  Also, he understands the part \emph{citi} denotes the verbal root \emph{cit-} and hence the \sakti\ to be the one expressed by the verbal root.  This observation is based on the sentence \emph{yathā pacyādayaḥ śaktayaḥ śaktimadapekṣāprāpaṇīyajanmāno na svayaṃ śaktayaḥ} that immediately follows the compound analysis above.  That \purusa\ is the root \emph{cit-} (following the understanding of the author of the YVi) is repeatedly expressed in the \YBh\ and in the YVi (first expressed in the \rYBh{1.02}{1.2} \emph{citiśaktir apariṇāminī\ldots} on which the YVi's compound analysis above is given).  The view implicit in such an interpretation is that \purusa\ is changeless because it is a verbal root \cit- and because the capacity (\sakti) of a verbal root is consistent.  Similarly \rYS{2.20}{2.20} has \emph{draṣṭā dṛśimātraḥ}, again identifying \purusa\ with a verbal root \emph{dṛś-}.  There it is more explicit that the root \emph{dṛś-} exists without a substance or a subject.  This is a solution to the problem that anything denoted by a substantive is subject to change.  This makes us think of \inidx{grammarian}s' influence on the formation of the \PYS.  It should be reminded that the beginning of the \PYS\ imitates that of the \VMBh\ and that the \PYS\ accepts the \sphota\ theory.  Note also that identification of \purusa\ (masculine) with \citisakti\ (feminine) is unique to the \PYS, absent in the \Samkhya\ tradition.  Another interesting point is that both in the YBh and the YVi the root \cit- is seen as special in that it does not require the locus. This is expressed in the YVi on \rYS{1.02}{1.2} \emph{citiḥ punaḥ svayam eva śaktir nārthāntaram apekṣate}, following the sentence \ldots\ \emph{svayaṃ śaktayaḥ} above.  The view that \purusa\ is nothing but the root \emph{cit-} is presupposed in the discussion in \rYBh{1.03}{1.3} starting with \emph{yadi citir eva puruṣaḥ\ldots}.  Such discussions are echoed in \Sankara's discussion on the meaning of the root \emph{upalabh-} in the \Upad\ 2.2.76i\,ff.\ where \Sankara\ argues that \atman/\emph{upalabdhṛ} is eternally only \emph{upalabdhi} (\emph{nityopalabdhimātra eva hy upalabdhā}), and that \atman\ does not change.} being established in its original state (\rYS{4.34}{4.34})'' is emancipation.\footnote{Although the last word of this sentence (\emph{kaivalyam}) is not enclosed in the quotes, it is necessary to complete the sentence meant to be referred. The whole \sutra\ 4.34 reads: \emph{puruṣārthaśūnyānāṃ guṇānāṃ pratiprasavaḥ kaivalyaṃ svarūpapratiṣṭhā vā citiśaktir iti}.  The word \emph{kaivalyam} is implied to be supplied in the second statement.  It is important for our author to have the word \kaivalya\ in this sentence to illustrate that similar expressions are used in the YS to describe the state in \asamprajnata\samadhi\ and \kaivalya.}
For, emancipation, [characterized as \citisakti's] being established in its original state, is the removal (\hana).\footnote{Two emendations are involved in this sentence.  The first is the reading \emph{sva\-rūpa\-prati\-ṣṭha\-tvaṃ} adopted from the 1952 edition instead of \emph{°prati\-ṣṭhaṃ} in the manuscripts.  The former reading is preferable because what we want here is a substantive rather than an adjective.  The other emendation is from \emph{kaivalyam āha} to \emph{kaivalyaṃ hānam}.  Considering that this paragraph elaborates the point at the end of the last paragraph (``in the state of immersion (\samadhi) there is no distinction to the ultimate removal (\hana),'') there should be a reference to \hana.  The reading, found in all the witnesses, is non-sensical, but the presence of the syllable \emph{ha} in that reading suggests that it is a corruption involving the word \hana.  Our author is pointing out the state in \asamprajnata\samadhi\ (defined in \rYS{1.02}{1.2}) is described in the same way as \kaivalya\ is described in the YS and therefore that they are not different.}  Therefore, in order to strengthen the
understanding of the purpose of the \sastra, the emancipation is revealed by [defining] the seedless immersion (\emph{nirbījasamādhi}) [in the next \sutra].

Accordingly,\refm{22+1}%\pratika{tasmād yad}
\pratikap{p01_25} it is not true that seedless immersion
(\emph{nirbīja\-samādhi}) is the means to establish
emancipation, as some say.\footnote{Cf.\ \Vacaspati\ (\Tattvavaisaradi) on \rYBh{1.01}{1.1} \emph{%
    niḥśreyasasya hetuḥ samādhir iti hi
    śruti\-smṛtī\-tihāsa\-purāṇe\-ṣu pra\-siddham\dd} \citep[2,4--5]{ys:anss}.  However, such a view that \nirbijasamadhi\ brings about the emancipation is a straight-forward interpretation of the beginning of the \PYS.  Our author's interpretation that (discriminative) knowledge is the means to attain emancipation is rather forced and particular.  This might be comparable to similar insistence on knowledge in \Sankara's works, such as the \BSBh, \BhGBh, etc.}  Instead, only knowledge is the means, through
cessation of \ignorance.  Because the state of being bound [to the cycle of rebirth] is caused by ignorance.
Therefore even if ``the instruction of yoga'' is placed in the [first]
\sutra, since the purpose of yoga is knowledge, [the \Sutrakara]
shows the relationship between knowledge and removal, too, with this
[\sutra].  Therefore [the \sutra\ is] appropriately composed.  Also, in order [to express] the relationship with the next \sutra, the [first] \sutra\ should contain [the words] ``the instruction of yoga.''\footnote{This paragraph concludes the discussion on the appropriateness of the first \sutra\ started in the paragraph \emph{nanu ca} \reftoeditionpar{p01_12} translated on p.~\pageref{JSTF}.  The argument in this paragraph is a little forced.  Still, one may note the author's insistence that knowledge is the means.}

\medskip
Even\refm{24}%\pratika{yady api}
\pratikap{p01_26} though there is also a teaching of the
means for those who seek a result (\emph{phala}) and have come
to expect the means to obtain the result, this (the means) is not the
objective.   Since everyone acts in view of a result, that (the
result) is the sole objective.\footnote{I have not given a literal translation of the word \emph{iti} at the end of this sentence.  The word \emph{iti} and even the whole sentence appear to be a little out of place.  Most likely the word \emph{iti} here signals the completion of a major segment of the text.  And the segment could be the preliminary discussion on the objective and the relationship of the whole \sastra.  However, the discussion on the objective and the relationship had completed (see page \pageref{PSCOMP}), and the mention of only \prayojana\ here is is not congruent with the immediately preceding discussions.  It is possible that this is a trace of a correction going wrong; this passage was originally omitted while copying but supplied in the margin with a marking in the body indicating where to insert the text; the insertion point was misunderstood by the next scribe. If this passage is indeed misplaced, the proper place would be at the end of the discussion on \prayojana, i.e., immediately before the passage starting with \emph{sambandho 'pi} \reftoeditionpar{p01_9} translated on p.~\pageref{SMBDHPI}. The presence of \emph{iti} at the end of this passage also becomes explicable if it indeed is misplaced. The word \emph{iti} would mark the end of the discussion on \prayojana.}

\subsubsection{Explanation of the \bhasya}
\label{YGNSSM} ``The%\pratika{yogānuśāsanam iti}
\pratikap{p01_27} instruction of yoga''---\refm{25} This is comparable to [the way] a student is instructed by means of 
determining specific actions (\emph{pravṛtti}) [that he should do] and withdrawal from actions (\emph{nivṛtti}) [that he should not do]. Likewise [the \sastra] is called the instruction of yoga, like the instruction to pupils  (\emph{antevāsyanuśāsana}), since it essentially resembles [the instruction to pupils] in that [the \sastra] determines a specific goal (\emph{sādhya}), means (\emph{sādhana}), and their divisions (\emph{aṅga}).\footnote{I suspect that the sentence, especially the part \emph{tathā viśiṣṭa\-sādhya\-sādhana\-tad\-aṃga\-niyama\-mātra\-sādṛśyāt}, is corrupt.  Given the somewhat similar discussion in the \TaiUBh\ (see the apparatus on p.\ \pageref{TAIUBH111REF} for the text), the discussion could have been longer.  However, I cannot offer a viable emendation.%  Note the expression \emph{°dvāreṇa} and the use of the word \emph{niyama} in the \TaiUBh.  Our author appears to be fond of these expressions in this part of the commentary.  They disappear in the rest of the work)
}  ``Instruction'' = instructing (\emph{anuśiṣṭi}).\footnote{This seemingly meaningless—particularly in translation—paraphrase is meant to reject the ambiguity that the word \emph{anuśāsana\/} with the \emph{-na\/} suffix can denote the means to instruct.  By paraphrasing \emph{anuśāsana\/} with \emph{anuśiṣṭi}, the author claims that the word \emph{anuśāsana\/} is meant to refer to the action of instruction, not the means of instruction.}
\bht{The \sastra\ is the instruction of yoga}\footnote{\label{VMBHEMU}I consider the sentence or the phrase \emph{yogānuśāsanaṃ śāstram} to be part of the \bhasya.  This is equivalent to the sentence \emph{yogānuśāsanaṃ śāstram adhikṛtaṃ veditavyam} in most transmission of the \YBh\ after the sentence \emph{athety ayam adhikārārthaḥ} (see below).  The sequence, starting from the \sutra, \emph{atha yogānuśāsanam\ddd\ athety ayam adhi\-kārā\-rthaḥ\dd\ yogānuśāsanam śāstram adhikṛtaṃ veditavyam} forms a closer parallel to the opening of the \VMBh: \emph{atha śa\-bdā\-nu\-śā\-sanam\dd\ athety ayaṃ śabdo ’dhikārārthaḥ prayujyate\dd\  śa\-bdā\-nu\-śā\-sa\-naṃ śā\-stram adhi\-kṛtaṃ ve\-di\-tavyam}.  In my reconstruction the \sutra\ and \bhasya, combined, reads \emph{atha yo\-gā\-nu\-śā\-sa\-nam\ddd\ yo\-gā\-nu\-śā\-sanam śā\-stram\dd\ athety ayam adhi\-kārā\-rthaḥ\dd}.  It is conceivable that the version of the \YBh\ our author used had the opening not exactly parallel to that of the \VMBh. It may, however, be noted that there are some uncertainties how well the text of the \VMBh\ is established in the editions. Cf.\ \citet[248--50]{cardona:1976} for the overview of the discussions on the transmission of the \VMBh. Cf.\ also \citet{witzel:1986}. Note that the reading of the quote from the \VMBh\ in the YVi below (see pp.\ \pageref{VMBH1138_E} for the edition and \pageref{VMBH1138_T} for the translation) is somewhat different, although close enough to call it a quote, from the readings found in editions of the \VMBh. Still, the parallelism is easily recognizable.  The current paragraph in our text could be easily seen to be an explanation of the word \emph{yogānuśāsanam} in the \sutra\ itself.  However, it is preferable to see that our author has started to gloss the \bhasya, and its first sentence is \emph{yogānuśāsanaṃ śāstram}.  Otherwise, the presence of the word \emph{śāstram} is inexplicable.  The sentence should be seen as \bhasya's gloss of the word \emph{yogānuśāsanam} in the \sutra.} by which or in which yoga is instructed.%\footnote{Even though it is not clear from the translation, the syntax in the YVi suggests that the author considers that in the paraphrase in the YBh (\emph{yogānuśāsanaṃ śāstram}) the word \emph{yogānuśāsanam\/} modifies the word \emph{śāstram\/} as a \emph{bahuvrīhi\/} compound.  According to this interpretation, the \Bhasyakara\ interprets the beginning of the YS as proclaiming the topic by introducing the word \sastra\ as implied in the \sutra.}  %[The compound] \emph{yogānuśāsana} means either that yoga is
%instructed by it, or in it, and it is the \sastra.

\bht{That%\pratika{athety ayam}
\pratikap{p01_28} word\refm{27} ``now'' means
  the declaration of the topic.}  %\footnote{Note that the first \sutra\ and this first
  %sentences of YBh is parallel to the beginning of MBh:
%
%\bigskip
%\begin{tabular}{p{.5\textwidth}p{.5\textwidth}}
%  \multicolumn{1}{c}{YS + YBh} & \multicolumn{1}{c}{MBh}  \\
%  \emph{atha yogā\-nu\-śā\-sa\-nam| athe\-tya\-ya\-m
%    a\-dhi\-kā\-rā\-rthaḥ| (yogā\-nuśā\-sa\-nam adhi\-kṛtaṃ
%    ve\-di\-ta\-vyam||)} This last sentence was not present in the YBh
%  commented by the author of the YVi.  See note\fnref{note:veditavyam} on
%  page~\pageref{note:veditavyam}.  & \emph{atha
%    śabdā\-nu\-śā\-sa\-nam| athe\-ty ayaṃ śa\-bdaḥ
%    adhi\-kā\-rā\-rthaḥ pra\-yu\-jya\-te|
%    śa\-bdā\-nu\-śā\-sa\-naṃ śā\-stra\-m adhi\-kṛ\-taṃ
%    ve\-di\-ta\-vyam|}    \\
%\end{tabular}
%\bigskip } 
``The declaration of the topic (\emph{adhikāra})'' = opening (\emph{%
ārambha}) = announcement (\emph{prastāva}), and ``meaning (\emph{%
artha})'' = sense (\abhidheya).  [These words are connected as
a \KD\ compound and then form a \BV\ compound.]\footnote{The sentence is a typical \emph{vigraha\/} to analyze a compound, \emph{adhikārārthaḥ} in this case.  The author of the YVi is clarifying the meanings of constituent words \emph{adhikāra\/} and \emph{artha} by listing synonyms because they both have divergent meanings.  Also, he is showing the relationship between the constituents (in the same case, viz., \KD) and the type of the compound syntactically.}  [This interpretation is]
based on the authority of a \smrti\/ (the Mahābhāṣya) by hands of a learned.\footnote{\label{PTJME}Thus the learned person is \Patanjali.  \Patanjali, the author of the \VMBh, is mentioned by name in the YVi in the commentary on \rYS{3.44}{3.44} \citep[299]{yvi}.  However, our author shows no awareness of \Patanjali\ being associated with the text he comments upon, viz., the \PYS.  Note that the beginning of the \PYS\ emulates that of the \VMBh.  See the parallel passage in the edition part.  See also note \ref{VMBHEMU}.}%For the relationship between the \emph{śiṣṭa\/}s and \emph{śāstra}, see the opening discussion of the Mahābhāṣya itself.  There Patañjali declares that ...}
  
%\footnote{Is he referring to the MBh?}

[Objection]\refm{28}%\pratika{nanu cā°}
\pratikap{p01_29}\footnote{Our author consistently starts the objection with the expression \nanuca\ as found here.  As noted by \citet[36]{wezler:yvi1}\index{Wezler, Albrecht}, this is fairly consistent in the YVi and might indeed be idiosyncratic\index{idiosyncrasy} to the author.  A notable text that uses the expression \nanuca\ frequently is the \VMBh\ that set the standard for Sanskrit commentaries.  In \Sankara's BSBh the use is significantly less (as far as we look at published editions).  Still, the Nirnaya Sagar Press edition of the BSBh has the expression in 2.3.19, 2.3.43, 2.3.47, 4.2.5, 4.2.6, and 4.2.20.  The concentration is a little intriguing.  It is nonetheless necessary to wait for a critical study of the \BSBh\ manuscripts before concluding anything from the use of this expression.  We cannot exclude the possibility that the \BSBh, being one of the most heavily and closely studied text, has gone through many generations of cleanup in style.} The learned ones teach in a \smrti\/ that the word \atha\ means ``immediately after.''  Namely,
[\Sabara] says:\label{SBRATH}
\begin{quote}
  It is experienced that [this word \emph{atha\/} at the
  beginning of the Mī\-māṃ\-sā\-sū\-tra (\JS)\footnote{\rJS{1.1.1}: \emph{athāto dharmajijñāsā}.}] has the meaning of introducing something
  that comes immediately after something has concluded.\footnote{\rSBh{1.1.1}.  The
    whole sentence runs as follows: \emph{tatra loke 'yam athaśabdo
      vṛttād a\-na\-nta\-ra\-sya prakriyārtho dṛṣṭaḥ\dd.} A similar discussion
    of the meaning of \atha{} at the beginning of \sutra\ texts
    appears also in the \rnnBSBh{1.1.01}{1.1.1} on \rBS{1.1.01}{1.1.1}: \emph{athāto brahmajijñāsā} (note the parallel structure of the \rJS{1.1.1} and the \rBS{1.1.01}{1.1.1}).  \Sankara\ holds that the meaning of \atha{} at the beginning of BS 1.1.1 is \emph{ānantarya}, just like Śabara, although he claims the intention to learn about Brahman (\emph{brahmajijñāsā}) does not necessarily come after the intention to learn about Dharma (\emph{dharmajijñāsā}).  The antagonist of \Sankara\ here holds the view that the \brahmajijnasa\ must come after \dharmajijnasa.  \Sankara\ was probably dealing with the traditional view that the two \sutra\ texts (the JS and the BS) formed a set and to be studied in succession.  (\Sankara\ considers the \JS\ and the \BS\ form one continuous text; he refers to \rJS{3.3.14}: \emph{śruti\-liṅga\-vākya\-prakaraṇa\-sthāna\-samākhyā\-nāṃ sama\-vā\-ye pāra\-daurbalyam artha\-vi\-pra\-karṣāt} as belonging to ``an earlier section (\emph{pūrvasmin kāṇḍe}).'') As \emph{dharmajijñāsā} is placed at the beginning of the JS, this interpretation of the meaning of \atha{} allows one to begin studying the \BS\ without learning the JS, or the discipline of Mīmāṃsā.  We may say that \Sankara\ and the author of the YVi are in agreement in that they both refer to the \Mimamsaka\ view of the interpretation of the word \atha\ and that they are both anti-\Mimamsaka\ when the word \atha\ is concerned.  I am unaware of any other author who is so concerned about how the \Mimamsaka s understand the meaning of the word \atha\ as \Sankara\ or the author of the YVi.  Note that the author of the YVi here ignores the more developed discussion by \Kumarila\ in his \SV\ while \Sankara\ appears to be aware of it.}
\end{quote}

[Answer]\refm{29} No.%\pratika{na\dd}
\pratikap{p01_30}  Because [the word
\atha{}] always has the meaning ``opening,'' and because the
meaning ``immediately after'' must be [secondarily] understood.  For
example---when we say ``[he is] a son,'' it is understood that [the man can be] a father.\footnote{The sentence is vague, but the idea behind this example should be the one that frequently appears in śāstric literature---that a man can be a son or a father or many other things.  In that vein, the text would read better if we replace \emph{putra} with \emph{puruṣa} and \emph{pitṛ} with \emph{putra}. (A man is always a son.) There is some suspicion that the original text read as such and the reading we have is a result of corruption. The idea is again already expressed in \rVMBh{1.1.66}{1.1.66}.  \Sankara\ mentions the idea in \rBSBh{2.2.17}{2.2.17}.  (See page \pageref{PTRPITR} for the text.)  There, \Sankara\ refers to the written symbol of the number one in the decimal notation system.  He says it is a line (\emph{rekhā}).  The passage is a paraphrase of another in \rYBh{3.13}{3.13}.  Here we might have a testimony to the writing system (a line representing the numeral 1) \Sankara\ and the author of the YBh were familiar with.}  But the meaning of ``son'' is not ``father.''  In the same
manner, here, only ``opening'' is being directly expressed.  On the other
hand, the meaning ``immediately after'' is [secondarily] understood.
For this very reason, [Śabara] says ``It is [often] experienced that
[this word \atha{} at the beginning of the JS] has the meaning of
introducing something that comes immediately after something has concluded.''  If the
meaning of \atha{} was``immediately after,'' then he would have
said [more explicitly], ``[This word \atha{} is used] to signify
that the beginning comes immediately after something has concluded.''
%[instead of saying ``it is experienced \ldots,'' in which case it is
%implied that that is not always the case.

Also,%\pratika{anantarasyeti}
\pratikap{p01_31} [the word] ``something that comes immediately after (\emph{anantara})''  signifies the thing[, a substance,] that comes immediately after (\emph{anantarabhāvin}).  %[This is the evidence that even \Sabara\ did not consider that the meaning of the word \atha\ was `immediately after' since \Sabara\ does not use the expression `immediately after (\emph{ānantaryārthe}),' but uses `something that comes immediately after (\emph{anantarasya}).']
% When I was trying to locate the usage of the simil putra-putṛ, I 
% encountered a reference to Vaiśeṣika-s in \rBSBh{2.2.17}{2.2.17}, where 
% śa"nkara lists six padārthas of V.  Was there V at the time of 
% S?  Also there is a mention to satkāryavādin that include Manu.  
% What is he referring?  Off course there is a mention to Sāṃkhya.  
% Who are the yogins for S?  Who are the Vedavādins?  Remember 
% Vedavādins in Jayamangalā? [This turns out to be Māṭharavṛtti 20080827]
% 
And,\refm{30} %%\pratika{api caivaṃ}
% \pratikap{p01_32}
a \smrti\/ (\VMBh) indeed says thus:
\begin{quote}
  The primary meaning of some indeclinables is nominal, of others it
  is primarily verbal.  Such words as ``upwards (\emph{uccair}),''
  ``downwards (\emph{nīcair}),'' etc., have the nominal meaning as
  their primary meaning.  Such words as ``near (\emph{hiruk}),''
  ``separately (\emph{pṛthak}),'' etc., have the verbal meaning as their
  primary meaning.  There is no other meaning of indeclinables apart from these. (\rVMBh{1.1.38 or 1.2.47}{1.1.38 or 1.2.47})\footnote{Note that the quote here is not exactly the same as what we find in edition of the \VMBh. Parallel passages appear twice in the \VMBh, and in both cases, the editions of the \VMBh\ have \emph{ādi} at the end of examples while the quote here does not. Perhaps more importantly, the whole part on \emph{taddhita} or on \emph{tibanta} is missing in the quote here in the YVi. Instead, it has a sentence to exclude the third category of \avyaya s.}\labelh{VMBH1138_T}
\end{quote}
Hence, if the meaning of the word \atha\ is ``immediately after'' [as you
claim], hearing a nominal suffix is reasonable, because its primary
meaning is a substance.  On the other hand, if [the meaning of \atha\ is] the action of beginning [as I claim], then it is reasonable not to hear
a nominal suffix, because its primary meaning is something that does not
exist[; and that is the case].\footnote{As seen in the critical apparatus, the text of the preceding two sentences has been transmitted in a confused state.  It is most conspicuously observed in the non-sensical reading \emph{ārambhakriyārthatve tu sa tu} recorded in both \Tm\ and L\@.  The reading was already in the common archetype.  Even more curious is the reading \emph{vibhaktyarthaśravaṇam} in all the Devanagari manuscripts.  The scribes of \Td\ and M inserted \emph{°rtha°} when their respective exemplars did not have the syllable.  I do not believe that there was communication between \Td\ and M\@.  What probably happened was that two scribes independently and subconsciously inserted the syllable \emph{°rtha°} to form a compound that sounded more natural to them.}

In%\pratika{tasmād atha°}
\pratikap{p01_33} conclusion,\refm{31} it is reasonable that
the word \emph{atha\/} only has the meaning of announcing [a topic].

The%\pratika{`athe'ty a\-dhi\-kā\-rā°}
\pratikap{p01_34} meaning\refm{32} of the word \emph{iti} in
the ``the meaning of \atha{} is the declaration of the topic (\emph{athety
  adhikārārthaḥ})'' is to show that [the word \atha{}] is being quoted.  For
example, ``He said `cow' (\emph{`gaur' ity āha}).''  Even though the meaning of the word \atha{} is well-known, it is mentioned [by the \Bhasyakara] in order to make a student pay attention to what is going to be started.  Indeed the following formula
is well-known in \emph{śruti}:
\begin{quote}
        I am going to explain to you.  Pay close attention while I am 
        explaining.  (\rBAU{2.4.04/4.5.5}{2.4.4/4.5.5})
\end{quote}

Being%\pratika{kaḥ punar}
\pratikap{p01_35}\labelh{KHPNR1} asked,\refm{33} ``Then what is yoga,
whose instruction is announced [as the topic]?'' [the \Bhasyakara] says---\bht{Yoga is \samadhi.} %\footnote{Note that the author of the YVi does not
%  comment on the sentence:
%\begin{quote}
%        \emph{yogānuśāsanam adhikṛtaṃ veditavyam\dd}
%\end{quote}
%which appears in vulgate YBh.\label{note:veditavyam}}
 Since [the \Bhasyakara\ uses] the word `\samadhi\/ (assimilation),'\footnote{\label{ASSMLT}The translation ``assimilation'' for the word \samadhi\ is chosen here because usual translations of the word, such as ``concentration,'' etc., do not fit in the discussions below.  In the following \samadhi\ is said to be ``a property of the mind common to all stages.''  It is explicitly mentioned in our text that \samadhi\ is the state of stability in all the five stages of the mind.  See the paragraph starting with \emph{nanu ca bhūmīnām\ldots} below \reftoeditionpar{p01_42} (translated on p.~\pageref{NNCBH1}).} [he] rejects the meaning ``to unite'' of the root \emph{yuj-}, which [appears in the \Dhatupatha] as \emph{yujir yoge} (7.7), instead, he accepts the meaning ``assimilation (\emph{samādhi}),'' [which appears in the \Dhatupatha\ as] \emph{yuja samādhau} (4.68).  Yoga is assimilating (\emph{samādhāna}).\footnote{It is interesting to note that the two entries in the \Dhatupatha\ appear to reflect the two major trends in defining yoga.  One is, as seen here, to define it as \samadhi\ (perhaps originally meant as concentration, but not necessarily as describing a mental process).  The other is to define it as \samyoga.  One example of the latter definition (from the \Vaisesikasutra) is dealt with later in our commentary.  See pp.~\pageref{samyoga}\,ff. The same definition was presupposed by the \Naiyayika s (NS/NBh 4.2.38); even the \PYS\ 2.53--55 may have had the same system in its background.  The system of that yoga probably had the \Nyaya/\Vaisesika\ type of view on how perception occurs: by the contact between objects, sense faculties, \manas\ and \atman.  Their yoga was a gradual process to detach the component of perception from the outermost one, viz., first the objects, then the sense faculties.  When \manas\ resides with \atman\ (\samyoga), it is yoga.  See also \Carakasamhita\ \Sarirasthana\ 2.138--139.}
%These two \sutra s from the \emph{Dhātupāṭha} may
%  be of importance in studying the history of yoga, as according to
%  Bronkhorst~\cite{bronkhorst:history}, they entered the \emph{
%    Dātupāṭha} after \Patanjali.  This fight between whether \emph{
%    yoga} is \emph{samādhi} or \emph{saṃyoga} is reflected in the
%  discussion below. (Pages \pageref{samyoga}ff.)}

[Objection]\refm{34}%\pratika{nanu ca}
\pratikap{p01_36} If the \sutra\ is explained up to ``yoga is
\samadhi,'' then the \bhasya\ that starts with ``That (\samadhi) is a [property of mind] common to all stages'' is seen as detached.\footnote{This question arises only from the assumption that yoga/\samadhi\ is something special, that should be achieved with efforts.  The statement that it is present in all the stages of mind appears betrays such an assumption and thus the following discussions on the stages of mind appears to be detached. This is also the reason I have difficulty translating the word \samadhi\ in this context.  (See note \ref{ASSMLT} above.) Thus I think this question was posed by our author himself, as a hermeneutical problem.  He felt that the transition was abrupt, but was compelled to defend the text he was commenting upon.}

[Answer]\refm{35}%\pratika{naiṣa}
\pratikap{p01_37} There is no such fault.  For, when it is said that [yoga is] \samadhi\/ here, since [\samadhi] presupposes something to be assimilated[, i.e., the subject,] and since there are many things that can be assimilated, there is a dispute among those who know much,%
  \footnote{As the editors of the 1952 edition emended to \emph{bahuvidhā pratipatteḥ}, the expression \emph{bahuvidāṃ vipratipatteḥ} might seem awkward.  However, the construction, genitive plural plus the word \emph{vipratipatti} is found in \Sankara's works.  See \rBSBh{2.1.11}{2.1.11}: \emph{prasiddha\-māhātmyā\-bhimatānām api tīrtha\-karāṇāṃ kapila\-kaṇabhu\-kprabhṛtī\-nāṃ paraspara\-vipratipatti\-darśanāt}; 21: \emph{tatraivaṃ sati samyagjñāne puruṣāṇāṃ vipratipattir anupapannā}; \rBAUBh{4.5.15}{4.5.15}: \emph{vipratipattiś ca śāstrārthapratipattṝṇāṃ bahuvidām api}.  Cf.\ also \rBSBh{2.3.18}{2.3.18}: \emph{sa kiṃ kāṇabhujānām ivāgantukacaitanyaḥ, svato 'cetanaḥ, āhosvit sāṃkhyānām iva nityacaitanyasvarūpa eva, iti vādipratipatteḥ saṃśayaḥ}.  Here the word \emph{bahuvid} is probably used somewhat sarcastically.%

  The alternatives (\emph{ātman}, \emph{śarīra}, \emph{indriya\/}s) listed as the subject of \emph{samādhi\/} call for some attention.  These three, the most conspicuously the term \atman, are not part of the \PYS's terminology (an indication that our author is thinking of a wide context).  This set of three appears to reflect our author's view of what constitute an individual.  What is missing here is the word \manas\ or \citta. (Our author does not distinguish \manas\ and \citta.  He does not clearly distinguish \manas/\citta\ and \buddhi, either.)  But he omitted it because that is the correct answer to the question what is assimilated.  The set of \atman, \manas\ (or \citta), \sarira\ (or \deha), and \indriya s is what our author generally thought of as the constituents of an individual.  This is a rather simple view, somewhat reminiscent of the \Naiyayika/\Vaisesika\ view of an individual.%
  } %
such as ``Is \atman\ assimilated?''\ ``Is the body assimilated?''\ or ``Are the senses assimilated?'' Also, since assimilation presupposes something that qualifies itself [as in ``what exactly is it?'']%
  \footnote{See the next paragraph for what is considered to be the qualifier (\emph{viśeṣaṇa}) of \samadhi.%
  } %
[it is appropriate that the \bhasya\ continues the discussion on \samadhi, and hence there is no fault that the \bhasya\ up to \emph{yogaḥ samādhiḥ} and the following are detached].%
  \footnote{The translation is based on the understanding that the sentence that starts with \emph{iha samādhir ity ukte\ldots} consists of two main clauses even though it has four phrases in the ablative case expressing reasons.  This understanding derives from the positions of the particle \emph{ca}.  The first two phrases that end in the ablative case are combined by the particle \emph{ca}.  And semantically those are two stages.  This should lead to the first major reason \emph{bahuvidāṃ vipratipatteḥ}.  The second reason \emph{svaviśeṣaṇakākāṅkṣatvāc ca samādheḥ} is again connected with first major reason by the particle \emph{ca}.  This understanding fits the following commentary as well.  The two components of the question \emph{kasyāyaṃ, kiṃviśeṣaṇakaḥ} at the beginning of the next paragraph refer to the two reasons of the current sentence.  The first (\emph{kasyāyam}) corresponds to the question of what is assimilated, and the second \emph{kiṃviśaṣaṇakaḥ} corresponds to similarly worded \emph{svaviśeṣaṇkākāṅkṣatvāc ca samādheḥ}.%
  }

In\refm{36}%\pratika{kasyāyaṃ}
\pratikap{p01_38} order to answer the inevitable question, ``To what does
it (\samadhi) belong, what qualifies it?''\ [the \Bhasyakara] states the following---\bht{And that is a property of the mind common to all stages.}  It is a property of mind, not of \atman, etc.  And the mind assimilates itself because it does not require another agent.%

Being\refm{37}%\pratika{kāḥ punar}
\pratikap{p01_39} asked, ``Then what are those stages?''\ [the \Bhasyakara] answers---\bht{The stages are fixated (\emph{kṣipta}) [mind], stupefied (\emph{mūḍha}) % diff. fr. suggested
[mind]}, etc.  The past passive suffix \emph{-ta} in `fixated (\emph{kṣipta}),' etc., is used in the meaning that the object of an action is the agent itself (\emph{karmakartṛ}).  Compare, ``The granary collapsed by itself.''%
  \footnote{See MBh 3.1.87, \citet[nn.~22--24]{wezler:yvi1}\index{Wezler, Albrecht}.%
  }  %
\bht{Fixated [mind]} means a benumbed [mind] being fixed with an undesirable object.  \bht{Stupefied [mind]} lacks the distinction [between different objects].  \bht{Scattered (\emph{vikṣipta}) [mind]} means being fixated at several objects [where the passive past participle is used] in the same sense [as in the case of fixated mind] in which the object of the action is the agent itself.  Also, it is not capable of distinguishing [different objects] for the very reason that it is scattered.  \bht{Focused [mind] (\emph{ekāgra})} means the [mind] holding a consistent notion.  \bht{Suppressed [mind] (\emph{niruddha})} is without any notion---[All the adjectives modify the word] the mind.%

[Objection]\refm{38}%\pratika{nanu ca}
\pratikap{p01_40} If [the \Bhasyakara] intends to talk about stages [of the mind], properties (\emph{dharma}), why is an property-bearer \emph{dharmin}) [that is the mind] discussed by ``fixated [mind],'' etc.?%
  \footnote{Note that all the terms to refer to the different stages of mind in the \bhasya\ are adjectives in the nominative case, neuter singular.  They naturally modify the word \emph{cittam}.  And this is the understanding our author holds and is the background of this objection.%
  }

[Answer]%\pratika{naiṣa}
\pratikap{p01_41}\refm{39} There is no such fault.  The property-bearer indicates exactly the property.  Because the properties belong to the property-bearer.  For
example: If one is asked ``What is the characteristic of being a
cow?''\ the property is indicated by the property-bearer, by saying, ``It has
horns and a hump, and is hairy toward its tail.''\footnote{According to
  \citet{wezler:yvi3}\index{Wezler, Albrecht}, this is a reference to \rVS{2.1.18}.  Wezler
  points out that the \sutra\ is referred to in support of the
  view held by the author of the YVi, and that the \emph{sūtra}, therefore,
  was favorable to him.  Of course the archetypal description of the cow appears already in the \VMBh.} Therefore [the \Bhasyakara] means to say that fixation, etc., are the stages, the properties of the mind.

% nanu ca bhuumiinaa.m 
[Objection]%\pratika{nanu ca}
\pratikap{p01_42}\labelh{NNCBH1}\refm{40} If the stages are properties of the mind and \samadhi, too, is [a property of the mind], then how do we distinguish \samadhi\/ that is said to be ``common to all the stages'' as a common denominator from the stages?  [If there is nothing special about \samadhi, it would not be necessary to introduce the term.]%
     \footnote{The whole exchange appears a little clumsy.  Answering to the question ``How (\emph{ka\-tha\-m})\ldots?'' with ``no'' does not appear to be usual.  I have interpreted and translated it as short for ``naiṣa doṣaḥ'' below.  Still, some text might have been lost.%
     }
%
% na| saamaanyabhuuta
[Answer]\refm{41} That is not correct.  Because \samadhi\/ is
universal, while the stages are particulars.  That is, in such phrases as ``[the mind] stays fixated,'' ``[the mind] stays stupefied,'' ``[the mind] stays scattered,'' and ``[the mind] stays focused,'' and ``[the mind] stays suppressed,''%
  \footnote{I have conjecturally added \emph{niruddham} after \emph{vikṣiptam} and \emph{ekāgram}, while it is missing in all the manuscript and the \EP.  Also, following the \EP, I have emended \emph{vikṣiptam} at the beginning of the series to \emph{kṣiptam} and \emph{kṣiptam} at the third position to \emph{vikṣiptam} so that the order of the stages follows that of the \bhasya.  I do not believe that the swapped positions of \emph{kṣipta} and \emph{vikṣipta} or the lack of \emph{niruddha} were intended by the author.  I judge rather that those are results of the damage to the text during the transmission.  Cf.\ the following for the state of the text preserved in manuscripts.  Frequent appearances of the same words (\emph{kṣiptam}, etc.) may have contributed to the corruption in this particular sentence.%
  } %
stasis (\emph{sthiti}) in the stages ``fixated,'' etc., accompanies universally.  And\labelh{station_immersion} \samadhi\/ is stasis.  Therefore [\samadhi] is commonly%
  \footnote{The reading \emph{sā\-dhā\-ra\-ṇye\-na} is a result of emendation.  Most manuscripts read \emph{prā\-dhā\-nye\-na}, except \Tm\ \emph{sā\-dhā\-nye\-na} before correction.  It is likely that the common archetype of \Tm\ and L indeed read \emph{prā\-dhā\-nye\-na}.  What probably happened was that the scribe of \Tm\ first wrote \emph{sā\-dhā\-nye\-na}, expecting the word \emph{sā\-dhā\-ra\-ṇye\-na} before realizing that he had written something not in his exemplar.  Nonetheless, I adopt the reading what he subconsciously expected; otherwise, the sentence is not coherent.%
  } %
present in all the stages in the form of a universal.%
  \footnote{The word \emph{iti} at the end of this paragraph marks the end of a topic, and hence no specific translation is given.%
  }

% sarvabhuumip.rthivii
% Vetter's trl.:
[According to \rAAP{5.1.41--42}{5.1.41--42}]%\pratika{`sarvabhūmi°}
\pratikap{p01_43} \emph{sarvabhūmipṛthivī}[\emph{bhyām aṇañau} and \emph{tasyeśvaraḥ}] the suffix \emph{a} (\emph{aṇ}) [added to \emph{sarvabhūmi}] gives the meaning ``governing [all stages].''  Since the rule \emph{anuśatikādī\emph{[}nāñ ca\emph{]}} (\rAAP{7.3.20}{7.3.20}) is applicable [because \emph{sarvabhūmi} is included in the \emph{gaṇa} \emph{anuśatikādi}], [like in \emph{ānuśātika}, the first vowel] in both constituents of the compound [\emph{sarvabhūmi}] is substituted by the longest modification[, thus resulting in the form \emph{sārvabhauma} as found in the \bhasya].

%[The\refm{42} word \emph{sārvabhauma} has] the \emph{aṇ} suffix
%(\emph{a}), because it is governed by [P\=a\d{n}ini's rule] ``\emph{%
%  sarvabhūmipṛthivī}'' (AA 5.1.41).  [The \emph{aṇ} is used] in
%the meaning of `governing'[, the meaning mentioned in the subsequent
%rule ``\emph{tasyeśvaraḥ}'' (AA 5.1.42)].  Thereupon the resulting
%compound is governed by the rule ``\emph{anuśatikādīnāñ ca}''
%(AA 7.3.20), resulting in the \emph{vṛddhi} [substituting the first
%vowel] of both constituents of the compound.

% anye punar
On%\pratika{anye}
\pratikap{p01_44}\refm{43}\labelh{ANYEPUNARtrl} the other hand, some people explain the stages as external and internal objects of \samyama\ meditation.%
  \footnote{The relation of \samyama\ with particular objects and its result is the main topic of the third {pāda}.  Here the view that the \emph{bhūmi\/}s in the context of \rYBh{1.01}{1.1} is the same \emph{bhūmi\/}s in the context of \emph{saṃyama\/}s in \rYS{3.06}{3.6} (\emph{tasya bhūmiṣu viniyogaḥ}) is criticized.  Someone held the view (according to the YVi) that they are the same, presumably because the same word \emph{bhūmi\/} is used in both places.  It should be noted that the author of the YVi is referring to an interpretation of the YBh here.  It is probable that he is using a sub-commentary on the YBh.  Note, however, that the interpretation is different from that of \Vacaspati.%  See section~\ref{vpmandyvi}.

  Philologically, our manuscripts indicate that their common ancestor consistently read \emph{samayama} instead of \samyama.  It is interesting that such an obviously faulty reading was consistently adopted and more or less faithfully reproduced.%
  } %
[But] for them there is no chance to place [a stage and the mind] in the same case (\emph{sāmānādhikaraṇya} relationship) [in the way the \Bhasyakara\ does:] ``When the mind is scattered\ldots\@,''%
  \footnote{This is a reference to the next sentence in the \bhasya.  See page~\pageref{ATRVKSP}\,ff.%
  } %
and they contradict their own doctrine.  How?  [According to their interpretation,] that on which [the mind] is fixated is ``fixated''; that with regard to which  [the mind is stupefied] is ``stupefied''; and that with regard to which [the mind is scattered] is ``scattered.''%
  \footnote{The last phrase \emph{vikṣpiyate 'sminn iti ca vikṣiptam} is missing in all the manuscripts and the \EP.  The context requires this last item to be mentioned as well.  One, albeit not very strong, indication that there is a textual corruption is the lack of the particle \emph{ca} in the preceding phrase regarding \emph{mūḍha}.%
  }  %
If this [interpretation] were true, then they could not be the object of \samyama\ meditation.  In that (in the state of \samyama) \emph{kṣepa}, etc., cannot occur.  For, \samyama\ meditation is possible at the state of focused (\emph{ekāgra}).

Furthermore,%\pratika{kiñ cānyat}
\pratikap{p01_45}\refm{44} since the root \emph{kṣip-} does not have the
meaning of ``steadiness,'' etc., it is not possible to have the \emph{%
  niṣṭhā} suffix in the meaning of locus (\emph{adhikaraṇa}) [as prescribed by {%
  \Panini}].\footnote{See \rAAP{3.4.76}{3.4.76}: \emph{kto ’dhikaraṇe ca dhrauvyagatipratyavasānārthebhyaḥ} (The past passive participle affix \emph{-ta\/} is used in the sense of the locus, too, when it follows verbal roots that have the meanings of steadiness, motion, or termination).  The argument here is still based on the analysis of the opponent's view: \emph{kṣipyate ’sminn iti kṣiptam}, etc.  Note the use of the locative case (\emph{yasmin}).  It represents the locus (\emph{adhikaraṇa}).  Our authors angle of attack here is that such a formulation, ``that on which something is thrown is `thrown'\thinspace,'' to analyze a past passive participle (\emph{kṣipta}, etc.)\ is grammatically impossible.  Pāṇini allows only verbs that have certain meanings to form a passive past participle in the meaning of the locus.  The root \emph{kṣip-} does not have any of the three meanings, \emph{dhrauvya}, \emph{gati}, or \emph{pratyavasāna}.  According to the \Dhatupatha, the root \emph{kṣip-} has the meaning \emph{preraṇa} (driving).}

Furthermore,%\pratika{kiñ ca}
\pratikap{p01_46}\refm{45} when the mind has been suppressed, \emph{saṃyama\/} meditation, too, will not exist because no activity exists.  Accordingly, the object of \samyama\ meditation does not exist.  For, if suppression was a specific object, the mind would not be suppressed.  Because at that time (when the mind is suppressed) [the mind] ceases to take an object.  Also, [there is a problem that] completely enumerating [the stages] is not possible.  Because the objects of \emph{saṃyama\/} meditation are not just five of \emph{kṣipta}, etc.  Because they are innumerable.

Objection:%\pratika{nanu ca}
\pratikap{p01_47}\labelh{samyoga} Some\refm{46} consider yoga differently.  Namely, they say---\footnote{\citet{wezler:yviandvs}\index{Wezler, Albrecht} deals with this
quotation from/reference to the \VS\ in detail.}
\begin{quotation}
        \bht{Pleasure%\pratika{indriya}
\pratikap{p01_48} and pain arise from the contact between senses, mind, and objects.  [Yoga] is a union (\emph{saṃyoga}),}%
		\footnote{All the manuscripts read \emph{saṃyogaḥ} at the end of the second \sutra.  \citet{wezler:yviandvs}\index{Wezler, Albrecht} proposes to emend this to \emph{sa yogaḥ}.  But this emendation does not seem necessary.%

		The commentary that is cited with these \sutra s below does not support \emph{sa yogaḥ} in the \sutra.  First, the word \emph{sa} in the commentary \emph{sa prāṇamanovinigrahāpekṣaḥ} towards the end of the quoted commentary corresponds to the preceding relative pronoun \emph{yaḥ} in \emph{yo ’sau}.  The syntax of the commentary should be viewed as ``\emph{yo ’sau\ldots saṃyogaḥ} (a word from the \sutra) \emph{sa prāṇamanovinigrahāpekṣaḥ} (a word from the \sutra) \ldots.''  Accordingly, the word \emph{sa} does not seem to be from the \sutra.  The parallel syntax in ``\emph{yo ’sau} \ldots \emph{sannikarṣaḥ} \ldots \emph{tasya}'' near the beginning of the commentary supports this view.  The commentary employs the similar syntax when explaining the relationship between words in the \sutra s. %
		Secondly, the commentary is explicit in that yoga is \emph{saṃyoga\/} and treats \emph{sukhaduḥkhābhāvaḥ} more like a condition (\emph{tasyām avasthāyām}).  [Cf.\ the Mithilā version of \rVS{5.2.14}: \emph{saśarīrasya sukhaduḥkhābhāvaḥ\dd} \citep[113]{vs} where the \sutra\ is terminated after \emph{sukhaduḥkhābhāvaḥ} and the word is not syntactically connected with the following \sutra: \emph{saṃyogaḥ\dd}; and \rnnNS{4.2.45}{4.2.45}/\rnnNBh{4.2.45} 4.2.45: \emph{\textbf{tadabhāvaś cāpavarge\ddd} tasya bu\-ddhi\-ni\-mittā\-śraya\-sya śa\-rīre\-ndri\-ya\-sya dha\-rmā\-dharmā\-bhāvād a\-bhā\-vo ’pa\-va\-rge\dd\ ta\-tra yad u\-kta\-m ``apa\-varge ’py evaṃ pra\-saṅga'' iti (4.2.43), tad a\-yu\-ktam\dd\ tasmāt sarva\-duḥkha\-vi\-mokṣo 'pa\-vargaḥ\dd\ yasmāt sarva\-duḥkha\-bījaṃ sa\-rva\-duḥkhā\-ya\-ta\-naṃ cā\-pa\-varge vi\-cchi\-dyate tasmāt sa\-rve\-ṇa duḥkhe\-na vi\-mu\-ktir apa\-vargaḥ\dd\ na nirbī\-jaṃ ni\-rā\-yatanaṃ ca duḥ\-kham utpa\-dya\-ta iti\ddd}.]  This makes it difficult for \emph{sukhaduḥkhābhāvaḥ} and \emph{yogaḥ} to be in the \emph{sāmānādhikaraṇya\/} relationship, at least according to the commentary. %
		Thirdly, the commentary refers to \emph{saṃyoga\/} twice. This insistance on \emph{yoga = saṃyoga\/} makes it easier to think that the word \emph{saṃyogaḥ} was already in the \sutra\ as part of the definition of yoga than that it was introduced from somewhere else. %
		Thus even if the word \emph{saṃyogaḥ} might not be original to the \VS, there seems to be little support that the author of the YVi \emph{did not} quote from the \VS\ that already contained the word.%

		Apart from the commentary quoted in the YVi below, evidence from the Vaiśeṣika side appears to support the presence of \emph{saṃyogaḥ\/} in the \sutra, too.  All the manuscripts of the VS used by Jambūvijaya have \emph{saṃyogaḥ} \citep[42]{vs}.%

		In addition, there appears to have been a great deal of controversy over the definition of yoga---whether it is \emph{saṃyoga\/} or \emph{samādhi}---before the authority of the YS/YBh prevailed.  This may be reflected even in two entries of the root \emph{yuj-\/} in the Dhātūpāṭha, namely \emph{yujir yoge} and \emph{yuja samādhau}.  The two entries of the Dhātupāṭha are said to be referred to in the \YBh. (See the paragraph \emph{kaḥ punar yogo\ldots} \reftoeditionpar{p01_35} translated on p.~\pageref{KHPNR1}.  Cf.\ \citet{bronkhorst:history}.)  Another reflection of this controversy may be seen between \rNS{4.2.38}{4.2.38} (\emph{samādhiviśeṣābhyāsāt}) and its interpretation in the NBh (\emph{sa tu pratyāhṛtasyendriyebhyo manaso dārakeṇa prayatnena dhāryamāṇasyātmanā saṃyogas tattvabubhutsāviśiṣṭaḥ} \ldots).  The \NBh\ uses \emph{saṃyoga\/} instead of \emph{samādhi\/} of the \sutra.  [The system of yoga adopted in the Nyāya and Vaiśeṣika systems is closely tied to their epistemology.  Cf.\ the criticism toward the Nyāya/Vaiśeṣika theory of the generation of knowledge through connection (\emph{saṃyoga}) between \emph{manas\/} and \emph{ātman} in the YVi on \PYS\ 1.2.% (\pageref{nvsamyoga}--\pageref{nvsamyogaend}).%  Close relation between yoga and the knowledge of the truth \emph{tattvajñāna\/} may also be reminded.  In the yoga system adopted in Nyāya/Vaiśeṣika systems, the goal of yoga is \emph{tattvajñāna\/} while in the system of Patañjali the goal is \emph{cittavṛttinirodha}.%
		]  Thus given the conflict between \emph{saṃyoga\/} and \emph{samādhi\/} surrounding the definition of yoga, it seems natural that a sūtra defining yoga included the other synonym for yoga.%

		And finally the criticism in the paragraph \emph{athāpi syāt---yogaḥ samādhir iti\ldots} \reftoeditionpar{p01_63} (translated on p.~\pageref{ATHPSY1}) by the author of the YVi presupposes a defining term (\emph{samādhi\/} or \emph{saṃyoga}) to be present in the \sutra.  Note that when he finally deals with the view \emph{yogaḥ samādhiḥ}, he has criticized most of the words in the opponent's sūtras:  \emph{indriyamano'rthasannikarṣāt} in two paragraphs \emph{kiñ cānyat\dd\ mano°\ldots} \reftoeditionpar{p01_55} (trl.\ p.~\pageref{KCMN1}); \emph{sukhaduḥkhe} in three paragraphs starting with \emph{kiñ cānyat---sa\-rva\-prāṇa°\ldots} \reftoeditionpar{p01_57} (trl.\ p.~\pageref{SRVPRNBH1}); \emph{tadanārambhaḥ} in the paragraph \emph{indriyasaṃyoge\ldots} \reftoeditionpar{p01_62} (trl.\ p.~\pageref{INDRYSMYG1}); \emph{ātmasthe manasi} in four paragraphs starting with these exact words \reftoeditionpar{p01_51} (trl.\ p.~\pageref{ATMSTH1}); \emph{saśarīrasya} in the paragraph \emph{kiñ cānyat---upātte 'pi\ldots} \reftoeditionpar{p01_60} (trl.\ p.~\pageref{UPTTP1}); \emph{sukhaduḥkhābhāvaḥ} in three paragraphs starting with \emph{kiñ cānyat---sarvaprāṇa°} \reftoeditionpar{p01_57} (p.~\pageref{SRVPRNBH1}); \emph{prāṇamanovinigrahāpekṣaḥ} in the paragraph that starts with \emph{kiñ cānyat---prāṇamano°\ldots} \reftoeditionpar{p01_61} (trl.\ p.~\pageref{KCPRMN1}).  There should be something to be criticized in the paragraph \emph{athāpi syāt---yogaḥ samādhir iti\ldots}.  Note again that in \rNBh{4.2.38} \emph{samādhi\/} in the \sutra\ is replaced by \emph{saṃyoga}.%

		From the above, it appears that there was the word \emph{saṃyogaḥ} in the \sutra\ in question.  Adding \emph{yogaḥ}, as suggested by the editors of the 1952 edition (p.~6, l.~10) might make the \sutra s complete as a definition of yoga.  The commentary of Candrānanda also seems to imply that there was the word \emph{yogaḥ} as the last element of the \sutra\ \citep[42,24]{vs}.  Despite my inclinations, both from those of the YVi and the \VS, I did not add the word in the \sutra.  The fact that the author of the YVi supplies the word \emph{yogaḥ} to both \emph{sukhaduḥkhābhāvaḥ} (here is a deviation from the interpretation given in the commentary accompanying the \sutra s) and to \emph{samādhiḥ} (the replacement I think for \emph{saṃyogaḥ}) suggests that the word \emph{yogaḥ} was provided by the context of surrounding other \sutra s.%  After all, the \sutra s in question are viewed as not original to the \Vaisesikasutra\ \citep{wezler:yviandvs}.  Perhaps we should not forget that there is even a possibility of interpolation (in the VS) and quotation (in the YVi) from the same source.
		} %
	\bht{which follows restraining of breaths and mind, and is [characterized as] its (of the contact) not happening, absence of pleasure and pain, and it belongs to someone who still has a body, under the condition that the mind stays with \atman.}%
        
    The%\pratika{tad iti}
\pratikap{p01_49} word ``\bht{Its}'' presupposes something that is being        discussed---the cause \bht{of pleasure and pain} = the \bht{contact between} \atman, \bht{senses, mind, and objects}.  Its \bht{not happening} = its not originating.  How does it occur?%
        \bht{Under the condition that the mind stays with \atman}---not when it stays with senses.  [Yoga] \bht{belongs to someone who still has body}---belongs to someone whose body has not perished. At that time, as the contact does not exist, \bht{pleasure and pain} do \bht{not exist}.   For it is said ``if the cause does not exist, its result does not exist (\rVS{1.2.1}).''  Under this circumstance there is a \bht{union} between omnipresent \bht{\atman} and the \bht{mind}.  That is yoga, which is a specific \bht{union} that \bht{follows restraining of breaths and the mind}.%
\end{quotation}

We%\pratika{atrocyate}
\pratikap{p01_50} answer---

[Answer]%\pratika{`ātmasthe}
\pratikap{p01_51}\labelh{ATMSTH1} To say\refm{47} that \bht{under the condition that the mind stays with \atman} is not rational.  Because the mind always stays with \atman.
  
If%\pratika{indriyādi°}
\pratikap{p01_52} you say, ``[It is said] with regards to [the condition] that contact between senses, etc., is not happening,'' then stating ``under the condition that the mind
stays with \atman'' would be meaningless since it is already established
just by saying ``its not happening.''

Furthermore,%\pratika{kiñ cānyat}
\pratikap{p01_53} even a liberated (\emph{mukta}) would have yoga since his mind stays with \atman, and he does not have contact with senses.\footnote{Here our author assumes that his opponent would not like the idea that a liberated person has yoga.  If the opponent holds the view that yoga is the direct means of liberation, this might not be a problem. This does not appear to be the case.  Liberation (\moksa) is defined in the Vaiśeṣikasūtra a few sūtras after the sūtras under consideration here: \emph{ta\-da\-bhā\-ve saṃyogābhāvo ’prādurbhāvaḥ sa mokṣaḥ} (5.2.20 in the numbering with Candrānanda's commentary).  Hence it is indeed probably undesirable that a liberated person has yoga.  It is, however, interesting to note that our author assumes that the ``liberated'' still has \manas.  The \emph{saṃyoga} in \emph{saṃyogābhāvaḥ} in the above sūtra is, according to Candrānanda, the union between \atman\ and \manas, and this interpretation is the most plausible.  (The word \emph{tad} in \emph{tadābhāve} refers to \emph{adṛṣṭa}s mentioned in the previous sūtra.)  Our author thus has a different understanding of what a ``liberated'' is.  The expression \emph{muktātman} or simply \emph{mukta} is used several times in this work, especially in relation to \Isvara.  For example, see p.~\pageref{LBRTDATM}.  There the assumption is that the liberated ātman or a liberated person does not have a body.  The use of the term is probably consistent with that of Kumārila.  See note \ref{KMRSAP} on p.~\pageref{KMRSAP}.}  For, since he, too, [being \atman,] is omnipresent, and the mind is eternal,\footnote{The \Vaisesika s hold that \manas\ is eternal.  See the references in the edition part.} [the mind] would indeed always stay with \atman.
% 
% Usage of `mukta' seems to be interesting because he is speaking of 
% a special being.
% See PDhS and p.84 (370):
% \begin{quotation}
%        yathāsmadādīnāṃ gavādiṣvaśvādibhyas
%       tulyākṛtiguṇakriyāvayavasaṃyoganimittā
%       pratyayavyāvṛttir dṛṣṭā gauḥ śuklaḥ śīghragatiḥ
%       pīnakakudmān mahāgaṇṭa iti. tathāsmadviśiṣṭānāṃ
%       yogināṃ nityeṣu tulyākṛtiguṇakariyeṣu paramāṇuṣu
%       muktātmamanaḥsu cānyanimittāsambhavād yebhyo
%       ^^^^^^^^^^^^^^^^^
%       nimittebhyaḥ pratyādhāraṃ vilakṣaṇo 'yaṃ vilakṣaṇo
%       'yam iti pratyayavyāvṛttiḥ, deśakālaviprakarṣe ca
%       paramāṇau sa evāyam iti pratyabhijñānaṃ ca bhavati, te
%       'tyā viśeṣāḥ.
% \end{quotation} , also cf. p.80 (359), p.84 (371)
% 
% And by the way rememver dāruyantra in Vivaraṇa?  PDhS has an
% entry---i.e. śarīraparigṛhīte vāyau vikṛtakarmadarśanād
% bhastrādhimāpayiteva.  nimeṣonmeṣakarmaṇā niyatena
% dāruyantraprayokteva.

Furthermore,%\pratika{kiñ cātmanaḥ}
\pratikap{p01_54}\refm{49} since there are no regions in \atman,\footnote{That \atman\ has no region is not part of the \Vaisesika\ system. Old \Vaisesika s take for granted that \atman\ has regions. This is observed in expressions such as \emph{antarhṛ\-da\-ye ni\-ri\-ndri\-ye ā\-tma\-pra\-de\-śe ni\-śca\-lam ma\-nas ti\-ṣṭha\-ti} \citep[183]{pdhs}, \emph{ha\-sta\-vaty ā\-tma\-pra\-de\-śe pra\-yatnaḥ saṃjā\-ya\-te} \citep[297]{pdhs}. Note that the expression \atmapradesa\ is used in the sense of “a region in \atman.” The \Naiyayika s, on the other hand, saw a problem in expressions such as \emph{ātmapradeśa}. \Nyayasutra\ \rnNS{2.2.17}{2.2.17} reads \emph{kāra\-ṇa\-dra\-vya\-sya pra\-deśa\-śabde\-nā\-bhi\-dhā\-nāt}. This is an answer to the criticism expressed in \rNS{2.2.14}{2.2.14} \emph{na, ghaṭā\-bhā\-va\-sā\-mā\-nya\-ni\-tya\-tvān ni\-tye\-ṣv apy a\-ni\-tya\-vad upa\-cā\-rāc ca}. The part \rNS{2.2.17}{2.2.17} is responding to is \emph{nityeṣv apy anityavad upacārāt}. This, in turn, is an objection to one of the reasons that the proponent gives to establish that \sabda\ (sound) is impermanent (\anitya) and to be created (\emph{utpadyate}) rather than revealed (\emph{abhivyajyate}). It is the part \emph{krṭakavad upacārāt} in \rNS{2.2.13}{2.2.13}. Following the explanation of the \NBh, this part of the discourse goes as follows: the sound is impermanent since like any other things that are made (\krtaka), i.e., impermanent, the sound gets the figurative expression such as intense or mild (2.2.13). The opponent argues that the reasoning has a deviation (\vyabhicara) since eternal things also get figurative expressions in the same ways as the things that are made do. Examples are a region of space (\akasa) or a region of \atman\ (2.2.14). The assumption is that only the things that are made (\krtaka) have parts, viz., regions, and they can be referred to. But we refer to regions of eternal, viz., non-made, things---presumably without regions---figuratively. Then the proponent replies that in the example the opponent gives, a region of eternal things (\akasa\ or \atman), the referent is not the eternal things but the fact that the union between the eternal thing and [spatially] limited substance does not fill the whole eternal thing (\emph{saṃyogasyāvyāpyavṛttitvam}). That is, the opponent's example does not annul the argument that the sound gets figurative expression just like any other things that are made. For, the expression that purportedly show eternal things being referred to in the same manner as impermanent things in fact does not refer to the eternal things. Rather, it is only a derivative (\emph{bhākta}) understanding that \akasa\ or \atman\ has a region.

Admittedly, it is hard to precisely understand what is meant by the root \emph{upa-car-}, which should be understood as ``to figuratively express,'' in this discourse. However, we can still glean that the \Naiyayika s were concerned about the expressions such as \emph{ākāśapradeśa} or \emph{ātmapradeśa} in their system adopted from the \Vaisesika s where \akasa\ or \atman\ should not have a region. It is also possible that their solution to the problem was that those are figurative or secondary expressions. This is seen in the next sentence in our text.

It is interesting that our author states that \atman\ has no regions matter-of-factly when there is no statement to the effect in old \Vaisesika\ works. \Sankara\ assumes the same. For example, in \rBSBh{2.2.17}{2.2.17} he simply states that \anu s, \atman, and \manas\ have no regions (\emph{aṇvātmamanasām apradeśatvāt}). In \rBSBh{2.3.53}{2.3.53} the topic is a region of omnipresent \atman\ delimited by a body. There he argues that, even assuming multiple \atman s, since \atman s are equally omnipresent (following the \Vaisesika\ doctrine) all the \atman s are inside all the bodies and thus it is impossible to postulate a region of ātman, even when it is delimited by a body. There is a leap of logic here. As he further discusses, it may be impossible to discern which \atman\ among all the \atman s is being affected by a body, yet it does not follow that the part of \atman\ inside a body and the rest cannot be distinguished. He seems to believe that omnipresent means having no region. He further qualifies the word \emph{ātman} with the adjective \nispradesa, and for such \atman, a region delimited by a body is ``made up (\emph{kālpanika})'' and incapable of determining a real effect. Along with the assumption that there is no regions in \atman, the mention of ``made up'' region of having no effect is similar to the discussion here. For the text of the BSBh, see p.~\pageref{BSBHTHREETHREE}} to say  ``under the condition that the mind stays with \atman'' is not rational.  Also, even if the union with a region of \atman\ is said figuratively, it will not be the cause of true yoga.  For something figurative is [ultimately] a false.

There%\pratika{kiñ cā\-nya\-t\dd\ ma\-no°}
\pratikap{p01_55}\labelh{KCMN1}\refm{50} is another [difficulty in your view].  The use of the word ``senses (\emph{indriya})'' would be pointless since [that contact with objects is] already established just by saying ``from the contact between the mind and objects.''

If%\pratika{athāpi}
\pratikap{p01_56} you\refm{51} may say---[Objection] But since the mind and objects
cannot have contact without senses, it would be irrational to negate
something that cannot be achieved, saying ``its not happening.''

It\refm{52} is not so.  Without senses, the contact between the mind and objects
does not exist.  Thus naturally, the person suppressing the contact should already know that the acquisition of objects is through senses.

There%\pratika{kiñ cānyat---sarva\-prāṇa°}
\pratikap{p01_57}\labelh{SRVPRNBH1}\refm{53} is another [difficulty in your view].  Because it is well
known that anything that breathes has pleasure and pain from the
contact between the senses, the mind and objects, and because the connection between \atman\/ and the mind is eternal, [what is meant] would be established by
merely saying ``yoga is the absense of pleasure and pain.''  Therefore the
rest of the \sutra\ would be meaningless.

There%\pratika{athāpi syād}
\pratikap{p01_58}\refm{54} might be the following objection: [Objection] If the \sutra\ says as much as ``[yoga is] the absence of pleasure and pain'' without
mentioning the contact between senses, the mind and objects, then,
since a liberated person does not have pleasure and pain, there would be an undesirable consequence of he becoming a yogin.  Therefore in order to avoid this [undesirable consequence], the rest of the \sutra\ is included.

[Answer]%\pratika{tac ca na}
\pratikap{p01_59}\refm{55} That again is not correct.  Because a liberated 
does not acquire pleasure and pain.  [And] because it is logical that
negation is done on something that has previously been acquired.  Someone
acquires pleasure and pain, and yoga is only the absence of pleasure
and pain with regard to him.  Therefore it is reasonable to
reject the rest of the \sutra.

There%\pratika{kiñ cānyat---upātte 'pi}
\pratikap{p01_60}\labelh{UPTTP1}\refm{56} is another [difficulty in the \sutra].  Even though the
word ``[that] belongs to someone who still has a body (\emph{%
  saśarīrasya})'' is adopted, it is impossible to exclude a liberated (\emph{%
  mukta}).  Because it lacks the purpose.  If yoga is the absence of
pleasure and pain, then the body does not mean anything, even when
it still exists, because it is does not produce any effect.  When we do not take the effect into consideration, there is absolutely no distinction between a liberated and not liberated.
                                % who is mukta?

There%\pratika{kiñ cā\-nya\-t---prā\-ṇa\-ma\-no°}
\pratikap{p01_61}\labelh{KCPRMN1}\refm{57} is another [problem in the \sutra].  The part ``[which]
follows restraining of breaths and mind'' is not appropriate, either.
Because the effort that is holding the restraining is inherent to
\atman.  Even though the mind does have activities, without being
connected with senses, it is impossible to restrain the mind and
breaths.
                                % manasaḥ kriyāvattva?
                                % prayatnasya ātmasamavetatva?

Also,%\pratika{indriyasaṃyoge}
\pratikap{p01_62}\labelh{INDRYSMYG1} if\refm{58} the mind is connected to senses, it is not reasonable to
say ``its (contact's) not beginning.''  Also, since the mind is immersed in
restraining breaths, it is impossible for it to not begin [contacting
with senses].  Also, if there is no contact between the senses and
the wind, there is no restraining of breaths.  Also, if breaths are 
restrained, there is no yoga.  Because the cause of restraining breaths
is the relation with the action in the mind through the holding effort
residing in \atman.

Then%\pratika{athāpi syāt---yogaḥ samādhir iti}
\pratikap{p01_63}\labelh{ATHPSY1}\refm{59} you might say---[Objection] Yoga is \samadhi.  [Answer] %%\pratika{tacca na}
% \pratikap{p01_64} 
That is not right.  Because [saying so] does not have any effect.  Again, \atman\ and manas are eternally assimilated.  Also, it is said that \samadhi\/ is stasis.\footnote{See page \pageref{station_immersion}.}

In%\pratika{tasmād}
\pratikap{p01_65} conclusion\refm{60}, it is appropriate to say ``And that (\samadhi) is
a property of the mind common to all stages.''

[Objection]%\pratika{yady evaṃ}
\pratikap{p01_66}\refm{61} If, as you say, yoga is \samadhi, and it is a
property of mind common to all stages, then it is [already]
established for all the living creatures without effort, because [it]
includes fixated [mind], etc.  Accordingly, the definition (\emph{lakṣaṇa})
is meaningless because it is established without effort, just
like inhaling and exhaling.  Thus the instruction of yoga (\emph{yogānuśāsana}) also becomes meaningless.

Answering%\pratika{atra vikṣipte}
\pratikap{p01_67}\refm{62}\labelh{ATRVKSP} this, [the \Bhasyakara]
states---\bht{among those, if the mind is scattered, the \samadhi\/ that
  is dominated by being scattered does not belong to the side of yoga.}
For, [\samadhi] that includes ``fixated [mind],'' and so on, is not intended to be on the side of yoga.  Namely, \samadhi\/ that includes [the condition of] being fixated, etc., is not capable of illuminating an object as is.  Because the action of being fixated is predominant in it.  By merely negating [the stage of] scattered (\emph{vikṣipta}), according to the maxim of defeating the representing
wrestler, the stages fixated (\emph{kṣipta}) and stupefied ({mūḍha}) are also negated.  \emph{Samādhi\/} of being scattered (\emph{vikṣepasamādhi}) represents [other types of \samadhi\/ that do not belong to the side of yoga] because it [still] has an adequate qualification [to be included on the side of
yoga].  Namely, the mind which is in the state of scattered is able
to acquire desirable object, for it has not fallen into [a
particular] side.  But stupefied [mind]---by means of attaching [it] to an object, for example---or fixated [mind]---by means of detaching [it] from undesirable object, for example---cannot be led to anywhere else.  [The \Bhasyakara] says ``the side of yoga''---Even though it is still
\samadhi, since it does not produce an effect [expected from yoga], he says ``the side.''  For example---even though a walking man stops at every step, since it
does not cause the effect of stopping, it is not called
standing.\footnote{Cf. {Nāgārjuna}\/'s famous discussion on going/walking (\gati)
  in {Madhyamikakārikā}s \citep[kk.~2.1--25]{mmk}.  Note that
  while {Nāgārjuna} tries to prove there is no such thing as
  ``going,'' the point the author of the YVi making is the
  opposite.}
  
With%\pratika{`vikṣepo°}
\pratikap{p01_68}\refm{63} the intention [of specifying] a particular kind of \samadhi\/ that is in the state of being scattered, [the \Bhasyakara\ refers to] the
quality of being ``dominated by being scattered.''  Because, on the other hand, [that yoga is \samadhi] primarily applies to all the stages equally, for [the \Bhasyakara] has stated ``common to all stages.''%. Since it is stated For, on the other hand, the property of being common to all the stages main
%[characteristic of immersion] is that it is ``common to all stages''.
%[It is main characteristic] because it exists at all stages.

If%\pratika{yady upa°}
\pratikap{p01_69}\refm{64} one might question---If being dominated is the cause of not belonging to the side of yoga, now, in the case of focused [mind] there would be the dominance of being focused.  Then [the \samadhi\/ in the stage of focused] does not belong to the side of yoga.  In order to answer this, [the \Bhasyakara]
states---\bht{But [\samadhi] when [the mind is] focused}. %\footnote{The above two
%  paragraphs may still contain errors in readings. See footntoe
%  \fnref{note:avrtti} on page \pageref{note:avrtti}.}
%
At\refm{65} the stage of being focused, there is no domination by the stage.  Why?  Because defects (\emph{kleśa}) and \karman\ are not powerful [at that stage].
Because when the power of defects and \karman\ are powerful, [the states of] diversion and so on arise.   \emph{Samādhi\/} at the stage of being focused is \bht{the one that illuminates} = makes \bht{the object} understood \bht{as is} = as
itself. [The \Bhasyakara] uses the expression ``as is (\emph{%
  bhūta}),'' because, if there is even a trace of non-exactness (\emph{a\-yathābhūta}), the knowledge is not for a yogin.%
  \footnote{%
	There is some suspicion that sentence is corrupt. The first word \emph{ayogyarthajñānam} is already suspect. One would rather expect \emph{ayogyārtha°}. Also, the appearance of the word \emph{jñāna} is somewhat confusing. Knowledge has not been discussed in relation to \cittabhumi s and therefore appears out of context. One could also conceive that the compound was in fact two words, \emph{ayogy(a/ā)rthañ jñānam}; the conjunct character for \emph{ñjña} is complicated and can easily look like \emph{jña}. Or it is even possible that the word \emph{jñānam} is a corruption from something entirely different. Furthermore, the word \emph{gandhita} is not well attested in the sense our author apparently intends. \Sankara\ uses the word \emph{āgandhita} or its negated form \emph{anāgandhita} in the sense ``not left with any trace of,'' and this is one of his idiosyncrasies. We could consider the possibility that instead of \emph{ayathābhūtatvagandhitam}, the text in fact read \emph{ayathābhūtānāgandhitam}. Even after considering these emendations, the text does not appear very coherent.%
  }  [\emph{Samādhi\/} at the stage of being focused] \bht{destroys}---makes the five fold \bht{defects}, \ignorance, etc.,\footnote{See \rYS{2.03}{2.3} (\emph{avidyāsmitārāgadveṣābhiniveṣāḥ kleśāḥ}).} destroyed---
  \bht{[It] loosens the bindings of \karman}---[the compound \emph{karmabandhanāni\/} is a \KD\ compound in which] \karman s are the bindings [connected in the \emph{sāmānādhikaraṇya\/} relationship]. The [\karman s] are good, bad, and mixture of good and bad; or, [the compound \emph{karmabandhanāni\/} is a \TP\ compound that means] those originated from \karman s and they are the birth, etc., i.e., the bindings---\bht{[and it] loosens} = makes something loose. %
%[``Loosens'' means] to make something loose.\footnote{This sentence first
%  seems awkward since there is ``{yadi karmajāti} in the
%  interpretation.  This, however, becomes clearer after consulting
%  relevant sections of the PYS.  First, for the first member of the
%  compound {karma-bandhana}, which is being interpreted, {%
%    karman}, see YS \ldots, YBh \ldots.  For {bandhana}, there
%  is a mention in YBh i.e.\ on \rYS{1.24---}{}.  The three bondages is
%  explained as ``{prākṛta-vaikṛta-dākṣiṇāani
%    bandhanāni}'' by the author of the YVi while interpreting the PYS
%  1.24.  For the misleading compound ``{%
%    dharmādharmavimiśrāṇi}, it is believed that it should be taken
%  as dvandva compound ({ekaśeṣa}).  The author of the YVi
%  holds an opinion that karmans are of three kinds.  This is seen at
%  the interpretation on 1.5 (page ??), 4.7--8.  On the other hand,
%  according to the \rYS{4.07@4.7}, karmans of people but yogins are of three
%  kind.  Anyways as its being in the same case and number with {%
%    karmāṇi} it should be taken as explanation of the word.  With
%  regards to a following word {janmādīni}, it seems natural to
%  think that the author of the YVi is talking about the karmans of
%  people but yogins.  As the word {karmajāti} also appears in
%  the YBh 4.7, and there is another kind of karman the reading {%
%    yadi karmajāti} can be understood as it is talking about the
%  condition under which {vipākas} such as the birth ({%
%    janman}), and so on are produced.  The effects are birth, length of
%  life, experiences, according to the \rYS{4.9}.}
 \bht{[It] directs [the yogin] to suppression (\nirodha)}---orients him toward
[suppression].  \bht{It is called ``conscious (\emph{samprajñāta})''} by
the master(s).\footnote{\label{no_yogah}Note that the author of the YVi reads the \YBh\ \emph{\ldots\ nirodham āmukhīkaroti\dd\ sa samprajñātaḥ}, while most versions of the YBh reads \emph{sa samprajñataḥ yogaḥ}.  See \citet[3]{maas:2006}.  As a result of this lack of the word \emph{yogaḥ}, the author of the YVi considers the subject of the YBh to be \emph{samādhi\/} throughout after \emph{yogaḥ samādhiḥ}.  Hence, the terms \emph{samprajñāta\/} and \emph{asamprajñāta\/} can only modify the word \emph{samādhi}.  Consequently he considers the next \sutra\ as the definition of \emph{asamprajñāta-samādhi}.  See the next paragraph.}

The%\pratika{sa ca vitarkā°}
\pratikap{p01_70}\refm{66} \bhasya\ continues---\bht{And that is accompanied by \emph{vitarka}, or by \emph{vic\=ara}}, and so on.  [The reason why the \bhasya\ mentions this] is because \emph{samprajñāta\/} occurs first.  On the other hand, when it comes to stating the definition, since  \emph{asamprajñāta\/} is more important, it is reasonable to state the definition of it here---\bht{this} \emph{samprajñāta} \bht{we will teach later}.%\footnote{There could be other interpretation
%  for this section.  Coupled with difficulty in distinguishing the
%  text from the \bhasya\ and reconstructing the text of the YVi
%  itself, the sentence translated here might still have problems.
%  Consult the following paragraph where the author of the YVi exhibits
%  almost identical interpretation in slightlymore elaborately.  The
%  place where the \bhasya\ is referring to teach {%
%    samprajñāta} is \rYS{1.17}.}

For%\pratika{itaś cā°}
\pratikap{p01_71}\refm{87} the following reason, too, the definition of the \emph{%
  asamprajñāta\/} should be given at this very place.  Without
having the \emph{samprajñāta} prerequisite, the \emph{asamprajñāta} can
be realized by either excellent detachment or practicing the notion
of termination.\footnote{See \rYS{1.12--18}{1.12--18}.}  In order to show this [the
definition of the \emph{asamprajñāta} comes first].  On the other
hand, if the definition of \emph{samprajñāta\/}
was stated here, and if the definition of \emph{%
  asamprajñāta} was stated later, there would be a suspicion that one
qualifies for \emph{asamprajñāta-samādhi\/} with \emph{%
  samprajñāta\/} as prerequisite.  Therefore [the \Bhasyakara] states \bht{we will teach later}.


%
% \chapter{The ending of the YVi}\label{ENDOFTML}
% \renewcommand{\rightmark}{\MakeUppercase{The ending of the YVi}}
% \thispagestyle{plain}
% % \section{The ending of the YVi}\label{ENDOFTML}
The ending of a text, especially the \colophon s in its manuscripts, can give us various information about the text itself, its author, provenance, transmission, etc. This applies also to the YVi. Its closing stanzas have already attracted some scholarly attention.\footnote{See \citet[420]{rukmani:yoga1}.}
Unfortunately, the ending of the YVi has not been reproduced properly in the 1952 edition. %What they presented was almost fictional. They altered the order of verses and introduced quite different wordings in some places.

%Having access to accurate data is prerequisite for investigation. %It is regrettable that such an important part as the ending of the text was not properly reported in the 1952 edition, possibly leading to arguments without basis. 
Here is an attempt to give a better picture of the ending of the YVi. I will first present transcripts from the manuscripts \Tm\ and L\@.\footnote{I am not concerned with the remaining manuscripts, \Td, M or A here. We know them to be direct descendants of either \Tm\ or L\@. (See pp.~\pageref{sec:stemma}\,ff.) Unlike the other parts of this volume where editions of the YVi was presented, the purpose of this section is to discuss the transmission history of the YVi. It was useful to include readings in those modern manuscripts when the goal was to reach the most plausible reading of the YVi (in its intended form). Readings resulting from the intelligence and knowledge of the modern copyists, particularly that of M, are worth consulting and reporting.} Then I will present an edition, accompanied by a translation and annotations, of the part that is common to both manuscripts. Then I will proceed with editions of additional endings (the \colophon\ or colophon stanzas) that are unique to \Tm\ and L each. A note on the presentation of the ending in the 1952 and discussions on the layers in the endings will follow. 

\section{Transcripts}
Here are transcripts of \Tm\ and L toward the end. The commentary proper ends with \ldots\ \emph{vedāraṇyakavad ity om iti} \citep[369]{yvi}.\footnote{After this, both manuscripts have the decorative sign that appears when the commentary on a \sutra\ ends. Here the sign is transcribed as \textdb{❋}. Cf.~p.~\pageref{punctuation}.} 

Conventions generally follow those that have been adopted to report readings in the critical apparatus.\footnote{See pp.~\pageref{sec:symbols}\,ff.\ for more complete explanations.} The ones used here are: 
  \ccs\ and \cce\ to represent canceled elements; 
  \ins\ and \ine\ to represent elements added later;
  [ and ] to represent physically lost portions (if signs are partially lost but recoverable with confidence, the brackets enclose the reading and if signs are completely lost, they are represented by center dots (\mist{1}) whose number represents the estimated number of lost signs);
  \{\ and \}\ to represent modifications to the text that do not necessarily involve simple deletion and insertion.
  
\newcommand{\HTG}{\textdb{❋}}
\subsection{The end of \Tm}
\Tm\ ends in folio 125 recto (line numbers are in parentheses):
\begin{quote}
	(1) ntarasvarūpatatvanirddhāraṇena tantrāntarīyakatvasva\-rū\-pa\-ni\-rā\-ka\-ra\ccs n\cce \ins ṇ\ine ena ca kaivalya\lost{4}vedi[tavyā\mist{1}]\-ta\-[tro\mist{3}maṃgalārtthaḥ] prayuktaḥ parameśvara\-nā\-māṃ\-ki\-ta\-ma\-sta\-ka\-tayā \ins ta\ine\ccs s ta\cce ntrapra\-ca\ins y\ine ārttham śāntyarttham ve\-dā\-ra\-ṇya\-ka\-va\-d ity om iti \HTG\ %
	(2) oṃkāro yasya vaktā samacarata phalaiḥ ka\-rmma yasmād aśeṣa\-n niṣkarmakleśapāko gha\-[ṭ]a\-ya\-[ti] sakalaṃ yaḥ phalen\{a ā $\rightarrow$ ā\}%
	\footnote{Due to the fact that in the Malayalam script, the diacritical sign for the vowel \emph{ā} is written as an independent letter, what happens here is that the diacritical sign for \emph{ā} has replaced what appears to be an independent vowel \emph{ā}. So, the manuscript originally had \emph{phalena ā°} and was corrected to \emph{phalenā°}.}%
	kri\-yāṇā\ccs m īśvaro\cce \ins m īś\ine ā\-na\-m īśvaro ya sthiti\-bhavanidhanapra\-kri\-yā\-ṇāṃ vi\ccs p\cce \ins dh\ine ā\-tā ddhyā\-ya\-n naś śu\-kli\-mā\-naṃ vya\-pa\-nu\-da\-tu ta\-rāṃ kṛṣṇimānaṃ sa kṛṣṇaḥ a\-ne\-ka\-pha\-% 
	(3)ṇa\-ra\-tnau\-gha\-vi\-dyo\-ti\ccs X\cce%
	\footnote{I cannot decipher the sign marked to be cancelled. It appears to have a \emph{d} element as first part of the conjunct letter. The rest of the conjunct letter could have a \emph{c}, \emph{ṣ}, and/or \emph{m} element. It could have the vowel \emph{i}, but the circle-like part of the letter is too small for the usual way of marking a syllable with the \emph{i} vowel.} %
	dyu\-na\-bho\-di\-śe yo\-gī\-ndrā\-ya pha\-ṇī\-ndrā\-ya tat pa\-ta\-ñja\-la\-ye namaḥ va\-da\-nā\-hi\-ta\-pū\-rṇṇa\-ca\-ndra\-kaṃ gurum īśānam a\-bhū\-ti\-bhū\-ṣa\-ṇa\-m pra\-ṇa\-mā\-my a\-bhu\-ñjaṃ\-ga\-saṃ\-gra\-ha\ccs m\cce\ins ḥ\ine\ bha\-ga\-va\-tpā\-da\-m a\-pū\-rvva\-śaṃ\-ka\-ram \HTG\ go\-vi\-nda\-bha\-ga\-vatpū\-jya% 
	(4)pā\-da\-śi\-ṣya\-sya pa\-ra\-ma\-haṃ\-sa\-pari\-vrā\-ja\-kā\-cā\-ryya\-sya śaṃ\-ka\-ra\-bha\-ga\-va\-taḥ kṛtau pā\-ta\-ñja\-la\-yoga\-śā\-stra\-bhā\-ṣya\-vi\-va\-ra\-ṇe catu\-rtthaḥ pādaḥ \HTG\ yā\-trā\-kṛ\-tas tri\-ja\-ga\-tāṃ ya\-da\-ci\-ntya\-śa\-kti\-le\-śān da\-śā\-hur iha mī\-na-mu\-khā\-va\-tā\-rā\-[\ccs t\cce\ins n\ine] kle\-śo\-% 
	(5)\-dbha\-va\-tri\-vi\-dha\-tā\-va\-va\-ra\-mpa\-rā\-rttās tan nā\-ga\-nā\-tha\-śa\-ya\-naṃ śa\-ra\-ṇaṃ vra\-jā\-maḥ yo\-ge\-na ci\-tta\-sya pa\-de\-na vā\-cā\ins ṃ\ine ma\-la\-śa\-rī\-ra\-sya ca vai\-dya\-ke\-na yo pā\-ka\-ro\ccs kta\cce\ins t ta\ine m pra\-va\-ra\-mu\-nī\-nām \ccs v\cce%
	\footnote{This \emph{va} is clearly marked to be cancelled but I do not find the obvious replacement, \emph{pa}, in the manuscript. Also, the diacritical sign for \emph{ā}, too, should be cancelled, but I do not see a trace of such  cancellation.}%
	āta\-ñja\-li\-m prā\-ñja\-lir \ccs X\cce ā%
	\footnote{What appears to be written here is the sign \emph{ra}, followed by what appears to be the diacritical sign to mark the following consonant sign with the vowel \emph{e} (which is cancelled), and another such sign to mark the preceding consonant with the vowel \emph{ā}. What probably happened here was that the scribe started to write \emph{rane} or \emph{rano}, but realized the mistake and wrote \emph{rā na}.}%
	na\-to smi \HTG\ \HTG\ 
	(6) śrīmat pa\-ta\-ñja\-li\-ma\-ha\-rṣi\-vara\-pra\-ṇī\-ta\-\ccs ṃ\cce śrī\-yoga\-bhā\-ṣya\-hṛ\-da\-ya\-pra\-ka\-ṭī\-kṛd etat śrī\-śāṃ\-ka\-ra\ccs ṇa\cce\-ṃ vi\-va\-ra\-ṇaṃ sa\-ma\-va\{tti $\rightarrow$ rtti\}śa\-tro \ccs [g $\rightarrow$ rgg]ā\cce \ins gā\ine rggyo%
	\footnote{\label{GATORGG}There are two stages of corrections. First the original scribe wrote \emph{°śatro gārggyo} as it should read. Someone altered \emph{gā} to \emph{rggā} by adding strokes to the original sign for \emph{ga}. This made the word \emph{samavartiśatro} in the vocative case to the \emph{samavartiśatroḥ} in the genitive case. Then again the whole of the modified \emph{rggā} was cancelled and in the margin \emph{gā} was inserted. These two stages of corrections may have been done by the same person. All these alterations are not inked.} %
	vya\-lī\-li\-kha\-m aha\ccs tvi\cce \ins n tva\ine\ccs ñca\cce da\-vā\-pti\-kā\-maḥ \HTG\ \HTG\ \ccs va\cce\ins pa\ine ra\-kro\-ḍaḥ \HTG\ \HTG\ \HTG
\end{quote}

\subsection{The end of L}
And here is a transcript of L (folio 138 recto):
\begin{quote}
(4) \ldots\ tatroṃkāro maṃgalārtthaḥ prayuktaḥ parame\-śva\-ra\-nā\-mā\-ṃki\-ta\-mastakatayā tantrapracayārtthaṃ śāntyartthaṃ vā ve\-%
(5)\-dā\-ra\-ṇya\-kavad ity om iti \HTG\ oṃkāro yasya vaktā samacarata phalaiḥ karmma yasmād aśeṣan niṣkarmmakleśapāko gha\-ṭa\-ya\-ti sakalaṃ yaḥ phalena kriyāṇām īśānām īśvaro ya sthi\-%
(6)\-ti\-bha\-va\-ni\-dha\-na\-pra\-kri\-yā\-ṇāṃ vi\-dhā\-tā dhyāyan naḥ śuklimānam vyapanudatu tarām kṛṣṇimānam sa kṛṣṇaḥ pha\-la\-ra\-tnau\-gha\-vi\-dyo\-ti [dyu]na\-bho\-di\-śe yogendrāya phaṇīndrāya [ta\mist{1}ta]ñja\-%
(7)\-la\-ye namaḥ padanāhitapū\-rṇṇa\-ca\-ndra\-kaṃ gurum īśānam abhūtibhūṣaṇam praṇamāmy abhu\-ñja\-ṃga\-saṃ\-gra\-ham bha\-ga\-va\-tpā\-da\-pū\-rvva\-śaṃ\-ka\-ram \HTG\ go\-vi\-nda\-bha\-ga\-va\-tpū\-jya\-pā\-da\-śi\-ṣya\-[sya \mist{1}]\-ra\-ma\-haṃ%
(8)sa\-parivrājakācā\-ryya\-sya śaṃ\-ka\-ra\-bha\-ga\-va\-taḥ kṛtau pā\-ta\-ñja\-la\-yo\-ga\-śā\-stra\-bhā\-ṣya\-vi\-va\-ra\-ṇe ca\-tu\-rttha\-ḥ pādaḥ \HTG\ sa\-mā\-pta\-ñ cedaṃ vivaraṇam \HTG\ yātrākṛtas tri\-ja\-ga\-tāṃ ya\-da\-[ci]\-ntya\-%
(9)\-śa\-kti\-leśān daśāhur iha mīnamukhā[\mist{3}n] kleśo\-t\-bha\-va\-tri\-vi\-dha\-tā\-pa\-pa\-ra\-mparā\-rttās tan nā\-ga\-tnā\-tha\-śa\-ya\-naṃ śa\-ra\-ṇam pra\-jā\-maḥ 
yo\-ge\-na ci\-tta\-sya pa\-de\-na vācāmmalāśarīrasya tu %
(10) vai\-dya\-ke\-na yo pākarot tam pra\ccs vara\cce munīnām pata\-ñja\-li[m] prāñjalir ā\-na\-to smi|\footnote{As I write on p.~\pageref{punctuation}, hardly any punctuation is used in either \Tm\ or L\@. Here, however, L has what appears to be a \danda.} jitam pata\-ñja\-li\-muninā yena śreyo\-rtthi\-sā\-rttha\-sa\-syā\-nām vihito harṣas tā\-pa\-tri\-ta\-ya\-hṛtā %
(138v1) dha\-rmma\-me\-ghe\-na| yas ta\-tā\-na bha\-va\-mārgga\-laṃ\-ghi\-nām kle\-śa\-ka\-rmma\-ma\-ya\-gha\-rmma\-tu\-tta\-ye dha\-rmma\-megha\-mu\-kha\-yogatoya śrī\-ke\-śa\-va\-pra\-kā\-bha\-ṭṭā\-ra\-kā\-ṇāṃ śiṣyasya sa\-tyā\-na\-nda\-sya iya\-m pādañjalaṭī%
(2)kā \HTG\ \HTG\ \HTG
\end{quote}

\section{An edition of the common part}
Quite clearly, the ending is shared between the two manuscripts up to some point; the shared part must have been in the common ancestor,\footnote{Note that this does not necessarily mean that \emph{only} the shared part was in the common ancestor. There are possibilities that the common ancestor had more and that some of the continuation made it to either or neither of \Tm\ or L\@. It is, however, highly unlikely that the text shared in both \Tm\ and L originated independently.} which I consider to have been the exemplar of \Tm.\footnote{See section ``\titleref{extm}'' on pp.~\pageref{extm}\,ff.}  Here is an edition, translation of the common part along with some comments.

\subsection{The end of the commentary}
The YVi as a commentary ends as follows:
\begin{quote}
{\dn तत्रौंकारो मंगलार्त्थः प्रयुक्तः। परमेश्वरनामांकितमस्तकतया तन्त्रप्रचयार्त्थं शा\-न्त्य\-र्त्थं वा वेदारण्यकवदित्योमिति॥}

There the syllable \OM\ that has the auspicious meaning is used. [It is used,] since it is the foremost of things bearing the name of \Isvara, so that the system (\tantra) [expounded in this \sastra] grows, or for the sake of tranquility, analogous to the \Veda s and the \Aranyaka s [ending with the syllable \OM]. Therefore [the \sastra\ ends with] ``\OM.''/Thus \OM. End.
\end{quote}
This concludes the commentary. It appears that the author loaded several meanings in the last words of his text. The last part \emph{om iti} can be seen as the quote of the end of the root text, but it is at the same time the conclusion of his own text, too. In that sense, readers understand that he, too, placed the syllable \OM\ at the end in the sense he has just explained. This is the same technique the author uses at the beginning of the commentary. He starts the commentary with \emph{athetyādi pātañjalayogaśāstram. tasya vivaraṇam ārabhyate}.\footnote{See p.~\pageref{FIRSTSTC} for the text and p.~\pageref{ATHAATHATRL} for the translation.} He uses the very first word of the root text as the first word of the commentary. Also, the very last word, \emph{iti} serves several functions. One is to signal the end of a quotation from the root text. Another is as part of the quote of the very last word of the root text; we naturally expect the word \emph{iti} at the end of the root text. The third is to mark the end of the commentary itself. The part \emph{parameśvaranāmāṅkitamastakatayā} alludes to the statement in the commentary on \rYS{1.27}{1.27} where the syllable \OM\ is said to be the favorite name of \Isvara.\footnote{See p.~\pageref{FAVNAMEED} for the edition and p.~\pageref{FAVNAMETRL} for the translation, and n.~\ref{FAVNAME} on p.~\pageref{FAVNAME} for parallel passages in the \BAUBh\ and the \ChUBh.}

\subsection{First common concluding stanza}
\label{krsnaverse}
This conclusion is followed by a stanza in the \sragdhara\ meter:
\begin{verse}
%oṃkāro yasya vaktā samacarata phalaiḥ karmma yasmād aśeṣan 
%  niṣkarmmakleśapāko ghaṭayati sakalaṃ yaḥ phalena kriyāṇām| 
%īśānām īśvaro yaḥ sthitibhavanidhanaprakriyāṇāṃ vidhātā 
%  dhyāyannaḥ śuklimānam vyapanudatu tarāṃ kṛṣṇimānaṃ sa kṛṣṇaḥ|
\textdn{ओंकारो यस्य वक्ता समचरत फलैः कर्म्म यस्मादशेषन्निष्कर्म्म\kle शपाको घटयति सकलं यः फलेन क्रियाणाम्। \\
ईशानामीश्वरो य स्थितिभवनिधनप्रक्रियाणां विधाता 
  ध्यायन्नः शुक्लिमानं व्यपनुदतु तरां कृष्णिमानं स कृष्णः॥}\\
{\footnotesize pāda b: phalena kri° L] phalena ākri° \Tmac, phalenākri° \Tmpc\\
pāda c: īśānām īśvaro L] īśvaro \Tmac, īśānam īśvaro \Tmpc; vidhātā L\Tmpc] vipātā \Tmac

}
The syllable \OM\ signifies him; from him all the actions appeared along with [their] fruits; \\ He, [himself] having no actions, impurities or [their] ripening, unites everyone with the fruit of actions; \\
  He is the \inidx{lord} of lords; he performs the procedures of stability, rise, and end.\\ May that \Krsna, thinking of my good, utterly remove [my] wrong.
\end{verse}
This stanza applies a similar playful composition technique to that just seen in the end of the commentary. The author finished the commentary with the syllable \OM. Now a closing stanza immediately starts with the same word. This is deliberate and the mind behind this choice appears to be the same as that which chose to finish the commentary with the syllable \OM. This stanza is thus probably integral to the bulk of the commentary.

Much of the content in this stanza also agrees with what is put forward in the \Isvara\ section of the commentary and the opening stanzas.\footnote{Both are dealt with in this volume. For the opening stanzas, see pp.~\pageref{edition01}\,ff.\ for the text and pp.~\pageref{OSTZ1}\,ff.\ for translation. The \Isvara\ section is the main part of this volume.} The first segment of the stanza refers back to \rYS{1.27}{1.27}. The second and third segment (about \Isvara's being the legislator and enforcer of the law of \karman s) strongly resonate with the first \emph{pāda} of the first stanza in the opening of the whole text (see p.~\pageref{FVFPEDN} for the text and p.~\pageref{FVFPTRL} for its translation). It is interesting that \Isvara's being the creator of the laws of \karman, rather than his being the author of the Vedas that prescribe them, is more or less clearly stated here. In the opening stanza and in the \Isvara\ section of the commentary, this point was ambiguous. Also notable is the similarity in wording to that of the \BhGBh. See n.~\ref{NONFSFP} on p.~\pageref{NONFSFP} for the wording in the \BhGBh. In that connection, the explicit naming of \Isvara\ as \Krsna\ is also noteworthy. The author of this commentary repeatedly hints that he considers \Isvara\ to be \Visnu, especially the \Bhagavat\ of the \BhG\ (see pp.~\pageref{narayana}, \pageref{gita1155}, \pageref{omvisnu}, \pageref{NRYNN25}, \pageref{NNRYN26}, \pageref{VISNU27}, and \pageref{VSNSM1}), who indeed is \Krsna. The last \emph{pāda} further employs a word-play, between black/bad (\emph{kṛṣṇiman}) and \Krsna, alluding to terminologies in the text that has been commented; \rYS{4.07}{4.7}, for example, refers to actions as white or black (and their mixture).

In the last pāda the author asks \Isvara, who is revealed to be \Krsna, for his forgiveness. Mentioning \Isvara\ or \Krsna\ as someone who can remove bad actions attracts attention. I am not aware of any passage in the YVi that expresses the view that \Isvara\ can remove bad \karman. Hence this last \emph{pāda} sounds a little like being addressed to a human king, a judicial authority who, in reality, can grant pardons. %The author sounds as though soliciting \Krsna\ to be tolerant in judging actions. 

The whole stanza gives the impression that it is about a governing personage. This impression comes from the lack of the mention of \Isvara's being omniscient or omnipotent---the two aspects the author of the YVi usually mentions along with his being the lord of all. The last pāda, by introducing the view not put forward in the text, reinforces that impression. As a result, this stanza can be read as a prayer to \Krsna\ the deity as well as a plea to \Krsna\ a human ruler.\footnote{This point catches my attention because I propose the theory that \Sankara, the author of the \BSBh, lived under the rule of the Rāṣṭrakūṭa king, Kṛṣṇa I \citep{harimoto:2006}.} The author has already shown some literary techniques in this stanza. It would not be surprising if he meant this stanza to be read in these two meanings.

\subsection{Second common concluding stanza}
\label{patanjaliverseone}
Then \Tm\ and L records a śloka:
\begin{verse}
%anekaphaṇaratnaughavidyoti dyunabhodiśe 
%yogendrāya phaṇīndrāya tat patañjalaye namaḥ||

\textdn{अनेकफणरत्नौघविद्योतिद्युनभोदिशे। \\
योगीन्द्राय फणीन्द्राय तत्पतञ्जलये नमः॥}\\
{\footnotesize 
pāda a: anekaphaṇa° \Tm] phala° L\\
pāda c: yogīndrāya \Tm] yogendrāya L \\
pāda d: tat patañjalaye \Tm] [ta\mist{1}ta]ñjalaye L

}
To him who irradiate the heaven, sky, and directions with multitude of jewels residing in [each of] his many hoods \\
Thus salutation to [him,] the \inidx{lord} of yogins, the lord of serpents, \Patanjali.
\end{verse}
The above translation is tentative. As presented, I do not think the Sanskrit text of the stanza is intelligible. Yet I cannot offer a better solution to make the text meaningful. Either the transmission of this stanza was bad or it was not well composed in the first place.

In pāda a, I have adopted the reading \emph{anekaphaṇaratnaugha°} found in \Tm. L does not have \emph{aneka°} and hence the first line is short of three syllables there. The 1952 edition,\labelh{EDITALT0} based primarily on the manuscript M (an \inidx{apograph} of L) supplied \emph{pṛthivī°} in the beginning of pāda b, resulting in the reading \emph{phaṇa\-ratnau\-gha\-vidyoti [pṛthivī]dyunabhodiśe} for the first line.\footnote{L, and hence M, read \emph{phalaratnau°} rather than \emph{phaṇaratnau°}. That the 1952 edition still reads \emph{phaṇaratnau°} is an emendation.} Although I do not think this was a successful emendation, it was a clever one. When one sees \emph{dyaus} and \emph{nabhas}, it is natural to expect something that means ``earth.'' I do not adopt this reading because we do have a reading from \Tm\ that satisfies metrical requirements. In addition, the word \emph{diś} in the compound might have meant the same thing, the earth. 

The reading adopted here, \emph{anekaphaṇaratnaugha°}, is still not completely satisfactory because \emph{aneka} and \emph{ogha} seem semantically redundant. I cannot eliminate the possibility that the word \emph{aneka} in the beginning of the line was an emendation on the part of the scribe of \Tm\ (\Tm\ is an \inidx{apograph} of the common ancestor of \Tm\ and L). 

In pāda c, while L reads \emph{yogendrāya}, \Tm\ reads \emph{yogīndrāya}. I have adopted more common word \yogindra\ of \Tm. If we applied the principle of \emph{lectio difficilior}, perhaps \emph{yogendrāya} would be more likely to have been in the common ancestor. Given how syllables with the vowel \emph{e} is written in the Malayalam script, it is hard to imagine how \emph{gī} or \emph{ge} becomes the other by misreading. If, then, the change was intentional, it would be from \emph{ge} to \emph{gī} (\yogindra\ being more common Sanskrit word). Note, however, that being in the common ancestor does not necessarily mean that the stanza was originally composed with that reading. I have decided to adopt the word that would have been used by an author who knew standard Sanskrit; that is, I have followed the probable judgement of the scribe of \Tm. Nonetheless, as a result of this decision, the constituted reading is more or less that of \Tm. This is quite a contrast to the previous stanza where the scribe of \Tm\ was apparently struggling to record correct readings (see the variants for the previous stanza). If there is a trend that L was more successful in reproducing readings in the common ancestor, more acceptable readings in \Tm\ might in fact be a result of conscious alterations. Note that \emph{phaṇa\-ratnau\-gha\-vidyoti} that starts the stanza in L is in itself a good pāda of \sloka. And there is a sign that I cannot decipher in \Tm\ after \emph{°vidyoti}. Perhaps, despite metrically more correct or using more common words, the readings adopted from \Tm\ may still be far from what was originally intended.

I have further problems with this stanza. One of them is the word \emph{tat} in the second line. Again, a pronoun that points to a missing entity. It cannot be the first member of a compound (or can it?) affixed to a proper name. As an independent word, what we need is \emph{tasmai} that does not fit metrically. As the last resort, I have translated it as an adverb, meaning ``therefore.'' It still is preferable if the preceding line had a relative pronoun such as \emph{yat}, \emph{yatas} or \emph{yasmāt}. We cannot easily supply it as an emendation in the previous line.

It to me appears that this stanza is a remnant of badly transmitted text. It could have constituted of more than just 4 pādas. One possibility is that the original text was a sort of \emph{nirukti} of the name \Patanjali, trying to explain that it constitutes of elements of the verb \emph{pat-} (to fall) and \emph{jalin} (possessing water), perhaps in the sense that he rains down gems that shine the universe. The association between snakes and rain is a universal topos. 

\subsection{Third common concluding stanza}
\label{THETHIRDST}
Next we have a stanza in the \viyogini\ meter, dedicated to the first \Sankara:

%vadanāhitapūrṇṇacandrakaṃ gurum īśānam abhūtibhūṣaṇam|
%praṇamāmy abhujaṃgasaṃgraham bhagavatpādam apūrvvaśaṃkaram|| (Viyoginī)

\begin{verse}\label{SANKARASIVASTANZA}
\textdn{वदनाहितपूर्ण्णचन्द्रकं गुरुमीशानमभूतिभूषणम्।\\
प्रणमाम्यभुजंगसंग्रहम् भगवत्पादमपूर्व्वशंकरम्॥}\\
{\footnotesize 
pāda a: vadanāhita° \Tm] padanāhita° L\\
pāda c: abhujaṃga° L] abhuñjaṃga° \Tm; °saṃgraham L\Tmac] °saṃgrahaḥ \Tmpc \\
pāda d: °pādam apūrvva° \Tm] °pādapūrvva° L

}
The one to whose face the full moon is attached, the master, the ruler, the one not adorned by ashes\\
I salute him, without a collection of snakes, the Bhagavatpāda, the unprecedented \Sankara.
\end{verse}
This stanza is interesting in that it emphatically distinguishes \Sankara\ the teacher from \Siva, one of whose popular epithets is \Sankara. The first attribute mentioned here is the full-moon like face of \Sankara, a common description of a beautiful face. Its wording \emph{vadanāhitapūrṇacandra}, literally ``to whose face the full moon is attached'' is a contrast to \Siva, to whose head the crescent moon is attached. The expression \emph{abhūtibhūṣaṇam} appears to be a straightforward negation of \Siva's attribute of wearing ashes. However, there can be more positive interpretations such as ``not wearing any ornament,'' or ``to whom not owning wealth is an ornament.'' % the last one is Haru's suggestion.
This should be a positive attribute for a renouncer. Another contrast to \Siva, \emph{abhujaṃgasaṃgraha}, may be a play on the two major meanings of the word \emph{bhujaṃga}, the snake and the lover. Since the word \emph{saṃgraha} has a wide range of meanings, it is difficult to construe it with the second meaning. One possibility is that the compound means something like ``the one who does not partake in intercourse as a lover.'' This would be a description of again a renouncer. Nonetheless, it appears safe to assume that the technique of \emph{śleṣa} is used in these three attributes.

%The author of this stanza might have had a slightly sectarian attitude toward \Saivism. The attributes he says \Sankara\ does not have (wearing ashes and surrounded by snakes) represent the fearful aspect of the god. %; they are typically mentioned in stories when someone has negative things to say about \Siva. %[version 2] % this does not seem justifiable. % or does it?
%This is quite a contrast to hagiographical accounts or the modern popular perception that strongly associate \Sankara\ with \Saivism. 
%The author of this stanza at least wanted to assure that when he paid homage to “\Sankara,” it was directed only to the teacher and not (additionally) to the god.

%This emphatic discrimination between the god \Siva\ and the human teacher \Sankara, despite sharing the name, is quite a contrast to the late or modern tradition (or tendency?)\ to associate \Sankara\ with \Siva. This stanza even has a slight negative tone toward \Siva. Some of the attributes that \Sankara\ is said not to have (wearing ashes, surrounded by snakes) in this stanza represent horrifying aspects of the god. It does not appear that the author of this stanza was very sympathetic toward \Siva\ worshippers. The origin of this stanza may go back to the period when the cult of \Sankara\ was not yet strongly associated with \Saivism. [version 1]

A few more things may be noted. One is that \Sankara\ is referred to as a \inidx{lord} (\emph{īśāna}). Īśāna is another popular name of \Siva\ and this one is not negated. The application of this name to \Sankara\ indicates that his authority as a spiritual leader had already been established when this stanza was composed. In addition, the author of this stanza used the expression \emph{apūrva} (the first), which implies that there were succeeding \Sankara s. This stanza presupposes the tradition of \Sankaracarya s already in place. We do not know when the \Sankara\ \matha s where their heads are called \Sankaracarya s were institutionalized, but their earliest possible establishment is during \Sankara's lifetime. Hence this stanza must be from at least a few generations after \Sankara's life time. %Yet, this order had not associated the founder to \Siva. The verse reads as though the author tried to avoid confusing the founder of the order with \Siva, particularly negating the fearful aspect of the god. Given that earlier followers, possibly including contemporary disciples, pay homage to \Visnu\ \citep[153]{hacker:relations}, the provenance of this verse might not be so far from them.

\subsection{Colophon of the work}
\label{WORKCOL}
A sub-colophon for the fourth pāda and the \colophon\ for the entire commentary follow:
\begin{quote}
{\dn गोविन्दभगवत्पूज्यपादशिष्यस्य परमहंसपरिव्राजकाचार्य्यस्य शंकरभगवतः कृतौ पातञ्जलयोगशास्त्रभाष्यविवरणे चतुर्त्थः पादः ॥ समाप्तञ्चेदं विवरणम् ॥}
\end{quote}
The first part, the sub-colophon, is a fairly standard \colophon\ to a work ascribed to the famous \Sankara.\footnote{See \citet{hacker:avsbhkl}.} At least, whether true or not, there is no ambiguity to whom the work is ascribed: not a random person called \Sankara, but the author of the \BSBh.\footnote{Another context in which this \colophon\ is significant is the discussion on the title of the work. For that, see \index{Wezler, Albrecht}\citet[17]{wezler:yvi1}} What troubles me slightly is the fact that there has been no \emph{iti} preceding these colophons. The end of the text is not clearly marked. This may or may not be a sign that some of the preceding concluding stanzas were not by the author of the bulk of the text itself.

\subsection{Post-colophon concluding stanza 1}
The \colophon\ is followed by two more stanzas in both \Tm\ and L\@. This first one is as follows:
\begin{verse}
%yātrākṛtas trijagatāṃ yadacintyaśaktileśān daśāhur iha mīnamukhāvatārān|
%kleśodbhavatrividhatāpaparamparārttās tan nāganāthaśayanaṃ śaraṇaṃ vrajāmaḥ||
{\dn यात्राकृतस्त्रिजगतां यदचिन्त्यशक्तिलेशान्दशाहुरिह मीनमुखावतारान्।\\
\kle शोद्भवत्रिविधतापपरम्परार्त्तास्तन्नागनाथशयनं शरणं व्रजामः॥}

{\footnotesize 
pāda a: yadacintya° \Tm] yada[ci]ntya° L\\
pāda b: °mukhāvatārān \Tmpc] °mukhāvatārāt \Tmac, °mukhā[\mist{3}n] L\\
pāda c: °tāpaparamparārttās L] °tāvavaramparārttās \Tm\ \\
pāda d: nāganātha° \Tm] nāgatnātha° L; vrajāmaḥ \Tm] prajāmah L

}
They say, in this world, that the ten \avatara s, starting with the fish, that keep the three worlds going, are fractions of his inconceivable power. \\%``The ten \avatara s, starting with the fish, in this world are fractions of his inconceivable power,'' say those who make the rounds of the three worlds,\\
We, being continuously afflicted by three kinds of sufferings originated from the impurity, seek refuge in him, the one who sleeps on the serpent king.
\end{verse}
This stanza in the \vasantatilaka\ meter has few textual problems (mostly exchanges between \emph{p} and \emph{v}, typical of the Malayalam script). It shows, in more or less a straightforward manner, an allegiance to \Visnu. 
%The only slightly tricky part in interpretation is the first word \emph{yātrākṛtaḥ}. Combined with \emph{trijagatām}, I take it to mean those who repeat the cycle of re-incarnations, i.e., ordinary beings. This goes best with the first word in the third pāda (\emph{kleśo\-dbhava\-tri\-vidha\-tāpa\-para\-mparā\-rttāḥ}), in the same case. 
It should be noted that \emph{tan} in the text in fact is \emph{tam}, an inflected word. The consonant \emph{m} is homogenized with the following \emph{n}.\footnote{An unlikely, yet not impossible, interpretation is that the \emph{tan-} is read as part of a compound \emph{ta\-nnāga\-nātha\-śayanam} (``his [\Visnu's] serpent king bed,'' i.e., \Sesa). If that (too) is meant, then the speaker of the stanza seeks refuge in \Sesa\ while praising \Visnu. Note that \Sesa\ and \Patanjali\ are sometimes identified.}
\subsection{Post-colophon concluding stanza 2}
\begin{verse}
{\dn योगेन चित्तस्य पदेन वाचाम्मलं शरीरस्य च वैद्यकेन।\\
यो ऽपाकरोत्तम्प्रवरम्मुनीनाम्पतञ्जलिम्प्राञ्जलिरानतो ऽस्मि॥} % upajāti

{\footnotesize 
pāda a: vācām \Tmpc (vācāṃ) L] vācā \Tmac\\
pāda b: malaṃ \emph{em.}]  mala° \Tm, malā° L; ca \Tm] tu L\\
pāda c: pākarot tam \Tmpc L] pākaroktam \Tmac; pravara \Tm\Lac] pra\ccs vara\cce\ L \\
pāda d: patañjalim L] \ccs v\cce ātañjalim \Tm, patañjali[m] L; °jalir ānato \Tmpc L] jalir \ccs X\cce ānato \Tm

}
The one who removed the impurity from the mind by yoga, from the speech by word, and from the body by medicine, \\
the best of the sages, \Patanjali, I bow to him, with my hands forming \emph{prāñjali}.
\end{verse}
This is, as \citet[(22)]{endoh:notes} points out, a stanza widely known among \inidx{grammarian}s.\footnote{See also \citet[xiv]{woods:hos} and \citet[503]{mbh1}. Cf.\ \citet[141--4]{meulenbeld:himia}} As such, it is doubtful that  this stanza was part of the original composition of the YVi.

This is the end of the shared text between \Tm\ and L\@. Following this stanza, \Tm\ has another \colophon\ stanza.

\section{An edition of the colophon stanza of T\kern -0.15em\raisebox{-.4ex}{\scriptsize m}}\label{TMCOL0405}
The \colophon\ stanza of \Tm\ involves many corrections. The text that was meant can be reconstructed as follows:
\begin{verse}
{\dn श्रीमत्पतञ्जलिमहर्षिवरप्रणीतश्रीयोगभाष्यहृदयप्रकटीकृदेतत्।\\
	श्रीशांकरं विवरणं समवर्त्तिशत्रो गार्ग्यो व्यलीलिखमहन्त्वदवाप्तिकामः॥} % vasantatilakā

{\footnotesize 
pāda a: °praṇīta° \Tmpc] °praṇītaṃ \Tmac\\
pāda c: śrīśāṃkaraṃ \Tmpc] śrīśāṃkaraṇaṃ \Tmac; 
        samavartti° \Tmpc] samavatti° \Tmac\\
pāda c-d: °śatro gārggyo \Tmac]  śatro \ccs [g $\rightarrow$ rgg]ā\cce \ins gā\ine rggyo \Tm\ (see note \ref{GATORGG})\\
pāda d: ahan tvad° \Tmpc] aha\ccs tvi\cce \ins n tva\ine\ccs ñca\cce d° \Tm

}

This [text] reveals the essence of the Yogabhāṣya composed by the venerable \Patanjali, the best of the great sages \\
I, \Gargya, have made it written, the commentary by venerable \Sankara, o slayer of \Samavartin, hoping to reach you.
\end{verse}
This stanza is followed by the very last word in the manuscript:
\begin{quote}
{\dn परक्रोडः}
\end{quote}
of which the first letter is originally written \emph{va} and corrected by a later hand (not inked) to \emph{pa}.

This \colophon\ informs us of several things. First, we learn the name, \Gargya, of the person who was one way or another responsible for the production of \Tm. He shows allegiance to \Samavartisatru, possibly \Visnu\ or \Siva: the slayer of \Samavartin\ (\Yama).\footnote{Dictionaries list names such as Yamāri, Yamaghna, Yamaripu as names of \Visnu. I am not aware of any work in which such names occur as such. On the other hand, a name similar in meaning, \Yamantaka\ is well-attested as a name of \Siva.} This much is certain.\footnote{Another potentially interesting point exists. The person who wrote this stanza considers \Patanjali\ to be the author of the \Yogabhasya.} 

However, there are many uncertainties. I am not certain if the name \Gargya\ belongs to the individual or it is his \gotra{} name. (For the possibility of it being a \gotra{} name, see below.) Nor am I certain in what capacity he was involved in the production of \Tm. The verb \emph{vi-likh-}, used in causative, usually means that the subject ordered the manuscript to be prepared. Most of the time, there is a mention of another name as the scribe.\footnote{These statements come from the observation of colophons of Nepalese manuscripts to which I have a good access to.}  But we do not have another name. We cannot take the name \Samavartisatru\ as that of the scribe. The word \emph{tvat-} “you” and the vocative case of \emph{samavartiśatro} compliment each other. Another curiosity is that the verb is in the first person. As far as the Nepalese manuscripts---to whose \colophon\ data I have an easier access---are concerned, when the verb \emph{vi-likh-} is used in the causative aorist, it is in the third person. The third person makes sense since the subject is not the person who was doing the actual writing. The stanza reads as though someone gave the scribe the text to be copied and the scribe reproduced it. I do not have enough experience to know if this was a common practice in Kerala.
%
%I cannot justifiably propose a better reading of the stanza to eliminate these oddities. %Yet, apart from some ambiguities, overall intention of the stanza presented here is more or less clear. 

It may further be noted that the constituted text of the stanza essentially follows the text after corrections in order to be more or less intelligible. The fact that there are so many corrections in the manuscript and the nature of the errors are also curious. 

Two of the errors originally made in the manuscript (\emph{praṇītaṃ} for \emph{praṇīta°} and \emph{śrīsāṃkaraṇaṃ} for \emph{śrīśāṃkaraṃ}) do not tell much of a story because they can occur in various circumstances;\footnote{They are surely mistakes. If the word \emph{śrīmatpatañjalimaharṣivarapraṇīta} ended there, then the \emph{etat} \ldots\ \emph{vivaraṇam} would be a work by \Patanjali. That would contradict with the adjective \emph{śrīśāṃkaram}. The word \emph{śrīśāṃkaraṇam} is unintelligible and would make the stanza unmetrical.} they are caused by mental expectations. Such writing mistakes could happen even when the person writing the manuscript was the composer himself.\footnote{As I see more colophons, I notice many mistakes are made in them even though the scribes were presumably writing their own words. I used to find it odd and suspected that colophons were being copied. However, it appears more likely that scribes are more used to copying than composing. Copying is a mechanical activity that one can train themselves to excel. However, composing and writing down at the same time is quite a different activity. I experience myself making more mistakes when writing something down especially by hand, at the same time thinking what to write. I find it not so rare to write my own name wrongly. I do not assume a \colophon\ being copied when mistakes, including the ones involving the scribe's own name, are found any more.} 
%
However, the mistake that was corrected to \emph{ahan tvad°} is rather inexplicable for any reason other than copying. The two signs of the original reading \emph{tvi} and \emph{ñca}, later replaced by \emph{ntva}, are completely unintelligible and hyper-metrical when read together. The only feature I find in these three signs is their graphical similarities. Unintelligible yet graphically similar signs appearing in a manuscript typically occurs when the scribe did not understand the nature of a correction. In this case, what happened was probably as follows: there was a manuscript where the part was first written \emph{vyalīlikham ahañ ca} and subsequently corrected to \emph{ahan tva°} by cancelling \emph{ñca} and inserting \emph{tva} in the margin; then a scribe (of \Tm) who was copying from this written text did not realize the cancellation and wrote down both, additionally misreading \emph{ntva} as \emph{tvi}. After \Tm\ was inked, someone further corrected it to the intended reading. The nature of errors being graphical, we can exclude the possibility that this error was produced by mishearing or mental associations. We can hence exclude the possibility that this \colophon\ stanza was being dictated (cf.\ the causative \emph{vyalīlikham} discussed above). 

Furthermore, I doubt that this stanza was present in the exemplar of \Tm. This is the very end of the manuscript where, in most cases, information relevant to the manuscript is recorded. There still is enough space on the side for about three lines of text and the side is the recto. It is not that the last folio was lost or that the scribe did not have enough space to write further. This stanza was intended to be the \colophon\ of the manuscript. The most likely scenario appears to be that it was prepared by the person who commissioned the manuscript (Gārgya) in written form and the scribe copied it.

I would now like to discuss the final word of the whole manuscript, \emph{pa\-ra\-kro\-ḍaḥ}.\footnote{I owe Yasutaka Muroya for calling my attention to the fact that \Parakroda\ is a name.} \Parakroda\ is the name of the family known to have produced \Jyesthadeva{} (active between 1500 and 1610)---the author of an astronomical work, \Yuktibhasa.\footnote{See \citet[187--188, and 207]{kvsarma:1991}. In addition to the post-colophon statements in a manuscript of a Malayalam commentary on the \Suryasiddhanta, Sarma cites Parayil Raman Namputiri, \emph{Nampūtirimār} (in Malayalam), Trichur 1918, p.~55, as a source of the information that \Jyesthadeva\ hailed from the family \Parakroda\ (Sanskritised name of Malayalam \emph{Paraṅṅoṭṭu}).} But \Parakroda\ may have been the name of the family deity as well.\footnote{This comes from the opening of a manuscript of the \Nyayabhasya\ to which Yasutaka Muroya drew my attention. The scribe of the manuscript shows allegiance to \emph{parakroḍadeva}. I do not have the details of the manuscript. Muroya kindly sent me a digital reproduction of the first side of the manuscript, and it starts with \emph{hari śrīgaṇapataye namaḥ śrīparakroḍadevāya namo namaḥ}. The first word here is \emph{hariḥ} whose \visarga\ has dropped before the combination of a sibilant and a semi-vowel.} I do not have a clear idea what function the word \emph{parakroḍaḥ} in the nominative singular at the end of \Tm\ serves. It would seem that it is there as an auspicious word.\footnote{This use would be analogous to the Nyāyabhāṣya manuscript mentioned in the previous note.} Yet again, the original scribe wrote \emph{varakroḍaḥ} instead of \emph{para°}. This is a typical error stemming from the graphical similarity between \emph{pa} and \emph{va} in the Malayalam script. As I have speculated above, the scribe might have been given what to write at the end of the manuscript and was copying it. He was probably not a \Parakroda\ himself; he was not familiar with the word.

It is in this connection that I cannot take the name \Gargya\ simply as a personal name. It can instead be the \gotra{} name of the person who commissioned \Tm. \Jyesthadeva\ mentioned above had some close relationship with another astronomer from Kerala, \NilakanthaSomayaji\ (born 1444 and composed the Tantrasaṃgraha, in 1500).\footnote{See \citet[187–8]{kvsarma:1991}. For the date of \NilakanthaSomayaji, see \citet[S 3]{sarmaandnarasimhan:tantrasangraha}.} They are said to have been fellow disciples of the same teacher and \Jyesthadeva\ based his \Yuktibhasa\ on the Tantrasaṃgraha. And this \NilakanthaSomayaji's \gotra{} was \Gargya.\footnote{See \citet[S 1]{sarmaandnarasimhan:tantrasangraha}.} Given the historically recorded relationship of a Gārgya and someone from the Parakroḍa family, it is hard not to consider the person referred to as \emph{gārgyaḥ} in \Tm's \colophon\ stanza as a \Gargya\ as well. On the other hand, it does not appear that \Jyesthadeva\ and \NilakanthaSomayaji\ came from the same family, and hence there is no need to consider that \Jyesthadeva's \gotra{}, too, was \Gargya, i.e., \Parakroda\ family's \gotra{} was \Gargya. I cannot deny the possibility that the nominative singular of the word \emph{parakroḍaḥ} might correspond to that of the word \emph{gārgyaḥ} in the \colophon\ stanza; but connecting these two words is not necessary.

There is too little information in the \colophon\ stanza to determine the provenance of \Tm\ precisely. The most we can say is that it may have had something to do with the circle of Kerala astronomers of the late 15th to early 16th century.


\section{An edition of colophon stanzas of L}\label{LCOL0405}
L has other dedicatory stanzas and a brief colophon. The text and a tentative translation of stanzas may be presented as follows:
%jitam pata\-ñja\-li\-muninā yena śreyortthisārtthasasyānām vihito harṣas tāpatritayahṛtā%
%(138v1)dharmmameghena| yas tatāna bhava\-mārgga\-laṃ\-ghi\-nām kleśakarmmamayagharmmatuttaye dharmma\-megha\-mu\-kha\-yogatoya śrīkeśavaprakābhaṭṭārakāṇāṃ śiṣyasya satyānandasya iya\-m pādañjalaṭī%
%(2)kā \HTG\ \HTG\ \HTG
 
%jitam patañjalimuninā yena śreyortthisārtthasasyānām| \\ % 12 18
%vihito harṣas tāpatritayahṛtā dharmmameghena|| % 12 15
\begin{verse}
{\dn जितम् पतञ्जलिमुनिना येन श्रेयोर्त्थिसार्त्थसस्यानाम्। \\
विहितो हर्षस्तापत्रितयहृता धर्म्ममेघेन॥}

Victorious is the sage \Patanjali, who caused joy by means of the cloud of Dharma (the \Dharmamegha [\samadhi]), the remover of heat (three torments), to the crops which are the multitude of those who seek better fortune.

%yas tatāna bhavamārggalaṃghinām % rathoddhatā
%kleśakarmmamayagharmmanuttaye|
%dharmmameghamukhayogatoyadaḥ
%śrī [patañjalimuniṃ taṃ astu śrīḥ]||


{\dn यस्ततान भवमार्ग्गलंघिनां
\kle शकर्म्ममयघर्म्मनुत्तये। \\
धर्म्ममेघमुखयोगतोय\ldots}

[[I salute \Patanjali\ldots]] who expanded [[\ldots\ that produces]] water that is yoga [pouring] from the \Dharmamegha\ (Dharma-cloud) in order to remove the heat, consisting of impurities and deeds, from those who are traversing the path of reincarnation.
\end{verse}
As can be seen, I did not make any emendation from the reading found in the manuscript. It does appear that the scribe wrote down what was meant to be written down. However, there are problems that raise questions about the circumstances under which these two stanzas were recorded in this manuscript.

The first and most glaring defect is that the second stanza is not complete. It lacks 12 syllables at the end. The stanza is interrupted in the middle of a word and the scribe simply continues to write the name of the owner (not certain; see below) after \emph{°yogatoya}. The scribe did not leave any indication that something was missing.

Another outward problem is a metrical defect in the first stanza. It starts with a \emph{ja}-\gana\ ({\applesymbols ⏑\,–\,⏑}). This is a metrical defect in the \arya\ meter.%meant to be in the \arya\ meter, has a metrical defect. It starts with syncopation, in other words, the first \gana\ is a \emph{ja}-\gana\ ({\applesymbols ⏑\,–\,⏑}). An \arya\ stanza should not have that rhythm in the odd numbered \gana s (the first to third syllables, 7th to 9th syllables, and so on).
\footnote{Harunaga Isaacson pointed this out to me.}
% The second stanza lacks of 12 syllables at the end. We could give the author(s) of the stanzas the benefit of the doubt: transmission errors (which I have just said is unlikely). Still, the stanzas are not well composed. 
%
This\labelh{EDITALT1} was probably the reason why the editors of the 1952 edition silently emended the first \pada\ of this stanza to \emph{sa jayati patañjalimuniḥ} \citep[370]{yvi}, which means the same thing, is metrically sound, and has the added benefit of having the corresponding pronoun to the relative pronoun \emph{yena}, the first word of the second \pada. 

The first stanza also appears clumsy for having so many words in the instrumental singular, masculine (\emph{patañjalinā}, \emph{yena}, \emph{dharmameghena}, and \emph{tāpatritayahṛtā}). The first three each occupy different syntactic functions and the last modifies the fourth. In particular, I find the composition ``\emph{yena \ldots\ harṣo vihito \ldots\ dharmamethena}'' (the words rearranged for clarity) to mean ``who caused the joy be means of \emph{dharmamegha}'' not clean. There the words \emph{yena} and \emph{dharmameghena} refer to two different things in the same clause: the agent of the action and an instrument, respectively.
%The first \arya\ stanza has three things in the instrumental singular, masculine, that serve different functions, viz., \emph{patañjalinā}, \emph{yena} and \emph{dharmameghena} (\emph{tāpatritayahṛtā} modifies this). The first, \emph{patañjalinā}, functions as the correlative of \emph{yena}. An explicit use of the correlative (\emph{tena}) is omitted, which in itself is fine, but \emph{patañjlinā} does not have to be in the instrumental case, somewhat confusingly. It can be in any case. This use of impersonal passive construction bothered the editors of the 1952 edition so much that they silently rewrote the first pāda into \emph{sa jayati patañjalimuniḥ} \citep[370]{yvi}. Furthermore, \emph{yena} and \emph{dharmameghena} refer to two different things in the same clause: the agent of the action and an instrument, respectively. Again, it is not hard to understand what is meant, but not an elegant composition. 

The compound \emph{śreyorthisārthasasyānām} also looks somewhat clumsy. What the author of the stanza wanted to convey was probably that \Patanjali\ promised joy to the \emph{śreyo\-rthin}s in the form of crops or the fruits achieved by the heat-reducing, i.e., rain producing, cloud, viz., the \Dharmamegha\ \samadhi. The author of the stanza is apparently playing on the word \emph{megha} (cloud) part of the \Dharmamegha\ and the meanings of the word \emph{tāpa} (heat and suffering). The topos behind this is the autumnal cloud that produces rain, relieving people of the summer heat, and contributing to  agriculture. However, with the compound \emph{śreyo\-rthi\-sārtha\-sasyānām} ending with the genitive plural, followed by \emph{harṣa}, the natural interpretation is that the people are compared to crops and they experience joy. This probably is not what was meant. I would think people rejoice due to the removal of heat and the crops, both given by the autumnal cloud. %This interpretation still results in a clumsy sense because it is unclear whether the joy comes from the success (crops) or the fact that people are relieved from the heat/sufferings. 
I would attribute the clumsiness to the inability of the author of this stanza to formulate his idea precisely.

The second stanza in the \rathoddhata\ meter is, as has been mentioned, short of one syllable and one pāda.\footnote{\citet[370]{yvi} supplies \emph{°dam [patañjalim ṛṣiṃ praṇato ‘smi]}. Although\labelh{EDITALT2} the editors indicate that only the last pāda was their reconstruction by brackets, they in fact added \emph{°dam} to have \emph{°toyadam}.} The parts enclosed in the double brackets in the translation are what I conjecture to have been in the missing part. Interestingly, the topos of the stanza significantly overlaps with that of the previous one. This stanza is about \Patanjali\ who rescues those who are suffering from the karmic existence by means of yoga; the suffering is compared to the (summer) heat; and the \Dharmamegha\ \samadhi\ is compared to the (autumnal) cloud that produces rain. Seeing the word \emph{tatāna} ``expanded,'' I would speculate that there was also an imagery of the serpent king \Patanjali\ who expands the hood to cause rain. However, as the editors of the 1952 edition conjectured, the most fitting syllable to supply at the end of the c-pāda indeed is \emph{da(m)}. There are not many words that have only one syllable after \emph{toya}. But then the whole compound (\emph{dharmameghamukhayogatoyada}) becomes rather unintelligible. The relative clause would then mean something like ``the one who spread the cloud that is yoga where the Dharmamegha is the foremost.'' It is not impossible but not good. In addition, there is no sign that the omission in the manuscript was a copying mistake. The best explanation for the lack of 12 syllables appears to be that the scribe (author of the stanzas himself?) just gave up finishing the stanza. That the last two stanzas in L convey more or less the same content contributes to this suspicion. It is as if the author of these stanzas were experimenting with different expressions to convey the same sense and gave up on finishing the second attempt.

The puzzling conclusion of L continues with the final statement. Here is the text without any corrections:\labelh{CFSGCLP1}

\begin{quote}
\emph{śrīkeśavaprakā\emph{[sic]}bhaṭṭārakāṇāṃ śiṣyasya satyānandasya i\emph{[sic]}yam pāda\emph{[sic]}\-ñja\-la\-ṭīkā} 

(This [manuscript of the] Pātañjaliṭīkā belongs to \Satyananda, a disciple of Keśavaprakāśabhaṭṭāraka).
\end{quote}
I note three problems (one is minor). First of all, the teacher's name should be \emph{śrīkeśavaprakāśa\-bhaṭṭāraka}. The manuscript omits \emph{śa} in \emph{°prakāśa°}. Also, it misspells \emph{pādañjala°} for \emph{pātañjala°}. This is most likely an influence from the scribe's native language (presumably Malayalam) where a consonant is voiced or unvoiced according to the position. Another minor issue is the lack of the \sandhi\ between \emph{satyānandasya} and \emph{iyam}; if the \sandhi\ is applied, the text would read \emph{°nandasyeyam}. I am not certain if \Satyananda\ was the scribe or the owner of L\@. (It is also possible that he was the scribe and the owner.) At least, we can gather that he was closely involved with L and his teacher was Keśavaprakāśa. I consider this \colophon\ to be genuine, as opposed to being copied from elsewhere. He was thinking of the text word by word (lack of the \sandhi), he wrote down what he would pronounce rather than what he saw (\emph{pādañjala} rather than \emph{pātañjala}). He might have been losing focus and was in a hurry with the prospect of finishing the manuscript (dropping a syllable from the teacher's name).

\section{Notes on the 1952 edition}
As has been noted, in the closing part of the YVi the editors of the 1952 edition introduced some changes to the text, mostly silently.\footnote{See pp.~\pageref{EDITALT0}, \pageref{EDITALT1} and \pageref{EDITALT2}.} 
%Without access to manuscripts, the reader would not know what they did. The editors 
They made another, more significant, change to the text without telling the reader: they moved around the stanzas and the colophon.

According to the two principle manuscripts, the end of the YVi consists of the following elements. First, both have:
\begin{enumerate}
\item a \sragdhara\ stanza, a prayer to \Krsna
\item a (corrupt) śloka(?) paying homage to \Patanjali
\item a \viyogini\ stanza paying homage to \Sankara
\item \colophon\ to the work
\item a \vasantatilaka\ stanza expressing allegiance to \Visnu, and 
\item an \upajati\ stanza paying homage to \Patanjali
\end{enumerate}
then \Tm\ has the \colophon\ stanza. L further has:
\begin{enumerate}
\setcounter{enumi}{6}
\item a defective \arya\ stanza saluting \Patanjali
\item an incomplete \rathoddhata\ stanza (presumably) saluting \Patanjali
\item the ownership statement
\end{enumerate}
The 1952 edition ends its text with most of the above but in a different order: 1, 5, 2, 6, 7, 8, 3, and 4 (the colophon).\footnote{It does not have 9; this is understandable because it pertains only to the manuscript L\@.} The edition derives from the manuscript M, which is a copy of L,\footnote{See pp.~\pageref{l2m2a}\,ff.} and M has the same ending as L\@. So, it was the editors of the 1952 edition who rearranged the text. The result was that the YVi ends with: two stanzas (1 and 5) dedicated to \Visnu, four stanzas (2, 6, 7 and 8) to \Patanjali, one stanza (3) to \Sankara, and the colophon (4). We can easily infer the motivation behind the change. Apparently the editors were not happy with the original arrangement because the objects of dedication went back and forth; usually, stanzas dedicated to one object appear in succession. The editors tried to \emph{fix} this perceived peculiarity. They did not, however, take into account if there is any reasonable explanation how the arrangement they imagined could have been mangled so badly. I doubt there is any. They in fact assumed that all the concluding stanzas were part of the text and found the arrangement strange. I, on the other hand, find the assumption not warranted. 
% WThis would be true if a same person was paying homage to various objects.  I certainly cannot explain the process and I do not think the rearrangement is necessary. Rather, it is not called for. %account for was how one could explain the peculiarity occurring under the same assumption. They should have reflected on their assumption.%but there is no need to assume that. Moreover, I cannot conceive how the ending in the 1952 edition would have changed to the one found in the manuscripts. 

\section{Layers in the ending of the YVi manuscripts}
I rather view the seemingly peculiar arrangement of dedicatory stanzas in the end of the YVi as a sign of the following: the ending grew as manuscripts were copied because different scribes added their own dedicatory/colophon material without (completely) removing their predecessors'. 
%I rather find the arrangement of stanzas in manuscripts to be a sign that the ending of the text grew in size by adding more elements to the end. 
There are several other signs that the ending of the YVi manuscripts consists of several layers.

%It is a norm that Sanskrit manuscripts convey more than the text for which they were produced. The additional material may be of two kinds: the one that pertains to the text (title, author, etc.) and the other that pertains to the manuscript (various information about the scribe, copying circumstances). may be divided roughly into the one about the text \ and the other about the manuscript . It is another norm that scribes want to start and end their manuscripts with auspicious symbols, phrases (both the beginning and the end) or even verses (in the end). Sometimes the borders between the text proper, colophons, remarks about the text itself and those about the manuscript and their provenances (e.g., did the author write the colophons between chapters or did someone else write them?), are not clear. It can happen that some portion of the additional material for a particular manuscript is coped in a subsequent manuscript. If such a process repeats, succeeding manuscripts of the same text may have a bigger and bigger ending.\footnote{This does not seem to apply to the beginning of manuscripts. The format for the beginning of manuscripts has little variations.} The ending of the YVi manuscript appears to be such a case. Luckily for us, the process left some marks.
%Primarily, what we find in what order at the end of the YVi manuscripts It is evident that the ending of the YVi manuscripts had layers that indicate the transmission process (adding more to the end) of the text.

%The outermost layer in the ending of the YVi manuscripts is relatively clearly visible because 

One immediately sees that the texts of \Tm\ and L deviate at one point. That is after the item numbered 6, using the list in the previous section. This tells us several things.
%This is an indication that anything that follows 6 in \Tm\ or L were introduced somewhere, in other words, they are not essential part of the text. 
The first is that items 1 to 6 were in the latest common ancestor of the two manuscripts, $\alpha$ (\Tm's exemplar).\footnote{See pp.~\pageref{extmandl}\,ff.}

As far as \Tm\ is concerned, it is most likely that the colophon stanza of \Tm\ is unique to it. Similarly, the ownership statement of L (item 9) is probably unique to L\@. The two faulty stanzas (7 and 8) of L may well belong to the same layer, judging from the curious quality (full of errors) shared by the stanzas and the ownership statement. Even if they were in $\alpha$,\footnote{There are various possibilities with regard to the origin of the two stanzas. It is not inconceivable that stanzas 7 and 8, in their entirety or partially, were in $\alpha$ and the scribe of \Tm\ dropped the part thinking that it pertained only to his exemplar and not something he had to reproduce. The faulty nature of those stanzas, if they were in the exemplar, might have contributed to the decision. Also, we cannot preclude that items 7 and 8 were added during the transmission between $\alpha$ and L since there is a possibility that there were intervening manuscripts between $\alpha$ and L (see p.~\pageref{LandALPHA}). Other possibilities include: that the colophon stanza found in \Tm\ was in $\alpha$ but it was replaced in L while it survived in \Tm; and that $\alpha$ had more text that was dropped both in \Tm\ and L\@.} 
%We should still be careful. The exact provenance of the last two stanzas could be a little complex. 
%
%We cannot, however, state that $\alpha$ had only up to item 6.   Nonetheless 
the scribe of \Tm\ felt them auxiliary enough to drop them. We can, therefore, draw a border after stanza 6.

%To summarize, items 7 and 8 in the above list may be introduced at the end of L by the scribe. Yet there are some possibilities that they were, at once or gradually, introduced to the end of the YVi text transmission in any of the ancestors of L, including $\alpha$. Even if any part of those two stanzas was in $\alpha$, . The colophon stanza of \Tm\ and the ownership statement of L (item 9) were introduced by the respective scribe.

Another border is after the colophon (4), revealing a layer consisting of stanzas 5 and 6. Unless there is a compelling reason to suggest otherwise, I consider anything after the colophon to have been added by a scribe. The scribe does not have to be the one who produced the manuscript at hand. Even if later scribes reproduced what an earlier scribe had added, it is still by a scribe. What we have here appears to be such a case. I find no reason to consider the two stanzas not to be scribal. Stanzas 5 and 6 express allegiance to \Visnu\ and pay homage to \Patanjali\ respectively. This is consistent with what a scribe would do. He was expressing allegiance to his tutelary deity and paying homage to the author of the \PYS\ on which the YVi is a commentary. %This is why the homages to them appear to be interrupted in the series of stanzas at the end of the YVi manuscripts. There is no reason to move these stanzas around.

%They were probably additions during the transmission by a scribe because they follow the colophon. In general, anything that comes after the phrase \emph{samāptam \ldots} in manuscripts is very fluid. I usually consider the text there only relevant to the manuscript unless there is evidence suggesting otherwise, in this case the fact that some of the following text is shared in two manuscripts. A scribe may add information about himself and/or the manuscript, or he might write auspicious stanzas. He may replace what was in his exemplar, or he may reproduce some part of it. This probably is such a case. Auspicious stanzas were reproduced. %In this case, however, at some point in the transmission, even as late as in $\alpha$ (\Tm's exemplar), the scribe copied the two stanzas that were meant to be unique in his exemplar.
%\footnote{There might have been more , but we have no means to tell.}
%That the homage to \Patanjali\ is interrupted by several stanzas, that the homage to \Visnu\ again (although the first one was to \Krsna) (after stanzas 1 and 2), is an indication that they were additions by a scribe. As the editors of the 1952 felt, an integral text is unlikely to go back and forth between \Visnu\ and \Patanjali\ as the object of homage. Furthermore, stanza 6 is not original but a widely circulated one. Such a stanza is fitting for a scribe to add at the end of his manuscript. This might indeed have been the last text in a manuscript, baring short statements about the owner, such as the one found in L\@. %A scribe may use such a cookie-cutter stanza but probably not an author, unless he was the first to compose the stanza that later became widely used. %Such a stanza can be easily added at the end of a manuscript by a scribe who wants to make his manuscript a bit more auspicious.

%Whether there were more layers before the colophon in the ending in the YVi transmission is unclear. 
Now we are left with three stanzas (items 1, 2, and 3) and the colophon. I do not see any clear borders in this part. Here I start examination from the first stanza.

The first stanza, a prayer to \Krsna, appears integral to the body of the text as has been discussed.\footnote{See pp.~\pageref{krsnaverse}--\pageref{patanjaliverseone}.} 

I do not see any strong positive or negative indications whether the second concluding stanza (numbered 2 above)\footnote{See p.~\pageref{patanjaliverseone}.} is part of the original composition. One thing that attracts my attention is that it pays homage to \Patanjali. The reason why \Patanjali\ should be paid homage to in the context of a commentary on the \PYS\ is that he is considered the author of it. However, while the author of the YVi mentions the name \Patanjali, he shows no awareness that \Patanjali\ is considered the author of the work he is commenting on. \Patanjali\ is referred to only as the author of the \VMBh, a \inidx{grammarian}.\footnote{See note~\ref{MLATHRS} on p.~\pageref{MLATHRS} and note \ref{PTJME} on p.~\pageref{PTJME}.} Thus the stanza gives a sense of incongruence with the rest of the work.\footnote{In this regard, all the stanzas at the end of the YVi manuscripts that pay homage to \Patanjali\ (2, 6, 7 and 8) are subject to the same doubt about their authenticity.} Moreover, as mentioned above (pp.~\pageref{patanjaliverseone}\,ff.), this stanza is, if it was a stanza at all, badly transmitted, showing contrast to the immediately preceding and following stanzas that were transmitted with few problems. This gives an impression of different origins. %I am slightly inclined toward the view that this stanza is not genuine part of the text.

 %This further implies that stanzas numbered 1, 2, and 3 all form separate layers.% Such a contrast was one of the things that became inaccessible in the 1952 edition because the orders of stanzas was shuffled.

%determine whether n other words, there is no decisive evidence whether any of the concluding stanzas were by the author himself or someone else.
%It to some degree depends on one's view on the authorship problem. If one considers that the YVi is not a work of the famous \Sankara, then the remaining elements, numbered 1--3 could be part of the original composition of the YVi.\footnote{I exclude number 4, the colophon, in this discussion for now. For, one could argue that the \colophon\ is always independent of the main body of the text. At least, matters are more complicated if we take the placement of the \colophon\ into account. See below.} On the other hand, if one's position is that the YVi is a work of the famous \Sankara, not everything up to 3 (dedicated clearly to the `first' \Sankara) is by himself. The point is that this part, particularly the stanza dedicated to \Sankara, cannot be evidence for or against \Sankara's authorship of the YVi. Rather, the authorship is a factor in determining whether there were more layers in the ending of the YVi.
%Yet I would like to point out that there are some doubts whether the second and third concluding stanzas are by the author of the YVi himself. 

The third stanza\footnote{See p.~\pageref{THETHIRDST}.} is in a similar situation: no clear boundaries around it, but some incongruence with the body of the work. It pays homage to \Sankara\ but in no place in the body of the YVi is he or his \BSBh\ mentioned. Additionally, if we assume that the text of the YVi continued up to this stanza, then the immediately following
%Another possibility is to assume a border just before the colophon and all the concluding stanzas are part of the work, i.e., written by the author of the YVi. In this case, 
colophon that ascribes the YVi to \Sankara\ will contradict what the stanza implies: the author is not \Sankara. The speaker of the stanza cannot be \Sankara\ because he praises \Sankara\ and is from generations later. Then the colophon ascribes the YVi, including this stanza, to \Sankara\ldots ?
%I also have some doubt about the next stanza that salutes \Sankara\ (item numbered 3) belonging to the same layer as the body of the text. 
%I already suspect a gap between the first concluding stanza and the second. Anything that follows is subject to the same degree of suspicion. %If the gap is real, then there are possibilities that the third stanza belongs to the same layer as the second stanza or even further outside layer.

Rather, the stanza and the colophon should be read together.
%On the other hand, the stanza and the colophon are in harmony if we don't consider the stanza to be part of the main text. %can be seen congruent with the following colophon. %, there is nothing odd about it.
%One advantage of viewing the third stanza not coming from the author himself is that we avoid conflict between it and the colophon (item numbered 4) that immediately follows. 
The stanza salutes \Sankara\ and the colophon attributes the work to him. The stanza can simply be understood as paying homage to the author of the text. If we do not wish to explain a contradiction, then the stanza should be placed outside the YVi, coming from the idea that \Sankara\ was the author of the YVi.
 %It is not possible to tell if these two small text portions were introduced by the same person or not, but 
%at least it was not \Sankara\ (no one salutes himself and the person who salutes is from generations later\footnote{See pp.~\pageref{THETHIRDST}\,ff.}) or the author of the YVi himself (the stanza/colophon author thinks the author is \Sankara, but he could not have written the stanza or the colophon). This is when we do not assume a border between the third concluding stanza and the colophon. Anyone can write them as long as the person is not \Sankara\ or the author of the YVi. Besides, It is still not a problem even if \Sankara\ was the author. 
%understanding them as coming from the same notion is more realistic. Note that if we consider that 
%
%The most natural way to read this stanza is alongside the colophon. 
%The third stanza and the colophon probably belong to the same layer in terms of the notion behind them. %Thus, again, my inclination is toward the stanza not being authorial.  %This scenario requires us to postulate a reason why this incongruence was brought about. (Did the mere mention of \Sankara in the stanza evoke the notion that the work was by him?) On the other hand, in the alternative scenario we have to assume a border between the text the author wrote and that which was added during the transmission of the text. This seems a little more likely because I already suspect such a border between the first concluding stanza and the second.

%In the 1952 edition, the order of the dedicatory stanzas was changed and some of the signs of the ending consisting of multiple layers were lost. That was unfortunate.

%has a problem with the \colophon\ that immediately follows . If the original author of the YVi wrote the concluding stanzas up to the \Sankara\ salutation and if the \colophon\ was added by someone else to satisfy convention, \footnote{I highly doubt the users of Sanskrit manuscripts, including those who copied them, considered colophons or sub-colophons authorial. They are often changed, omitted, expanded, or abbreviated more or less by scribes. This is observable in many manuscripts of the same text or sometimes even in apographs. There was little fidelity among scribes when it comes to colophons and sub-colophons.} there is an incongruence between the stanza and the colophon. That is, the stanza salutes \Sankara\ (the author of the stanza is not \Sankara) and the \colophon\ ascribes the text (up to the salutation) to \Sankara. This amounts to the \colophon\ ignoring what the preceding stanza indicates. It is more natural if the \Sankara\ salutation and the \colophon\ belong to the same layer and that they stem from the same notion that the first \Sankara\ was the author of the work.


%Another difficulty I have with seeing stanzas 1--3 as by the same author is the relationship between stanza 3 (dedicated to \Sankara) and the colophon. Let us for the moment assume that the stanza was written by the author of the whole YVi. He salutes \Sankara. Then the colophon ascribes the work (including the salutation) to \Sankara. This is not congruous. Something inexplicable must happen for the colophon (attributing the work to \Sankara) to replace the original colophon that recorded the real author. (There must have been another colophon for the current colophon to be there, and it must be another one for a colophon to be meaningful.) Rather, it is far more plausible that the stanza and the colophon forms a set. The stanza pays homage to the first \Sankara\ in the stanza and the colophon attributes the work to the same \Sankara. It is natural that the \Sankara\ salutation was done because someone thought that the work was by \Sankara. Again, whether the work was truly by \Sankara\ or not, the presence of \Sankara\ salutation is due to the ascription of the work to him. 

Finally the colophon.\footnote{See p.~\pageref{WORKCOL}.} It is about the text (YVi) and is not part of the text. We cannot treat it as in the same level as the body of the text. It is outside of the YVi. 
%The original author might have had some sort of colophon, but what we have is conventional.
%It is, including the phrase \emph{samāptam} \ldots, is just conventional. 
%The format, \emph{govindabhagavatpūjyapādaśiṣyasya} \ldots, is a standard one that ascribes a work to the famous \Sankara. 
%We rarely treat the colophon as by the author of the work. The colophon refers to the text and therefore is by nature from the third party perspective. In fact, as just has been discussed, it appears that the preceding stanza and the colophon come from the same notion that the YVi was by (\emph{that} venerable) \Sankara. This is from a third party perspective. 
%There might have been some sort of colophon in the first ever copy of the YVi, but 
%Regardless of the authorship, it is hard to imagine that the original author had this standard colophon already. %it already had this standard format. The colophon, in its current format, is the least likely part to have been in the original text. %thus from the third party, regardless if it is true. To be thorough, let us consider that the third stanza was in fact by the author of the YVi himself. In that case, the author of the colophon ignored or somehow mistook what the stanza implied and attributed the work to \Sankara. In that case, too, the colophon is not part of the YVi.\footnote{There is another possibility: the colophon, as well as the preceding stanza, was part of the original composition. In this case, the author deliberately attributed his commentary to \Sankara\ by saluting him and attributing the work to \Sankara. This is extremely unlikely. He should have realized that he was saluting himself.} %It is not inconceivable that someone wrote the whole YVi and deliberately attributed it to \Sankara\ by means of the colophon, . I do not think the colophon (item 4), including the phrase ``\emph{samāptam\ldots}'' belongs to the same layer as the main body of the text. The presumed speaker of those statements in the colophon, ``This is the work by \ldots. Here ends the [manuscript of] \ldots,'' should be the scribe. Furthermore, if the author is \Sankara, it is hardly conceivable that he calls himself \emph{bhagavat}. If the author is not \Sankara, then, the colophon is telling un
 
In the end, the first stanza (numbered 1 above) following the conclusion of the commentary is the only stanza that I feel compelled to attribute to the same author who composed the body of the commentary. For the rest, I have some reservations. The likelihood of the subsequent elements being additions during the transmission of the YVi increases as we go down the list. 

Here I would like to call attention to a lack of the word \emph{iti} that was last seen at the end of the commentary and before the first stanza. The word would  clearly mark the end of the text. One of the reasons why there are uncertainties about the extent of the original text, and why the editors of the 1952 edition struggled to make sense out of the arrangement of stanzas, is the lack of that word. The lack may indicate that there was more tampering in the ending than has been discussed here.\footnote{Although I do not have any statistics, that scribal additions being preserved in subsequent copies to such an extent as seen here may already be a rare case. On the one hand, this indicates a complex text transmission process. On the other hand, it might suggest the possibility that even the body of the text itself consists of several layers. If this is the case, it will be much harder to reveal layers than to do so in the ending. Note that at least for the part edited in this volume, I did not consider the possibility of text being tampered by later hands. See pp.~\pageref{what_text}\,ff.} Someone who was adding his own material might have inadvertently removed the clear marking of the end of the text, the word \emph{iti}.
% while adding his material, 
% be a sign of the convoluted situation in the end of the manuscripts may testify to 
%overall impression I get from examining the ending of the YVi is 
%a long period of text transmission of the YVi.
% that jumbled up the ending of the text. I do have sympathies for the editors of the 1952 who felt it necessary to move the stanzas and colophons around.

%\subsection{}


%
% \chapter{Materials for the Authorship Problem}
% \renewcommand{\rightmark}{\MakeUppercase{Materials for the Authorship Problem}}
% \thispagestyle{plain}
% \input{thework}
%
\backmatter
\renewcommand{\rightmark}{\MakeUppercase{Bibliography}}
% \renewcommand{\leftmark}{\MakeUppercase{Bibliography}}
\thispagestyle{plain}

\bibliography{mybibli} 

\thispagestyle{plain}
% \hfuzz=40pt
\renewcommand{\leftmark}{\MakeUppercase{INDEX}}
\renewcommand{\rightmark}{\MakeUppercase{INDEX}}
\printindex
\end{document}
